\documentclass[10pt,a4paper]{article}

\usepackage[T1]{fontenc}
\usepackage[utf8]{inputenc}
\usepackage[english]{babel}
\usepackage[sc]{mathpazo}
\usepackage{microtype}
\usepackage[a4paper,margin=2.1cm,top=1.9cm,bottom=1.9cm]{geometry}
\usepackage{amsmath,amssymb}
\usepackage{graphicx}
\usepackage{booktabs}
\usepackage{enumitem}
\usepackage[font=footnotesize,labelfont=bf,labelsep=colon,skip=3pt]{caption}
\usepackage[hidelinks]{hyperref}
\usepackage{xcolor}
\usepackage{float}
\usepackage{fontawesome5}
\usepackage[most]{tcolorbox}
\usepackage{adjustbox}

\graphicspath{{../out/d120_eh500_ez300_th45_g120_q0.1/}}
\newcommand{\G}{\mathrm{G}}
\newcommand{\A}{\mathrm{A}}
\newcommand{\Sl}{\mathrm{S}}
\newcommand{\Hloc}{\alpha}  % the code and the macros still call it H_loc
\newcommand{\Hk}{\Hloc(1000~\mathrm{s})}
\setlist{itemsep=4pt, topsep=6pt, parsep=0pt, leftmargin=1.3em}
\setlength{\parskip}{6pt}
\linespread{1.08}
\setlength{\parindent}{0pt}
\newcommand{\AlphaStepBeforeMax}{0.003}
\newcommand{\AlphaStepHourMax}{0.009}
\newcommand{\AlphaStepHourMin}{0.007}
\newcommand{\BandZ}{1.645}
\newcommand{\BelowEvery}{17}
\newcommand{\BelowFixed}{19}
\newcommand{\BelowLongest}{26}
\newcommand{\CarryAllFar}{$+$0.010}
\newcommand{\CarryAllFarBand}{[$+$0.006, $+$0.015]}
\newcommand{\CarryAllNear}{$+$0.001}
\newcommand{\CarryAllNearBand}{[$-$0.003, $+$0.004]}
\newcommand{\CarryBins}{4}
\newcommand{\CarryFarLabel}{$\ge80$~min}
\newcommand{\CarryHighMtnFar}{$+$0.026}
\newcommand{\CarryHighMtnFarBand}{[$+$0.019, $+$0.032]}
\newcommand{\CarryHighMtnNear}{$+$0.020}
\newcommand{\CarryHighMtnNearBand}{[$+$0.014, $+$0.026]}
\newcommand{\CarryHillsFar}{$+$0.024}
\newcommand{\CarryHillsFarBand}{[$+$0.011, $+$0.037]}
\newcommand{\CarryHillsNear}{$+$0.010}
\newcommand{\CarryHillsNearBand}{[$+$0.005, $+$0.016]}
\newcommand{\CarryLowMtnFar}{$+$0.027}
\newcommand{\CarryLowMtnFarBand}{[$+$0.021, $+$0.033]}
\newcommand{\CarryLowMtnNear}{$+$0.013}
\newcommand{\CarryLowMtnNearBand}{[$+$0.008, $+$0.018]}
\newcommand{\CarryNearLabel}{0--20~min}
\newcommand{\CarryPlainsFar}{$+$0.012}
\newcommand{\CarryPlainsFarBand}{[$+$0.008, $+$0.017]}
\newcommand{\CarryPlainsNear}{$-$0.000}
\newcommand{\CarryPlainsNearBand}{[$-$0.004, $+$0.004]}
\newcommand{\CarryWindowsAllFar}{3\,894}
\newcommand{\CarryWindowsAllNear}{6\,965}
\newcommand{\CarryWindowsHighMtnFar}{1\,199}
\newcommand{\CarryWindowsHighMtnNear}{1\,802}
\newcommand{\CarryWindowsHillsFar}{174}
\newcommand{\CarryWindowsHillsNear}{637}
\newcommand{\CarryWindowsLowMtnFar}{2\,002}
\newcommand{\CarryWindowsLowMtnNear}{3\,467}
\newcommand{\CarryWindowsPlainsFar}{418}
\newcommand{\CarryWindowsPlainsNear}{889}
\newcommand{\CarryWindowsSAll}{88\,654}
\newcommand{\CarryWindowsSHighMtn}{15\,286}
\newcommand{\CarryWindowsSHills}{16\,554}
\newcommand{\CarryWindowsSLowMtn}{34\,372}
\newcommand{\CarryWindowsSPlains}{15\,779}
\newcommand{\ChainAFirstPieces}{16\,736}
\newcommand{\ChainALastLag}{17\,190}
\newcommand{\ChainALastPieces}{131}
\newcommand{\ChainAMinusCohortMax}{0.007}
\newcommand{\ChainASHundred}{$+$0.003}
\newcommand{\ChainASHundredBand}{[$+$0.002, $+$0.004]}
\newcommand{\ChainASLong}{$-$0.005}
\newcommand{\ChainASLongBand}{[$-$0.010, $+$0.001]}
\newcommand{\ChainASNegative}{870}
\newcommand{\ChainASTen}{$+$0.007}
\newcommand{\ChainASTenBand}{[$+$0.006, $+$0.009]}
\newcommand{\ChainASThousand}{$-$0.003}
\newcommand{\ChainASThousandBand}{[$-$0.006, $-$0.000]}
\newcommand{\ChainASThreeTwenty}{$+$0.003}
\newcommand{\ChainASThreeTwentyBand}{[$+$0.001, $+$0.005]}
\newcommand{\ChainGAGrow}{170}
\newcommand{\ChainGAHundred}{$-$0.001}
\newcommand{\ChainGAHundredBand}{[$-$0.002, $-$0.000]}
\newcommand{\ChainGALong}{$+$0.026}
\newcommand{\ChainGALongBand}{[$+$0.021, $+$0.032]}
\newcommand{\ChainGATen}{$-$0.009}
\newcommand{\ChainGATenBand}{[$-$0.010, $-$0.008]}
\newcommand{\ChainGAThousand}{$+$0.027}
\newcommand{\ChainGAThousandBand}{[$+$0.025, $+$0.030]}
\newcommand{\ChainGAThreeTwenty}{$+$0.011}
\newcommand{\ChainGAThreeTwentyBand}{[$+$0.009, $+$0.012]}
\newcommand{\ChainGAZero}{150}
\newcommand{\ChainGFirstPieces}{17\,240}
\newcommand{\ChainGLastLag}{11\,450}
\newcommand{\ChainGLastPieces}{153}
\newcommand{\ChainGMinusCohortMax}{0.006}
\newcommand{\ChainGSGrow}{210}
\newcommand{\ChainGSHundred}{$+$0.002}
\newcommand{\ChainGSHundredBand}{[$+$0.001, $+$0.003]}
\newcommand{\ChainGSLong}{$+$0.022}
\newcommand{\ChainGSLongBand}{[$+$0.015, $+$0.028]}
\newcommand{\ChainGSSmall}{0.002}
\newcommand{\ChainGSTen}{$-$0.002}
\newcommand{\ChainGSTenBand}{[$-$0.004, $-$0.000]}
\newcommand{\ChainGSThousand}{$+$0.024}
\newcommand{\ChainGSThousandBand}{[$+$0.021, $+$0.027]}
\newcommand{\ChainGSThreeTwenty}{$+$0.014}
\newcommand{\ChainGSThreeTwentyBand}{[$+$0.012, $+$0.015]}
\newcommand{\ChainSFirstPieces}{52\,812}
\newcommand{\ChainSLastLag}{21\,060}
\newcommand{\ChainSLastPieces}{394}
\newcommand{\ChainSMinusCohortMax}{0.004}
\newcommand{\ChainSelFree}{3\,600}
\newcommand{\ChainSolidEnd}{2\,250}
\newcommand{\CohortPiecesFirst}{17\,240}
\newcommand{\CohortPiecesLast}{195}
\newcommand{\CovDraws}{400}
\newcommand{\CovNewMax}{99}
\newcommand{\CovNewMin}{88}
\newcommand{\CovNewSdMax}{1.43}
\newcommand{\CovNewSdMin}{0.93}
\newcommand{\CovOldMatchedMin}{99}
\newcommand{\CovOldMatchedSdMax}{2.7}
\newcommand{\CovOldMatchedSdMin}{1.5}
\newcommand{\CovOldRawMin}{59}
\newcommand{\CovReplicates}{400}
\newcommand{\CovReuse}{6.7}
\newcommand{\DecControlFlightsF}{1\,404}
\newcommand{\DecControlFlightsNone}{20\,405}
\newcommand{\DecControlsF}{8\,094}
\newcommand{\DecControlsNone}{135\,850}
\newcommand{\DecControlsR}{40\,637}
\newcommand{\DecCtlClimbF}{26.2}
\newcommand{\DecCtlClimbNone}{38.7}
\newcommand{\DecCtlClimbR}{27.0}
\newcommand{\DecCtlClimbRT}{26.5}
\newcommand{\DecCtlClimbRW}{26.5}
\newcommand{\DecCtlClimbT}{39.5}
\newcommand{\DecCtlClimbTW}{38.5}
\newcommand{\DecCtlClimbW}{37.9}
\newcommand{\DecCtlCphF}{6.83}
\newcommand{\DecCtlCphNone}{6.94}
\newcommand{\DecCtlCphR}{6.97}
\newcommand{\DecCtlCphRT}{6.89}
\newcommand{\DecCtlCphRW}{6.91}
\newcommand{\DecCtlCphT}{7.10}
\newcommand{\DecCtlCphTW}{7.09}
\newcommand{\DecCtlCphW}{6.91}
\newcommand{\DecCtlFoundF}{46.6}
\newcommand{\DecCtlFoundNone}{45.5}
\newcommand{\DecCtlFoundR}{47.2}
\newcommand{\DecCtlFoundRT}{47.4}
\newcommand{\DecCtlFoundRW}{46.4}
\newcommand{\DecCtlFoundT}{46.4}
\newcommand{\DecCtlFoundTW}{45.6}
\newcommand{\DecCtlFoundW}{44.5}
\newcommand{\DecCtlGlideF}{30.7}
\newcommand{\DecCtlGlideNone}{35.2}
\newcommand{\DecCtlGlideR}{30.9}
\newcommand{\DecCtlGlideRT}{31.0}
\newcommand{\DecCtlGlideRW}{30.6}
\newcommand{\DecCtlGlideT}{35.7}
\newcommand{\DecCtlGlideTW}{33.8}
\newcommand{\DecCtlGlideW}{33.7}
\newcommand{\DecCtlHKF}{0.897}
\newcommand{\DecCtlHKNone}{0.908}
\newcommand{\DecCtlHKR}{0.895}
\newcommand{\DecCtlHKRT}{0.896}
\newcommand{\DecCtlHKRW}{0.896}
\newcommand{\DecCtlHKT}{0.910}
\newcommand{\DecCtlHKTW}{0.902}
\newcommand{\DecCtlHKW}{0.902}
\newcommand{\DecCtlHQF}{0.868}
\newcommand{\DecCtlHQNone}{0.928}
\newcommand{\DecCtlHQR}{0.863}
\newcommand{\DecCtlHQRT}{0.865}
\newcommand{\DecCtlHQRW}{0.868}
\newcommand{\DecCtlHQT}{0.930}
\newcommand{\DecCtlHQTW}{0.920}
\newcommand{\DecCtlHQW}{0.919}
\newcommand{\DecCtlRateF}{1.22}
\newcommand{\DecCtlRateNone}{1.10}
\newcommand{\DecCtlRateR}{1.20}
\newcommand{\DecCtlRateRT}{1.21}
\newcommand{\DecCtlRateRW}{1.21}
\newcommand{\DecCtlRateT}{1.13}
\newcommand{\DecCtlRateTW}{1.13}
\newcommand{\DecCtlRateW}{1.10}
\newcommand{\DecCtlSearchF}{13.9}
\newcommand{\DecCtlSearchNone}{9.7}
\newcommand{\DecCtlSearchR}{13.4}
\newcommand{\DecCtlSearchRT}{13.7}
\newcommand{\DecCtlSearchRW}{13.6}
\newcommand{\DecCtlSearchT}{9.5}
\newcommand{\DecCtlSearchTW}{10.3}
\newcommand{\DecCtlSearchW}{10.4}
\newcommand{\DecCtlSinkF}{62}
\newcommand{\DecCtlSinkNone}{96}
\newcommand{\DecCtlSinkR}{63}
\newcommand{\DecCtlSinkRT}{62}
\newcommand{\DecCtlSinkRW}{62}
\newcommand{\DecCtlSinkT}{96}
\newcommand{\DecCtlSinkTW}{100}
\newcommand{\DecCtlSinkW}{99}
\newcommand{\DecDeltaClimbF}{$+$3.0}
\newcommand{\DecDeltaClimbFBand}{[$+$1.4, $+$4.5]}
\newcommand{\DecDeltaClimbNone}{$-$9.5}
\newcommand{\DecDeltaClimbNoneBand}{[$-$10.4, $-$8.7]}
\newcommand{\DecDeltaClimbR}{$+$2.2}
\newcommand{\DecDeltaClimbRBand}{[$+$0.9, $+$3.5]}
\newcommand{\DecDeltaClimbRT}{$+$2.6}
\newcommand{\DecDeltaClimbRTBand}{[$+$1.2, $+$4.1]}
\newcommand{\DecDeltaClimbRW}{$+$2.6}
\newcommand{\DecDeltaClimbRWBand}{[$+$1.2, $+$4.1]}
\newcommand{\DecDeltaClimbT}{$-$10.4}
\newcommand{\DecDeltaClimbTBand}{[$-$11.2, $-$9.5]}
\newcommand{\DecDeltaClimbTW}{$-$9.3}
\newcommand{\DecDeltaClimbTWBand}{[$-$10.1, $-$8.4]}
\newcommand{\DecDeltaClimbW}{$-$8.7}
\newcommand{\DecDeltaClimbWBand}{[$-$9.5, $-$7.8]}
\newcommand{\DecDeltaCphF}{$+$0.73}
\newcommand{\DecDeltaCphFBand}{[$+$0.37, $+$1.10]}
\newcommand{\DecDeltaCphNone}{$+$0.63}
\newcommand{\DecDeltaCphNoneBand}{[$+$0.48, $+$0.77]}
\newcommand{\DecDeltaCphR}{$+$0.59}
\newcommand{\DecDeltaCphRBand}{[$+$0.30, $+$0.89]}
\newcommand{\DecDeltaCphRT}{$+$0.67}
\newcommand{\DecDeltaCphRTBand}{[$+$0.34, $+$1.00]}
\newcommand{\DecDeltaCphRW}{$+$0.66}
\newcommand{\DecDeltaCphRWBand}{[$+$0.32, $+$0.99]}
\newcommand{\DecDeltaCphT}{$+$0.46}
\newcommand{\DecDeltaCphTBand}{[$+$0.31, $+$0.61]}
\newcommand{\DecDeltaCphTW}{$+$0.47}
\newcommand{\DecDeltaCphTWBand}{[$+$0.32, $+$0.62]}
\newcommand{\DecDeltaCphW}{$+$0.65}
\newcommand{\DecDeltaCphWBand}{[$+$0.50, $+$0.80]}
\newcommand{\DecDeltaFoundF}{$+$2.3}
\newcommand{\DecDeltaFoundFBand}{[$-$0.2, $+$4.9]}
\newcommand{\DecDeltaFoundNone}{$+$3.4}
\newcommand{\DecDeltaFoundNoneBand}{[$+$2.1, $+$4.7]}
\newcommand{\DecDeltaFoundR}{$+$1.7}
\newcommand{\DecDeltaFoundRBand}{[$-$0.4, $+$3.9]}
\newcommand{\DecDeltaFoundRT}{$+$1.6}
\newcommand{\DecDeltaFoundRTBand}{[$-$0.8, $+$3.9]}
\newcommand{\DecDeltaFoundRW}{$+$2.5}
\newcommand{\DecDeltaFoundRWBand}{[$+$0.1, $+$4.9]}
\newcommand{\DecDeltaFoundT}{$+$2.5}
\newcommand{\DecDeltaFoundTBand}{[$+$1.2, $+$3.8]}
\newcommand{\DecDeltaFoundTW}{$+$3.4}
\newcommand{\DecDeltaFoundTWBand}{[$+$2.1, $+$4.7]}
\newcommand{\DecDeltaFoundW}{$+$4.4}
\newcommand{\DecDeltaFoundWBand}{[$+$3.1, $+$5.7]}
\newcommand{\DecDeltaGlideF}{$-$0.9}
\newcommand{\DecDeltaGlideFBand}{[$-$1.7, $-$0.1]}
\newcommand{\DecDeltaGlideNone}{$-$5.4}
\newcommand{\DecDeltaGlideNoneBand}{[$-$6.0, $-$4.8]}
\newcommand{\DecDeltaGlideR}{$-$1.1}
\newcommand{\DecDeltaGlideRBand}{[$-$1.8, $-$0.4]}
\newcommand{\DecDeltaGlideRT}{$-$1.2}
\newcommand{\DecDeltaGlideRTBand}{[$-$2.0, $-$0.4]}
\newcommand{\DecDeltaGlideRW}{$-$0.8}
\newcommand{\DecDeltaGlideRWBand}{[$-$1.6, $-$0.1]}
\newcommand{\DecDeltaGlideT}{$-$5.9}
\newcommand{\DecDeltaGlideTBand}{[$-$6.4, $-$5.3]}
\newcommand{\DecDeltaGlideTW}{$-$4.0}
\newcommand{\DecDeltaGlideTWBand}{[$-$4.5, $-$3.6]}
\newcommand{\DecDeltaGlideW}{$-$3.9}
\newcommand{\DecDeltaGlideWBand}{[$-$4.4, $-$3.4]}
\newcommand{\DecDeltaHKF}{$+$0.010}
\newcommand{\DecDeltaHKFBand}{[$+$0.004, $+$0.017]}
\newcommand{\DecDeltaHKNone}{$-$0.001}
\newcommand{\DecDeltaHKNoneBand}{[$-$0.006, $+$0.004]}
\newcommand{\DecDeltaHKR}{$+$0.012}
\newcommand{\DecDeltaHKRBand}{[$+$0.007, $+$0.017]}
\newcommand{\DecDeltaHKRT}{$+$0.011}
\newcommand{\DecDeltaHKRTBand}{[$+$0.006, $+$0.017]}
\newcommand{\DecDeltaHKRW}{$+$0.011}
\newcommand{\DecDeltaHKRWBand}{[$+$0.005, $+$0.017]}
\newcommand{\DecDeltaHKT}{$-$0.003}
\newcommand{\DecDeltaHKTBand}{[$-$0.008, $+$0.002]}
\newcommand{\DecDeltaHKTW}{$+$0.005}
\newcommand{\DecDeltaHKTWBand}{[$+$0.000, $+$0.010]}
\newcommand{\DecDeltaHKW}{$+$0.005}
\newcommand{\DecDeltaHKWBand}{[$+$0.000, $+$0.010]}
\newcommand{\DecDeltaHQF}{$+$0.021}
\newcommand{\DecDeltaHQFBand}{[$+$0.008, $+$0.034]}
\newcommand{\DecDeltaHQNone}{$-$0.039}
\newcommand{\DecDeltaHQNoneBand}{[$-$0.052, $-$0.026]}
\newcommand{\DecDeltaHQR}{$+$0.026}
\newcommand{\DecDeltaHQRBand}{[$+$0.016, $+$0.037]}
\newcommand{\DecDeltaHQRT}{$+$0.024}
\newcommand{\DecDeltaHQRTBand}{[$+$0.012, $+$0.036]}
\newcommand{\DecDeltaHQRW}{$+$0.021}
\newcommand{\DecDeltaHQRWBand}{[$+$0.009, $+$0.034]}
\newcommand{\DecDeltaHQT}{$-$0.041}
\newcommand{\DecDeltaHQTBand}{[$-$0.054, $-$0.029]}
\newcommand{\DecDeltaHQTW}{$-$0.031}
\newcommand{\DecDeltaHQTWBand}{[$-$0.044, $-$0.018]}
\newcommand{\DecDeltaHQW}{$-$0.030}
\newcommand{\DecDeltaHQWBand}{[$-$0.042, $-$0.018]}
\newcommand{\DecDeltaRateF}{$-$0.044}
\newcommand{\DecDeltaRateFBand}{[$-$0.085, $-$0.002]}
\newcommand{\DecDeltaRateNone}{$+$0.077}
\newcommand{\DecDeltaRateNoneBand}{[$+$0.051, $+$0.102]}
\newcommand{\DecDeltaRateR}{$-$0.032}
\newcommand{\DecDeltaRateRBand}{[$-$0.069, $+$0.005]}
\newcommand{\DecDeltaRateRT}{$-$0.039}
\newcommand{\DecDeltaRateRTBand}{[$-$0.080, $+$0.003]}
\newcommand{\DecDeltaRateRW}{$-$0.039}
\newcommand{\DecDeltaRateRWBand}{[$-$0.079, $-$0.000]}
\newcommand{\DecDeltaRateT}{$+$0.041}
\newcommand{\DecDeltaRateTBand}{[$+$0.016, $+$0.066]}
\newcommand{\DecDeltaRateTW}{$+$0.043}
\newcommand{\DecDeltaRateTWBand}{[$+$0.018, $+$0.069]}
\newcommand{\DecDeltaRateW}{$+$0.076}
\newcommand{\DecDeltaRateWBand}{[$+$0.049, $+$0.102]}
\newcommand{\DecDeltaSearchF}{$+$0.8}
\newcommand{\DecDeltaSearchFBand}{[$-$0.5, $+$2.1]}
\newcommand{\DecDeltaSearchNone}{$+$5.0}
\newcommand{\DecDeltaSearchNoneBand}{[$+$4.2, $+$5.8]}
\newcommand{\DecDeltaSearchR}{$+$1.3}
\newcommand{\DecDeltaSearchRBand}{[$+$0.2, $+$2.4]}
\newcommand{\DecDeltaSearchRT}{$+$1.0}
\newcommand{\DecDeltaSearchRTBand}{[$-$0.3, $+$2.2]}
\newcommand{\DecDeltaSearchRW}{$+$1.1}
\newcommand{\DecDeltaSearchRWBand}{[$-$0.1, $+$2.3]}
\newcommand{\DecDeltaSearchT}{$+$5.2}
\newcommand{\DecDeltaSearchTBand}{[$+$4.4, $+$6.0]}
\newcommand{\DecDeltaSearchTW}{$+$4.4}
\newcommand{\DecDeltaSearchTWBand}{[$+$3.6, $+$5.1]}
\newcommand{\DecDeltaSearchW}{$+$4.3}
\newcommand{\DecDeltaSearchWBand}{[$+$3.6, $+$5.1]}
\newcommand{\DecDeltaSinkF}{$+$6}
\newcommand{\DecDeltaSinkFBand}{[$+$2, $+$10]}
\newcommand{\DecDeltaSinkNone}{$-$28}
\newcommand{\DecDeltaSinkNoneBand}{[$-$30, $-$25]}
\newcommand{\DecDeltaSinkR}{$+$5}
\newcommand{\DecDeltaSinkRBand}{[$+$2, $+$8]}
\newcommand{\DecDeltaSinkRT}{$+$6}
\newcommand{\DecDeltaSinkRTBand}{[$+$2, $+$9]}
\newcommand{\DecDeltaSinkRW}{$+$6}
\newcommand{\DecDeltaSinkRWBand}{[$+$2, $+$10]}
\newcommand{\DecDeltaSinkT}{$-$28}
\newcommand{\DecDeltaSinkTBand}{[$-$30, $-$26]}
\newcommand{\DecDeltaSinkTW}{$-$32}
\newcommand{\DecDeltaSinkTWBand}{[$-$34, $-$29]}
\newcommand{\DecDeltaSinkW}{$-$31}
\newcommand{\DecDeltaSinkWBand}{[$-$33, $-$29]}
\newcommand{\DecGroupClimb}{29.2}
\newcommand{\DecGroupCph}{7.56}
\newcommand{\DecGroupFound}{48.9}
\newcommand{\DecGroupGlide}{29.8}
\newcommand{\DecGroupHK}{0.907}
\newcommand{\DecGroupHQ}{0.889}
\newcommand{\DecGroupRate}{1.17}
\newcommand{\DecGroupSearch}{14.7}
\newcommand{\DecGroupSink}{68}
\newcommand{\DecK}{50}
\newcommand{\DecNoneClimb}{$-$9.5}
\newcommand{\DecNoneClimbBand}{[$-$10.4, $-$8.7]}
\newcommand{\DecNoneCph}{$+$0.63}
\newcommand{\DecNoneCphBand}{[$+$0.48, $+$0.77]}
\newcommand{\DecNoneFound}{$+$3.4}
\newcommand{\DecNoneFoundBand}{[$+$2.1, $+$4.7]}
\newcommand{\DecNoneGlide}{$-$5.4}
\newcommand{\DecNoneGlideBand}{[$-$6.0, $-$4.8]}
\newcommand{\DecNoneHK}{$-$0.001}
\newcommand{\DecNoneHKBand}{[$-$0.006, $+$0.004]}
\newcommand{\DecNoneHQ}{$-$0.039}
\newcommand{\DecNoneHQBand}{[$-$0.052, $-$0.026]}
\newcommand{\DecNoneRate}{$+$0.077}
\newcommand{\DecNoneRateBand}{[$+$0.051, $+$0.102]}
\newcommand{\DecNoneSearch}{$+$5.0}
\newcommand{\DecNoneSearchBand}{[$+$4.2, $+$5.8]}
\newcommand{\DecNoneSink}{$-$28}
\newcommand{\DecNoneSinkBand}{[$-$30, $-$25]}
\newcommand{\DecPhaseControlsF}{3\,210}
\newcommand{\DecPhaseControlsNone}{46\,267}
\newcommand{\DecPhiClimbR}{$-$12.0}
\newcommand{\DecPhiClimbRBand}{[$-$13.4, $-$10.7]}
\newcommand{\DecPhiClimbT}{$+$0.2}
\newcommand{\DecPhiClimbTBand}{[$-$0.1, $+$0.5]}
\newcommand{\DecPhiClimbW}{$-$0.7}
\newcommand{\DecPhiClimbWBand}{[$-$1.0, $-$0.3]}
\newcommand{\DecPhiCphR}{$-$0.11}
\newcommand{\DecPhiCphRBand}{[$-$0.40, $+$0.18]}
\newcommand{\DecPhiCphT}{$+$0.05}
\newcommand{\DecPhiCphTBand}{[$-$0.02, $+$0.11]}
\newcommand{\DecPhiCphW}{$-$0.04}
\newcommand{\DecPhiCphWBand}{[$-$0.13, $+$0.05]}
\newcommand{\DecPhiFoundR}{$+$1.4}
\newcommand{\DecPhiFoundRBand}{[$-$0.4, $+$3.2]}
\newcommand{\DecPhiFoundT}{$+$0.6}
\newcommand{\DecPhiFoundTBand}{[$+$0.0, $+$1.1]}
\newcommand{\DecPhiFoundW}{$-$0.9}
\newcommand{\DecPhiFoundWBand}{[$-$1.6, $-$0.1]}
\newcommand{\DecPhiGlideR}{$-$3.8}
\newcommand{\DecPhiGlideRBand}{[$-$4.5, $-$3.0]}
\newcommand{\DecPhiGlideT}{$+$0.2}
\newcommand{\DecPhiGlideTBand}{[$+$0.1, $+$0.4]}
\newcommand{\DecPhiGlideW}{$-$1.0}
\newcommand{\DecPhiGlideWBand}{[$-$1.2, $-$0.7]}
\newcommand{\DecPhiHKR}{$-$0.010}
\newcommand{\DecPhiHKRBand}{[$-$0.015, $-$0.004]}
\newcommand{\DecPhiHKT}{$+$0.001}
\newcommand{\DecPhiHKTBand}{[$-$0.000, $+$0.002]}
\newcommand{\DecPhiHKW}{$-$0.003}
\newcommand{\DecPhiHKWBand}{[$-$0.005, $-$0.001]}
\newcommand{\DecPhiHQR}{$-$0.059}
\newcommand{\DecPhiHQRBand}{[$-$0.073, $-$0.045]}
\newcommand{\DecPhiHQT}{$+$0.001}
\newcommand{\DecPhiHQTBand}{[$-$0.001, $+$0.004]}
\newcommand{\DecPhiHQW}{$-$0.003}
\newcommand{\DecPhiHQWBand}{[$-$0.007, $+$0.001]}
\newcommand{\DecPhiRateR}{$+$0.098}
\newcommand{\DecPhiRateRBand}{[$+$0.062, $+$0.133]}
\newcommand{\DecPhiRateT}{$+$0.020}
\newcommand{\DecPhiRateTBand}{[$+$0.009, $+$0.030]}
\newcommand{\DecPhiRateW}{$+$0.003}
\newcommand{\DecPhiRateWBand}{[$-$0.008, $+$0.013]}
\newcommand{\DecPhiSearchR}{$+$3.6}
\newcommand{\DecPhiSearchRBand}{[$+$2.6, $+$4.7]}
\newcommand{\DecPhiSearchT}{$+$0.1}
\newcommand{\DecPhiSearchTBand}{[$-$0.2, $+$0.3]}
\newcommand{\DecPhiSearchW}{$+$0.4}
\newcommand{\DecPhiSearchWBand}{[$+$0.1, $+$0.8]}
\newcommand{\DecPhiSinkR}{$-$35}
\newcommand{\DecPhiSinkRBand}{[$-$39, $-$32]}
\newcommand{\DecPhiSinkT}{$+$0}
\newcommand{\DecPhiSinkTBand}{[$-$1, $+$1]}
\newcommand{\DecPhiSinkW}{$+$2}
\newcommand{\DecPhiSinkWBand}{[$+$1, $+$3]}
\newcommand{\DecPiecesK}{2\,717}
\newcommand{\DecPiecesPhase}{1\,422}
\newcommand{\DecPiecesQ}{700}
\newcommand{\DecShareClimbF}{$-$31}
\newcommand{\DecShareClimbR}{126}
\newcommand{\DecShareClimbT}{$-$2}
\newcommand{\DecShareClimbW}{7}
\newcommand{\DecShareCphF}{116}
\newcommand{\DecShareCphR}{$-$18}
\newcommand{\DecShareCphT}{8}
\newcommand{\DecShareCphW}{$-$6}
\newcommand{\DecShareFoundF}{68}
\newcommand{\DecShareFoundR}{41}
\newcommand{\DecShareFoundT}{17}
\newcommand{\DecShareFoundW}{$-$25}
\newcommand{\DecShareGlideF}{17}
\newcommand{\DecShareGlideR}{70}
\newcommand{\DecShareGlideT}{$-$4}
\newcommand{\DecShareGlideW}{18}
\newcommand{\DecShareHKF}{$-$773}
\newcommand{\DecShareHKR}{719}
\newcommand{\DecShareHKT}{$-$61}
\newcommand{\DecShareHKW}{215}
\newcommand{\DecShareHQF}{$-$53}
\newcommand{\DecShareHQR}{149}
\newcommand{\DecShareHQT}{$-$4}
\newcommand{\DecShareHQW}{7}
\newcommand{\DecShareRateF}{$-$57}
\newcommand{\DecShareRateR}{127}
\newcommand{\DecShareRateT}{26}
\newcommand{\DecShareRateW}{4}
\newcommand{\DecShareSearchF}{17}
\newcommand{\DecShareSearchR}{73}
\newcommand{\DecShareSearchT}{1}
\newcommand{\DecShareSearchW}{9}
\newcommand{\DecShareSinkF}{$-$21}
\newcommand{\DecShareSinkR}{127}
\newcommand{\DecShareSinkT}{$-$0}
\newcommand{\DecShareSinkW}{$-$6}
\newcommand{\DeclFixedFitA}{0.914}
\newcommand{\DeclFixedFitG}{0.920}
\newcommand{\DeclFixedFitS}{0.907}
\newcommand{\DeclFixedFlightsA}{7\,827}
\newcommand{\DeclFixedFlightsG}{7\,448}
\newcommand{\DeclFixedFlightsS}{39\,610}
\newcommand{\DeclFixedHA}{0.886}
\newcommand{\DeclFixedHABand}{[0.880, 0.892]}
\newcommand{\DeclFixedHG}{0.908}
\newcommand{\DeclFixedHGBand}{[0.902, 0.914]}
\newcommand{\DeclFixedHS}{0.886}
\newcommand{\DeclFixedHSBand}{[0.882, 0.890]}
\newcommand{\DeclFixedHThousandA}{0.887}
\newcommand{\DeclFixedHThousandG}{0.911}
\newcommand{\DeclFixedHThousandS}{0.883}
\newcommand{\DeclLongestFitA}{0.920}
\newcommand{\DeclLongestFitG}{0.931}
\newcommand{\DeclLongestFitS}{0.912}
\newcommand{\DeclLongestFlightsA}{820}
\newcommand{\DeclLongestFlightsG}{153}
\newcommand{\DeclLongestFlightsS}{5\,492}
\newcommand{\DeclLongestHA}{0.808}
\newcommand{\DeclLongestHABand}{[0.776, 0.839]}
\newcommand{\DeclLongestHG}{0.890}
\newcommand{\DeclLongestHGBand}{[0.839, 0.941]}
\newcommand{\DeclLongestHS}{0.827}
\newcommand{\DeclLongestHSBand}{[0.811, 0.843]}
\newcommand{\DeclLongestHThousandA}{0.907}
\newcommand{\DeclLongestHThousandG}{0.923}
\newcommand{\DeclLongestHThousandS}{0.901}
\newcommand{\DeclMinFlights}{100}
\newcommand{\DeclTauFixed}{2\,760}
\newcommand{\DeclTauLongest}{11\,450}
\newcommand{\DtDroppedPercent}{3.7}
\newcommand{\DtFlights}{63\,433}
\newcommand{\DtMax}{5}
\newcommand{\DurationLongFlights}{32\,106}
\newcommand{\DurationLongKeep}{77}
\newcommand{\DurationLongKeepN}{24\,794}
\newcommand{\DurationLongMedian}{0.25}
\newcommand{\DurationLongStrict}{24}
\newcommand{\DurationLongStrictN}{7\,548}
\newcommand{\DurationS}{9\,000}
\newcommand{\DurationShortFlights}{32\,185}
\newcommand{\DurationShortKeep}{90}
\newcommand{\DurationShortKeepN}{28\,822}
\newcommand{\DurationShortMedian}{0.37}
\newcommand{\DurationShortStrict}{29}
\newcommand{\DurationShortStrictN}{9\,270}
\newcommand{\EnterAtDeparture}{751}
\newcommand{\EnterBaseClimbAll}{28.5}
\newcommand{\EnterBaseClimbAllBand}{[28.0, 29.0]}
\newcommand{\EnterBaseClimbHighMtn}{26.5}
\newcommand{\EnterBaseClimbHighMtnBand}{[25.9, 27.0]}
\newcommand{\EnterBaseClimbHills}{34.6}
\newcommand{\EnterBaseClimbHillsBand}{[33.5, 35.7]}
\newcommand{\EnterBaseClimbLowMtn}{25.4}
\newcommand{\EnterBaseClimbLowMtnBand}{[24.9, 25.8]}
\newcommand{\EnterBaseClimbPlains}{36.9}
\newcommand{\EnterBaseClimbPlainsBand}{[36.1, 37.7]}
\newcommand{\EnterBaseGlideAll}{57.0}
\newcommand{\EnterBaseGlideAllBand}{[56.6, 57.5]}
\newcommand{\EnterBaseGlideHighMtn}{59.8}
\newcommand{\EnterBaseGlideHighMtnBand}{[59.3, 60.3]}
\newcommand{\EnterBaseGlideHills}{52.2}
\newcommand{\EnterBaseGlideHillsBand}{[51.1, 53.3]}
\newcommand{\EnterBaseGlideLowMtn}{59.3}
\newcommand{\EnterBaseGlideLowMtnBand}{[58.9, 59.8]}
\newcommand{\EnterBaseGlidePlains}{48.8}
\newcommand{\EnterBaseGlidePlainsBand}{[47.8, 49.8]}
\newcommand{\EnterBaseSearchAll}{14.5}
\newcommand{\EnterBaseSearchAllBand}{[14.2, 14.8]}
\newcommand{\EnterBaseSearchHighMtn}{13.8}
\newcommand{\EnterBaseSearchHighMtnBand}{[13.3, 14.2]}
\newcommand{\EnterBaseSearchHills}{13.2}
\newcommand{\EnterBaseSearchHillsBand}{[12.4, 13.9]}
\newcommand{\EnterBaseSearchLowMtn}{15.3}
\newcommand{\EnterBaseSearchLowMtnBand}{[14.9, 15.7]}
\newcommand{\EnterBaseSearchPlains}{14.3}
\newcommand{\EnterBaseSearchPlainsBand}{[13.6, 15.1]}
\newcommand{\EnterCountAll}{16\,105}
\newcommand{\EnterCountHighMtn}{3\,944}
\newcommand{\EnterCountHills}{1\,767}
\newcommand{\EnterCountLowMtn}{7\,612}
\newcommand{\EnterCountPlains}{2\,441}
\newcommand{\EnterRatioClimbAll}{1.04}
\newcommand{\EnterRatioClimbAllBand}{[1.01, 1.06]}
\newcommand{\EnterRatioClimbHighMtn}{1.11}
\newcommand{\EnterRatioClimbHighMtnBand}{[1.06, 1.16]}
\newcommand{\EnterRatioClimbHills}{0.98}
\newcommand{\EnterRatioClimbHillsBand}{[0.92, 1.05]}
\newcommand{\EnterRatioClimbLowMtn}{0.99}
\newcommand{\EnterRatioClimbLowMtnBand}{[0.96, 1.02]}
\newcommand{\EnterRatioClimbPlains}{1.10}
\newcommand{\EnterRatioClimbPlainsBand}{[1.04, 1.15]}
\newcommand{\EnterRatioGlideAll}{1.03}
\newcommand{\EnterRatioGlideAllBand}{[1.02, 1.04]}
\newcommand{\EnterRatioGlideHighMtn}{0.99}
\newcommand{\EnterRatioGlideHighMtnBand}{[0.97, 1.01]}
\newcommand{\EnterRatioGlideHills}{1.09}
\newcommand{\EnterRatioGlideHillsBand}{[1.05, 1.13]}
\newcommand{\EnterRatioGlideLowMtn}{1.05}
\newcommand{\EnterRatioGlideLowMtnBand}{[1.03, 1.07]}
\newcommand{\EnterRatioGlidePlains}{0.99}
\newcommand{\EnterRatioGlidePlainsBand}{[0.95, 1.04]}
\newcommand{\EnterRatioSearchAll}{0.81}
\newcommand{\EnterRatioSearchAllBand}{[0.77, 0.84]}
\newcommand{\EnterRatioSearchHighMtn}{0.83}
\newcommand{\EnterRatioSearchHighMtnBand}{[0.77, 0.89]}
\newcommand{\EnterRatioSearchHills}{0.69}
\newcommand{\EnterRatioSearchHillsBand}{[0.61, 0.78]}
\newcommand{\EnterRatioSearchLowMtn}{0.83}
\newcommand{\EnterRatioSearchLowMtnBand}{[0.78, 0.87]}
\newcommand{\EnterRatioSearchPlains}{0.78}
\newcommand{\EnterRatioSearchPlainsBand}{[0.70, 0.86]}
\newcommand{\EnterShareClimbAll}{29.5}
\newcommand{\EnterShareClimbAllBand}{[28.8, 30.2]}
\newcommand{\EnterShareClimbHighMtn}{29.4}
\newcommand{\EnterShareClimbHighMtnBand}{[28.1, 30.6]}
\newcommand{\EnterShareClimbHills}{34.1}
\newcommand{\EnterShareClimbHillsBand}{[31.9, 36.3]}
\newcommand{\EnterShareClimbLowMtn}{25.1}
\newcommand{\EnterShareClimbLowMtnBand}{[24.2, 25.9]}
\newcommand{\EnterShareClimbPlains}{40.4}
\newcommand{\EnterShareClimbPlainsBand}{[38.4, 42.4]}
\newcommand{\EnterShareGlideAll}{58.8}
\newcommand{\EnterShareGlideAllBand}{[58.1, 59.6]}
\newcommand{\EnterShareGlideHighMtn}{59.2}
\newcommand{\EnterShareGlideHighMtnBand}{[57.9, 60.6]}
\newcommand{\EnterShareGlideHills}{56.8}
\newcommand{\EnterShareGlideHillsBand}{[54.4, 59.3]}
\newcommand{\EnterShareGlideLowMtn}{62.3}
\newcommand{\EnterShareGlideLowMtnBand}{[61.3, 63.2]}
\newcommand{\EnterShareGlidePlains}{48.5}
\newcommand{\EnterShareGlidePlainsBand}{[46.3, 50.6]}
\newcommand{\EnterShareSearchAll}{11.6}
\newcommand{\EnterShareSearchAllBand}{[11.2, 12.1]}
\newcommand{\EnterShareSearchHighMtn}{11.4}
\newcommand{\EnterShareSearchHighMtnBand}{[10.5, 12.3]}
\newcommand{\EnterShareSearchHills}{9.1}
\newcommand{\EnterShareSearchHillsBand}{[8.0, 10.2]}
\newcommand{\EnterShareSearchLowMtn}{12.7}
\newcommand{\EnterShareSearchLowMtnBand}{[12.0, 13.4]}
\newcommand{\EnterShareSearchPlains}{11.1}
\newcommand{\EnterShareSearchPlainsBand}{[9.9, 12.4]}
\newcommand{\EventsNoPhase}{303}
\newcommand{\FlightsAll}{154\,598}
\newcommand{\FlightsClosed}{87\,484}
\newcommand{\FlightsEvery}{64\,291}
\newcommand{\FlightsFixed}{51\,017}
\newcommand{\FlightsLongest}{6\,464}
\newcommand{\FlightsNeither}{269}
\newcommand{\FlightsOpen}{66\,845}
\newcommand{\GapAllASHundred}{$+$0.009}
\newcommand{\GapAllASHundredBand}{[$+$0.008, $+$0.010]}
\newcommand{\GapAllASLong}{$+$0.000}
\newcommand{\GapAllASLongBand}{[$-$0.005, $+$0.006]}
\newcommand{\GapAllASTen}{$+$0.015}
\newcommand{\GapAllASTenBand}{[$+$0.013, $+$0.016]}
\newcommand{\GapAllASThousand}{$+$0.006}
\newcommand{\GapAllASThousandBand}{[$+$0.003, $+$0.009]}
\newcommand{\GapAllASThreeTwenty}{$+$0.010}
\newcommand{\GapAllASThreeTwentyBand}{[$+$0.009, $+$0.012]}
\newcommand{\GapAllGAFrom}{210}
\newcommand{\GapAllGAHundred}{$-$0.002}
\newcommand{\GapAllGAHundredBand}{[$-$0.003, $-$0.000]}
\newcommand{\GapAllGALong}{$+$0.026}
\newcommand{\GapAllGALongBand}{[$+$0.019, $+$0.032]}
\newcommand{\GapAllGATen}{$-$0.010}
\newcommand{\GapAllGATenBand}{[$-$0.012, $-$0.008]}
\newcommand{\GapAllGAThousand}{$+$0.024}
\newcommand{\GapAllGAThousandBand}{[$+$0.020, $+$0.027]}
\newcommand{\GapAllGAThreeTwenty}{$+$0.007}
\newcommand{\GapAllGAThreeTwentyBand}{[$+$0.005, $+$0.009]}
\newcommand{\GapAllGSFrom}{10}
\newcommand{\GapAllGSHundred}{$+$0.008}
\newcommand{\GapAllGSHundredBand}{[$+$0.006, $+$0.009]}
\newcommand{\GapAllGSLong}{$+$0.026}
\newcommand{\GapAllGSLongBand}{[$+$0.019, $+$0.032]}
\newcommand{\GapAllGSTen}{$+$0.005}
\newcommand{\GapAllGSTenBand}{[$+$0.002, $+$0.007]}
\newcommand{\GapAllGSThousand}{$+$0.030}
\newcommand{\GapAllGSThousandBand}{[$+$0.026, $+$0.034]}
\newcommand{\GapAllGSThreeTwenty}{$+$0.017}
\newcommand{\GapAllGSThreeTwentyBand}{[$+$0.014, $+$0.019]}
\newcommand{\GapCohortTau}{3\,600}
\newcommand{\GapMinASHundred}{$+$0.007}
\newcommand{\GapMinASHundredBand}{[$+$0.006, $+$0.008]}
\newcommand{\GapMinASLong}{$-$0.004}
\newcommand{\GapMinASLongBand}{[$-$0.009, $+$0.001]}
\newcommand{\GapMinASTen}{$+$0.014}
\newcommand{\GapMinASTenBand}{[$+$0.012, $+$0.015]}
\newcommand{\GapMinASThousand}{$+$0.001}
\newcommand{\GapMinASThousandBand}{[$-$0.003, $+$0.004]}
\newcommand{\GapMinASThreeTwenty}{$+$0.007}
\newcommand{\GapMinASThreeTwentyBand}{[$+$0.005, $+$0.009]}
\newcommand{\GapMinGAFrom}{260}
\newcommand{\GapMinGAHundred}{$-$0.002}
\newcommand{\GapMinGAHundredBand}{[$-$0.003, $-$0.001]}
\newcommand{\GapMinGALong}{$+$0.022}
\newcommand{\GapMinGALongBand}{[$+$0.017, $+$0.028]}
\newcommand{\GapMinGATen}{$-$0.011}
\newcommand{\GapMinGATenBand}{[$-$0.013, $-$0.008]}
\newcommand{\GapMinGAThousand}{$+$0.021}
\newcommand{\GapMinGAThousandBand}{[$+$0.018, $+$0.024]}
\newcommand{\GapMinGAThreeTwenty}{$+$0.005}
\newcommand{\GapMinGAThreeTwentyBand}{[$+$0.003, $+$0.007]}
\newcommand{\GapMinGSFrom}{10}
\newcommand{\GapMinGSHundred}{$+$0.005}
\newcommand{\GapMinGSHundredBand}{[$+$0.003, $+$0.006]}
\newcommand{\GapMinGSLong}{$+$0.019}
\newcommand{\GapMinGSLongBand}{[$+$0.013, $+$0.024]}
\newcommand{\GapMinGSTen}{$+$0.003}
\newcommand{\GapMinGSTenBand}{[$+$0.001, $+$0.006]}
\newcommand{\GapMinGSThousand}{$+$0.022}
\newcommand{\GapMinGSThousandBand}{[$+$0.018, $+$0.025]}
\newcommand{\GapMinGSThreeTwenty}{$+$0.012}
\newcommand{\GapMinGSThreeTwentyBand}{[$+$0.009, $+$0.014]}
\newcommand{\GapSpreadGA}{0.007}
\newcommand{\GapSpreadHighAS}{0.012}
\newcommand{\GapSpreadHighGA}{0.016}
\newcommand{\GapSpreadHighGS}{0.016}
\newcommand{\GapSpreadLag}{1\,000}
\newcommand{\GapSpreadLowAS}{0.003}
\newcommand{\GapSpreadLowGA}{0.008}
\newcommand{\GapSpreadLowGS}{0.006}
\newcommand{\GapStrictASHundred}{$+$0.002}
\newcommand{\GapStrictASHundredBand}{[$+$0.001, $+$0.003]}
\newcommand{\GapStrictASLong}{$-$0.000}
\newcommand{\GapStrictASLongBand}{[$-$0.003, $+$0.002]}
\newcommand{\GapStrictASTen}{$+$0.008}
\newcommand{\GapStrictASTenBand}{[$+$0.005, $+$0.010]}
\newcommand{\GapStrictASThousand}{$-$0.000}
\newcommand{\GapStrictASThousandBand}{[$-$0.003, $+$0.003]}
\newcommand{\GapStrictASThreeTwenty}{$+$0.001}
\newcommand{\GapStrictASThreeTwentyBand}{[$-$0.001, $+$0.003]}
\newcommand{\GapStrictGAFrom}{580}
\newcommand{\GapStrictGAHundred}{$-$0.005}
\newcommand{\GapStrictGAHundredBand}{[$-$0.006, $-$0.003]}
\newcommand{\GapStrictGALong}{$+$0.007}
\newcommand{\GapStrictGALongBand}{[$+$0.004, $+$0.011]}
\newcommand{\GapStrictGATen}{$-$0.012}
\newcommand{\GapStrictGATenBand}{[$-$0.015, $-$0.008]}
\newcommand{\GapStrictGAThousand}{$+$0.008}
\newcommand{\GapStrictGAThousandBand}{[$+$0.005, $+$0.011]}
\newcommand{\GapStrictGAThreeTwenty}{$-$0.001}
\newcommand{\GapStrictGAThreeTwentyBand}{[$-$0.003, $+$0.002]}
\newcommand{\GapStrictGSFrom}{470}
\newcommand{\GapStrictGSHundred}{$-$0.003}
\newcommand{\GapStrictGSHundredBand}{[$-$0.004, $-$0.001]}
\newcommand{\GapStrictGSLong}{$+$0.007}
\newcommand{\GapStrictGSLongBand}{[$+$0.004, $+$0.010]}
\newcommand{\GapStrictGSTen}{$-$0.004}
\newcommand{\GapStrictGSTenBand}{[$-$0.008, $+$0.000]}
\newcommand{\GapStrictGSThousand}{$+$0.008}
\newcommand{\GapStrictGSThousandBand}{[$+$0.005, $+$0.011]}
\newcommand{\GapStrictGSThreeTwenty}{$+$0.000}
\newcommand{\GapStrictGSThreeTwentyBand}{[$-$0.002, $+$0.003]}
\newcommand{\GridS}{10}
\newcommand{\GroundCell}{200}
\newcommand{\GroundChanged}{7.1}
\newcommand{\GroundEdgeHigh}{1\,500}
\newcommand{\GroundEdgeLow}{300}
\newcommand{\GroundEdgeMid}{800}
\newcommand{\GroundHighMtnASThousand}{$+$0.022}
\newcommand{\GroundHighMtnASThousandBand}{[$+$0.017, $+$0.027]}
\newcommand{\GroundHighMtnASThreeTwenty}{$+$0.012}
\newcommand{\GroundHighMtnASThreeTwentyBand}{[$+$0.010, $+$0.015]}
\newcommand{\GroundHighMtnFlightsA}{2\,800}
\newcommand{\GroundHighMtnFlightsG}{3\,112}
\newcommand{\GroundHighMtnFlightsS}{8\,028}
\newcommand{\GroundHighMtnGAPositive}{170}
\newcommand{\GroundHighMtnGAThousand}{$+$0.033}
\newcommand{\GroundHighMtnGAThousandBand}{[$+$0.028, $+$0.038]}
\newcommand{\GroundHighMtnGAThreeTwenty}{$+$0.014}
\newcommand{\GroundHighMtnGAThreeTwentyBand}{[$+$0.012, $+$0.017]}
\newcommand{\GroundHighMtnGSThousand}{$+$0.055}
\newcommand{\GroundHighMtnGSThousandBand}{[$+$0.050, $+$0.059]}
\newcommand{\GroundHighMtnGSThreeTwenty}{$+$0.027}
\newcommand{\GroundHighMtnGSThreeTwentyBand}{[$+$0.024, $+$0.029]}
\newcommand{\GroundHighMtnPeakGA}{0.23}
\newcommand{\GroundHighMtnPeakGALag}{870}
\newcommand{\GroundHighMtnPeakGS}{0.33}
\newcommand{\GroundHighMtnPeakGSLag}{1\,230}
\newcommand{\GroundHillsASThousand}{$+$0.012}
\newcommand{\GroundHillsASThousandBand}{[$+$0.007, $+$0.018]}
\newcommand{\GroundHillsASThreeTwenty}{$+$0.010}
\newcommand{\GroundHillsASThreeTwentyBand}{[$+$0.006, $+$0.014]}
\newcommand{\GroundHillsFlightsA}{1\,120}
\newcommand{\GroundHillsFlightsG}{1\,550}
\newcommand{\GroundHillsFlightsS}{9\,815}
\newcommand{\GroundHillsGAPositive}{580}
\newcommand{\GroundHillsGAThousand}{$+$0.012}
\newcommand{\GroundHillsGAThousandBand}{[$+$0.007, $+$0.018]}
\newcommand{\GroundHillsGAThreeTwenty}{$+$0.001}
\newcommand{\GroundHillsGAThreeTwentyBand}{[$-$0.003, $+$0.005]}
\newcommand{\GroundHillsGSThousand}{$+$0.025}
\newcommand{\GroundHillsGSThousandBand}{[$+$0.020, $+$0.030]}
\newcommand{\GroundHillsGSThreeTwenty}{$+$0.011}
\newcommand{\GroundHillsGSThreeTwentyBand}{[$+$0.008, $+$0.015]}
\newcommand{\GroundHillsPeakGA}{0.22}
\newcommand{\GroundHillsPeakGALag}{2\,250}
\newcommand{\GroundHillsPeakGS}{0.39}
\newcommand{\GroundHillsPeakGSLag}{2\,250}
\newcommand{\GroundLowMtnASThousand}{$+$0.011}
\newcommand{\GroundLowMtnASThousandBand}{[$+$0.007, $+$0.015]}
\newcommand{\GroundLowMtnASThreeTwenty}{$+$0.012}
\newcommand{\GroundLowMtnASThreeTwentyBand}{[$+$0.010, $+$0.014]}
\newcommand{\GroundLowMtnFlightsA}{5\,277}
\newcommand{\GroundLowMtnFlightsG}{6\,160}
\newcommand{\GroundLowMtnFlightsS}{18\,902}
\newcommand{\GroundLowMtnGAPositive}{210}
\newcommand{\GroundLowMtnGAThousand}{$+$0.031}
\newcommand{\GroundLowMtnGAThousandBand}{[$+$0.027, $+$0.035]}
\newcommand{\GroundLowMtnGAThreeTwenty}{$+$0.013}
\newcommand{\GroundLowMtnGAThreeTwentyBand}{[$+$0.011, $+$0.015]}
\newcommand{\GroundLowMtnGSThousand}{$+$0.042}
\newcommand{\GroundLowMtnGSThousandBand}{[$+$0.038, $+$0.046]}
\newcommand{\GroundLowMtnGSThreeTwenty}{$+$0.025}
\newcommand{\GroundLowMtnGSThreeTwentyBand}{[$+$0.022, $+$0.027]}
\newcommand{\GroundLowMtnPeakGA}{0.22}
\newcommand{\GroundLowMtnPeakGALag}{710}
\newcommand{\GroundLowMtnPeakGS}{0.28}
\newcommand{\GroundLowMtnPeakGSLag}{710}
\newcommand{\GroundLowlandA}{25}
\newcommand{\GroundLowlandG}{28}
\newcommand{\GroundLowlandS}{42}
\newcommand{\GroundMinOnDem}{50}
\newcommand{\GroundMountainA}{75}
\newcommand{\GroundMountainG}{72}
\newcommand{\GroundMountainS}{58}
\newcommand{\GroundPlainsASThousand}{$+$0.005}
\newcommand{\GroundPlainsASThousandBand}{[$+$0.002, $+$0.008]}
\newcommand{\GroundPlainsASThreeTwenty}{$+$0.004}
\newcommand{\GroundPlainsASThreeTwentyBand}{[$+$0.002, $+$0.007]}
\newcommand{\GroundPlainsFlightsA}{1\,548}
\newcommand{\GroundPlainsFlightsG}{2\,124}
\newcommand{\GroundPlainsFlightsS}{9\,763}
\newcommand{\GroundPlainsGAPositive}{580}
\newcommand{\GroundPlainsGAThousand}{$+$0.007}
\newcommand{\GroundPlainsGAThousandBand}{[$+$0.005, $+$0.010]}
\newcommand{\GroundPlainsGAThreeTwenty}{$-$0.002}
\newcommand{\GroundPlainsGAThreeTwentyBand}{[$-$0.004, $+$0.001]}
\newcommand{\GroundPlainsGSThousand}{$+$0.013}
\newcommand{\GroundPlainsGSThousandBand}{[$+$0.010, $+$0.016]}
\newcommand{\GroundPlainsGSThreeTwenty}{$+$0.003}
\newcommand{\GroundPlainsGSThreeTwentyBand}{[$-$0.000, $+$0.006]}
\newcommand{\GroundPlainsPeakGA}{0.21}
\newcommand{\GroundPlainsPeakGALag}{1\,500}
\newcommand{\GroundPlainsPeakGS}{0.28}
\newcommand{\GroundPlainsPeakGSLag}{2\,250}
\newcommand{\GroundTrim}{10}
\newcommand{\GroundUnknown}{6.2}
\newcommand{\GroupCheckDays}{20}
\newcommand{\GroupCheckFlightsChange}{0.0}
\newcommand{\GroupCheckShared}{99.6}
\newcommand{\GroupCheckTimeChange}{$-$0.07}
\newcommand{\HalfDecade}{0.25}
\newcommand{\HqHighest}{3\,272}
\newcommand{\HqLowest}{1\,035}
\newcommand{\JoinFlightsAll}{13\,188}
\newcommand{\JoinFlightsHighMtn}{3\,113}
\newcommand{\JoinFlightsHills}{1\,545}
\newcommand{\JoinFlightsLowMtn}{6\,131}
\newcommand{\JoinFlightsPlains}{2\,113}
\newcommand{\Kept}{53\,616}
\newcommand{\KeptShare}{83}
\newcommand{\KindGroupDuration}{13\,042}
\newcommand{\KindGroupFlights}{13\,252}
\newcommand{\KindGroupKeep}{81}
\newcommand{\KindGroupMedian}{0.28}
\newcommand{\KindGroupMedianHigh}{0.22}
\newcommand{\KindGroupMedianLong}{0.25}
\newcommand{\KindGroupMedianLow}{0.43}
\newcommand{\KindGroupMedianShort}{0.35}
\newcommand{\KindGroupMountain}{55}
\newcommand{\KindGroupStrict}{24}
\newcommand{\KindSoloDuration}{8\,106}
\newcommand{\KindSoloFlights}{51\,039}
\newcommand{\KindSoloKeep}{84}
\newcommand{\KindSoloMedian}{0.33}
\newcommand{\KindSoloMedianHigh}{0.24}
\newcommand{\KindSoloMedianLong}{0.25}
\newcommand{\KindSoloMedianLow}{0.45}
\newcommand{\KindSoloMedianShort}{0.37}
\newcommand{\KindSoloMountain}{53}
\newcommand{\KindSoloStrict}{27}
\newcommand{\LagCount}{37}
\newcommand{\LagDecades}{3}
\newcommand{\LagFactor}{1\,000}
\newcommand{\LagList}{10, 20, 30, 40, 50, 60, 80, 90, 100, 110, 140, 170, 210, 260, 320, 390, 470, 580, 710, 870,\\ 1\,000, 1\,230, 1\,500, 1\,840, 2\,250, 2\,760, 3\,380, 4\,150, 5\,080, 6\,220, 7\,630, 9\,350, 11\,450, 14\,030, 17\,190, 21\,060, 25\,810}
\newcommand{\LagMax}{25\,810}
\newcommand{\LagMin}{10}
\newcommand{\LagRatio}{1.225}
\newcommand{\LagSteps}{34}
\newcommand{\LagsPerDecade}{11}
\newcommand{\LeaveAtLanding}{1\,071}
\newcommand{\LeaveBaseClimbAll}{28.3}
\newcommand{\LeaveBaseClimbAllBand}{[27.8, 28.7]}
\newcommand{\LeaveBaseClimbHighMtn}{26.3}
\newcommand{\LeaveBaseClimbHighMtnBand}{[25.9, 26.8]}
\newcommand{\LeaveBaseClimbHills}{34.5}
\newcommand{\LeaveBaseClimbHillsBand}{[33.5, 35.5]}
\newcommand{\LeaveBaseClimbLowMtn}{25.4}
\newcommand{\LeaveBaseClimbLowMtnBand}{[24.9, 25.8]}
\newcommand{\LeaveBaseClimbPlains}{35.7}
\newcommand{\LeaveBaseClimbPlainsBand}{[35.0, 36.4]}
\newcommand{\LeaveBaseGlideAll}{56.9}
\newcommand{\LeaveBaseGlideAllBand}{[56.5, 57.4]}
\newcommand{\LeaveBaseGlideHighMtn}{59.7}
\newcommand{\LeaveBaseGlideHighMtnBand}{[59.1, 60.2]}
\newcommand{\LeaveBaseGlideHills}{51.6}
\newcommand{\LeaveBaseGlideHillsBand}{[50.5, 52.7]}
\newcommand{\LeaveBaseGlideLowMtn}{59.1}
\newcommand{\LeaveBaseGlideLowMtnBand}{[58.7, 59.6]}
\newcommand{\LeaveBaseGlidePlains}{49.2}
\newcommand{\LeaveBaseGlidePlainsBand}{[48.3, 50.1]}
\newcommand{\LeaveBaseSearchAll}{14.8}
\newcommand{\LeaveBaseSearchAllBand}{[14.5, 15.1]}
\newcommand{\LeaveBaseSearchHighMtn}{14.0}
\newcommand{\LeaveBaseSearchHighMtnBand}{[13.5, 14.4]}
\newcommand{\LeaveBaseSearchHills}{13.9}
\newcommand{\LeaveBaseSearchHillsBand}{[13.1, 14.7]}
\newcommand{\LeaveBaseSearchLowMtn}{15.5}
\newcommand{\LeaveBaseSearchLowMtnBand}{[15.1, 15.9]}
\newcommand{\LeaveBaseSearchPlains}{15.1}
\newcommand{\LeaveBaseSearchPlainsBand}{[14.4, 15.8]}
\newcommand{\LeaveCountAll}{15\,784}
\newcommand{\LeaveCountHighMtn}{3\,921}
\newcommand{\LeaveCountHills}{1\,667}
\newcommand{\LeaveCountLowMtn}{7\,492}
\newcommand{\LeaveCountPlains}{2\,368}
\newcommand{\LeaveRatioClimbAll}{0.98}
\newcommand{\LeaveRatioClimbAllBand}{[0.96, 1.00]}
\newcommand{\LeaveRatioClimbHighMtn}{1.08}
\newcommand{\LeaveRatioClimbHighMtnBand}{[1.03, 1.12]}
\newcommand{\LeaveRatioClimbHills}{0.87}
\newcommand{\LeaveRatioClimbHillsBand}{[0.81, 0.92]}
\newcommand{\LeaveRatioClimbLowMtn}{1.02}
\newcommand{\LeaveRatioClimbLowMtnBand}{[0.99, 1.05]}
\newcommand{\LeaveRatioClimbPlains}{0.85}
\newcommand{\LeaveRatioClimbPlainsBand}{[0.79, 0.90]}
\newcommand{\LeaveRatioGlideAll}{1.02}
\newcommand{\LeaveRatioGlideAllBand}{[1.01, 1.04]}
\newcommand{\LeaveRatioGlideHighMtn}{0.97}
\newcommand{\LeaveRatioGlideHighMtnBand}{[0.95, 1.00]}
\newcommand{\LeaveRatioGlideHills}{1.13}
\newcommand{\LeaveRatioGlideHillsBand}{[1.09, 1.17]}
\newcommand{\LeaveRatioGlideLowMtn}{1.00}
\newcommand{\LeaveRatioGlideLowMtnBand}{[0.98, 1.01]}
\newcommand{\LeaveRatioGlidePlains}{1.15}
\newcommand{\LeaveRatioGlidePlainsBand}{[1.10, 1.19]}
\newcommand{\LeaveRatioSearchAll}{0.94}
\newcommand{\LeaveRatioSearchAllBand}{[0.91, 0.97]}
\newcommand{\LeaveRatioSearchHighMtn}{0.97}
\newcommand{\LeaveRatioSearchHighMtnBand}{[0.90, 1.03]}
\newcommand{\LeaveRatioSearchHills}{0.83}
\newcommand{\LeaveRatioSearchHillsBand}{[0.75, 0.92]}
\newcommand{\LeaveRatioSearchLowMtn}{0.97}
\newcommand{\LeaveRatioSearchLowMtnBand}{[0.93, 1.02]}
\newcommand{\LeaveRatioSearchPlains}{0.89}
\newcommand{\LeaveRatioSearchPlainsBand}{[0.81, 0.96]}
\newcommand{\LeaveShareClimbAll}{27.7}
\newcommand{\LeaveShareClimbAllBand}{[27.1, 28.4]}
\newcommand{\LeaveShareClimbHighMtn}{28.3}
\newcommand{\LeaveShareClimbHighMtnBand}{[27.0, 29.6]}
\newcommand{\LeaveShareClimbHills}{29.9}
\newcommand{\LeaveShareClimbHillsBand}{[27.8, 32.0]}
\newcommand{\LeaveShareClimbLowMtn}{25.9}
\newcommand{\LeaveShareClimbLowMtnBand}{[25.0, 26.7]}
\newcommand{\LeaveShareClimbPlains}{30.2}
\newcommand{\LeaveShareClimbPlainsBand}{[28.3, 32.2]}
\newcommand{\LeaveShareGlideAll}{58.3}
\newcommand{\LeaveShareGlideAllBand}{[57.6, 59.1]}
\newcommand{\LeaveShareGlideHighMtn}{58.1}
\newcommand{\LeaveShareGlideHighMtnBand}{[56.8, 59.5]}
\newcommand{\LeaveShareGlideHills}{58.5}
\newcommand{\LeaveShareGlideHillsBand}{[56.3, 60.8]}
\newcommand{\LeaveShareGlideLowMtn}{59.0}
\newcommand{\LeaveShareGlideLowMtnBand}{[58.0, 60.1]}
\newcommand{\LeaveShareGlidePlains}{56.4}
\newcommand{\LeaveShareGlidePlainsBand}{[54.2, 58.6]}
\newcommand{\LeaveShareSearchAll}{14.0}
\newcommand{\LeaveShareSearchAllBand}{[13.5, 14.5]}
\newcommand{\LeaveShareSearchHighMtn}{13.5}
\newcommand{\LeaveShareSearchHighMtnBand}{[12.6, 14.5]}
\newcommand{\LeaveShareSearchHills}{11.6}
\newcommand{\LeaveShareSearchHillsBand}{[10.3, 12.8]}
\newcommand{\LeaveShareSearchLowMtn}{15.1}
\newcommand{\LeaveShareSearchLowMtnBand}{[14.3, 15.8]}
\newcommand{\LeaveShareSearchPlains}{13.4}
\newcommand{\LeaveShareSearchPlainsBand}{[12.1, 14.7]}
\newcommand{\MatChosenText}{month $\pm2$, hour $\pm2$, same class}
\newcommand{\MatCover}{0.8}
\newcommand{\MatDeltaClimbA}{$+$0.5}
\newcommand{\MatDeltaClimbABand}{[$-$1.2, $+$2.3]}
\newcommand{\MatDeltaClimbAYears}{[$-$0.1, $+$1.0]}
\newcommand{\MatDeltaClimbD}{$+$2.4}
\newcommand{\MatDeltaClimbDBand}{[$+$0.8, $+$4.0]}
\newcommand{\MatDeltaClimbDYears}{[$+$1.8, $+$2.8]}
\newcommand{\MatDeltaClimbG}{$+$3.0}
\newcommand{\MatDeltaClimbGBand}{[$+$1.4, $+$4.5]}
\newcommand{\MatDeltaClimbGYears}{[$+$2.3, $+$3.5]}
\newcommand{\MatDeltaCphA}{$+$0.39}
\newcommand{\MatDeltaCphABand}{[$-$0.07, $+$0.84]}
\newcommand{\MatDeltaCphAYears}{[$+$0.21, $+$0.56]}
\newcommand{\MatDeltaCphD}{$+$0.35}
\newcommand{\MatDeltaCphDBand}{[$-$0.09, $+$0.78]}
\newcommand{\MatDeltaCphDYears}{[$+$0.16, $+$0.45]}
\newcommand{\MatDeltaCphG}{$+$0.73}
\newcommand{\MatDeltaCphGBand}{[$+$0.37, $+$1.10]}
\newcommand{\MatDeltaCphGYears}{[$+$0.64, $+$0.85]}
\newcommand{\MatDeltaFitKA}{$-$0.000}
\newcommand{\MatDeltaFitKABand}{[$-$0.004, $+$0.004]}
\newcommand{\MatDeltaFitKAYears}{[$-$0.001, $+$0.001]}
\newcommand{\MatDeltaFitKD}{$-$0.005}
\newcommand{\MatDeltaFitKDBand}{[$-$0.009, $-$0.001]}
\newcommand{\MatDeltaFitKDYears}{[$-$0.006, $-$0.004]}
\newcommand{\MatDeltaFitKG}{$-$0.005}
\newcommand{\MatDeltaFitKGBand}{[$-$0.008, $-$0.002]}
\newcommand{\MatDeltaFitKGYears}{[$-$0.006, $-$0.004]}
\newcommand{\MatDeltaFitTA}{$+$0.006}
\newcommand{\MatDeltaFitTABand}{[$+$0.000, $+$0.011]}
\newcommand{\MatDeltaFitTAYears}{[$+$0.005, $+$0.007]}
\newcommand{\MatDeltaFitTD}{$+$0.002}
\newcommand{\MatDeltaFitTDBand}{[$-$0.004, $+$0.009]}
\newcommand{\MatDeltaFitTDYears}{[$+$0.001, $+$0.003]}
\newcommand{\MatDeltaFitTG}{$+$0.008}
\newcommand{\MatDeltaFitTGBand}{[$+$0.004, $+$0.012]}
\newcommand{\MatDeltaFitTGYears}{[$+$0.007, $+$0.009]}
\newcommand{\MatDeltaFoundA}{$+$1.2}
\newcommand{\MatDeltaFoundABand}{[$-$1.7, $+$4.1]}
\newcommand{\MatDeltaFoundAYears}{[$+$0.2, $+$2.1]}
\newcommand{\MatDeltaFoundD}{$+$1.1}
\newcommand{\MatDeltaFoundDBand}{[$-$2.4, $+$4.7]}
\newcommand{\MatDeltaFoundDYears}{[$+$0.1, $+$2.4]}
\newcommand{\MatDeltaFoundG}{$+$2.3}
\newcommand{\MatDeltaFoundGBand}{[$-$0.2, $+$4.9]}
\newcommand{\MatDeltaFoundGYears}{[$+$1.7, $+$3.0]}
\newcommand{\MatDeltaGlideA}{$-$1.0}
\newcommand{\MatDeltaGlideABand}{[$-$1.8, $-$0.2]}
\newcommand{\MatDeltaGlideAYears}{[$-$1.1, $-$0.7]}
\newcommand{\MatDeltaGlideD}{$+$0.1}
\newcommand{\MatDeltaGlideDBand}{[$-$0.9, $+$1.0]}
\newcommand{\MatDeltaGlideDYears}{[$-$0.2, $+$0.3]}
\newcommand{\MatDeltaGlideG}{$-$0.9}
\newcommand{\MatDeltaGlideGBand}{[$-$1.7, $-$0.1]}
\newcommand{\MatDeltaGlideGYears}{[$-$1.2, $-$0.8]}
\newcommand{\MatDeltaHKA}{$+$0.006}
\newcommand{\MatDeltaHKABand}{[$-$0.003, $+$0.016]}
\newcommand{\MatDeltaHKAYears}{[$+$0.004, $+$0.009]}
\newcommand{\MatDeltaHKD}{$+$0.004}
\newcommand{\MatDeltaHKDBand}{[$-$0.007, $+$0.015]}
\newcommand{\MatDeltaHKDYears}{[$+$0.002, $+$0.005]}
\newcommand{\MatDeltaHKG}{$+$0.010}
\newcommand{\MatDeltaHKGBand}{[$+$0.004, $+$0.017]}
\newcommand{\MatDeltaHKGYears}{[$+$0.008, $+$0.012]}
\newcommand{\MatDeltaHKLongA}{$+$0.013}
\newcommand{\MatDeltaHKLongABand}{[$+$0.001, $+$0.026]}
\newcommand{\MatDeltaHKLongG}{$+$0.024}
\newcommand{\MatDeltaHKLongGBand}{[$+$0.014, $+$0.034]}
\newcommand{\MatDeltaHQA}{$+$0.027}
\newcommand{\MatDeltaHQABand}{[$+$0.008, $+$0.046]}
\newcommand{\MatDeltaHQAYears}{[$+$0.023, $+$0.032]}
\newcommand{\MatDeltaHQD}{$-$0.006}
\newcommand{\MatDeltaHQDBand}{[$-$0.028, $+$0.016]}
\newcommand{\MatDeltaHQDYears}{[$-$0.009, $-$0.002]}
\newcommand{\MatDeltaHQG}{$+$0.021}
\newcommand{\MatDeltaHQGBand}{[$+$0.008, $+$0.034]}
\newcommand{\MatDeltaHQGYears}{[$+$0.017, $+$0.024]}
\newcommand{\MatDeltaRateA}{$-$0.074}
\newcommand{\MatDeltaRateABand}{[$-$0.126, $-$0.023]}
\newcommand{\MatDeltaRateAYears}{[$-$0.093, $-$0.043]}
\newcommand{\MatDeltaRateD}{$+$0.031}
\newcommand{\MatDeltaRateDBand}{[$-$0.021, $+$0.083]}
\newcommand{\MatDeltaRateDYears}{[$+$0.004, $+$0.047]}
\newcommand{\MatDeltaRateG}{$-$0.044}
\newcommand{\MatDeltaRateGBand}{[$-$0.085, $-$0.002]}
\newcommand{\MatDeltaRateGYears}{[$-$0.055, $-$0.031]}
\newcommand{\MatDeltaSearchA}{$+$1.9}
\newcommand{\MatDeltaSearchABand}{[$+$0.3, $+$3.6]}
\newcommand{\MatDeltaSearchAYears}{[$+$1.4, $+$2.6]}
\newcommand{\MatDeltaSearchD}{$-$1.1}
\newcommand{\MatDeltaSearchDBand}{[$-$2.6, $+$0.4]}
\newcommand{\MatDeltaSearchDYears}{[$-$1.7, $-$0.6]}
\newcommand{\MatDeltaSearchG}{$+$0.8}
\newcommand{\MatDeltaSearchGBand}{[$-$0.5, $+$2.1]}
\newcommand{\MatDeltaSearchGYears}{[$+$0.6, $+$1.3]}
\newcommand{\MatDeltaSinkA}{$+$3}
\newcommand{\MatDeltaSinkABand}{[$-$3, $+$8]}
\newcommand{\MatDeltaSinkAYears}{[$+$0, $+$4]}
\newcommand{\MatDeltaSinkD}{$+$3}
\newcommand{\MatDeltaSinkDBand}{[$-$1, $+$8]}
\newcommand{\MatDeltaSinkDYears}{[$+$2, $+$4]}
\newcommand{\MatDeltaSinkG}{$+$6}
\newcommand{\MatDeltaSinkGBand}{[$+$2, $+$10]}
\newcommand{\MatDeltaSinkGYears}{[$+$4, $+$7]}
\newcommand{\MatDeltaXcA}{$-$0.0}
\newcommand{\MatDeltaXcABand}{[$-$0.8, $+$0.8]}
\newcommand{\MatDeltaXcAYears}{[$-$0.2, $+$0.4]}
\newcommand{\MatDeltaXcD}{$+$1.3}
\newcommand{\MatDeltaXcDBand}{[$+$0.4, $+$2.2]}
\newcommand{\MatDeltaXcDYears}{[$+$0.7, $+$1.5]}
\newcommand{\MatDeltaXcG}{$+$1.2}
\newcommand{\MatDeltaXcGBand}{[$+$0.7, $+$1.8]}
\newcommand{\MatDeltaXcGYears}{[$+$1.1, $+$1.4]}
\newcommand{\MatHour}{2}
\newcommand{\MatLevelAllAHK}{$+$0.006}
\newcommand{\MatLevelAllDHK}{$+$0.004}
\newcommand{\MatLevelAllGHK}{$+$0.010}
\newcommand{\MatLevelBeginnersAHK}{$+$0.005}
\newcommand{\MatLevelBeginnersDHK}{$+$0.003}
\newcommand{\MatLevelBeginnersGHK}{$+$0.008}
\newcommand{\MatLevelExperiencedAHK}{$+$0.007}
\newcommand{\MatLevelExperiencedDHK}{$+$0.005}
\newcommand{\MatLevelExperiencedGHK}{$+$0.011}
\newcommand{\MatLongT}{2\,760}
\newcommand{\MatLongTAll}{2\,760}
\newcommand{\MatLongTBeginners}{1\,840}
\newcommand{\MatLongTExperienced}{2\,760}
\newcommand{\MatMatches}{96\,462}
\newcommand{\MatMin}{200}
\newcommand{\MatMinBeg}{279}
\newcommand{\MatMinExp}{421}
\newcommand{\MatMonth}{2}
\newcommand{\MatNAControls}{4\,522}
\newcommand{\MatNAControlsWithPhases}{1\,680}
\newcommand{\MatNAFlightsASide}{1\,504}
\newcommand{\MatNAPieces}{1\,607}
\newcommand{\MatNAPiecesOfTheLongSeries}{338}
\newcommand{\MatNAPiecesWithPhases}{831}
\newcommand{\MatNAPilotsASide}{990}
\newcommand{\MatNAPilotsControlSide}{766}
\newcommand{\MatNAPoolFlights}{1\,191}
\newcommand{\MatNGControls}{8\,094}
\newcommand{\MatNGControlsWithPhases}{3\,210}
\newcommand{\MatNGFlightsGSide}{2\,474}
\newcommand{\MatNGPieces}{2\,717}
\newcommand{\MatNGPiecesOfTheLongSeries}{700}
\newcommand{\MatNGPiecesWithPhases}{1\,422}
\newcommand{\MatNGPilotsControlSide}{887}
\newcommand{\MatNGPilotsGSide}{1\,378}
\newcommand{\MatNGPoolFlights}{1\,404}
\newcommand{\MatNQuant}{9}
\newcommand{\MatPool}{4\,518}
\newcommand{\MatRho}{500}
\newcommand{\MatRobHighHK}{$+$0.012}
\newcommand{\MatRobHighHQ}{$+$0.022}
\newcommand{\MatRobIsolationPieces}{1\,773}
\newcommand{\MatRobLines}{8}
\newcommand{\MatRobLinesHK}{8}
\newcommand{\MatRobLinesHQ}{5}
\newcommand{\MatRobLowHK}{$+$0.007}
\newcommand{\MatRobLowHQ}{$+$0.017}
\newcommand{\MatRobPiecesHigh}{4\,833}
\newcommand{\MatRobPiecesLow}{1\,321}
\newcommand{\MatRobShortHighHQ}{$+$0.034}
\newcommand{\MatRobShortLinesHQ}{3}
\newcommand{\MatRobShortLowHQ}{$+$0.031}
\newcommand{\MatRobShortTHigh}{2\,250}
\newcommand{\MatRobShortTLow}{1\,840}
\newcommand{\MatRobZeroHK}{0}
\newcommand{\MatRobZeroHQ}{0}
\newcommand{\MatRunsClimbA}{4\,972}
\newcommand{\MatRunsClimbG}{10\,063}
\newcommand{\MatRunsClimbSA}{8\,495}
\newcommand{\MatRunsClimbSG}{18\,254}
\newcommand{\MatRunsGlideA}{5\,484}
\newcommand{\MatRunsGlideG}{10\,648}
\newcommand{\MatRunsGlideSA}{9\,706}
\newcommand{\MatRunsGlideSG}{20\,353}
\newcommand{\MatRunsSearchA}{2\,170}
\newcommand{\MatRunsSearchG}{4\,200}
\newcommand{\MatRunsSearchSA}{3\,563}
\newcommand{\MatRunsSearchSG}{8\,008}
\newcommand{\MatSideClimbA}{27.4}
\newcommand{\MatSideClimbABand}{[26.2, 28.6]}
\newcommand{\MatSideClimbG}{29.2}
\newcommand{\MatSideClimbGBand}{[28.3, 30.1]}
\newcommand{\MatSideClimbSA}{26.8}
\newcommand{\MatSideClimbSABand}{[25.2, 28.4]}
\newcommand{\MatSideClimbSG}{26.2}
\newcommand{\MatSideClimbSGBand}{[24.8, 27.7]}
\newcommand{\MatSideCphA}{7.27}
\newcommand{\MatSideCphABand}{[6.99, 7.55]}
\newcommand{\MatSideCphG}{7.56}
\newcommand{\MatSideCphGBand}{[7.39, 7.74]}
\newcommand{\MatSideCphSA}{6.88}
\newcommand{\MatSideCphSABand}{[6.49, 7.27]}
\newcommand{\MatSideCphSG}{6.83}
\newcommand{\MatSideCphSGBand}{[6.50, 7.16]}
\newcommand{\MatSideFitKA}{0.900}
\newcommand{\MatSideFitKABand}{[0.896, 0.904]}
\newcommand{\MatSideFitKG}{0.911}
\newcommand{\MatSideFitKGBand}{[0.907, 0.914]}
\newcommand{\MatSideFitKSA}{0.900}
\newcommand{\MatSideFitKSABand}{[0.896, 0.905]}
\newcommand{\MatSideFitKSG}{0.916}
\newcommand{\MatSideFitKSGBand}{[0.913, 0.919]}
\newcommand{\MatSideFitTA}{0.904}
\newcommand{\MatSideFitTABand}{[0.897, 0.911]}
\newcommand{\MatSideFitTG}{0.911}
\newcommand{\MatSideFitTGBand}{[0.905, 0.916]}
\newcommand{\MatSideFitTSA}{0.899}
\newcommand{\MatSideFitTSABand}{[0.892, 0.905]}
\newcommand{\MatSideFitTSG}{0.903}
\newcommand{\MatSideFitTSGBand}{[0.897, 0.909]}
\newcommand{\MatSideFoundA}{46.7}
\newcommand{\MatSideFoundABand}{[44.7, 48.7]}
\newcommand{\MatSideFoundG}{48.9}
\newcommand{\MatSideFoundGBand}{[47.4, 50.4]}
\newcommand{\MatSideFoundSA}{45.5}
\newcommand{\MatSideFoundSABand}{[43.1, 47.9]}
\newcommand{\MatSideFoundSG}{46.6}
\newcommand{\MatSideFoundSGBand}{[44.6, 48.6]}
\newcommand{\MatSideGlideA}{28.9}
\newcommand{\MatSideGlideABand}{[28.2, 29.5]}
\newcommand{\MatSideGlideG}{29.8}
\newcommand{\MatSideGlideGBand}{[29.3, 30.3]}
\newcommand{\MatSideGlideSA}{29.8}
\newcommand{\MatSideGlideSABand}{[29.1, 30.6]}
\newcommand{\MatSideGlideSG}{30.7}
\newcommand{\MatSideGlideSGBand}{[29.9, 31.5]}
\newcommand{\MatSideHKA}{0.870}
\newcommand{\MatSideHKABand}{[0.861, 0.878]}
\newcommand{\MatSideHKG}{0.907}
\newcommand{\MatSideHKGBand}{[0.902, 0.912]}
\newcommand{\MatSideHKSA}{0.863}
\newcommand{\MatSideHKSABand}{[0.853, 0.874]}
\newcommand{\MatSideHKSG}{0.897}
\newcommand{\MatSideHKSGBand}{[0.890, 0.904]}
\newcommand{\MatSideHQA}{0.866}
\newcommand{\MatSideHQABand}{[0.844, 0.888]}
\newcommand{\MatSideHQG}{0.889}
\newcommand{\MatSideHQGBand}{[0.872, 0.906]}
\newcommand{\MatSideHQSA}{0.839}
\newcommand{\MatSideHQSABand}{[0.811, 0.867]}
\newcommand{\MatSideHQSG}{0.868}
\newcommand{\MatSideHQSGBand}{[0.849, 0.887]}
\newcommand{\MatSidePiecesA}{1\,607}
\newcommand{\MatSidePiecesG}{2\,717}
\newcommand{\MatSidePiecesSA}{4\,522}
\newcommand{\MatSidePiecesSG}{8\,094}
\newcommand{\MatSideRateA}{1.12}
\newcommand{\MatSideRateABand}{[1.09, 1.16]}
\newcommand{\MatSideRateG}{1.17}
\newcommand{\MatSideRateGBand}{[1.14, 1.20]}
\newcommand{\MatSideRateSA}{1.20}
\newcommand{\MatSideRateSABand}{[1.15, 1.25]}
\newcommand{\MatSideRateSG}{1.22}
\newcommand{\MatSideRateSGBand}{[1.18, 1.25]}
\newcommand{\MatSideSearchA}{16.2}
\newcommand{\MatSideSearchABand}{[15.0, 17.3]}
\newcommand{\MatSideSearchG}{14.7}
\newcommand{\MatSideSearchGBand}{[13.8, 15.6]}
\newcommand{\MatSideSearchSA}{14.2}
\newcommand{\MatSideSearchSABand}{[12.7, 15.8]}
\newcommand{\MatSideSearchSG}{13.9}
\newcommand{\MatSideSearchSGBand}{[12.6, 15.1]}
\newcommand{\MatSideSinkA}{71}
\newcommand{\MatSideSinkABand}{[66, 75]}
\newcommand{\MatSideSinkG}{68}
\newcommand{\MatSideSinkGBand}{[65, 71]}
\newcommand{\MatSideSinkSA}{68}
\newcommand{\MatSideSinkSABand}{[62, 74]}
\newcommand{\MatSideSinkSG}{62}
\newcommand{\MatSideSinkSGBand}{[57, 67]}
\newcommand{\MatSideXcA}{16.5}
\newcommand{\MatSideXcABand}{[15.5, 17.5]}
\newcommand{\MatSideXcG}{17.7}
\newcommand{\MatSideXcGBand}{[17.2, 18.3]}
\newcommand{\MatSideXcSA}{16.5}
\newcommand{\MatSideXcSABand}{[15.6, 17.4]}
\newcommand{\MatSideXcSG}{16.5}
\newcommand{\MatSideXcSGBand}{[15.9, 17.1]}
\newcommand{\MatWithMatchA}{3\,851}
\newcommand{\MatWithMatchG}{6\,018}
\newcommand{\MatXcControlsA}{588}
\newcommand{\MatXcControlsG}{1\,976}
\newcommand{\MatXcPiecesA}{368}
\newcommand{\MatXcPiecesG}{1\,127}
\newcommand{\MatZeroA}{5}
\newcommand{\MatZeroD}{8}
\newcommand{\MatZeroG}{2}
\newcommand{\MedianEvery}{0.32}
\newcommand{\MedianFixed}{0.29}
\newcommand{\MedianFixedA}{0.23}
\newcommand{\MedianFixedG}{0.29}
\newcommand{\MedianFixedS}{0.29}
\newcommand{\MedianLongest}{0.22}
\newcommand{\MedianLongestA}{0.19}
\newcommand{\MedianLongestG}{0.28}
\newcommand{\MedianLongestS}{0.23}
\newcommand{\MinOpening}{0.1}
\newcommand{\MountainM}{1\,000}
\newcommand{\NBoot}{1\,000}
\newcommand{\OpenFixedFitA}{0.920}
\newcommand{\OpenFixedFitG}{0.924}
\newcommand{\OpenFixedFitS}{0.915}
\newcommand{\OpenFixedFlightsA}{6\,235}
\newcommand{\OpenFixedFlightsG}{6\,297}
\newcommand{\OpenFixedFlightsS}{33\,613}
\newcommand{\OpenFixedHA}{0.905}
\newcommand{\OpenFixedHABand}{[0.900, 0.911]}
\newcommand{\OpenFixedHG}{0.923}
\newcommand{\OpenFixedHGBand}{[0.918, 0.929]}
\newcommand{\OpenFixedHS}{0.908}
\newcommand{\OpenFixedHSBand}{[0.905, 0.912]}
\newcommand{\OpenFixedHThousandA}{0.898}
\newcommand{\OpenFixedHThousandG}{0.918}
\newcommand{\OpenFixedHThousandS}{0.898}
\newcommand{\OpenLongestFitA}{0.929}
\newcommand{\OpenLongestFitG}{0.943}
\newcommand{\OpenLongestFitS}{0.923}
\newcommand{\OpenLongestFlightsA}{602}
\newcommand{\OpenLongestFlightsG}{115}
\newcommand{\OpenLongestFlightsS}{4\,309}
\newcommand{\OpenLongestHA}{0.872}
\newcommand{\OpenLongestHABand}{[0.846, 0.897]}
\newcommand{\OpenLongestHG}{0.945}
\newcommand{\OpenLongestHGBand}{[0.919, 0.971]}
\newcommand{\OpenLongestHS}{0.878}
\newcommand{\OpenLongestHSBand}{[0.865, 0.890]}
\newcommand{\OpenLongestHThousandA}{0.917}
\newcommand{\OpenLongestHThousandG}{0.936}
\newcommand{\OpenLongestHThousandS}{0.912}
\newcommand{\OpenLowlandAll}{39}
\newcommand{\OpenLowlandStrict}{78}
\newcommand{\OpenMinFlights}{100}
\newcommand{\OpenTauFixed}{2\,760}
\newcommand{\OpenTauLongest}{11\,450}
\newcommand{\OpeningAboveOne}{11}
\newcommand{\PerLagAFirstPieces}{16\,736}
\newcommand{\PerLagALastLag}{11\,450}
\newcommand{\PerLagALastPieces}{131}
\newcommand{\PerLagAMinusCohortMax}{0.006}
\newcommand{\PerLagASHundred}{$+$0.003}
\newcommand{\PerLagASHundredBand}{[$+$0.002, $+$0.004]}
\newcommand{\PerLagASLong}{$+$0.001}
\newcommand{\PerLagASLongBand}{[$-$0.005, $+$0.007]}
\newcommand{\PerLagASTen}{$+$0.007}
\newcommand{\PerLagASTenBand}{[$+$0.006, $+$0.009]}
\newcommand{\PerLagASThousand}{$-$0.001}
\newcommand{\PerLagASThousandBand}{[$-$0.004, $+$0.002]}
\newcommand{\PerLagASThreeTwenty}{$+$0.003}
\newcommand{\PerLagASThreeTwentyBand}{[$+$0.001, $+$0.005]}
\newcommand{\PerLagCoincideEnd}{580}
\newcommand{\PerLagExample}{2\,250}
\newcommand{\PerLagGAHundred}{$-$0.001}
\newcommand{\PerLagGAHundredBand}{[$-$0.002, $-$0.000]}
\newcommand{\PerLagGALong}{$+$0.026}
\newcommand{\PerLagGALongBand}{[$+$0.019, $+$0.033]}
\newcommand{\PerLagGATen}{$-$0.009}
\newcommand{\PerLagGATenBand}{[$-$0.010, $-$0.008]}
\newcommand{\PerLagGAThousand}{$+$0.026}
\newcommand{\PerLagGAThousandBand}{[$+$0.024, $+$0.028]}
\newcommand{\PerLagGAThreeTwenty}{$+$0.011}
\newcommand{\PerLagGAThreeTwentyBand}{[$+$0.009, $+$0.012]}
\newcommand{\PerLagGFirstPieces}{17\,240}
\newcommand{\PerLagGLastLag}{7\,630}
\newcommand{\PerLagGLastPieces}{153}
\newcommand{\PerLagGMinusCohortMax}{0.004}
\newcommand{\PerLagGSHundred}{$+$0.002}
\newcommand{\PerLagGSHundredBand}{[$+$0.001, $+$0.003]}
\newcommand{\PerLagGSLong}{$+$0.027}
\newcommand{\PerLagGSLongBand}{[$+$0.019, $+$0.034]}
\newcommand{\PerLagGSTen}{$-$0.002}
\newcommand{\PerLagGSTenBand}{[$-$0.004, $-$0.000]}
\newcommand{\PerLagGSThousand}{$+$0.025}
\newcommand{\PerLagGSThousandBand}{[$+$0.022, $+$0.028]}
\newcommand{\PerLagGSThreeTwenty}{$+$0.014}
\newcommand{\PerLagGSThreeTwentyBand}{[$+$0.012, $+$0.015]}
\newcommand{\PerLagPoolRunsAwayA}{1\,200}
\newcommand{\PerLagPoolRunsAwayG}{1\,500}
\newcommand{\PerLagPoolRunsAwayS}{3\,400}
\newcommand{\PerLagSFirstPieces}{52\,812}
\newcommand{\PerLagSLastLag}{14\,030}
\newcommand{\PerLagSLastPieces}{394}
\newcommand{\PerLagSMinusCohortMax}{0.004}
\newcommand{\PerLagSelFree}{3\,600}
\newcommand{\PerLagSolidEnd}{2\,250}
\newcommand{\PhAboveClimbAAll}{97}
\newcommand{\PhAboveClimbAHighMtn}{324}
\newcommand{\PhAboveClimbAHills}{22}
\newcommand{\PhAboveClimbALowMtn}{273}
\newcommand{\PhAboveClimbAPlains}{24}
\newcommand{\PhAboveClimbGAll}{67}
\newcommand{\PhAboveClimbGHighMtn}{228}
\newcommand{\PhAboveClimbGHills}{15}
\newcommand{\PhAboveClimbGLowMtn}{177}
\newcommand{\PhAboveClimbGPlains}{26}
\newcommand{\PhAboveClimbSAll}{495}
\newcommand{\PhAboveClimbSHighMtn}{1\,077}
\newcommand{\PhAboveClimbSHills}{168}
\newcommand{\PhAboveClimbSLowMtn}{1\,104}
\newcommand{\PhAboveClimbSPlains}{141}
\newcommand{\PhAboveGlideAAll}{888}
\newcommand{\PhAboveGlideAHighMtn}{444}
\newcommand{\PhAboveGlideAHills}{168}
\newcommand{\PhAboveGlideALowMtn}{566}
\newcommand{\PhAboveGlideAPlains}{127}
\newcommand{\PhAboveGlideGAll}{522}
\newcommand{\PhAboveGlideGHighMtn}{237}
\newcommand{\PhAboveGlideGHills}{114}
\newcommand{\PhAboveGlideGLowMtn}{325}
\newcommand{\PhAboveGlideGPlains}{76}
\newcommand{\PhAboveGlideSAll}{2\,573}
\newcommand{\PhAboveGlideSHighMtn}{1\,253}
\newcommand{\PhAboveGlideSHills}{1\,396}
\newcommand{\PhAboveGlideSLowMtn}{1\,314}
\newcommand{\PhAboveGlideSPlains}{645}
\newcommand{\PhAboveSearchAAll}{592}
\newcommand{\PhAboveSearchAHighMtn}{115}
\newcommand{\PhAboveSearchAHills}{33}
\newcommand{\PhAboveSearchALowMtn}{144}
\newcommand{\PhAboveSearchAPlains}{362}
\newcommand{\PhAboveSearchGAll}{360}
\newcommand{\PhAboveSearchGHighMtn}{60}
\newcommand{\PhAboveSearchGHills}{30}
\newcommand{\PhAboveSearchGLowMtn}{92}
\newcommand{\PhAboveSearchGPlains}{364}
\newcommand{\PhAboveSearchSAll}{2\,168}
\newcommand{\PhAboveSearchSHighMtn}{403}
\newcommand{\PhAboveSearchSHills}{527}
\newcommand{\PhAboveSearchSLowMtn}{489}
\newcommand{\PhAboveSearchSPlains}{2\,205}
\newcommand{\PhChordAAll}{501}
\newcommand{\PhChordAAllBand}{[484, 517]}
\newcommand{\PhChordAHighMtn}{347}
\newcommand{\PhChordAHighMtnBand}{[339, 354]}
\newcommand{\PhChordAHills}{723}
\newcommand{\PhChordAHillsBand}{[694, 753]}
\newcommand{\PhChordALowMtn}{377}
\newcommand{\PhChordALowMtnBand}{[366, 389]}
\newcommand{\PhChordAPlains}{971}
\newcommand{\PhChordAPlainsBand}{[942, 1\,001]}
\newcommand{\PhChordASAll}{$-$39.6}
\newcommand{\PhChordASAllBand}{[$-$53.2, $-$26.0]}
\newcommand{\PhChordASHighMtn}{$+$0.737}
\newcommand{\PhChordASHighMtnBand}{[$-$6.732, $+$8.206]}
\newcommand{\PhChordASHills}{$+$79.6}
\newcommand{\PhChordASHillsBand}{[$+$51.9, $+$107.3]}
\newcommand{\PhChordASLowMtn}{$-$18.5}
\newcommand{\PhChordASLowMtnBand}{[$-$29.0, $-$8.1]}
\newcommand{\PhChordASPlains}{$+$92.3}
\newcommand{\PhChordASPlainsBand}{[$+$62.5, $+$122.1]}
\newcommand{\PhChordGAAll}{$+$6.01}
\newcommand{\PhChordGAAllBand}{[$-$0.71, $+$12.73]}
\newcommand{\PhChordGAHighMtn}{$-$8.86}
\newcommand{\PhChordGAHighMtnBand}{[$-$14.79, $-$2.93]}
\newcommand{\PhChordGAHills}{$+$29.8}
\newcommand{\PhChordGAHillsBand}{[$+$7.2, $+$52.4]}
\newcommand{\PhChordGALowMtn}{$-$10.2}
\newcommand{\PhChordGALowMtnBand}{[$-$16.1, $-$4.4]}
\newcommand{\PhChordGAPlains}{$-$24.7}
\newcommand{\PhChordGAPlainsBand}{[$-$46.2, $-$3.1]}
\newcommand{\PhChordGAll}{507}
\newcommand{\PhChordGAllBand}{[489, 524]}
\newcommand{\PhChordGHighMtn}{338}
\newcommand{\PhChordGHighMtnBand}{[331, 345]}
\newcommand{\PhChordGHills}{753}
\newcommand{\PhChordGHillsBand}{[727, 780]}
\newcommand{\PhChordGLowMtn}{367}
\newcommand{\PhChordGLowMtnBand}{[356, 378]}
\newcommand{\PhChordGPlains}{947}
\newcommand{\PhChordGPlainsBand}{[923, 970]}
\newcommand{\PhChordGSAll}{$-$33.6}
\newcommand{\PhChordGSAllBand}{[$-$48.2, $-$19.1]}
\newcommand{\PhChordGSHighMtn}{$-$8.12}
\newcommand{\PhChordGSHighMtnBand}{[$-$15.44, $-$0.80]}
\newcommand{\PhChordGSHills}{$+$109}
\newcommand{\PhChordGSHillsBand}{[$+$84, $+$135]}
\newcommand{\PhChordGSLowMtn}{$-$28.8}
\newcommand{\PhChordGSLowMtnBand}{[$-$39.1, $-$18.4]}
\newcommand{\PhChordGSPlains}{$+$67.6}
\newcommand{\PhChordGSPlainsBand}{[$+$42.4, $+$92.9]}
\newcommand{\PhChordSAll}{540}
\newcommand{\PhChordSAllBand}{[529, 552]}
\newcommand{\PhChordSHighMtn}{346}
\newcommand{\PhChordSHighMtnBand}{[342, 350]}
\newcommand{\PhChordSHills}{644}
\newcommand{\PhChordSHillsBand}{[630, 658]}
\newcommand{\PhChordSLowMtn}{396}
\newcommand{\PhChordSLowMtnBand}{[388, 404]}
\newcommand{\PhChordSPlains}{879}
\newcommand{\PhChordSPlainsBand}{[860, 898]}
\newcommand{\PhClimbMeanAAll}{167}
\newcommand{\PhClimbMeanAAllBand}{[165, 170]}
\newcommand{\PhClimbMeanAHighMtn}{155}
\newcommand{\PhClimbMeanAHighMtnBand}{[152, 157]}
\newcommand{\PhClimbMeanAHills}{197}
\newcommand{\PhClimbMeanAHillsBand}{[191, 202]}
\newcommand{\PhClimbMeanALowMtn}{155}
\newcommand{\PhClimbMeanALowMtnBand}{[153, 157]}
\newcommand{\PhClimbMeanAPlains}{206}
\newcommand{\PhClimbMeanAPlainsBand}{[201, 210]}
\newcommand{\PhClimbMeanASAll}{$-$1.11}
\newcommand{\PhClimbMeanASAllBand}{[$-$2.97, $+$0.74]}
\newcommand{\PhClimbMeanASHighMtn}{$+$0.577}
\newcommand{\PhClimbMeanASHighMtnBand}{[$-$1.785, $+$2.939]}
\newcommand{\PhClimbMeanASHills}{$+$12.4}
\newcommand{\PhClimbMeanASHillsBand}{[$+$7.2, $+$17.5]}
\newcommand{\PhClimbMeanASLowMtn}{$+$1.50}
\newcommand{\PhClimbMeanASLowMtnBand}{[$-$0.83, $+$3.82]}
\newcommand{\PhClimbMeanASPlains}{$+$13.2}
\newcommand{\PhClimbMeanASPlainsBand}{[$+$9.2, $+$17.2]}
\newcommand{\PhClimbMeanClassBound}{3}
\newcommand{\PhClimbMeanGAAll}{$+$1.12}
\newcommand{\PhClimbMeanGAAllBand}{[$-$0.26, $+$2.51]}
\newcommand{\PhClimbMeanGAHighMtn}{$+$3.68}
\newcommand{\PhClimbMeanGAHighMtnBand}{[$+$1.40, $+$5.96]}
\newcommand{\PhClimbMeanGAHills}{$+$4.39}
\newcommand{\PhClimbMeanGAHillsBand}{[$-$0.07, $+$8.86]}
\newcommand{\PhClimbMeanGALowMtn}{$-$0.876}
\newcommand{\PhClimbMeanGALowMtnBand}{[$-$2.649, $+$0.897]}
\newcommand{\PhClimbMeanGAPlains}{$-$5.26}
\newcommand{\PhClimbMeanGAPlainsBand}{[$-$8.97, $-$1.56]}
\newcommand{\PhClimbMeanGAll}{169}
\newcommand{\PhClimbMeanGAllBand}{[166, 171]}
\newcommand{\PhClimbMeanGHighMtn}{158}
\newcommand{\PhClimbMeanGHighMtnBand}{[156, 161]}
\newcommand{\PhClimbMeanGHills}{201}
\newcommand{\PhClimbMeanGHillsBand}{[197, 206]}
\newcommand{\PhClimbMeanGLowMtn}{154}
\newcommand{\PhClimbMeanGLowMtnBand}{[152, 157]}
\newcommand{\PhClimbMeanGPlains}{200}
\newcommand{\PhClimbMeanGPlainsBand}{[196, 205]}
\newcommand{\PhClimbMeanGSAll}{$+$0.0107}
\newcommand{\PhClimbMeanGSAllBand}{[$-$2.0200, $+$2.0415]}
\newcommand{\PhClimbMeanGSHighMtn}{$+$4.25}
\newcommand{\PhClimbMeanGSHighMtnBand}{[$+$1.64, $+$6.87]}
\newcommand{\PhClimbMeanGSHills}{$+$16.8}
\newcommand{\PhClimbMeanGSHillsBand}{[$+$12.3, $+$21.2]}
\newcommand{\PhClimbMeanGSLowMtn}{$+$0.622}
\newcommand{\PhClimbMeanGSLowMtnBand}{[$-$1.802, $+$3.046]}
\newcommand{\PhClimbMeanGSPlains}{$+$7.94}
\newcommand{\PhClimbMeanGSPlainsBand}{[$+$3.52, $+$12.36]}
\newcommand{\PhClimbMeanSAll}{169}
\newcommand{\PhClimbMeanSAllBand}{[167, 170]}
\newcommand{\PhClimbMeanSHighMtn}{154}
\newcommand{\PhClimbMeanSHighMtnBand}{[152, 155]}
\newcommand{\PhClimbMeanSHills}{184}
\newcommand{\PhClimbMeanSHillsBand}{[182, 187]}
\newcommand{\PhClimbMeanSLowMtn}{154}
\newcommand{\PhClimbMeanSLowMtnBand}{[152, 155]}
\newcommand{\PhClimbMeanSPlains}{192}
\newcommand{\PhClimbMeanSPlainsBand}{[189, 196]}
\newcommand{\PhCphAAll}{6.56}
\newcommand{\PhCphAAllBand}{[6.51, 6.61]}
\newcommand{\PhCphAHighMtn}{6.40}
\newcommand{\PhCphAHighMtnBand}{[6.31, 6.49]}
\newcommand{\PhCphAHills}{7.10}
\newcommand{\PhCphAHillsBand}{[6.96, 7.24]}
\newcommand{\PhCphALowMtn}{6.37}
\newcommand{\PhCphALowMtnBand}{[6.30, 6.43]}
\newcommand{\PhCphAPlains}{7.00}
\newcommand{\PhCphAPlainsBand}{[6.88, 7.12]}
\newcommand{\PhCphASAll}{$-$0.199}
\newcommand{\PhCphASAllBand}{[$-$0.249, $-$0.150]}
\newcommand{\PhCphASHighMtn}{$-$0.130}
\newcommand{\PhCphASHighMtnBand}{[$-$0.227, $-$0.034]}
\newcommand{\PhCphASHills}{$+$0.268}
\newcommand{\PhCphASHillsBand}{[$+$0.122, $+$0.414]}
\newcommand{\PhCphASLowMtn}{$-$0.222}
\newcommand{\PhCphASLowMtnBand}{[$-$0.293, $-$0.152]}
\newcommand{\PhCphASPlains}{$-$0.0933}
\newcommand{\PhCphASPlainsBand}{[$-$0.2128, $+$0.0262]}
\newcommand{\PhCphGAAll}{$+$0.907}
\newcommand{\PhCphGAAllBand}{[$+$0.854, $+$0.959]}
\newcommand{\PhCphGAHighMtn}{$+$1.00}
\newcommand{\PhCphGAHighMtnBand}{[$+$0.89, $+$1.11]}
\newcommand{\PhCphGAHills}{$+$0.460}
\newcommand{\PhCphGAHillsBand}{[$+$0.301, $+$0.619]}
\newcommand{\PhCphGALowMtn}{$+$0.969}
\newcommand{\PhCphGALowMtnBand}{[$+$0.891, $+$1.047]}
\newcommand{\PhCphGAPlains}{$+$0.798}
\newcommand{\PhCphGAPlainsBand}{[$+$0.669, $+$0.927]}
\newcommand{\PhCphGAll}{7.46}
\newcommand{\PhCphGAllBand}{[7.41, 7.51]}
\newcommand{\PhCphGHighMtn}{7.40}
\newcommand{\PhCphGHighMtnBand}{[7.30, 7.50]}
\newcommand{\PhCphGHills}{7.56}
\newcommand{\PhCphGHillsBand}{[7.43, 7.69]}
\newcommand{\PhCphGLowMtn}{7.34}
\newcommand{\PhCphGLowMtnBand}{[7.26, 7.41]}
\newcommand{\PhCphGPlains}{7.80}
\newcommand{\PhCphGPlainsBand}{[7.68, 7.91]}
\newcommand{\PhCphGSAll}{$+$0.707}
\newcommand{\PhCphGSAllBand}{[$+$0.653, $+$0.761]}
\newcommand{\PhCphGSHighMtn}{$+$0.871}
\newcommand{\PhCphGSHighMtnBand}{[$+$0.771, $+$0.970]}
\newcommand{\PhCphGSHills}{$+$0.728}
\newcommand{\PhCphGSHillsBand}{[$+$0.589, $+$0.867]}
\newcommand{\PhCphGSLowMtn}{$+$0.747}
\newcommand{\PhCphGSLowMtnBand}{[$+$0.671, $+$0.823]}
\newcommand{\PhCphGSPlains}{$+$0.704}
\newcommand{\PhCphGSPlainsBand}{[$+$0.591, $+$0.818]}
\newcommand{\PhCphSAll}{6.75}
\newcommand{\PhCphSAllBand}{[6.72, 6.79]}
\newcommand{\PhCphSHighMtn}{6.53}
\newcommand{\PhCphSHighMtnBand}{[6.48, 6.58]}
\newcommand{\PhCphSHills}{6.83}
\newcommand{\PhCphSHillsBand}{[6.77, 6.90]}
\newcommand{\PhCphSLowMtn}{6.59}
\newcommand{\PhCphSLowMtnBand}{[6.54, 6.63]}
\newcommand{\PhCphSPlains}{7.09}
\newcommand{\PhCphSPlainsBand}{[7.03, 7.16]}
\newcommand{\PhCtxExit}{26.2}
\newcommand{\PhCtxExitAll}{26.2}
\newcommand{\PhCtxExitHighMtn}{23.8}
\newcommand{\PhCtxExitHills}{29.8}
\newcommand{\PhCtxExitLowMtn}{24.0}
\newcommand{\PhCtxExitPlains}{31.5}
\newcommand{\PhCtxFailed}{24.3}
\newcommand{\PhCtxFailedAll}{24.3}
\newcommand{\PhCtxFailedHighMtn}{25.4}
\newcommand{\PhCtxFailedHills}{21.4}
\newcommand{\PhCtxFailedLowMtn}{26.9}
\newcommand{\PhCtxFailedPlains}{18.7}
\newcommand{\PhCtxFound}{21.8}
\newcommand{\PhCtxFoundAll}{21.8}
\newcommand{\PhCtxFoundHighMtn}{23.5}
\newcommand{\PhCtxFoundHills}{19.7}
\newcommand{\PhCtxFoundLowMtn}{22.7}
\newcommand{\PhCtxFoundPlains}{19.0}
\newcommand{\PhCtxMaxDiff}{2.0}
\newcommand{\PhCtxRefound}{23.1}
\newcommand{\PhCtxRefoundAll}{23.1}
\newcommand{\PhCtxRefoundHighMtn}{22.8}
\newcommand{\PhCtxRefoundHills}{23.8}
\newcommand{\PhCtxRefoundLowMtn}{21.7}
\newcommand{\PhCtxRefoundPlains}{26.8}
\newcommand{\PhCycleAAll}{9.15}
\newcommand{\PhCycleAAllBand}{[9.08, 9.22]}
\newcommand{\PhCycleAHighMtn}{9.38}
\newcommand{\PhCycleAHighMtnBand}{[9.24, 9.51]}
\newcommand{\PhCycleAHills}{8.45}
\newcommand{\PhCycleAHillsBand}{[8.29, 8.62]}
\newcommand{\PhCycleALowMtn}{9.42}
\newcommand{\PhCycleALowMtnBand}{[9.33, 9.52]}
\newcommand{\PhCycleAPlains}{8.57}
\newcommand{\PhCycleAPlainsBand}{[8.43, 8.72]}
\newcommand{\PhCycleASAll}{$+$0.270}
\newcommand{\PhCycleASAllBand}{[$+$0.202, $+$0.338]}
\newcommand{\PhCycleASHighMtn}{$+$0.187}
\newcommand{\PhCycleASHighMtnBand}{[$+$0.047, $+$0.328]}
\newcommand{\PhCycleASHills}{$-$0.332}
\newcommand{\PhCycleASHillsBand}{[$-$0.508, $-$0.156]}
\newcommand{\PhCycleASLowMtn}{$+$0.318}
\newcommand{\PhCycleASLowMtnBand}{[$+$0.215, $+$0.421]}
\newcommand{\PhCycleASPlains}{$+$0.113}
\newcommand{\PhCycleASPlainsBand}{[$-$0.033, $+$0.258]}
\newcommand{\PhCycleGAAll}{$-$1.11}
\newcommand{\PhCycleGAAllBand}{[$-$1.18, $-$1.05]}
\newcommand{\PhCycleGAHighMtn}{$-$1.27}
\newcommand{\PhCycleGAHighMtnBand}{[$-$1.41, $-$1.12]}
\newcommand{\PhCycleGAHills}{$-$0.514}
\newcommand{\PhCycleGAHillsBand}{[$-$0.694, $-$0.334]}
\newcommand{\PhCycleGALowMtn}{$-$1.25}
\newcommand{\PhCycleGALowMtnBand}{[$-$1.35, $-$1.14]}
\newcommand{\PhCycleGAPlains}{$-$0.877}
\newcommand{\PhCycleGAPlainsBand}{[$-$1.023, $-$0.731]}
\newcommand{\PhCycleGAll}{8.04}
\newcommand{\PhCycleGAllBand}{[7.98, 8.10]}
\newcommand{\PhCycleGHighMtn}{8.11}
\newcommand{\PhCycleGHighMtnBand}{[8.00, 8.22]}
\newcommand{\PhCycleGHills}{7.94}
\newcommand{\PhCycleGHillsBand}{[7.80, 8.07]}
\newcommand{\PhCycleGLowMtn}{8.18}
\newcommand{\PhCycleGLowMtnBand}{[8.10, 8.26]}
\newcommand{\PhCycleGPlains}{7.69}
\newcommand{\PhCycleGPlainsBand}{[7.58, 7.81]}
\newcommand{\PhCycleGSAll}{$-$0.842}
\newcommand{\PhCycleGSAllBand}{[$-$0.904, $-$0.780]}
\newcommand{\PhCycleGSHighMtn}{$-$1.08}
\newcommand{\PhCycleGSHighMtnBand}{[$-$1.20, $-$0.97]}
\newcommand{\PhCycleGSHills}{$-$0.846}
\newcommand{\PhCycleGSHillsBand}{[$-$0.999, $-$0.694]}
\newcommand{\PhCycleGSLowMtn}{$-$0.927}
\newcommand{\PhCycleGSLowMtnBand}{[$-$1.017, $-$0.837]}
\newcommand{\PhCycleGSPlains}{$-$0.764}
\newcommand{\PhCycleGSPlainsBand}{[$-$0.880, $-$0.648]}
\newcommand{\PhCycleMax}{9.4}
\newcommand{\PhCycleMin}{7.7}
\newcommand{\PhCycleRatioMax}{4.6}
\newcommand{\PhCycleRatioMin}{1.3}
\newcommand{\PhCycleSAll}{8.88}
\newcommand{\PhCycleSAllBand}{[8.84, 8.93]}
\newcommand{\PhCycleSHighMtn}{9.19}
\newcommand{\PhCycleSHighMtnBand}{[9.12, 9.26]}
\newcommand{\PhCycleSHills}{8.78}
\newcommand{\PhCycleSHillsBand}{[8.70, 8.87]}
\newcommand{\PhCycleSLowMtn}{9.11}
\newcommand{\PhCycleSLowMtnBand}{[9.04, 9.17]}
\newcommand{\PhCycleSPlains}{8.46}
\newcommand{\PhCycleSPlainsBand}{[8.38, 8.53]}
\newcommand{\PhCycleShortAll}{1.11}
\newcommand{\PhCycleShortMax}{1.3}
\newcommand{\PhCycleShortMin}{0.5}
\newcommand{\PhDriftAAll}{3.00}
\newcommand{\PhDriftAAllBand}{[2.93, 3.07]}
\newcommand{\PhDriftAHighMtn}{2.26}
\newcommand{\PhDriftAHighMtnBand}{[2.22, 2.30]}
\newcommand{\PhDriftAHills}{3.68}
\newcommand{\PhDriftAHillsBand}{[3.56, 3.80]}
\newcommand{\PhDriftALowMtn}{2.45}
\newcommand{\PhDriftALowMtnBand}{[2.40, 2.50]}
\newcommand{\PhDriftAPlains}{4.72}
\newcommand{\PhDriftAPlainsBand}{[4.58, 4.86]}
\newcommand{\PhDriftASAll}{$-$0.226}
\newcommand{\PhDriftASAllBand}{[$-$0.287, $-$0.164]}
\newcommand{\PhDriftASHighMtn}{$-$0.0069}
\newcommand{\PhDriftASHighMtnBand}{[$-$0.0453, $+$0.0315]}
\newcommand{\PhDriftASHills}{$+$0.166}
\newcommand{\PhDriftASHillsBand}{[$+$0.054, $+$0.278]}
\newcommand{\PhDriftASLowMtn}{$-$0.147}
\newcommand{\PhDriftASLowMtnBand}{[$-$0.193, $-$0.101]}
\newcommand{\PhDriftASPlains}{$+$0.134}
\newcommand{\PhDriftASPlainsBand}{[$+$0.007, $+$0.261]}
\newcommand{\PhDriftGAAll}{$+$0.0109}
\newcommand{\PhDriftGAAllBand}{[$-$0.0164, $+$0.0382]}
\newcommand{\PhDriftGAHighMtn}{$-$0.111}
\newcommand{\PhDriftGAHighMtnBand}{[$-$0.141, $-$0.081]}
\newcommand{\PhDriftGAHills}{$+$0.0539}
\newcommand{\PhDriftGAHillsBand}{[$-$0.0164, $+$0.1243]}
\newcommand{\PhDriftGALowMtn}{$-$0.0580}
\newcommand{\PhDriftGALowMtnBand}{[$-$0.0849, $-$0.0311]}
\newcommand{\PhDriftGAPlains}{$+$0.0148}
\newcommand{\PhDriftGAPlainsBand}{[$-$0.0423, $+$0.0719]}
\newcommand{\PhDriftGAll}{3.01}
\newcommand{\PhDriftGAllBand}{[2.93, 3.09]}
\newcommand{\PhDriftGHighMtn}{2.15}
\newcommand{\PhDriftGHighMtnBand}{[2.11, 2.18]}
\newcommand{\PhDriftGHills}{3.73}
\newcommand{\PhDriftGHillsBand}{[3.63, 3.84]}
\newcommand{\PhDriftGLowMtn}{2.39}
\newcommand{\PhDriftGLowMtnBand}{[2.34, 2.44]}
\newcommand{\PhDriftGPlains}{4.73}
\newcommand{\PhDriftGPlainsBand}{[4.61, 4.86]}
\newcommand{\PhDriftGSAll}{$-$0.215}
\newcommand{\PhDriftGSAllBand}{[$-$0.280, $-$0.149]}
\newcommand{\PhDriftGSHighMtn}{$-$0.118}
\newcommand{\PhDriftGSHighMtnBand}{[$-$0.153, $-$0.082]}
\newcommand{\PhDriftGSHills}{$+$0.220}
\newcommand{\PhDriftGSHillsBand}{[$+$0.119, $+$0.321]}
\newcommand{\PhDriftGSLowMtn}{$-$0.205}
\newcommand{\PhDriftGSLowMtnBand}{[$-$0.252, $-$0.158]}
\newcommand{\PhDriftGSPlains}{$+$0.148}
\newcommand{\PhDriftGSPlainsBand}{[$+$0.035, $+$0.262]}
\newcommand{\PhDriftSAll}{3.23}
\newcommand{\PhDriftSAllBand}{[3.18, 3.27]}
\newcommand{\PhDriftSHighMtn}{2.26}
\newcommand{\PhDriftSHighMtnBand}{[2.24, 2.29]}
\newcommand{\PhDriftSHills}{3.51}
\newcommand{\PhDriftSHillsBand}{[3.46, 3.57]}
\newcommand{\PhDriftSLowMtn}{2.59}
\newcommand{\PhDriftSLowMtnBand}{[2.56, 2.63]}
\newcommand{\PhDriftSPlains}{4.58}
\newcommand{\PhDriftSPlainsBand}{[4.52, 4.65]}
\newcommand{\PhDurLimitAll}{4\,609}
\newcommand{\PhDurMeanClimbAAll}{171}
\newcommand{\PhDurMeanClimbAAllBand}{[169, 173]}
\newcommand{\PhDurMeanClimbAHighMtn}{162}
\newcommand{\PhDurMeanClimbAHighMtnBand}{[160, 164]}
\newcommand{\PhDurMeanClimbAHills}{206}
\newcommand{\PhDurMeanClimbAHillsBand}{[201, 211]}
\newcommand{\PhDurMeanClimbALowMtn}{161}
\newcommand{\PhDurMeanClimbALowMtnBand}{[159, 164]}
\newcommand{\PhDurMeanClimbAPlains}{211}
\newcommand{\PhDurMeanClimbAPlainsBand}{[206, 215]}
\newcommand{\PhDurMeanClimbASAll}{$-$4.08}
\newcommand{\PhDurMeanClimbASAllBand}{[$-$5.82, $-$2.35]}
\newcommand{\PhDurMeanClimbASHighMtn}{$+$1.11}
\newcommand{\PhDurMeanClimbASHighMtnBand}{[$-$1.02, $+$3.24]}
\newcommand{\PhDurMeanClimbASHills}{$+$9.39}
\newcommand{\PhDurMeanClimbASHillsBand}{[$+$4.33, $+$14.45]}
\newcommand{\PhDurMeanClimbASLowMtn}{$-$0.386}
\newcommand{\PhDurMeanClimbASLowMtnBand}{[$-$2.550, $+$1.778]}
\newcommand{\PhDurMeanClimbASPlains}{$+$5.09}
\newcommand{\PhDurMeanClimbASPlainsBand}{[$+$1.19, $+$8.98]}
\newcommand{\PhDurMeanClimbGAAll}{$+$2.65}
\newcommand{\PhDurMeanClimbGAAllBand}{[$+$0.83, $+$4.47]}
\newcommand{\PhDurMeanClimbGAHighMtn}{$+$1.19}
\newcommand{\PhDurMeanClimbGAHighMtnBand}{[$-$1.11, $+$3.48]}
\newcommand{\PhDurMeanClimbGAHills}{$+$2.55}
\newcommand{\PhDurMeanClimbGAHillsBand}{[$-$1.81, $+$6.91]}
\newcommand{\PhDurMeanClimbGALowMtn}{$-$2.20}
\newcommand{\PhDurMeanClimbGALowMtnBand}{[$-$4.31, $-$0.09]}
\newcommand{\PhDurMeanClimbGAPlains}{$-$5.90}
\newcommand{\PhDurMeanClimbGAPlainsBand}{[$-$10.19, $-$1.61]}
\newcommand{\PhDurMeanClimbGAll}{174}
\newcommand{\PhDurMeanClimbGAllBand}{[172, 177]}
\newcommand{\PhDurMeanClimbGHighMtn}{163}
\newcommand{\PhDurMeanClimbGHighMtnBand}{[161, 166]}
\newcommand{\PhDurMeanClimbGHills}{209}
\newcommand{\PhDurMeanClimbGHillsBand}{[204, 213]}
\newcommand{\PhDurMeanClimbGLowMtn}{159}
\newcommand{\PhDurMeanClimbGLowMtnBand}{[156, 162]}
\newcommand{\PhDurMeanClimbGPlains}{205}
\newcommand{\PhDurMeanClimbGPlainsBand}{[199, 210]}
\newcommand{\PhDurMeanClimbGSAll}{$-$1.43}
\newcommand{\PhDurMeanClimbGSAllBand}{[$-$3.82, $+$0.96]}
\newcommand{\PhDurMeanClimbGSHighMtn}{$+$2.29}
\newcommand{\PhDurMeanClimbGSHighMtnBand}{[$-$0.31, $+$4.90]}
\newcommand{\PhDurMeanClimbGSHills}{$+$11.9}
\newcommand{\PhDurMeanClimbGSHillsBand}{[$+$7.2, $+$16.6]}
\newcommand{\PhDurMeanClimbGSLowMtn}{$-$2.59}
\newcommand{\PhDurMeanClimbGSLowMtnBand}{[$-$5.36, $+$0.18]}
\newcommand{\PhDurMeanClimbGSPlains}{$-$0.815}
\newcommand{\PhDurMeanClimbGSPlainsBand}{[$-$6.080, $+$4.449]}
\newcommand{\PhDurMeanClimbSAll}{176}
\newcommand{\PhDurMeanClimbSAllBand}{[174, 177]}
\newcommand{\PhDurMeanClimbSHighMtn}{161}
\newcommand{\PhDurMeanClimbSHighMtnBand}{[159, 162]}
\newcommand{\PhDurMeanClimbSHills}{197}
\newcommand{\PhDurMeanClimbSHillsBand}{[194, 199]}
\newcommand{\PhDurMeanClimbSLowMtn}{162}
\newcommand{\PhDurMeanClimbSLowMtnBand}{[160, 163]}
\newcommand{\PhDurMeanClimbSPlains}{205}
\newcommand{\PhDurMeanClimbSPlainsBand}{[202, 209]}
\newcommand{\PhDurMeanGlideAAll}{300}
\newcommand{\PhDurMeanGlideAAllBand}{[296, 304]}
\newcommand{\PhDurMeanGlideAHighMtn}{315}
\newcommand{\PhDurMeanGlideAHighMtnBand}{[310, 319]}
\newcommand{\PhDurMeanGlideAHills}{265}
\newcommand{\PhDurMeanGlideAHillsBand}{[259, 272]}
\newcommand{\PhDurMeanGlideALowMtn}{306}
\newcommand{\PhDurMeanGlideALowMtnBand}{[302, 311]}
\newcommand{\PhDurMeanGlideAPlains}{258}
\newcommand{\PhDurMeanGlideAPlainsBand}{[252, 264]}
\newcommand{\PhDurMeanGlideASAll}{$+$26.2}
\newcommand{\PhDurMeanGlideASAllBand}{[$+$22.9, $+$29.5]}
\newcommand{\PhDurMeanGlideASHighMtn}{$+$18.5}
\newcommand{\PhDurMeanGlideASHighMtnBand}{[$+$13.6, $+$23.3]}
\newcommand{\PhDurMeanGlideASHills}{$+$6.40}
\newcommand{\PhDurMeanGlideASHillsBand}{[$+$0.24, $+$12.56]}
\newcommand{\PhDurMeanGlideASLowMtn}{$+$26.1}
\newcommand{\PhDurMeanGlideASLowMtnBand}{[$+$21.4, $+$30.8]}
\newcommand{\PhDurMeanGlideASPlains}{$+$11.6}
\newcommand{\PhDurMeanGlideASPlainsBand}{[$+$5.9, $+$17.2]}
\newcommand{\PhDurMeanGlideGAAll}{$-$31.6}
\newcommand{\PhDurMeanGlideGAAllBand}{[$-$34.6, $-$28.7]}
\newcommand{\PhDurMeanGlideGAHighMtn}{$-$27.1}
\newcommand{\PhDurMeanGlideGAHighMtnBand}{[$-$32.4, $-$21.8]}
\newcommand{\PhDurMeanGlideGAHills}{$-$20.7}
\newcommand{\PhDurMeanGlideGAHillsBand}{[$-$27.3, $-$14.1]}
\newcommand{\PhDurMeanGlideGALowMtn}{$-$28.4}
\newcommand{\PhDurMeanGlideGALowMtnBand}{[$-$33.0, $-$23.8]}
\newcommand{\PhDurMeanGlideGAPlains}{$-$28.9}
\newcommand{\PhDurMeanGlideGAPlainsBand}{[$-$34.6, $-$23.2]}
\newcommand{\PhDurMeanGlideGAll}{268}
\newcommand{\PhDurMeanGlideGAllBand}{[265, 272]}
\newcommand{\PhDurMeanGlideGHighMtn}{288}
\newcommand{\PhDurMeanGlideGHighMtnBand}{[282, 294]}
\newcommand{\PhDurMeanGlideGHills}{245}
\newcommand{\PhDurMeanGlideGHillsBand}{[238, 251]}
\newcommand{\PhDurMeanGlideGLowMtn}{278}
\newcommand{\PhDurMeanGlideGLowMtnBand}{[273, 283]}
\newcommand{\PhDurMeanGlideGPlains}{229}
\newcommand{\PhDurMeanGlideGPlainsBand}{[223, 235]}
\newcommand{\PhDurMeanGlideGSAll}{$-$5.42}
\newcommand{\PhDurMeanGlideGSAllBand}{[$-$8.81, $-$2.04]}
\newcommand{\PhDurMeanGlideGSHighMtn}{$-$8.63}
\newcommand{\PhDurMeanGlideGSHighMtnBand}{[$-$14.58, $-$2.67]}
\newcommand{\PhDurMeanGlideGSHills}{$-$14.3}
\newcommand{\PhDurMeanGlideGSHillsBand}{[$-$20.8, $-$7.8]}
\newcommand{\PhDurMeanGlideGSLowMtn}{$-$2.24}
\newcommand{\PhDurMeanGlideGSLowMtnBand}{[$-$6.93, $+$2.45]}
\newcommand{\PhDurMeanGlideGSPlains}{$-$17.4}
\newcommand{\PhDurMeanGlideGSPlainsBand}{[$-$23.7, $-$11.0]}
\newcommand{\PhDurMeanGlideSAll}{274}
\newcommand{\PhDurMeanGlideSAllBand}{[272, 276]}
\newcommand{\PhDurMeanGlideSHighMtn}{296}
\newcommand{\PhDurMeanGlideSHighMtnBand}{[293, 299]}
\newcommand{\PhDurMeanGlideSHills}{259}
\newcommand{\PhDurMeanGlideSHillsBand}{[256, 262]}
\newcommand{\PhDurMeanGlideSLowMtn}{280}
\newcommand{\PhDurMeanGlideSLowMtnBand}{[278, 283]}
\newcommand{\PhDurMeanGlideSPlains}{247}
\newcommand{\PhDurMeanGlideSPlainsBand}{[244, 250]}
\newcommand{\PhDurMeanSearchAAll}{149}
\newcommand{\PhDurMeanSearchAAllBand}{[141, 157]}
\newcommand{\PhDurMeanSearchASAll}{$+$4.77}
\newcommand{\PhDurMeanSearchASAllBand}{[$-$4.08, $+$13.62]}
\newcommand{\PhDurMeanSearchGAAll}{$-$29.4}
\newcommand{\PhDurMeanSearchGAAllBand}{[$-$37.7, $-$21.1]}
\newcommand{\PhDurMeanSearchGAll}{120}
\newcommand{\PhDurMeanSearchGAllBand}{[116, 123]}
\newcommand{\PhDurMeanSearchSAll}{144}
\newcommand{\PhDurMeanSearchSAllBand}{[140, 149]}
\newcommand{\PhDurMedianClimbAAll}{129}
\newcommand{\PhDurMedianClimbAAllBand}{[127, 131]}
\newcommand{\PhDurMedianClimbAHighMtn}{124}
\newcommand{\PhDurMedianClimbAHighMtnBand}{[122, 126]}
\newcommand{\PhDurMedianClimbAHills}{155}
\newcommand{\PhDurMedianClimbAHillsBand}{[150, 160]}
\newcommand{\PhDurMedianClimbALowMtn}{121}
\newcommand{\PhDurMedianClimbALowMtnBand}{[119, 123]}
\newcommand{\PhDurMedianClimbAPlains}{158}
\newcommand{\PhDurMedianClimbAPlainsBand}{[153, 163]}
\newcommand{\PhDurMedianClimbASAll}{$-$1.00}
\newcommand{\PhDurMedianClimbASAllBand}{[$-$2.59, $+$0.59]}
\newcommand{\PhDurMedianClimbASHighMtn}{$+$2.00}
\newcommand{\PhDurMedianClimbASHighMtnBand}{[$-$0.26, $+$4.26]}
\newcommand{\PhDurMedianClimbASHills}{$+$10.0}
\newcommand{\PhDurMedianClimbASHillsBand}{[$+$5.1, $+$14.9]}
\newcommand{\PhDurMedianClimbASLowMtn}{$+$1.00}
\newcommand{\PhDurMedianClimbASLowMtnBand}{[$-$0.83, $+$2.83]}
\newcommand{\PhDurMedianClimbASPlains}{$+$5.00}
\newcommand{\PhDurMedianClimbASPlainsBand}{[$+$0.95, $+$9.05]}
\newcommand{\PhDurMedianClimbGAAll}{$+$3.00}
\newcommand{\PhDurMedianClimbGAAllBand}{[$+$1.07, $+$4.93]}
\newcommand{\PhDurMedianClimbGAHighMtn}{$+$2.00}
\newcommand{\PhDurMedianClimbGAHighMtnBand}{[$-$0.35, $+$4.35]}
\newcommand{\PhDurMedianClimbGAHills}{$+$7.00}
\newcommand{\PhDurMedianClimbGAHillsBand}{[$+$2.75, $+$11.25]}
\newcommand{\PhDurMedianClimbGALowMtn}{$-$1.00}
\newcommand{\PhDurMedianClimbGALowMtnBand}{[$-$2.91, $+$0.91]}
\newcommand{\PhDurMedianClimbGAPlains}{$+$2.00}
\newcommand{\PhDurMedianClimbGAPlainsBand}{[$-$2.59, $+$6.59]}
\newcommand{\PhDurMedianClimbGAll}{132}
\newcommand{\PhDurMedianClimbGAllBand}{[130, 134]}
\newcommand{\PhDurMedianClimbGHighMtn}{126}
\newcommand{\PhDurMedianClimbGHighMtnBand}{[123, 129]}
\newcommand{\PhDurMedianClimbGHills}{162}
\newcommand{\PhDurMedianClimbGHillsBand}{[157, 167]}
\newcommand{\PhDurMedianClimbGLowMtn}{120}
\newcommand{\PhDurMedianClimbGLowMtnBand}{[118, 122]}
\newcommand{\PhDurMedianClimbGPlains}{160}
\newcommand{\PhDurMedianClimbGPlainsBand}{[155, 165]}
\newcommand{\PhDurMedianClimbGSAll}{$+$2.00}
\newcommand{\PhDurMedianClimbGSAllBand}{[$-$0.26, $+$4.26]}
\newcommand{\PhDurMedianClimbGSHighMtn}{$+$4.00}
\newcommand{\PhDurMedianClimbGSHighMtnBand}{[$+$1.37, $+$6.63]}
\newcommand{\PhDurMedianClimbGSHills}{$+$17.0}
\newcommand{\PhDurMedianClimbGSHillsBand}{[$+$12.1, $+$21.9]}
\newcommand{\PhDurMedianClimbGSLowMtn}{0}
\newcommand{\PhDurMedianClimbGSLowMtnBand}{[$-$2, $+$2]}
\newcommand{\PhDurMedianClimbGSPlains}{$+$7.00}
\newcommand{\PhDurMedianClimbGSPlainsBand}{[$+$1.82, $+$12.18]}
\newcommand{\PhDurMedianClimbSAll}{130}
\newcommand{\PhDurMedianClimbSAllBand}{[129, 131]}
\newcommand{\PhDurMedianClimbSHighMtn}{122}
\newcommand{\PhDurMedianClimbSHighMtnBand}{[120, 124]}
\newcommand{\PhDurMedianClimbSHills}{145}
\newcommand{\PhDurMedianClimbSHillsBand}{[142, 148]}
\newcommand{\PhDurMedianClimbSLowMtn}{120}
\newcommand{\PhDurMedianClimbSLowMtnBand}{[119, 121]}
\newcommand{\PhDurMedianClimbSPlains}{153}
\newcommand{\PhDurMedianClimbSPlainsBand}{[149, 157]}
\newcommand{\PhDurMedianGlideAAll}{198}
\newcommand{\PhDurMedianGlideAAllBand}{[195, 201]}
\newcommand{\PhDurMedianGlideAHighMtn}{206}
\newcommand{\PhDurMedianGlideAHighMtnBand}{[202, 210]}
\newcommand{\PhDurMedianGlideAHills}{195}
\newcommand{\PhDurMedianGlideAHillsBand}{[189, 201]}
\newcommand{\PhDurMedianGlideALowMtn}{192}
\newcommand{\PhDurMedianGlideALowMtnBand}{[189, 195]}
\newcommand{\PhDurMedianGlideAPlains}{201}
\newcommand{\PhDurMedianGlideAPlainsBand}{[195, 207]}
\newcommand{\PhDurMedianGlideASAll}{$+$11.0}
\newcommand{\PhDurMedianGlideASAllBand}{[$+$8.3, $+$13.7]}
\newcommand{\PhDurMedianGlideASHighMtn}{$+$15.0}
\newcommand{\PhDurMedianGlideASHighMtnBand}{[$+$10.3, $+$19.7]}
\newcommand{\PhDurMedianGlideASHills}{$+$5.00}
\newcommand{\PhDurMedianGlideASHillsBand}{[$-$0.66, $+$10.66]}
\newcommand{\PhDurMedianGlideASLowMtn}{$+$9.00}
\newcommand{\PhDurMedianGlideASLowMtnBand}{[$+$5.42, $+$12.58]}
\newcommand{\PhDurMedianGlideASPlains}{$+$13.0}
\newcommand{\PhDurMedianGlideASPlainsBand}{[$+$6.9, $+$19.1]}
\newcommand{\PhDurMedianGlideGAAll}{$-$15.0}
\newcommand{\PhDurMedianGlideGAAllBand}{[$-$17.7, $-$12.3]}
\newcommand{\PhDurMedianGlideGAHighMtn}{$-$16.0}
\newcommand{\PhDurMedianGlideGAHighMtnBand}{[$-$20.9, $-$11.1]}
\newcommand{\PhDurMedianGlideGAHills}{$-$7.00}
\newcommand{\PhDurMedianGlideGAHillsBand}{[$-$14.03, $+$0.03]}
\newcommand{\PhDurMedianGlideGALowMtn}{$-$12.0}
\newcommand{\PhDurMedianGlideGALowMtnBand}{[$-$15.9, $-$8.1]}
\newcommand{\PhDurMedianGlideGAPlains}{$-$21.0}
\newcommand{\PhDurMedianGlideGAPlainsBand}{[$-$26.7, $-$15.3]}
\newcommand{\PhDurMedianGlideGAll}{183}
\newcommand{\PhDurMedianGlideGAllBand}{[180, 186]}
\newcommand{\PhDurMedianGlideGHighMtn}{190}
\newcommand{\PhDurMedianGlideGHighMtnBand}{[185, 195]}
\newcommand{\PhDurMedianGlideGHills}{188}
\newcommand{\PhDurMedianGlideGHillsBand}{[182, 194]}
\newcommand{\PhDurMedianGlideGLowMtn}{180}
\newcommand{\PhDurMedianGlideGLowMtnBand}{[176, 184]}
\newcommand{\PhDurMedianGlideGPlains}{180}
\newcommand{\PhDurMedianGlideGPlainsBand}{[175, 185]}
\newcommand{\PhDurMedianGlideGSAll}{$-$4.00}
\newcommand{\PhDurMedianGlideGSAllBand}{[$-$7.05, $-$0.95]}
\newcommand{\PhDurMedianGlideGSHighMtn}{$-$1.00}
\newcommand{\PhDurMedianGlideGSHighMtnBand}{[$-$6.57, $+$4.57]}
\newcommand{\PhDurMedianGlideGSHills}{$-$2.00}
\newcommand{\PhDurMedianGlideGSHillsBand}{[$-$8.69, $+$4.69]}
\newcommand{\PhDurMedianGlideGSLowMtn}{$-$3.00}
\newcommand{\PhDurMedianGlideGSLowMtnBand}{[$-$6.98, $+$0.98]}
\newcommand{\PhDurMedianGlideGSPlains}{$-$8.00}
\newcommand{\PhDurMedianGlideGSPlainsBand}{[$-$13.83, $-$2.17]}
\newcommand{\PhDurMedianGlideSAll}{187}
\newcommand{\PhDurMedianGlideSAllBand}{[185, 189]}
\newcommand{\PhDurMedianGlideSHighMtn}{191}
\newcommand{\PhDurMedianGlideSHighMtnBand}{[188, 194]}
\newcommand{\PhDurMedianGlideSHills}{190}
\newcommand{\PhDurMedianGlideSHillsBand}{[187, 193]}
\newcommand{\PhDurMedianGlideSLowMtn}{183}
\newcommand{\PhDurMedianGlideSLowMtnBand}{[181, 185]}
\newcommand{\PhDurMedianGlideSPlains}{188}
\newcommand{\PhDurMedianGlideSPlainsBand}{[185, 191]}
\newcommand{\PhDurMedianSearchAAll}{80.0}
\newcommand{\PhDurMedianSearchAAllBand}{[79.0, 81.0]}
\newcommand{\PhDurMedianSearchASAll}{$+$1.00}
\newcommand{\PhDurMedianSearchASAllBand}{[$-$0.12, $+$2.12]}
\newcommand{\PhDurMedianSearchGAAll}{$-$6.00}
\newcommand{\PhDurMedianSearchGAAllBand}{[$-$7.16, $-$4.84]}
\newcommand{\PhDurMedianSearchGAll}{74.0}
\newcommand{\PhDurMedianSearchGAllBand}{[72.8, 75.2]}
\newcommand{\PhDurMedianSearchSAll}{79.0}
\newcommand{\PhDurMedianSearchSAllBand}{[78.0, 80.0]}
\newcommand{\PhExitDurAAll}{50.0}
\newcommand{\PhExitDurAAllBand}{[49.4, 50.6]}
\newcommand{\PhExitDurAHighMtn}{49.0}
\newcommand{\PhExitDurAHighMtnBand}{[47.8, 50.2]}
\newcommand{\PhExitDurAHills}{50.0}
\newcommand{\PhExitDurAHillsBand}{[48.4, 51.6]}
\newcommand{\PhExitDurALowMtn}{50.0}
\newcommand{\PhExitDurALowMtnBand}{[49.1, 50.9]}
\newcommand{\PhExitDurAPlains}{54.0}
\newcommand{\PhExitDurAPlainsBand}{[52.5, 55.5]}
\newcommand{\PhExitDurASAll}{0}
\newcommand{\PhExitDurASAllBand}{[$-$1, $+$1]}
\newcommand{\PhExitDurASHighMtn}{$-$1.00}
\newcommand{\PhExitDurASHighMtnBand}{[$-$2.38, $+$0.38]}
\newcommand{\PhExitDurASHills}{0}
\newcommand{\PhExitDurASHillsBand}{[$-$2, $+$2]}
\newcommand{\PhExitDurASLowMtn}{0}
\newcommand{\PhExitDurASLowMtnBand}{[$-$1, $+$1]}
\newcommand{\PhExitDurASPlains}{$-$1.00}
\newcommand{\PhExitDurASPlainsBand}{[$-$2.57, $+$0.57]}
\newcommand{\PhExitDurGAAll}{$+$2.00}
\newcommand{\PhExitDurGAAllBand}{[$+$1.16, $+$2.84]}
\newcommand{\PhExitDurGAHighMtn}{$+$1.00}
\newcommand{\PhExitDurGAHighMtnBand}{[$-$0.60, $+$2.60]}
\newcommand{\PhExitDurGAHills}{$+$2.00}
\newcommand{\PhExitDurGAHillsBand}{[$-$0.04, $+$4.04]}
\newcommand{\PhExitDurGALowMtn}{$+$1.00}
\newcommand{\PhExitDurGALowMtnBand}{[$-$0.35, $+$2.35]}
\newcommand{\PhExitDurGAPlains}{$+$3.00}
\newcommand{\PhExitDurGAPlainsBand}{[$+$0.98, $+$5.02]}
\newcommand{\PhExitDurGAll}{52.0}
\newcommand{\PhExitDurGAllBand}{[51.3, 52.7]}
\newcommand{\PhExitDurGHighMtn}{50.0}
\newcommand{\PhExitDurGHighMtnBand}{[48.8, 51.2]}
\newcommand{\PhExitDurGHills}{52.0}
\newcommand{\PhExitDurGHillsBand}{[50.2, 53.8]}
\newcommand{\PhExitDurGLowMtn}{51.0}
\newcommand{\PhExitDurGLowMtnBand}{[49.8, 52.2]}
\newcommand{\PhExitDurGPlains}{57.0}
\newcommand{\PhExitDurGPlainsBand}{[55.5, 58.5]}
\newcommand{\PhExitDurGSAll}{$+$2.00}
\newcommand{\PhExitDurGSAllBand}{[$+$1.26, $+$2.74]}
\newcommand{\PhExitDurGSHighMtn}{0}
\newcommand{\PhExitDurGSHighMtnBand}{[$-$1, $+$1]}
\newcommand{\PhExitDurGSHills}{$+$2.00}
\newcommand{\PhExitDurGSHillsBand}{[$+$0.15, $+$3.85]}
\newcommand{\PhExitDurGSLowMtn}{$+$1.00}
\newcommand{\PhExitDurGSLowMtnBand}{[$-$0.32, $+$2.32]}
\newcommand{\PhExitDurGSPlains}{$+$2.00}
\newcommand{\PhExitDurGSPlainsBand}{[$+$0.44, $+$3.56]}
\newcommand{\PhExitDurSAll}{50.0}
\newcommand{\PhExitDurSAllBand}{[49.7, 50.3]}
\newcommand{\PhExitDurSHighMtn}{50.0}
\newcommand{\PhExitDurSHighMtnBand}{[49.2, 50.8]}
\newcommand{\PhExitDurSHills}{50.0}
\newcommand{\PhExitDurSHillsBand}{[49.5, 50.5]}
\newcommand{\PhExitDurSLowMtn}{50.0}
\newcommand{\PhExitDurSLowMtnBand}{[49.3, 50.7]}
\newcommand{\PhExitDurSPlains}{55.0}
\newcommand{\PhExitDurSPlainsBand}{[54.2, 55.8]}
\newcommand{\PhExitMeanAAll}{68.5}
\newcommand{\PhExitMeanAAllBand}{[67.4, 69.5]}
\newcommand{\PhExitMeanAHighMtn}{68.9}
\newcommand{\PhExitMeanAHighMtnBand}{[67.0, 70.8]}
\newcommand{\PhExitMeanAHills}{61.8}
\newcommand{\PhExitMeanAHillsBand}{[59.4, 64.2]}
\newcommand{\PhExitMeanALowMtn}{70.5}
\newcommand{\PhExitMeanALowMtnBand}{[68.8, 72.1]}
\newcommand{\PhExitMeanAPlains}{68.7}
\newcommand{\PhExitMeanAPlainsBand}{[66.6, 70.8]}
\newcommand{\PhExitMeanASAll}{$+$0.729}
\newcommand{\PhExitMeanASAllBand}{[$-$0.355, $+$1.813]}
\newcommand{\PhExitMeanASHighMtn}{$-$0.212}
\newcommand{\PhExitMeanASHighMtnBand}{[$-$2.331, $+$1.907]}
\newcommand{\PhExitMeanASHills}{$-$3.04}
\newcommand{\PhExitMeanASHillsBand}{[$-$5.58, $-$0.51]}
\newcommand{\PhExitMeanASLowMtn}{$+$1.89}
\newcommand{\PhExitMeanASLowMtnBand}{[$+$0.18, $+$3.59]}
\newcommand{\PhExitMeanASPlains}{$-$1.99}
\newcommand{\PhExitMeanASPlainsBand}{[$-$4.25, $+$0.26]}
\newcommand{\PhExitMeanGAAll}{$-$0.522}
\newcommand{\PhExitMeanGAAllBand}{[$-$1.696, $+$0.653]}
\newcommand{\PhExitMeanGAHighMtn}{$-$1.72}
\newcommand{\PhExitMeanGAHighMtnBand}{[$-$4.25, $+$0.80]}
\newcommand{\PhExitMeanGAHills}{$+$2.70}
\newcommand{\PhExitMeanGAHillsBand}{[$-$0.05, $+$5.45]}
\newcommand{\PhExitMeanGALowMtn}{$-$1.77}
\newcommand{\PhExitMeanGALowMtnBand}{[$-$3.80, $+$0.27]}
\newcommand{\PhExitMeanGAPlains}{$+$2.07}
\newcommand{\PhExitMeanGAPlainsBand}{[$-$0.50, $+$4.63]}
\newcommand{\PhExitMeanGAll}{67.9}
\newcommand{\PhExitMeanGAllBand}{[67.0, 68.9]}
\newcommand{\PhExitMeanGHighMtn}{67.1}
\newcommand{\PhExitMeanGHighMtnBand}{[65.2, 69.1]}
\newcommand{\PhExitMeanGHills}{64.5}
\newcommand{\PhExitMeanGHillsBand}{[62.0, 67.0]}
\newcommand{\PhExitMeanGLowMtn}{68.7}
\newcommand{\PhExitMeanGLowMtnBand}{[67.2, 70.3]}
\newcommand{\PhExitMeanGPlains}{70.8}
\newcommand{\PhExitMeanGPlainsBand}{[69.0, 72.6]}
\newcommand{\PhExitMeanGSAll}{$+$0.208}
\newcommand{\PhExitMeanGSAllBand}{[$-$0.784, $+$1.199]}
\newcommand{\PhExitMeanGSHighMtn}{$-$1.94}
\newcommand{\PhExitMeanGSHighMtnBand}{[$-$4.13, $+$0.25]}
\newcommand{\PhExitMeanGSHills}{$-$0.347}
\newcommand{\PhExitMeanGSHillsBand}{[$-$2.951, $+$2.257]}
\newcommand{\PhExitMeanGSLowMtn}{$+$0.121}
\newcommand{\PhExitMeanGSLowMtnBand}{[$-$1.519, $+$1.762]}
\newcommand{\PhExitMeanGSPlains}{$+$0.0724}
\newcommand{\PhExitMeanGSPlainsBand}{[$-$1.8207, $+$1.9655]}
\newcommand{\PhExitMeanSAll}{67.7}
\newcommand{\PhExitMeanSAllBand}{[67.2, 68.3]}
\newcommand{\PhExitMeanSHighMtn}{69.1}
\newcommand{\PhExitMeanSHighMtnBand}{[68.1, 70.0]}
\newcommand{\PhExitMeanSHills}{64.8}
\newcommand{\PhExitMeanSHillsBand}{[63.7, 65.9]}
\newcommand{\PhExitMeanSLowMtn}{68.6}
\newcommand{\PhExitMeanSLowMtnBand}{[67.8, 69.4]}
\newcommand{\PhExitMeanSPlains}{70.7}
\newcommand{\PhExitMeanSPlainsBand}{[69.6, 71.9]}
\newcommand{\PhExitPerHourAAll}{1.54}
\newcommand{\PhExitPerHourAAllBand}{[1.52, 1.57]}
\newcommand{\PhExitPerHourAHighMtn}{1.38}
\newcommand{\PhExitPerHourAHighMtnBand}{[1.34, 1.42]}
\newcommand{\PhExitPerHourAHills}{1.83}
\newcommand{\PhExitPerHourAHillsBand}{[1.74, 1.92]}
\newcommand{\PhExitPerHourALowMtn}{1.42}
\newcommand{\PhExitPerHourALowMtnBand}{[1.39, 1.46]}
\newcommand{\PhExitPerHourAPlains}{2.05}
\newcommand{\PhExitPerHourAPlainsBand}{[1.97, 2.13]}
\newcommand{\PhExitPerHourASAll}{$-$0.204}
\newcommand{\PhExitPerHourASAllBand}{[$-$0.232, $-$0.177]}
\newcommand{\PhExitPerHourASHighMtn}{$-$0.0865}
\newcommand{\PhExitPerHourASHighMtnBand}{[$-$0.1301, $-$0.0429]}
\newcommand{\PhExitPerHourASHills}{$-$0.0869}
\newcommand{\PhExitPerHourASHillsBand}{[$-$0.1734, $-$0.0004]}
\newcommand{\PhExitPerHourASLowMtn}{$-$0.151}
\newcommand{\PhExitPerHourASLowMtnBand}{[$-$0.185, $-$0.116]}
\newcommand{\PhExitPerHourASPlains}{$-$0.152}
\newcommand{\PhExitPerHourASPlainsBand}{[$-$0.231, $-$0.073]}
\newcommand{\PhExitPerHourGAAll}{$+$0.163}
\newcommand{\PhExitPerHourGAAllBand}{[$+$0.135, $+$0.192]}
\newcommand{\PhExitPerHourGAHighMtn}{$+$0.117}
\newcommand{\PhExitPerHourGAHighMtnBand}{[$+$0.067, $+$0.167]}
\newcommand{\PhExitPerHourGAHills}{$+$0.210}
\newcommand{\PhExitPerHourGAHillsBand}{[$+$0.113, $+$0.307]}
\newcommand{\PhExitPerHourGALowMtn}{$+$0.127}
\newcommand{\PhExitPerHourGALowMtnBand}{[$+$0.090, $+$0.165]}
\newcommand{\PhExitPerHourGAPlains}{$+$0.174}
\newcommand{\PhExitPerHourGAPlainsBand}{[$+$0.094, $+$0.255]}
\newcommand{\PhExitPerHourGAll}{1.71}
\newcommand{\PhExitPerHourGAllBand}{[1.67, 1.74]}
\newcommand{\PhExitPerHourGHighMtn}{1.50}
\newcommand{\PhExitPerHourGHighMtnBand}{[1.45, 1.55]}
\newcommand{\PhExitPerHourGHills}{2.04}
\newcommand{\PhExitPerHourGHillsBand}{[1.95, 2.13]}
\newcommand{\PhExitPerHourGLowMtn}{1.55}
\newcommand{\PhExitPerHourGLowMtnBand}{[1.51, 1.59]}
\newcommand{\PhExitPerHourGPlains}{2.23}
\newcommand{\PhExitPerHourGPlainsBand}{[2.14, 2.31]}
\newcommand{\PhExitPerHourGSAll}{$-$0.0407}
\newcommand{\PhExitPerHourGSAllBand}{[$-$0.0743, $-$0.0072]}
\newcommand{\PhExitPerHourGSHighMtn}{$+$0.0307}
\newcommand{\PhExitPerHourGSHighMtnBand}{[$-$0.0194, $+$0.0809]}
\newcommand{\PhExitPerHourGSHills}{$+$0.123}
\newcommand{\PhExitPerHourGSHillsBand}{[$+$0.032, $+$0.215]}
\newcommand{\PhExitPerHourGSLowMtn}{$-$0.0233}
\newcommand{\PhExitPerHourGSLowMtnBand}{[$-$0.0642, $+$0.0176]}
\newcommand{\PhExitPerHourGSPlains}{$+$0.0223}
\newcommand{\PhExitPerHourGSPlainsBand}{[$-$0.0599, $+$0.1044]}
\newcommand{\PhExitPerHourSAll}{1.75}
\newcommand{\PhExitPerHourSAllBand}{[1.73, 1.77]}
\newcommand{\PhExitPerHourSHighMtn}{1.47}
\newcommand{\PhExitPerHourSHighMtnBand}{[1.44, 1.50]}
\newcommand{\PhExitPerHourSHills}{1.92}
\newcommand{\PhExitPerHourSHillsBand}{[1.88, 1.96]}
\newcommand{\PhExitPerHourSLowMtn}{1.57}
\newcommand{\PhExitPerHourSLowMtnBand}{[1.55, 1.60]}
\newcommand{\PhExitPerHourSPlains}{2.20}
\newcommand{\PhExitPerHourSPlainsBand}{[2.16, 2.25]}
\newcommand{\PhFlightsAAll}{10\,916}
\newcommand{\PhFlightsAHighMtn}{2\,786}
\newcommand{\PhFlightsAHills}{1\,113}
\newcommand{\PhFlightsALowMtn}{5\,241}
\newcommand{\PhFlightsAPlains}{1\,525}
\newcommand{\PhFlightsGAll}{13\,133}
\newcommand{\PhFlightsGHighMtn}{3\,097}
\newcommand{\PhFlightsGHills}{1\,543}
\newcommand{\PhFlightsGLowMtn}{6\,120}
\newcommand{\PhFlightsGPlains}{2\,097}
\newcommand{\PhFlightsSAll}{49\,781}
\newcommand{\PhFlightsSHighMtn}{7\,952}
\newcommand{\PhFlightsSHills}{9\,731}
\newcommand{\PhFlightsSLowMtn}{18\,796}
\newcommand{\PhFlightsSPlains}{9\,634}
\newcommand{\PhFoundAAll}{62.6}
\newcommand{\PhFoundAAllBand}{[62.0, 63.1]}
\newcommand{\PhFoundAHighMtn}{61.8}
\newcommand{\PhFoundAHighMtnBand}{[60.9, 62.8]}
\newcommand{\PhFoundAHills}{68.7}
\newcommand{\PhFoundAHillsBand}{[67.3, 70.1]}
\newcommand{\PhFoundALowMtn}{59.7}
\newcommand{\PhFoundALowMtnBand}{[58.9, 60.4]}
\newcommand{\PhFoundAPlains}{70.1}
\newcommand{\PhFoundAPlainsBand}{[68.8, 71.5]}
\newcommand{\PhFoundASAll}{$+$0.991}
\newcommand{\PhFoundASAllBand}{[$+$0.447, $+$1.535]}
\newcommand{\PhFoundASHighMtn}{$+$1.09}
\newcommand{\PhFoundASHighMtnBand}{[$+$0.07, $+$2.11]}
\newcommand{\PhFoundASHills}{$+$5.37}
\newcommand{\PhFoundASHillsBand}{[$+$3.88, $+$6.86]}
\newcommand{\PhFoundASLowMtn}{$+$1.36}
\newcommand{\PhFoundASLowMtnBand}{[$+$0.54, $+$2.19]}
\newcommand{\PhFoundASPlains}{$+$2.41}
\newcommand{\PhFoundASPlainsBand}{[$+$1.00, $+$3.82]}
\newcommand{\PhFoundGAAll}{$+$5.80}
\newcommand{\PhFoundGAAllBand}{[$+$5.19, $+$6.40]}
\newcommand{\PhFoundGAHighMtn}{$+$6.83}
\newcommand{\PhFoundGAHighMtnBand}{[$+$5.75, $+$7.92]}
\newcommand{\PhFoundGAHills}{$+$3.81}
\newcommand{\PhFoundGAHillsBand}{[$+$2.20, $+$5.42]}
\newcommand{\PhFoundGALowMtn}{$+$6.08}
\newcommand{\PhFoundGALowMtnBand}{[$+$5.19, $+$6.96]}
\newcommand{\PhFoundGAPlains}{$+$3.69}
\newcommand{\PhFoundGAPlainsBand}{[$+$2.12, $+$5.27]}
\newcommand{\PhFoundGAll}{68.4}
\newcommand{\PhFoundGAllBand}{[67.8, 68.9]}
\newcommand{\PhFoundGHighMtn}{68.7}
\newcommand{\PhFoundGHighMtnBand}{[67.8, 69.6]}
\newcommand{\PhFoundGHills}{72.5}
\newcommand{\PhFoundGHillsBand}{[71.2, 73.8]}
\newcommand{\PhFoundGLowMtn}{65.7}
\newcommand{\PhFoundGLowMtnBand}{[65.0, 66.5]}
\newcommand{\PhFoundGPlains}{73.8}
\newcommand{\PhFoundGPlainsBand}{[72.9, 74.7]}
\newcommand{\PhFoundGSAll}{$+$6.79}
\newcommand{\PhFoundGSAllBand}{[$+$6.24, $+$7.34]}
\newcommand{\PhFoundGSHighMtn}{$+$7.92}
\newcommand{\PhFoundGSHighMtnBand}{[$+$6.97, $+$8.87]}
\newcommand{\PhFoundGSHills}{$+$9.18}
\newcommand{\PhFoundGSHillsBand}{[$+$7.77, $+$10.59]}
\newcommand{\PhFoundGSLowMtn}{$+$7.44}
\newcommand{\PhFoundGSLowMtnBand}{[$+$6.66, $+$8.22]}
\newcommand{\PhFoundGSPlains}{$+$6.10}
\newcommand{\PhFoundGSPlainsBand}{[$+$5.05, $+$7.15]}
\newcommand{\PhFoundSAll}{61.6}
\newcommand{\PhFoundSAllBand}{[61.2, 62.0]}
\newcommand{\PhFoundSHighMtn}{60.8}
\newcommand{\PhFoundSHighMtnBand}{[60.2, 61.3]}
\newcommand{\PhFoundSHills}{63.3}
\newcommand{\PhFoundSHillsBand}{[62.6, 64.1]}
\newcommand{\PhFoundSLowMtn}{58.3}
\newcommand{\PhFoundSLowMtnBand}{[57.8, 58.8]}
\newcommand{\PhFoundSPlains}{67.7}
\newcommand{\PhFoundSPlainsBand}{[67.1, 68.3]}
\newcommand{\PhGainAAll}{197}
\newcommand{\PhGainAAllBand}{[194, 200]}
\newcommand{\PhGainAHighMtn}{199}
\newcommand{\PhGainAHighMtnBand}{[195, 204]}
\newcommand{\PhGainAHills}{214}
\newcommand{\PhGainAHillsBand}{[207, 221]}
\newcommand{\PhGainALowMtn}{189}
\newcommand{\PhGainALowMtnBand}{[185, 194]}
\newcommand{\PhGainAPlains}{202}
\newcommand{\PhGainAPlainsBand}{[197, 207]}
\newcommand{\PhGainASAll}{$+$11.0}
\newcommand{\PhGainASAllBand}{[$+$8.2, $+$13.8]}
\newcommand{\PhGainASHighMtn}{$+$4.16}
\newcommand{\PhGainASHighMtnBand}{[$-$0.13, $+$8.45]}
\newcommand{\PhGainASHills}{$+$19.3}
\newcommand{\PhGainASHillsBand}{[$+$12.4, $+$26.3]}
\newcommand{\PhGainASLowMtn}{$+$9.84}
\newcommand{\PhGainASLowMtnBand}{[$+$5.83, $+$13.85]}
\newcommand{\PhGainASPlains}{$+$15.7}
\newcommand{\PhGainASPlainsBand}{[$+$10.7, $+$20.7]}
\newcommand{\PhGainGAAll}{$+$0.136}
\newcommand{\PhGainGAAllBand}{[$-$2.009, $+$2.282]}
\newcommand{\PhGainGAHighMtn}{$+$7.46}
\newcommand{\PhGainGAHighMtnBand}{[$+$3.44, $+$11.48]}
\newcommand{\PhGainGAHills}{$+$4.61}
\newcommand{\PhGainGAHillsBand}{[$-$2.09, $+$11.32]}
\newcommand{\PhGainGALowMtn}{$-$0.480}
\newcommand{\PhGainGALowMtnBand}{[$-$3.537, $+$2.578]}
\newcommand{\PhGainGAPlains}{$-$12.1}
\newcommand{\PhGainGAPlainsBand}{[$-$17.3, $-$6.9]}
\newcommand{\PhGainGAll}{197}
\newcommand{\PhGainGAllBand}{[194, 200]}
\newcommand{\PhGainGHighMtn}{207}
\newcommand{\PhGainGHighMtnBand}{[202, 211]}
\newcommand{\PhGainGHills}{219}
\newcommand{\PhGainGHillsBand}{[212, 225]}
\newcommand{\PhGainGLowMtn}{189}
\newcommand{\PhGainGLowMtnBand}{[185, 193]}
\newcommand{\PhGainGPlains}{190}
\newcommand{\PhGainGPlainsBand}{[185, 195]}
\newcommand{\PhGainGSAll}{$+$11.2}
\newcommand{\PhGainGSAllBand}{[$+$8.2, $+$14.1]}
\newcommand{\PhGainGSHighMtn}{$+$11.6}
\newcommand{\PhGainGSHighMtnBand}{[$+$7.0, $+$16.3]}
\newcommand{\PhGainGSHills}{$+$24.0}
\newcommand{\PhGainGSHillsBand}{[$+$17.0, $+$30.9]}
\newcommand{\PhGainGSLowMtn}{$+$9.36}
\newcommand{\PhGainGSLowMtnBand}{[$+$5.34, $+$13.38]}
\newcommand{\PhGainGSPlains}{$+$3.64}
\newcommand{\PhGainGSPlainsBand}{[$-$1.22, $+$8.50]}
\newcommand{\PhGainSAll}{186}
\newcommand{\PhGainSAllBand}{[184, 188]}
\newcommand{\PhGainSHighMtn}{195}
\newcommand{\PhGainSHighMtnBand}{[193, 198]}
\newcommand{\PhGainSHills}{195}
\newcommand{\PhGainSHillsBand}{[191, 198]}
\newcommand{\PhGainSLowMtn}{180}
\newcommand{\PhGainSLowMtnBand}{[177, 182]}
\newcommand{\PhGainSPlains}{186}
\newcommand{\PhGainSPlainsBand}{[183, 189]}
\newcommand{\PhGlideMeanAAll}{285}
\newcommand{\PhGlideMeanAAllBand}{[282, 288]}
\newcommand{\PhGlideMeanAHighMtn}{304}
\newcommand{\PhGlideMeanAHighMtnBand}{[298, 310]}
\newcommand{\PhGlideMeanAHills}{256}
\newcommand{\PhGlideMeanAHillsBand}{[250, 262]}
\newcommand{\PhGlideMeanALowMtn}{291}
\newcommand{\PhGlideMeanALowMtnBand}{[286, 295]}
\newcommand{\PhGlideMeanAPlains}{251}
\newcommand{\PhGlideMeanAPlainsBand}{[245, 256]}
\newcommand{\PhGlideMeanASAll}{$+$18.3}
\newcommand{\PhGlideMeanASAllBand}{[$+$15.3, $+$21.2]}
\newcommand{\PhGlideMeanASHighMtn}{$+$15.4}
\newcommand{\PhGlideMeanASHighMtnBand}{[$+$9.5, $+$21.2]}
\newcommand{\PhGlideMeanASHills}{$-$0.478}
\newcommand{\PhGlideMeanASHillsBand}{[$-$6.334, $+$5.377]}
\newcommand{\PhGlideMeanASLowMtn}{$+$15.5}
\newcommand{\PhGlideMeanASLowMtnBand}{[$+$11.2, $+$19.9]}
\newcommand{\PhGlideMeanASPlains}{$+$10.2}
\newcommand{\PhGlideMeanASPlainsBand}{[$+$4.7, $+$15.6]}
\newcommand{\PhGlideMeanGAAll}{$-$27.6}
\newcommand{\PhGlideMeanGAAllBand}{[$-$30.1, $-$25.1]}
\newcommand{\PhGlideMeanGAHighMtn}{$-$31.4}
\newcommand{\PhGlideMeanGAHighMtnBand}{[$-$36.8, $-$25.9]}
\newcommand{\PhGlideMeanGAHills}{$-$19.2}
\newcommand{\PhGlideMeanGAHillsBand}{[$-$25.3, $-$13.1]}
\newcommand{\PhGlideMeanGALowMtn}{$-$24.3}
\newcommand{\PhGlideMeanGALowMtnBand}{[$-$28.2, $-$20.4]}
\newcommand{\PhGlideMeanGAPlains}{$-$28.2}
\newcommand{\PhGlideMeanGAPlainsBand}{[$-$33.6, $-$22.8]}
\newcommand{\PhGlideMeanGAll}{257}
\newcommand{\PhGlideMeanGAllBand}{[254, 260]}
\newcommand{\PhGlideMeanGHighMtn}{272}
\newcommand{\PhGlideMeanGHighMtnBand}{[267, 278]}
\newcommand{\PhGlideMeanGHills}{237}
\newcommand{\PhGlideMeanGHillsBand}{[231, 243]}
\newcommand{\PhGlideMeanGLowMtn}{266}
\newcommand{\PhGlideMeanGLowMtnBand}{[262, 270]}
\newcommand{\PhGlideMeanGPlains}{223}
\newcommand{\PhGlideMeanGPlainsBand}{[218, 227]}
\newcommand{\PhGlideMeanGSAll}{$-$9.34}
\newcommand{\PhGlideMeanGSAllBand}{[$-$12.35, $-$6.34]}
\newcommand{\PhGlideMeanGSHighMtn}{$-$16.0}
\newcommand{\PhGlideMeanGSHighMtnBand}{[$-$21.7, $-$10.2]}
\newcommand{\PhGlideMeanGSHills}{$-$19.7}
\newcommand{\PhGlideMeanGSHillsBand}{[$-$25.8, $-$13.5]}
\newcommand{\PhGlideMeanGSLowMtn}{$-$8.83}
\newcommand{\PhGlideMeanGSLowMtnBand}{[$-$13.08, $-$4.59]}
\newcommand{\PhGlideMeanGSPlains}{$-$18.1}
\newcommand{\PhGlideMeanGSPlainsBand}{[$-$22.9, $-$13.2]}
\newcommand{\PhGlideMeanSAll}{266}
\newcommand{\PhGlideMeanSAllBand}{[265, 268]}
\newcommand{\PhGlideMeanSHighMtn}{288}
\newcommand{\PhGlideMeanSHighMtnBand}{[285, 292]}
\newcommand{\PhGlideMeanSHills}{257}
\newcommand{\PhGlideMeanSHillsBand}{[253, 260]}
\newcommand{\PhGlideMeanSLowMtn}{275}
\newcommand{\PhGlideMeanSLowMtnBand}{[273, 278]}
\newcommand{\PhGlideMeanSPlains}{241}
\newcommand{\PhGlideMeanSPlainsBand}{[237, 244]}
\newcommand{\PhGlidePerHourAAll}{7.19}
\newcommand{\PhGlidePerHourAAllBand}{[7.13, 7.25]}
\newcommand{\PhGlidePerHourAHighMtn}{7.12}
\newcommand{\PhGlidePerHourAHighMtnBand}{[7.01, 7.23]}
\newcommand{\PhGlidePerHourAHills}{7.15}
\newcommand{\PhGlidePerHourAHillsBand}{[7.04, 7.26]}
\newcommand{\PhGlidePerHourALowMtn}{7.32}
\newcommand{\PhGlidePerHourALowMtnBand}{[7.23, 7.41]}
\newcommand{\PhGlidePerHourAPlains}{6.89}
\newcommand{\PhGlidePerHourAPlainsBand}{[6.78, 7.00]}
\newcommand{\PhGlidePerHourASAll}{$-$0.290}
\newcommand{\PhGlidePerHourASAllBand}{[$-$0.346, $-$0.234]}
\newcommand{\PhGlidePerHourASHighMtn}{$-$0.206}
\newcommand{\PhGlidePerHourASHighMtnBand}{[$-$0.312, $-$0.100]}
\newcommand{\PhGlidePerHourASHills}{$-$0.283}
\newcommand{\PhGlidePerHourASHillsBand}{[$-$0.398, $-$0.167]}
\newcommand{\PhGlidePerHourASLowMtn}{$-$0.300}
\newcommand{\PhGlidePerHourASLowMtnBand}{[$-$0.389, $-$0.210]}
\newcommand{\PhGlidePerHourASPlains}{$-$0.414}
\newcommand{\PhGlidePerHourASPlainsBand}{[$-$0.526, $-$0.302]}
\newcommand{\PhGlidePerHourGAAll}{$+$0.165}
\newcommand{\PhGlidePerHourGAAllBand}{[$+$0.116, $+$0.214]}
\newcommand{\PhGlidePerHourGAHighMtn}{$+$0.192}
\newcommand{\PhGlidePerHourGAHighMtnBand}{[$+$0.102, $+$0.282]}
\newcommand{\PhGlidePerHourGAHills}{$-$0.0383}
\newcommand{\PhGlidePerHourGAHillsBand}{[$-$0.1692, $+$0.0925]}
\newcommand{\PhGlidePerHourGALowMtn}{$+$0.160}
\newcommand{\PhGlidePerHourGALowMtnBand}{[$+$0.087, $+$0.233]}
\newcommand{\PhGlidePerHourGAPlains}{$+$0.288}
\newcommand{\PhGlidePerHourGAPlainsBand}{[$+$0.149, $+$0.427]}
\newcommand{\PhGlidePerHourGAll}{7.36}
\newcommand{\PhGlidePerHourGAllBand}{[7.30, 7.42]}
\newcommand{\PhGlidePerHourGHighMtn}{7.31}
\newcommand{\PhGlidePerHourGHighMtnBand}{[7.19, 7.42]}
\newcommand{\PhGlidePerHourGHills}{7.11}
\newcommand{\PhGlidePerHourGHillsBand}{[7.00, 7.23]}
\newcommand{\PhGlidePerHourGLowMtn}{7.48}
\newcommand{\PhGlidePerHourGLowMtnBand}{[7.39, 7.56]}
\newcommand{\PhGlidePerHourGPlains}{7.18}
\newcommand{\PhGlidePerHourGPlainsBand}{[7.06, 7.29]}
\newcommand{\PhGlidePerHourGSAll}{$-$0.125}
\newcommand{\PhGlidePerHourGSAllBand}{[$-$0.183, $-$0.067]}
\newcommand{\PhGlidePerHourGSHighMtn}{$-$0.0138}
\newcommand{\PhGlidePerHourGSHighMtnBand}{[$-$0.1257, $+$0.0981]}
\newcommand{\PhGlidePerHourGSHills}{$-$0.321}
\newcommand{\PhGlidePerHourGSHillsBand}{[$-$0.445, $-$0.197]}
\newcommand{\PhGlidePerHourGSLowMtn}{$-$0.139}
\newcommand{\PhGlidePerHourGSLowMtnBand}{[$-$0.223, $-$0.055]}
\newcommand{\PhGlidePerHourGSPlains}{$-$0.126}
\newcommand{\PhGlidePerHourGSPlainsBand}{[$-$0.250, $-$0.002]}
\newcommand{\PhGlidePerHourSAll}{7.48}
\newcommand{\PhGlidePerHourSAllBand}{[7.45, 7.52]}
\newcommand{\PhGlidePerHourSHighMtn}{7.32}
\newcommand{\PhGlidePerHourSHighMtnBand}{[7.26, 7.39]}
\newcommand{\PhGlidePerHourSHills}{7.43}
\newcommand{\PhGlidePerHourSHillsBand}{[7.37, 7.50]}
\newcommand{\PhGlidePerHourSLowMtn}{7.62}
\newcommand{\PhGlidePerHourSLowMtnBand}{[7.57, 7.67]}
\newcommand{\PhGlidePerHourSPlains}{7.30}
\newcommand{\PhGlidePerHourSPlainsBand}{[7.23, 7.37]}
\newcommand{\PhHoursAAll}{24\,741}
\newcommand{\PhHoursAHighMtn}{6\,968}
\newcommand{\PhHoursAHills}{1\,788}
\newcommand{\PhHoursALowMtn}{12\,525}
\newcommand{\PhHoursAPlains}{2\,823}
\newcommand{\PhHoursGAll}{20\,210}
\newcommand{\PhHoursGHighMtn}{4\,734}
\newcommand{\PhHoursGHills}{2\,308}
\newcommand{\PhHoursGLowMtn}{9\,383}
\newcommand{\PhHoursGPlains}{3\,412}
\newcommand{\PhHoursSAll}{107\,423}
\newcommand{\PhHoursSHighMtn}{22\,031}
\newcommand{\PhHoursSHills}{16\,852}
\newcommand{\PhHoursSLowMtn}{43\,005}
\newcommand{\PhHoursSPlains}{17\,179}
\newcommand{\PhKindExitRunsAAll}{26.7}
\newcommand{\PhKindExitRunsAAllBand}{[26.3, 27.1]}
\newcommand{\PhKindExitRunsAHighMtn}{24.8}
\newcommand{\PhKindExitRunsAHighMtnBand}{[24.3, 25.4]}
\newcommand{\PhKindExitRunsAHills}{32.2}
\newcommand{\PhKindExitRunsAHillsBand}{[30.9, 33.5]}
\newcommand{\PhKindExitRunsALowMtn}{24.4}
\newcommand{\PhKindExitRunsALowMtnBand}{[24.0, 24.8]}
\newcommand{\PhKindExitRunsAPlains}{33.3}
\newcommand{\PhKindExitRunsAPlainsBand}{[32.3, 34.2]}
\newcommand{\PhKindExitRunsASAll}{$-$1.34}
\newcommand{\PhKindExitRunsASAllBand}{[$-$1.70, $-$0.97]}
\newcommand{\PhKindExitRunsASHighMtn}{$+$0.40}
\newcommand{\PhKindExitRunsASHighMtnBand}{[$-$0.18, $+$0.99]}
\newcommand{\PhKindExitRunsASHills}{$+$1.22}
\newcommand{\PhKindExitRunsASHillsBand}{[$-$0.09, $+$2.52]}
\newcommand{\PhKindExitRunsASLowMtn}{$-$0.89}
\newcommand{\PhKindExitRunsASLowMtnBand}{[$-$1.35, $-$0.43]}
\newcommand{\PhKindExitRunsASPlains}{$+$0.58}
\newcommand{\PhKindExitRunsASPlainsBand}{[$-$0.30, $+$1.46]}
\newcommand{\PhKindExitRunsGAAll}{$+$1.51}
\newcommand{\PhKindExitRunsGAAllBand}{[$+$1.14, $+$1.87]}
\newcommand{\PhKindExitRunsGAHighMtn}{$+$1.10}
\newcommand{\PhKindExitRunsGAHighMtnBand}{[$+$0.40, $+$1.80]}
\newcommand{\PhKindExitRunsGAHills}{$+$1.61}
\newcommand{\PhKindExitRunsGAHillsBand}{[$+$0.32, $+$2.91]}
\newcommand{\PhKindExitRunsGALowMtn}{$+$1.13}
\newcommand{\PhKindExitRunsGALowMtnBand}{[$+$0.64, $+$1.63]}
\newcommand{\PhKindExitRunsGAPlains}{$+$1.28}
\newcommand{\PhKindExitRunsGAPlainsBand}{[$+$0.34, $+$2.22]}
\newcommand{\PhKindExitRunsGAll}{28.2}
\newcommand{\PhKindExitRunsGAllBand}{[27.8, 28.6]}
\newcommand{\PhKindExitRunsGHighMtn}{25.9}
\newcommand{\PhKindExitRunsGHighMtnBand}{[25.4, 26.5]}
\newcommand{\PhKindExitRunsGHills}{33.8}
\newcommand{\PhKindExitRunsGHillsBand}{[32.8, 34.8]}
\newcommand{\PhKindExitRunsGLowMtn}{25.5}
\newcommand{\PhKindExitRunsGLowMtnBand}{[25.1, 26.0]}
\newcommand{\PhKindExitRunsGPlains}{34.6}
\newcommand{\PhKindExitRunsGPlainsBand}{[33.9, 35.2]}
\newcommand{\PhKindExitRunsGSAll}{$+$0.17}
\newcommand{\PhKindExitRunsGSAllBand}{[$-$0.23, $+$0.57]}
\newcommand{\PhKindExitRunsGSHighMtn}{$+$1.50}
\newcommand{\PhKindExitRunsGSHighMtnBand}{[$+$0.91, $+$2.09]}
\newcommand{\PhKindExitRunsGSHills}{$+$2.83}
\newcommand{\PhKindExitRunsGSHillsBand}{[$+$1.82, $+$3.84]}
\newcommand{\PhKindExitRunsGSLowMtn}{$+$0.24}
\newcommand{\PhKindExitRunsGSLowMtnBand}{[$-$0.25, $+$0.74]}
\newcommand{\PhKindExitRunsGSPlains}{$+$1.86}
\newcommand{\PhKindExitRunsGSPlainsBand}{[$+$1.16, $+$2.56]}
\newcommand{\PhKindExitRunsSAll}{28.0}
\newcommand{\PhKindExitRunsSAllBand}{[27.8, 28.3]}
\newcommand{\PhKindExitRunsSHighMtn}{24.4}
\newcommand{\PhKindExitRunsSHighMtnBand}{[24.1, 24.7]}
\newcommand{\PhKindExitRunsSHills}{31.0}
\newcommand{\PhKindExitRunsSHillsBand}{[30.5, 31.5]}
\newcommand{\PhKindExitRunsSLowMtn}{25.3}
\newcommand{\PhKindExitRunsSLowMtnBand}{[25.0, 25.6]}
\newcommand{\PhKindExitRunsSPlains}{32.7}
\newcommand{\PhKindExitRunsSPlainsBand}{[32.2, 33.2]}
\newcommand{\PhKindExitTimeAAll}{23.3}
\newcommand{\PhKindExitTimeAAllBand}{[22.8, 23.7]}
\newcommand{\PhKindExitTimeAHighMtn}{21.2}
\newcommand{\PhKindExitTimeAHighMtnBand}{[20.5, 21.9]}
\newcommand{\PhKindExitTimeAHills}{30.3}
\newcommand{\PhKindExitTimeAHillsBand}{[28.9, 31.8]}
\newcommand{\PhKindExitTimeALowMtn}{20.6}
\newcommand{\PhKindExitTimeALowMtnBand}{[20.1, 21.2]}
\newcommand{\PhKindExitTimeAPlains}{32.5}
\newcommand{\PhKindExitTimeAPlainsBand}{[31.3, 33.8]}
\newcommand{\PhKindExitTimeASAll}{$-$2.06}
\newcommand{\PhKindExitTimeASAllBand}{[$-$2.53, $-$1.59]}
\newcommand{\PhKindExitTimeASHighMtn}{$+$0.18}
\newcommand{\PhKindExitTimeASHighMtnBand}{[$-$0.57, $+$0.94]}
\newcommand{\PhKindExitTimeASHills}{$+$1.70}
\newcommand{\PhKindExitTimeASHillsBand}{[$+$0.24, $+$3.16]}
\newcommand{\PhKindExitTimeASLowMtn}{$-$1.28}
\newcommand{\PhKindExitTimeASLowMtnBand}{[$-$1.88, $-$0.69]}
\newcommand{\PhKindExitTimeASPlains}{$-$0.05}
\newcommand{\PhKindExitTimeASPlainsBand}{[$-$1.30, $+$1.21]}
\newcommand{\PhKindExitTimeGAAll}{$+$2.50}
\newcommand{\PhKindExitTimeGAAllBand}{[$+$1.98, $+$3.02]}
\newcommand{\PhKindExitTimeGAHighMtn}{$+$1.79}
\newcommand{\PhKindExitTimeGAHighMtnBand}{[$+$0.80, $+$2.78]}
\newcommand{\PhKindExitTimeGAHills}{$+$2.94}
\newcommand{\PhKindExitTimeGAHillsBand}{[$+$1.34, $+$4.54]}
\newcommand{\PhKindExitTimeGALowMtn}{$+$1.67}
\newcommand{\PhKindExitTimeGALowMtnBand}{[$+$0.99, $+$2.35]}
\newcommand{\PhKindExitTimeGAPlains}{$+$3.23}
\newcommand{\PhKindExitTimeGAPlainsBand}{[$+$1.87, $+$4.59]}
\newcommand{\PhKindExitTimeGAll}{25.8}
\newcommand{\PhKindExitTimeGAllBand}{[25.2, 26.3]}
\newcommand{\PhKindExitTimeGHighMtn}{23.0}
\newcommand{\PhKindExitTimeGHighMtnBand}{[22.2, 23.8]}
\newcommand{\PhKindExitTimeGHills}{33.3}
\newcommand{\PhKindExitTimeGHillsBand}{[32.0, 34.6]}
\newcommand{\PhKindExitTimeGLowMtn}{22.3}
\newcommand{\PhKindExitTimeGLowMtnBand}{[21.7, 22.9]}
\newcommand{\PhKindExitTimeGPlains}{35.8}
\newcommand{\PhKindExitTimeGPlainsBand}{[34.8, 36.8]}
\newcommand{\PhKindExitTimeGSAll}{$+$0.44}
\newcommand{\PhKindExitTimeGSAllBand}{[$-$0.07, $+$0.94]}
\newcommand{\PhKindExitTimeGSHighMtn}{$+$1.97}
\newcommand{\PhKindExitTimeGSHighMtnBand}{[$+$1.14, $+$2.81]}
\newcommand{\PhKindExitTimeGSHills}{$+$4.64}
\newcommand{\PhKindExitTimeGSHillsBand}{[$+$3.31, $+$5.97]}
\newcommand{\PhKindExitTimeGSLowMtn}{$+$0.39}
\newcommand{\PhKindExitTimeGSLowMtnBand}{[$-$0.20, $+$0.98]}
\newcommand{\PhKindExitTimeGSPlains}{$+$3.18}
\newcommand{\PhKindExitTimeGSPlainsBand}{[$+$2.18, $+$4.18]}
\newcommand{\PhKindExitTimeSAll}{25.3}
\newcommand{\PhKindExitTimeSAllBand}{[25.0, 25.7]}
\newcommand{\PhKindExitTimeSHighMtn}{21.0}
\newcommand{\PhKindExitTimeSHighMtnBand}{[20.7, 21.4]}
\newcommand{\PhKindExitTimeSHills}{28.6}
\newcommand{\PhKindExitTimeSHillsBand}{[28.0, 29.2]}
\newcommand{\PhKindExitTimeSLowMtn}{21.9}
\newcommand{\PhKindExitTimeSLowMtnBand}{[21.5, 22.3]}
\newcommand{\PhKindExitTimeSPlains}{32.6}
\newcommand{\PhKindExitTimeSPlainsBand}{[32.0, 33.2]}
\newcommand{\PhKindFailedRunsAAll}{27.4}
\newcommand{\PhKindFailedRunsAAllBand}{[26.9, 27.9]}
\newcommand{\PhKindFailedRunsAHighMtn}{28.7}
\newcommand{\PhKindFailedRunsAHighMtnBand}{[27.9, 29.5]}
\newcommand{\PhKindFailedRunsAHills}{21.2}
\newcommand{\PhKindFailedRunsAHillsBand}{[20.1, 22.3]}
\newcommand{\PhKindFailedRunsALowMtn}{30.5}
\newcommand{\PhKindFailedRunsALowMtnBand}{[29.8, 31.2]}
\newcommand{\PhKindFailedRunsAPlains}{19.9}
\newcommand{\PhKindFailedRunsAPlainsBand}{[19.0, 20.9]}
\newcommand{\PhKindFailedRunsASAll}{$-$0.21}
\newcommand{\PhKindFailedRunsASAllBand}{[$-$0.68, $+$0.25]}
\newcommand{\PhKindFailedRunsASHighMtn}{$-$0.98}
\newcommand{\PhKindFailedRunsASHighMtnBand}{[$-$1.83, $-$0.12]}
\newcommand{\PhKindFailedRunsASHills}{$-$4.09}
\newcommand{\PhKindFailedRunsASHillsBand}{[$-$5.19, $-$2.98]}
\newcommand{\PhKindFailedRunsASLowMtn}{$-$0.66}
\newcommand{\PhKindFailedRunsASLowMtnBand}{[$-$1.35, $+$0.04]}
\newcommand{\PhKindFailedRunsASPlains}{$-$1.79}
\newcommand{\PhKindFailedRunsASPlainsBand}{[$-$2.81, $-$0.78]}
\newcommand{\PhKindFailedRunsGAAll}{$-$4.73}
\newcommand{\PhKindFailedRunsGAAllBand}{[$-$5.21, $-$4.24]}
\newcommand{\PhKindFailedRunsGAHighMtn}{$-$5.48}
\newcommand{\PhKindFailedRunsGAHighMtnBand}{[$-$6.35, $-$4.60]}
\newcommand{\PhKindFailedRunsGAHills}{$-$3.03}
\newcommand{\PhKindFailedRunsGAHillsBand}{[$-$4.19, $-$1.86]}
\newcommand{\PhKindFailedRunsGALowMtn}{$-$4.98}
\newcommand{\PhKindFailedRunsGALowMtnBand}{[$-$5.71, $-$4.26]}
\newcommand{\PhKindFailedRunsGAPlains}{$-$2.80}
\newcommand{\PhKindFailedRunsGAPlainsBand}{[$-$3.92, $-$1.68]}
\newcommand{\PhKindFailedRunsGAll}{22.7}
\newcommand{\PhKindFailedRunsGAllBand}{[22.3, 23.1]}
\newcommand{\PhKindFailedRunsGHighMtn}{23.2}
\newcommand{\PhKindFailedRunsGHighMtnBand}{[22.5, 23.9]}
\newcommand{\PhKindFailedRunsGHills}{18.2}
\newcommand{\PhKindFailedRunsGHillsBand}{[17.3, 19.1]}
\newcommand{\PhKindFailedRunsGLowMtn}{25.5}
\newcommand{\PhKindFailedRunsGLowMtnBand}{[24.9, 26.1]}
\newcommand{\PhKindFailedRunsGPlains}{17.1}
\newcommand{\PhKindFailedRunsGPlainsBand}{[16.5, 17.8]}
\newcommand{\PhKindFailedRunsGSAll}{$-$4.94}
\newcommand{\PhKindFailedRunsGSAllBand}{[$-$5.40, $-$4.48]}
\newcommand{\PhKindFailedRunsGSHighMtn}{$-$6.45}
\newcommand{\PhKindFailedRunsGSHighMtnBand}{[$-$7.21, $-$5.69]}
\newcommand{\PhKindFailedRunsGSHills}{$-$7.11}
\newcommand{\PhKindFailedRunsGSHillsBand}{[$-$8.15, $-$6.08]}
\newcommand{\PhKindFailedRunsGSLowMtn}{$-$5.64}
\newcommand{\PhKindFailedRunsGSLowMtnBand}{[$-$6.29, $-$4.99]}
\newcommand{\PhKindFailedRunsGSPlains}{$-$4.59}
\newcommand{\PhKindFailedRunsGSPlainsBand}{[$-$5.35, $-$3.83]}
\newcommand{\PhKindFailedRunsSAll}{27.6}
\newcommand{\PhKindFailedRunsSAllBand}{[27.3, 28.0]}
\newcommand{\PhKindFailedRunsSHighMtn}{29.7}
\newcommand{\PhKindFailedRunsSHighMtnBand}{[29.2, 30.1]}
\newcommand{\PhKindFailedRunsSHills}{25.3}
\newcommand{\PhKindFailedRunsSHillsBand}{[24.7, 25.9]}
\newcommand{\PhKindFailedRunsSLowMtn}{31.1}
\newcommand{\PhKindFailedRunsSLowMtnBand}{[30.7, 31.6]}
\newcommand{\PhKindFailedRunsSPlains}{21.7}
\newcommand{\PhKindFailedRunsSPlainsBand}{[21.2, 22.2]}
\newcommand{\PhKindFailedTimeAAll}{31.7}
\newcommand{\PhKindFailedTimeAAllBand}{[31.0, 32.3]}
\newcommand{\PhKindFailedTimeAHighMtn}{33.4}
\newcommand{\PhKindFailedTimeAHighMtnBand}{[32.2, 34.5]}
\newcommand{\PhKindFailedTimeAHills}{24.9}
\newcommand{\PhKindFailedTimeAHillsBand}{[23.3, 26.5]}
\newcommand{\PhKindFailedTimeALowMtn}{34.4}
\newcommand{\PhKindFailedTimeALowMtnBand}{[33.5, 35.3]}
\newcommand{\PhKindFailedTimeAPlains}{22.1}
\newcommand{\PhKindFailedTimeAPlainsBand}{[20.6, 23.6]}
\newcommand{\PhKindFailedTimeASAll}{$-$0.23}
\newcommand{\PhKindFailedTimeASAllBand}{[$-$0.89, $+$0.42]}
\newcommand{\PhKindFailedTimeASHighMtn}{$-$0.35}
\newcommand{\PhKindFailedTimeASHighMtnBand}{[$-$1.58, $+$0.88]}
\newcommand{\PhKindFailedTimeASHills}{$-$4.83}
\newcommand{\PhKindFailedTimeASHillsBand}{[$-$6.47, $-$3.18]}
\newcommand{\PhKindFailedTimeASLowMtn}{$-$1.26}
\newcommand{\PhKindFailedTimeASLowMtnBand}{[$-$2.21, $-$0.31]}
\newcommand{\PhKindFailedTimeASPlains}{$-$2.06}
\newcommand{\PhKindFailedTimeASPlainsBand}{[$-$3.62, $-$0.49]}
\newcommand{\PhKindFailedTimeGAAll}{$-$6.48}
\newcommand{\PhKindFailedTimeGAAllBand}{[$-$7.20, $-$5.76]}
\newcommand{\PhKindFailedTimeGAHighMtn}{$-$7.81}
\newcommand{\PhKindFailedTimeGAHighMtnBand}{[$-$9.19, $-$6.43]}
\newcommand{\PhKindFailedTimeGAHills}{$-$4.46}
\newcommand{\PhKindFailedTimeGAHillsBand}{[$-$6.22, $-$2.69]}
\newcommand{\PhKindFailedTimeGALowMtn}{$-$6.27}
\newcommand{\PhKindFailedTimeGALowMtnBand}{[$-$7.38, $-$5.17]}
\newcommand{\PhKindFailedTimeGAPlains}{$-$4.26}
\newcommand{\PhKindFailedTimeGAPlainsBand}{[$-$6.03, $-$2.48]}
\newcommand{\PhKindFailedTimeGAll}{25.2}
\newcommand{\PhKindFailedTimeGAllBand}{[24.6, 25.8]}
\newcommand{\PhKindFailedTimeGHighMtn}{25.6}
\newcommand{\PhKindFailedTimeGHighMtnBand}{[24.6, 26.6]}
\newcommand{\PhKindFailedTimeGHills}{20.5}
\newcommand{\PhKindFailedTimeGHillsBand}{[19.0, 21.9]}
\newcommand{\PhKindFailedTimeGLowMtn}{28.1}
\newcommand{\PhKindFailedTimeGLowMtnBand}{[27.2, 29.0]}
\newcommand{\PhKindFailedTimeGPlains}{17.9}
\newcommand{\PhKindFailedTimeGPlainsBand}{[17.0, 18.7]}
\newcommand{\PhKindFailedTimeGSAll}{$-$6.71}
\newcommand{\PhKindFailedTimeGSAllBand}{[$-$7.35, $-$6.08]}
\newcommand{\PhKindFailedTimeGSHighMtn}{$-$8.15}
\newcommand{\PhKindFailedTimeGSHighMtnBand}{[$-$9.26, $-$7.05]}
\newcommand{\PhKindFailedTimeGSHills}{$-$9.28}
\newcommand{\PhKindFailedTimeGSHillsBand}{[$-$10.86, $-$7.71]}
\newcommand{\PhKindFailedTimeGSLowMtn}{$-$7.53}
\newcommand{\PhKindFailedTimeGSLowMtnBand}{[$-$8.48, $-$6.59]}
\newcommand{\PhKindFailedTimeGSPlains}{$-$6.31}
\newcommand{\PhKindFailedTimeGSPlainsBand}{[$-$7.29, $-$5.34]}
\newcommand{\PhKindFailedTimeSAll}{31.9}
\newcommand{\PhKindFailedTimeSAllBand}{[31.5, 32.3]}
\newcommand{\PhKindFailedTimeSHighMtn}{33.7}
\newcommand{\PhKindFailedTimeSHighMtnBand}{[33.1, 34.3]}
\newcommand{\PhKindFailedTimeSHills}{29.8}
\newcommand{\PhKindFailedTimeSHillsBand}{[28.9, 30.6]}
\newcommand{\PhKindFailedTimeSLowMtn}{35.6}
\newcommand{\PhKindFailedTimeSLowMtnBand}{[35.1, 36.2]}
\newcommand{\PhKindFailedTimeSPlains}{24.2}
\newcommand{\PhKindFailedTimeSPlainsBand}{[23.6, 24.8]}
\newcommand{\PhKindFoundRunsAAll}{22.4}
\newcommand{\PhKindFoundRunsAAllBand}{[22.1, 22.6]}
\newcommand{\PhKindFoundRunsAHighMtn}{23.2}
\newcommand{\PhKindFoundRunsAHighMtnBand}{[22.7, 23.6]}
\newcommand{\PhKindFoundRunsAHills}{19.8}
\newcommand{\PhKindFoundRunsAHillsBand}{[18.9, 20.8]}
\newcommand{\PhKindFoundRunsALowMtn}{23.5}
\newcommand{\PhKindFoundRunsALowMtnBand}{[23.1, 23.9]}
\newcommand{\PhKindFoundRunsAPlains}{19.0}
\newcommand{\PhKindFoundRunsAPlainsBand}{[18.4, 19.7]}
\newcommand{\PhKindFoundRunsASAll}{$+$1.10}
\newcommand{\PhKindFoundRunsASAllBand}{[$+$0.82, $+$1.38]}
\newcommand{\PhKindFoundRunsASHighMtn}{$-$0.44}
\newcommand{\PhKindFoundRunsASHighMtnBand}{[$-$0.95, $+$0.07]}
\newcommand{\PhKindFoundRunsASHills}{$-$0.31}
\newcommand{\PhKindFoundRunsASHillsBand}{[$-$1.24, $+$0.62]}
\newcommand{\PhKindFoundRunsASLowMtn}{$+$1.36}
\newcommand{\PhKindFoundRunsASLowMtnBand}{[$+$0.95, $+$1.77]}
\newcommand{\PhKindFoundRunsASPlains}{$-$0.03}
\newcommand{\PhKindFoundRunsASPlainsBand}{[$-$0.67, $+$0.61]}
\newcommand{\PhKindFoundRunsGAAll}{$+$0.88}
\newcommand{\PhKindFoundRunsGAAllBand}{[$+$0.57, $+$1.18]}
\newcommand{\PhKindFoundRunsGAHighMtn}{$+$1.83}
\newcommand{\PhKindFoundRunsGAHighMtnBand}{[$+$1.21, $+$2.44]}
\newcommand{\PhKindFoundRunsGAHills}{$+$0.93}
\newcommand{\PhKindFoundRunsGAHillsBand}{[$-$0.01, $+$1.88]}
\newcommand{\PhKindFoundRunsGALowMtn}{$+$1.16}
\newcommand{\PhKindFoundRunsGALowMtnBand}{[$+$0.71, $+$1.62]}
\newcommand{\PhKindFoundRunsGAPlains}{$-$0.16}
\newcommand{\PhKindFoundRunsGAPlainsBand}{[$-$0.91, $+$0.59]}
\newcommand{\PhKindFoundRunsGAll}{23.3}
\newcommand{\PhKindFoundRunsGAllBand}{[22.9, 23.6]}
\newcommand{\PhKindFoundRunsGHighMtn}{25.0}
\newcommand{\PhKindFoundRunsGHighMtnBand}{[24.5, 25.5]}
\newcommand{\PhKindFoundRunsGHills}{20.8}
\newcommand{\PhKindFoundRunsGHillsBand}{[20.1, 21.5]}
\newcommand{\PhKindFoundRunsGLowMtn}{24.7}
\newcommand{\PhKindFoundRunsGLowMtnBand}{[24.3, 25.1]}
\newcommand{\PhKindFoundRunsGPlains}{18.9}
\newcommand{\PhKindFoundRunsGPlainsBand}{[18.2, 19.5]}
\newcommand{\PhKindFoundRunsGSAll}{$+$1.98}
\newcommand{\PhKindFoundRunsGSAllBand}{[$+$1.66, $+$2.29]}
\newcommand{\PhKindFoundRunsGSHighMtn}{$+$1.38}
\newcommand{\PhKindFoundRunsGSHighMtnBand}{[$+$0.83, $+$1.94]}
\newcommand{\PhKindFoundRunsGSHills}{$+$0.62}
\newcommand{\PhKindFoundRunsGSHillsBand}{[$-$0.15, $+$1.40]}
\newcommand{\PhKindFoundRunsGSLowMtn}{$+$2.52}
\newcommand{\PhKindFoundRunsGSLowMtnBand}{[$+$2.12, $+$2.92]}
\newcommand{\PhKindFoundRunsGSPlains}{$-$0.19}
\newcommand{\PhKindFoundRunsGSPlainsBand}{[$-$0.82, $+$0.44]}
\newcommand{\PhKindFoundRunsSAll}{21.3}
\newcommand{\PhKindFoundRunsSAllBand}{[21.1, 21.5]}
\newcommand{\PhKindFoundRunsSHighMtn}{23.6}
\newcommand{\PhKindFoundRunsSHighMtnBand}{[23.3, 23.9]}
\newcommand{\PhKindFoundRunsSHills}{20.2}
\newcommand{\PhKindFoundRunsSHillsBand}{[19.8, 20.5]}
\newcommand{\PhKindFoundRunsSLowMtn}{22.2}
\newcommand{\PhKindFoundRunsSLowMtnBand}{[21.9, 22.4]}
\newcommand{\PhKindFoundRunsSPlains}{19.1}
\newcommand{\PhKindFoundRunsSPlainsBand}{[18.7, 19.5]}
\newcommand{\PhKindFoundTimeAAll}{24.4}
\newcommand{\PhKindFoundTimeAAllBand}{[23.9, 24.8]}
\newcommand{\PhKindFoundTimeAHighMtn}{24.9}
\newcommand{\PhKindFoundTimeAHighMtnBand}{[24.2, 25.6]}
\newcommand{\PhKindFoundTimeAHills}{20.7}
\newcommand{\PhKindFoundTimeAHillsBand}{[19.4, 22.1]}
\newcommand{\PhKindFoundTimeALowMtn}{26.2}
\newcommand{\PhKindFoundTimeALowMtnBand}{[25.6, 26.8]}
\newcommand{\PhKindFoundTimeAPlains}{19.1}
\newcommand{\PhKindFoundTimeAPlainsBand}{[18.1, 20.1]}
\newcommand{\PhKindFoundTimeASAll}{$+$1.73}
\newcommand{\PhKindFoundTimeASAllBand}{[$+$1.27, $+$2.19]}
\newcommand{\PhKindFoundTimeASHighMtn}{$-$1.00}
\newcommand{\PhKindFoundTimeASHighMtnBand}{[$-$1.81, $-$0.19]}
\newcommand{\PhKindFoundTimeASHills}{$+$0.02}
\newcommand{\PhKindFoundTimeASHillsBand}{[$-$1.41, $+$1.44]}
\newcommand{\PhKindFoundTimeASLowMtn}{$+$1.90}
\newcommand{\PhKindFoundTimeASLowMtnBand}{[$+$1.26, $+$2.54]}
\newcommand{\PhKindFoundTimeASPlains}{$+$0.32}
\newcommand{\PhKindFoundTimeASPlainsBand}{[$-$0.77, $+$1.40]}
\newcommand{\PhKindFoundTimeGAAll}{$+$0.90}
\newcommand{\PhKindFoundTimeGAAllBand}{[$+$0.37, $+$1.43]}
\newcommand{\PhKindFoundTimeGAHighMtn}{$+$2.68}
\newcommand{\PhKindFoundTimeGAHighMtnBand}{[$+$1.63, $+$3.73]}
\newcommand{\PhKindFoundTimeGAHills}{$-$0.49}
\newcommand{\PhKindFoundTimeGAHillsBand}{[$-$1.85, $+$0.87]}
\newcommand{\PhKindFoundTimeGALowMtn}{$+$1.50}
\newcommand{\PhKindFoundTimeGALowMtnBand}{[$+$0.67, $+$2.32]}
\newcommand{\PhKindFoundTimeGAPlains}{$-$1.12}
\newcommand{\PhKindFoundTimeGAPlainsBand}{[$-$2.38, $+$0.14]}
\newcommand{\PhKindFoundTimeGAll}{25.3}
\newcommand{\PhKindFoundTimeGAllBand}{[24.8, 25.8]}
\newcommand{\PhKindFoundTimeGHighMtn}{27.6}
\newcommand{\PhKindFoundTimeGHighMtnBand}{[26.7, 28.5]}
\newcommand{\PhKindFoundTimeGHills}{20.3}
\newcommand{\PhKindFoundTimeGHillsBand}{[19.2, 21.4]}
\newcommand{\PhKindFoundTimeGLowMtn}{27.7}
\newcommand{\PhKindFoundTimeGLowMtnBand}{[27.0, 28.3]}
\newcommand{\PhKindFoundTimeGPlains}{18.0}
\newcommand{\PhKindFoundTimeGPlainsBand}{[17.2, 18.8]}
\newcommand{\PhKindFoundTimeGSAll}{$+$2.63}
\newcommand{\PhKindFoundTimeGSAllBand}{[$+$2.13, $+$3.13]}
\newcommand{\PhKindFoundTimeGSHighMtn}{$+$1.68}
\newcommand{\PhKindFoundTimeGSHighMtnBand}{[$+$0.75, $+$2.61]}
\newcommand{\PhKindFoundTimeGSHills}{$-$0.47}
\newcommand{\PhKindFoundTimeGSHillsBand}{[$-$1.64, $+$0.70]}
\newcommand{\PhKindFoundTimeGSLowMtn}{$+$3.40}
\newcommand{\PhKindFoundTimeGSLowMtnBand}{[$+$2.72, $+$4.08]}
\newcommand{\PhKindFoundTimeGSPlains}{$-$0.80}
\newcommand{\PhKindFoundTimeGSPlainsBand}{[$-$1.66, $+$0.06]}
\newcommand{\PhKindFoundTimeSAll}{22.7}
\newcommand{\PhKindFoundTimeSAllBand}{[22.4, 22.9]}
\newcommand{\PhKindFoundTimeSHighMtn}{25.9}
\newcommand{\PhKindFoundTimeSHighMtnBand}{[25.5, 26.3]}
\newcommand{\PhKindFoundTimeSHills}{20.7}
\newcommand{\PhKindFoundTimeSHillsBand}{[20.2, 21.2]}
\newcommand{\PhKindFoundTimeSLowMtn}{24.3}
\newcommand{\PhKindFoundTimeSLowMtnBand}{[23.9, 24.6]}
\newcommand{\PhKindFoundTimeSPlains}{18.8}
\newcommand{\PhKindFoundTimeSPlainsBand}{[18.3, 19.3]}
\newcommand{\PhKindRefoundRunsAAll}{23.5}
\newcommand{\PhKindRefoundRunsAAllBand}{[23.1, 23.8]}
\newcommand{\PhKindRefoundRunsAHighMtn}{23.3}
\newcommand{\PhKindRefoundRunsAHighMtnBand}{[22.7, 23.9]}
\newcommand{\PhKindRefoundRunsAHills}{26.7}
\newcommand{\PhKindRefoundRunsAHillsBand}{[25.6, 27.8]}
\newcommand{\PhKindRefoundRunsALowMtn}{21.6}
\newcommand{\PhKindRefoundRunsALowMtnBand}{[21.1, 22.0]}
\newcommand{\PhKindRefoundRunsAPlains}{27.8}
\newcommand{\PhKindRefoundRunsAPlainsBand}{[26.8, 28.7]}
\newcommand{\PhKindRefoundRunsASAll}{$+$0.45}
\newcommand{\PhKindRefoundRunsASAllBand}{[$+$0.11, $+$0.79]}
\newcommand{\PhKindRefoundRunsASHighMtn}{$+$1.01}
\newcommand{\PhKindRefoundRunsASHighMtnBand}{[$+$0.36, $+$1.66]}
\newcommand{\PhKindRefoundRunsASHills}{$+$3.18}
\newcommand{\PhKindRefoundRunsASHillsBand}{[$+$2.02, $+$4.34]}
\newcommand{\PhKindRefoundRunsASLowMtn}{$+$0.19}
\newcommand{\PhKindRefoundRunsASLowMtnBand}{[$-$0.29, $+$0.67]}
\newcommand{\PhKindRefoundRunsASPlains}{$+$1.25}
\newcommand{\PhKindRefoundRunsASPlainsBand}{[$+$0.25, $+$2.24]}
\newcommand{\PhKindRefoundRunsGAAll}{$+$2.34}
\newcommand{\PhKindRefoundRunsGAAllBand}{[$+$1.93, $+$2.75]}
\newcommand{\PhKindRefoundRunsGAHighMtn}{$+$2.56}
\newcommand{\PhKindRefoundRunsGAHighMtnBand}{[$+$1.72, $+$3.39]}
\newcommand{\PhKindRefoundRunsGAHills}{$+$0.48}
\newcommand{\PhKindRefoundRunsGAHillsBand}{[$-$0.90, $+$1.85]}
\newcommand{\PhKindRefoundRunsGALowMtn}{$+$2.69}
\newcommand{\PhKindRefoundRunsGALowMtnBand}{[$+$2.13, $+$3.24]}
\newcommand{\PhKindRefoundRunsGAPlains}{$+$1.68}
\newcommand{\PhKindRefoundRunsGAPlainsBand}{[$+$0.53, $+$2.83]}
\newcommand{\PhKindRefoundRunsGAll}{25.8}
\newcommand{\PhKindRefoundRunsGAllBand}{[25.5, 26.2]}
\newcommand{\PhKindRefoundRunsGHighMtn}{25.9}
\newcommand{\PhKindRefoundRunsGHighMtnBand}{[25.2, 26.5]}
\newcommand{\PhKindRefoundRunsGHills}{27.2}
\newcommand{\PhKindRefoundRunsGHillsBand}{[26.3, 28.1]}
\newcommand{\PhKindRefoundRunsGLowMtn}{24.3}
\newcommand{\PhKindRefoundRunsGLowMtnBand}{[23.8, 24.8]}
\newcommand{\PhKindRefoundRunsGPlains}{29.4}
\newcommand{\PhKindRefoundRunsGPlainsBand}{[28.7, 30.2]}
\newcommand{\PhKindRefoundRunsGSAll}{$+$2.79}
\newcommand{\PhKindRefoundRunsGSAllBand}{[$+$2.41, $+$3.17]}
\newcommand{\PhKindRefoundRunsGSHighMtn}{$+$3.57}
\newcommand{\PhKindRefoundRunsGSHighMtnBand}{[$+$2.85, $+$4.29]}
\newcommand{\PhKindRefoundRunsGSHills}{$+$3.66}
\newcommand{\PhKindRefoundRunsGSHillsBand}{[$+$2.65, $+$4.67]}
\newcommand{\PhKindRefoundRunsGSLowMtn}{$+$2.87}
\newcommand{\PhKindRefoundRunsGSLowMtnBand}{[$+$2.37, $+$3.37]}
\newcommand{\PhKindRefoundRunsGSPlains}{$+$2.92}
\newcommand{\PhKindRefoundRunsGSPlainsBand}{[$+$2.09, $+$3.76]}
\newcommand{\PhKindRefoundRunsSAll}{23.0}
\newcommand{\PhKindRefoundRunsSAllBand}{[22.8, 23.3]}
\newcommand{\PhKindRefoundRunsSHighMtn}{22.3}
\newcommand{\PhKindRefoundRunsSHighMtnBand}{[21.9, 22.6]}
\newcommand{\PhKindRefoundRunsSHills}{23.5}
\newcommand{\PhKindRefoundRunsSHillsBand}{[23.1, 23.9]}
\newcommand{\PhKindRefoundRunsSLowMtn}{21.4}
\newcommand{\PhKindRefoundRunsSLowMtnBand}{[21.1, 21.7]}
\newcommand{\PhKindRefoundRunsSPlains}{26.5}
\newcommand{\PhKindRefoundRunsSPlainsBand}{[26.1, 26.9]}
\newcommand{\PhKindRefoundTimeAAll}{20.7}
\newcommand{\PhKindRefoundTimeAAllBand}{[20.3, 21.1]}
\newcommand{\PhKindRefoundTimeAHighMtn}{20.5}
\newcommand{\PhKindRefoundTimeAHighMtnBand}{[19.7, 21.3]}
\newcommand{\PhKindRefoundTimeAHills}{24.0}
\newcommand{\PhKindRefoundTimeAHillsBand}{[22.6, 25.4]}
\newcommand{\PhKindRefoundTimeALowMtn}{18.8}
\newcommand{\PhKindRefoundTimeALowMtnBand}{[18.3, 19.4]}
\newcommand{\PhKindRefoundTimeAPlains}{26.2}
\newcommand{\PhKindRefoundTimeAPlainsBand}{[24.9, 27.5]}
\newcommand{\PhKindRefoundTimeASAll}{$+$0.56}
\newcommand{\PhKindRefoundTimeASAllBand}{[$+$0.13, $+$1.00]}
\newcommand{\PhKindRefoundTimeASHighMtn}{$+$1.16}
\newcommand{\PhKindRefoundTimeASHighMtnBand}{[$+$0.32, $+$2.01]}
\newcommand{\PhKindRefoundTimeASHills}{$+$3.11}
\newcommand{\PhKindRefoundTimeASHillsBand}{[$+$1.68, $+$4.54]}
\newcommand{\PhKindRefoundTimeASLowMtn}{$+$0.64}
\newcommand{\PhKindRefoundTimeASLowMtnBand}{[$+$0.04, $+$1.23]}
\newcommand{\PhKindRefoundTimeASPlains}{$+$1.79}
\newcommand{\PhKindRefoundTimeASPlainsBand}{[$+$0.47, $+$3.10]}
\newcommand{\PhKindRefoundTimeGAAll}{$+$3.08}
\newcommand{\PhKindRefoundTimeGAAllBand}{[$+$2.54, $+$3.63]}
\newcommand{\PhKindRefoundTimeGAHighMtn}{$+$3.34}
\newcommand{\PhKindRefoundTimeGAHighMtnBand}{[$+$2.20, $+$4.48]}
\newcommand{\PhKindRefoundTimeGAHills}{$+$2.01}
\newcommand{\PhKindRefoundTimeGAHillsBand}{[$+$0.23, $+$3.79]}
\newcommand{\PhKindRefoundTimeGALowMtn}{$+$3.11}
\newcommand{\PhKindRefoundTimeGALowMtnBand}{[$+$2.36, $+$3.86]}
\newcommand{\PhKindRefoundTimeGAPlains}{$+$2.14}
\newcommand{\PhKindRefoundTimeGAPlainsBand}{[$+$0.62, $+$3.67]}
\newcommand{\PhKindRefoundTimeGAll}{23.8}
\newcommand{\PhKindRefoundTimeGAllBand}{[23.3, 24.2]}
\newcommand{\PhKindRefoundTimeGHighMtn}{23.8}
\newcommand{\PhKindRefoundTimeGHighMtnBand}{[22.9, 24.8]}
\newcommand{\PhKindRefoundTimeGHills}{26.0}
\newcommand{\PhKindRefoundTimeGHillsBand}{[24.8, 27.2]}
\newcommand{\PhKindRefoundTimeGLowMtn}{21.9}
\newcommand{\PhKindRefoundTimeGLowMtnBand}{[21.3, 22.6]}
\newcommand{\PhKindRefoundTimeGPlains}{28.3}
\newcommand{\PhKindRefoundTimeGPlainsBand}{[27.4, 29.2]}
\newcommand{\PhKindRefoundTimeGSAll}{$+$3.65}
\newcommand{\PhKindRefoundTimeGSAllBand}{[$+$3.15, $+$4.14]}
\newcommand{\PhKindRefoundTimeGSHighMtn}{$+$4.50}
\newcommand{\PhKindRefoundTimeGSHighMtnBand}{[$+$3.53, $+$5.47]}
\newcommand{\PhKindRefoundTimeGSHills}{$+$5.12}
\newcommand{\PhKindRefoundTimeGSHillsBand}{[$+$3.82, $+$6.41]}
\newcommand{\PhKindRefoundTimeGSLowMtn}{$+$3.74}
\newcommand{\PhKindRefoundTimeGSLowMtnBand}{[$+$3.10, $+$4.38]}
\newcommand{\PhKindRefoundTimeGSPlains}{$+$3.93}
\newcommand{\PhKindRefoundTimeGSPlainsBand}{[$+$2.90, $+$4.96]}
\newcommand{\PhKindRefoundTimeSAll}{20.1}
\newcommand{\PhKindRefoundTimeSAllBand}{[19.9, 20.4]}
\newcommand{\PhKindRefoundTimeSHighMtn}{19.3}
\newcommand{\PhKindRefoundTimeSHighMtnBand}{[18.9, 19.7]}
\newcommand{\PhKindRefoundTimeSHills}{20.9}
\newcommand{\PhKindRefoundTimeSHillsBand}{[20.4, 21.4]}
\newcommand{\PhKindRefoundTimeSLowMtn}{18.2}
\newcommand{\PhKindRefoundTimeSLowMtnBand}{[17.9, 18.5]}
\newcommand{\PhKindRefoundTimeSPlains}{24.4}
\newcommand{\PhKindRefoundTimeSPlainsBand}{[24.0, 24.9]}
\newcommand{\PhMuClimbAAll}{10.1}
\newcommand{\PhMuClimbAAllBand}{[7.0, 13.2]}
\newcommand{\PhMuClimbAHighMtn}{6.49}
\newcommand{\PhMuClimbAHighMtnBand}{[5.06, 7.93]}
\newcommand{\PhMuClimbAHills}{8.38}
\newcommand{\PhMuClimbAHillsBand}{[5.88, 10.88]}
\newcommand{\PhMuClimbALowMtn}{6.94}
\newcommand{\PhMuClimbALowMtnBand}{[4.96, 8.92]}
\newcommand{\PhMuClimbAPlains}{9.77}
\newcommand{\PhMuClimbAPlainsBand}{[7.41, 12.14]}
\newcommand{\PhMuClimbASAll}{$+$0.258}
\newcommand{\PhMuClimbASAllBand}{[$-$1.824, $+$2.341]}
\newcommand{\PhMuClimbASHighMtn}{$+$0.0260}
\newcommand{\PhMuClimbASHighMtnBand}{[$-$0.7933, $+$0.8452]}
\newcommand{\PhMuClimbASHills}{$-$1.15}
\newcommand{\PhMuClimbASHillsBand}{[$-$3.40, $+$1.09]}
\newcommand{\PhMuClimbASLowMtn}{$-$0.194}
\newcommand{\PhMuClimbASLowMtnBand}{[$-$1.098, $+$0.711]}
\newcommand{\PhMuClimbASPlains}{$-$0.768}
\newcommand{\PhMuClimbASPlainsBand}{[$-$2.620, $+$1.084]}
\newcommand{\PhMuClimbGAAll}{$-$0.600}
\newcommand{\PhMuClimbGAAllBand}{[$-$3.547, $+$2.348]}
\newcommand{\PhMuClimbGAHighMtn}{$-$0.107}
\newcommand{\PhMuClimbGAHighMtnBand}{[$-$1.141, $+$0.926]}
\newcommand{\PhMuClimbGAHills}{$+$7.91}
\newcommand{\PhMuClimbGAHillsBand}{[$+$1.43, $+$14.39]}
\newcommand{\PhMuClimbGALowMtn}{$+$1.33}
\newcommand{\PhMuClimbGALowMtnBand}{[$-$0.40, $+$3.06]}
\newcommand{\PhMuClimbGAPlains}{$-$1.78}
\newcommand{\PhMuClimbGAPlainsBand}{[$-$3.77, $+$0.22]}
\newcommand{\PhMuClimbGAll}{9.50}
\newcommand{\PhMuClimbGAllBand}{[7.41, 11.60]}
\newcommand{\PhMuClimbGHighMtn}{6.39}
\newcommand{\PhMuClimbGHighMtnBand}{[5.19, 7.58]}
\newcommand{\PhMuClimbGHills}{16.3}
\newcommand{\PhMuClimbGHillsBand}{[9.0, 23.6]}
\newcommand{\PhMuClimbGLowMtn}{8.27}
\newcommand{\PhMuClimbGLowMtnBand}{[5.40, 11.14]}
\newcommand{\PhMuClimbGPlains}{8.00}
\newcommand{\PhMuClimbGPlainsBand}{[6.27, 9.73]}
\newcommand{\PhMuClimbGSAll}{$-$0.341}
\newcommand{\PhMuClimbGSAllBand}{[$-$2.302, $+$1.620]}
\newcommand{\PhMuClimbGSHighMtn}{$-$0.0812}
\newcommand{\PhMuClimbGSHighMtnBand}{[$-$0.8817, $+$0.7193]}
\newcommand{\PhMuClimbGSHills}{$+$6.76}
\newcommand{\PhMuClimbGSHillsBand}{[$+$1.04, $+$12.47]}
\newcommand{\PhMuClimbGSLowMtn}{$+$1.14}
\newcommand{\PhMuClimbGSLowMtnBand}{[$-$0.63, $+$2.91]}
\newcommand{\PhMuClimbGSPlains}{$-$2.54}
\newcommand{\PhMuClimbGSPlainsBand}{[$-$4.64, $-$0.45]}
\newcommand{\PhMuClimbSAll}{9.84}
\newcommand{\PhMuClimbSAllBand}{[7.84, 11.84]}
\newcommand{\PhMuClimbSHighMtn}{6.47}
\newcommand{\PhMuClimbSHighMtnBand}{[5.48, 7.46]}
\newcommand{\PhMuClimbSHills}{9.54}
\newcommand{\PhMuClimbSHillsBand}{[7.00, 12.08]}
\newcommand{\PhMuClimbSLowMtn}{7.13}
\newcommand{\PhMuClimbSLowMtnBand}{[5.22, 9.04]}
\newcommand{\PhMuClimbSPlains}{10.5}
\newcommand{\PhMuClimbSPlainsBand}{[7.7, 13.4]}
\newcommand{\PhMuGlideAAll}{4.34}
\newcommand{\PhMuGlideAAllBand}{[4.02, 4.66]}
\newcommand{\PhMuGlideAHighMtn}{4.38}
\newcommand{\PhMuGlideAHighMtnBand}{[3.97, 4.80]}
\newcommand{\PhMuGlideAHills}{5.12}
\newcommand{\PhMuGlideAHillsBand}{[4.19, 6.05]}
\newcommand{\PhMuGlideALowMtn}{4.34}
\newcommand{\PhMuGlideALowMtnBand}{[4.02, 4.67]}
\newcommand{\PhMuGlideAPlains}{5.26}
\newcommand{\PhMuGlideAPlainsBand}{[4.33, 6.20]}
\newcommand{\PhMuGlideASAll}{$-$0.388}
\newcommand{\PhMuGlideASAllBand}{[$-$0.792, $+$0.016]}
\newcommand{\PhMuGlideASHighMtn}{$-$0.208}
\newcommand{\PhMuGlideASHighMtnBand}{[$-$0.746, $+$0.330]}
\newcommand{\PhMuGlideASHills}{$+$0.749}
\newcommand{\PhMuGlideASHillsBand}{[$-$0.041, $+$1.538]}
\newcommand{\PhMuGlideASLowMtn}{$-$0.286}
\newcommand{\PhMuGlideASLowMtnBand}{[$-$0.683, $+$0.111]}
\newcommand{\PhMuGlideASPlains}{$-$0.632}
\newcommand{\PhMuGlideASPlainsBand}{[$-$1.479, $+$0.215]}
\newcommand{\PhMuGlideGAAll}{$-$0.883}
\newcommand{\PhMuGlideGAAllBand}{[$-$1.311, $-$0.455]}
\newcommand{\PhMuGlideGAHighMtn}{$-$1.20}
\newcommand{\PhMuGlideGAHighMtnBand}{[$-$1.82, $-$0.58]}
\newcommand{\PhMuGlideGAHills}{$+$0.240}
\newcommand{\PhMuGlideGAHillsBand}{[$-$0.773, $+$1.253]}
\newcommand{\PhMuGlideGALowMtn}{$-$0.861}
\newcommand{\PhMuGlideGALowMtnBand}{[$-$1.304, $-$0.419]}
\newcommand{\PhMuGlideGAPlains}{$+$0.127}
\newcommand{\PhMuGlideGAPlainsBand}{[$-$1.232, $+$1.485]}
\newcommand{\PhMuGlideGAll}{3.46}
\newcommand{\PhMuGlideGAllBand}{[3.14, 3.77]}
\newcommand{\PhMuGlideGHighMtn}{3.19}
\newcommand{\PhMuGlideGHighMtnBand}{[2.77, 3.60]}
\newcommand{\PhMuGlideGHills}{5.36}
\newcommand{\PhMuGlideGHillsBand}{[4.28, 6.44]}
\newcommand{\PhMuGlideGLowMtn}{3.48}
\newcommand{\PhMuGlideGLowMtnBand}{[3.14, 3.83]}
\newcommand{\PhMuGlideGPlains}{5.39}
\newcommand{\PhMuGlideGPlainsBand}{[4.09, 6.69]}
\newcommand{\PhMuGlideGSAll}{$-$1.27}
\newcommand{\PhMuGlideGSAllBand}{[$-$1.72, $-$0.82]}
\newcommand{\PhMuGlideGSHighMtn}{$-$1.41}
\newcommand{\PhMuGlideGSHighMtnBand}{[$-$2.04, $-$0.77]}
\newcommand{\PhMuGlideGSHills}{$+$0.989}
\newcommand{\PhMuGlideGSHillsBand}{[$+$0.042, $+$1.936]}
\newcommand{\PhMuGlideGSLowMtn}{$-$1.15}
\newcommand{\PhMuGlideGSLowMtnBand}{[$-$1.65, $-$0.64]}
\newcommand{\PhMuGlideGSPlains}{$-$0.505}
\newcommand{\PhMuGlideGSPlainsBand}{[$-$1.781, $+$0.771]}
\newcommand{\PhMuGlideSAll}{4.73}
\newcommand{\PhMuGlideSAllBand}{[4.46, 4.99]}
\newcommand{\PhMuGlideSHighMtn}{4.59}
\newcommand{\PhMuGlideSHighMtnBand}{[4.21, 4.98]}
\newcommand{\PhMuGlideSHills}{4.37}
\newcommand{\PhMuGlideSHillsBand}{[4.06, 4.68]}
\newcommand{\PhMuGlideSLowMtn}{4.63}
\newcommand{\PhMuGlideSLowMtnBand}{[4.31, 4.95]}
\newcommand{\PhMuGlideSPlains}{5.89}
\newcommand{\PhMuGlideSPlainsBand}{[5.05, 6.73]}
\newcommand{\PhMuSearchAAll}{3.93}
\newcommand{\PhMuSearchAAllBand}{[3.53, 4.33]}
\newcommand{\PhMuSearchAHighMtn}{3.84}
\newcommand{\PhMuSearchAHighMtnBand}{[3.22, 4.46]}
\newcommand{\PhMuSearchAHills}{4.31}
\newcommand{\PhMuSearchAHillsBand}{[2.62, 6.01]}
\newcommand{\PhMuSearchALowMtn}{4.35}
\newcommand{\PhMuSearchALowMtnBand}{[3.78, 4.92]}
\newcommand{\PhMuSearchAPlains}{3.40}
\newcommand{\PhMuSearchAPlainsBand}{[1.73, 5.07]}
\newcommand{\PhMuSearchASAll}{$+$0.131}
\newcommand{\PhMuSearchASAllBand}{[$-$0.257, $+$0.520]}
\newcommand{\PhMuSearchASHighMtn}{$-$0.138}
\newcommand{\PhMuSearchASHighMtnBand}{[$-$0.619, $+$0.342]}
\newcommand{\PhMuSearchASHills}{$+$0.562}
\newcommand{\PhMuSearchASHillsBand}{[$-$0.767, $+$1.890]}
\newcommand{\PhMuSearchASLowMtn}{$+$0.427}
\newcommand{\PhMuSearchASLowMtnBand}{[$-$0.092, $+$0.947]}
\newcommand{\PhMuSearchASPlains}{$-$0.0191}
\newcommand{\PhMuSearchASPlainsBand}{[$-$1.0263, $+$0.9882]}
\newcommand{\PhMuSearchGAAll}{$+$0.250}
\newcommand{\PhMuSearchGAAllBand}{[$-$0.368, $+$0.867]}
\newcommand{\PhMuSearchGAHighMtn}{$+$0.858}
\newcommand{\PhMuSearchGAHighMtnBand}{[$+$0.125, $+$1.592]}
\newcommand{\PhMuSearchGAHills}{$+$0.106}
\newcommand{\PhMuSearchGAHillsBand}{[$-$1.575, $+$1.788]}
\newcommand{\PhMuSearchGALowMtn}{$-$0.0209}
\newcommand{\PhMuSearchGALowMtnBand}{[$-$0.7458, $+$0.7040]}
\newcommand{\PhMuSearchGAPlains}{$+$0.407}
\newcommand{\PhMuSearchGAPlainsBand}{[$-$0.636, $+$1.450]}
\newcommand{\PhMuSearchGAll}{4.18}
\newcommand{\PhMuSearchGAllBand}{[3.64, 4.72]}
\newcommand{\PhMuSearchGHighMtn}{4.70}
\newcommand{\PhMuSearchGHighMtnBand}{[3.68, 5.71]}
\newcommand{\PhMuSearchGHills}{4.42}
\newcommand{\PhMuSearchGHillsBand}{[2.95, 5.89]}
\newcommand{\PhMuSearchGLowMtn}{4.33}
\newcommand{\PhMuSearchGLowMtnBand}{[3.75, 4.91]}
\newcommand{\PhMuSearchGPlains}{3.81}
\newcommand{\PhMuSearchGPlainsBand}{[2.45, 5.17]}
\newcommand{\PhMuSearchGSAll}{$+$0.381}
\newcommand{\PhMuSearchGSAllBand}{[$-$0.110, $+$0.872]}
\newcommand{\PhMuSearchGSHighMtn}{$+$0.720}
\newcommand{\PhMuSearchGSHighMtnBand}{[$+$0.040, $+$1.401]}
\newcommand{\PhMuSearchGSHills}{$+$0.668}
\newcommand{\PhMuSearchGSHillsBand}{[$-$0.464, $+$1.800]}
\newcommand{\PhMuSearchGSLowMtn}{$+$0.406}
\newcommand{\PhMuSearchGSLowMtnBand}{[$-$0.135, $+$0.948]}
\newcommand{\PhMuSearchGSPlains}{$+$0.388}
\newcommand{\PhMuSearchGSPlainsBand}{[$-$0.366, $+$1.142]}
\newcommand{\PhMuSearchSAll}{3.80}
\newcommand{\PhMuSearchSAllBand}{[3.54, 4.06]}
\newcommand{\PhMuSearchSHighMtn}{3.98}
\newcommand{\PhMuSearchSHighMtnBand}{[3.42, 4.53]}
\newcommand{\PhMuSearchSHills}{3.75}
\newcommand{\PhMuSearchSHillsBand}{[3.07, 4.43]}
\newcommand{\PhMuSearchSLowMtn}{3.92}
\newcommand{\PhMuSearchSLowMtnBand}{[3.57, 4.28]}
\newcommand{\PhMuSearchSPlains}{3.42}
\newcommand{\PhMuSearchSPlainsBand}{[2.48, 4.36]}
\newcommand{\PhPeakGALagAll}{1\,230}
\newcommand{\PhPeakGALagHighMtn}{870}
\newcommand{\PhPeakGALagHills}{2\,250}
\newcommand{\PhPeakGALagLowMtn}{710}
\newcommand{\PhPeakGALagPlains}{1\,500}
\newcommand{\PhPeakGSLagAll}{870}
\newcommand{\PhPeakGSLagHighMtn}{1\,230}
\newcommand{\PhPeakGSLagHills}{2\,250}
\newcommand{\PhPeakGSLagLowMtn}{710}
\newcommand{\PhPeakGSLagPlains}{2\,250}
\newcommand{\PhRateAAll}{1.18}
\newcommand{\PhRateAAllBand}{[1.17, 1.19]}
\newcommand{\PhRateAHighMtn}{1.30}
\newcommand{\PhRateAHighMtnBand}{[1.28, 1.31]}
\newcommand{\PhRateAHills}{1.09}
\newcommand{\PhRateAHillsBand}{[1.06, 1.11]}
\newcommand{\PhRateALowMtn}{1.23}
\newcommand{\PhRateALowMtnBand}{[1.21, 1.24]}
\newcommand{\PhRateAPlains}{0.980}
\newcommand{\PhRateAPlainsBand}{[0.959, 1.001]}
\newcommand{\PhRateASAll}{$+$0.0701}
\newcommand{\PhRateASAllBand}{[$+$0.0588, $+$0.0814]}
\newcommand{\PhRateASHighMtn}{$+$0.0205}
\newcommand{\PhRateASHighMtnBand}{[$+$0.0045, $+$0.0365]}
\newcommand{\PhRateASHills}{$+$0.0262}
\newcommand{\PhRateASHillsBand}{[$+$0.0010, $+$0.0514]}
\newcommand{\PhRateASLowMtn}{$+$0.0517}
\newcommand{\PhRateASLowMtnBand}{[$+$0.0364, $+$0.0669]}
\newcommand{\PhRateASPlains}{$+$0.0098}
\newcommand{\PhRateASPlainsBand}{[$-$0.0080, $+$0.0276]}
\newcommand{\PhRateGAAll}{$-$0.0089}
\newcommand{\PhRateGAAllBand}{[$-$0.0164, $-$0.0014]}
\newcommand{\PhRateGAHighMtn}{$+$0.0161}
\newcommand{\PhRateGAHighMtnBand}{[$+$0.0024, $+$0.0298]}
\newcommand{\PhRateGAHills}{$-$0.0049}
\newcommand{\PhRateGAHillsBand}{[$-$0.0233, $+$0.0136]}
\newcommand{\PhRateGALowMtn}{$+$0.0012}
\newcommand{\PhRateGALowMtnBand}{[$-$0.0095, $+$0.0120]}
\newcommand{\PhRateGAPlains}{$-$0.0317}
\newcommand{\PhRateGAPlainsBand}{[$-$0.0474, $-$0.0160]}
\newcommand{\PhRateGAll}{1.17}
\newcommand{\PhRateGAllBand}{[1.16, 1.18]}
\newcommand{\PhRateGHighMtn}{1.31}
\newcommand{\PhRateGHighMtnBand}{[1.30, 1.33]}
\newcommand{\PhRateGHills}{1.08}
\newcommand{\PhRateGHillsBand}{[1.06, 1.11]}
\newcommand{\PhRateGLowMtn}{1.23}
\newcommand{\PhRateGLowMtnBand}{[1.22, 1.25]}
\newcommand{\PhRateGPlains}{0.948}
\newcommand{\PhRateGPlainsBand}{[0.930, 0.967]}
\newcommand{\PhRateGSAll}{$+$0.0612}
\newcommand{\PhRateGSAllBand}{[$+$0.0492, $+$0.0732]}
\newcommand{\PhRateGSHighMtn}{$+$0.0366}
\newcommand{\PhRateGSHighMtnBand}{[$+$0.0197, $+$0.0535]}
\newcommand{\PhRateGSHills}{$+$0.0213}
\newcommand{\PhRateGSHillsBand}{[$-$0.0030, $+$0.0456]}
\newcommand{\PhRateGSLowMtn}{$+$0.0529}
\newcommand{\PhRateGSLowMtnBand}{[$+$0.0379, $+$0.0679]}
\newcommand{\PhRateGSPlains}{$-$0.0219}
\newcommand{\PhRateGSPlainsBand}{[$-$0.0379, $-$0.0059]}
\newcommand{\PhRateSAll}{1.11}
\newcommand{\PhRateSAllBand}{[1.10, 1.12]}
\newcommand{\PhRateSHighMtn}{1.28}
\newcommand{\PhRateSHighMtnBand}{[1.27, 1.29]}
\newcommand{\PhRateSHills}{1.06}
\newcommand{\PhRateSHillsBand}{[1.05, 1.07]}
\newcommand{\PhRateSLowMtn}{1.18}
\newcommand{\PhRateSLowMtnBand}{[1.17, 1.19]}
\newcommand{\PhRateSPlains}{0.970}
\newcommand{\PhRateSPlainsBand}{[0.960, 0.981]}
\newcommand{\PhRefoundMeanAAll}{69.3}
\newcommand{\PhRefoundMeanAAllBand}{[68.0, 70.5]}
\newcommand{\PhRefoundMeanAHighMtn}{71.0}
\newcommand{\PhRefoundMeanAHighMtnBand}{[68.8, 73.1]}
\newcommand{\PhRefoundMeanAHills}{58.9}
\newcommand{\PhRefoundMeanAHillsBand}{[55.8, 62.1]}
\newcommand{\PhRefoundMeanALowMtn}{72.8}
\newcommand{\PhRefoundMeanALowMtnBand}{[71.0, 74.6]}
\newcommand{\PhRefoundMeanAPlains}{66.3}
\newcommand{\PhRefoundMeanAPlainsBand}{[64.0, 68.7]}
\newcommand{\PhRefoundMeanASAll}{$+$3.70}
\newcommand{\PhRefoundMeanASAllBand}{[$+$2.46, $+$4.93]}
\newcommand{\PhRefoundMeanASHighMtn}{$+$1.29}
\newcommand{\PhRefoundMeanASHighMtnBand}{[$-$1.02, $+$3.61]}
\newcommand{\PhRefoundMeanASHills}{$-$3.36}
\newcommand{\PhRefoundMeanASHillsBand}{[$-$6.59, $-$0.13]}
\newcommand{\PhRefoundMeanASLowMtn}{$+$5.42}
\newcommand{\PhRefoundMeanASLowMtnBand}{[$+$3.56, $+$7.27]}
\newcommand{\PhRefoundMeanASPlains}{$+$1.00}
\newcommand{\PhRefoundMeanASPlainsBand}{[$-$1.50, $+$3.51]}
\newcommand{\PhRefoundMeanGAAll}{$-$0.746}
\newcommand{\PhRefoundMeanGAAllBand}{[$-$2.057, $+$0.565]}
\newcommand{\PhRefoundMeanGAHighMtn}{$-$1.15}
\newcommand{\PhRefoundMeanGAHighMtnBand}{[$-$3.82, $+$1.52]}
\newcommand{\PhRefoundMeanGAHills}{$+$3.75}
\newcommand{\PhRefoundMeanGAHillsBand}{[$+$0.59, $+$6.91]}
\newcommand{\PhRefoundMeanGALowMtn}{$-$1.62}
\newcommand{\PhRefoundMeanGALowMtnBand}{[$-$3.73, $+$0.49]}
\newcommand{\PhRefoundMeanGAPlains}{$-$0.485}
\newcommand{\PhRefoundMeanGAPlainsBand}{[$-$3.213, $+$2.243]}
\newcommand{\PhRefoundMeanGAll}{68.5}
\newcommand{\PhRefoundMeanGAllBand}{[67.5, 69.5]}
\newcommand{\PhRefoundMeanGHighMtn}{69.8}
\newcommand{\PhRefoundMeanGHighMtnBand}{[67.6, 72.0]}
\newcommand{\PhRefoundMeanGHills}{62.7}
\newcommand{\PhRefoundMeanGHillsBand}{[60.4, 65.0]}
\newcommand{\PhRefoundMeanGLowMtn}{71.2}
\newcommand{\PhRefoundMeanGLowMtnBand}{[69.6, 72.7]}
\newcommand{\PhRefoundMeanGPlains}{65.8}
\newcommand{\PhRefoundMeanGPlainsBand}{[64.0, 67.7]}
\newcommand{\PhRefoundMeanGSAll}{$+$2.95}
\newcommand{\PhRefoundMeanGSAllBand}{[$+$1.87, $+$4.03]}
\newcommand{\PhRefoundMeanGSHighMtn}{$+$0.145}
\newcommand{\PhRefoundMeanGSHighMtnBand}{[$-$2.211, $+$2.501]}
\newcommand{\PhRefoundMeanGSHills}{$+$0.395}
\newcommand{\PhRefoundMeanGSHillsBand}{[$-$1.963, $+$2.752]}
\newcommand{\PhRefoundMeanGSLowMtn}{$+$3.80}
\newcommand{\PhRefoundMeanGSLowMtnBand}{[$+$2.10, $+$5.49]}
\newcommand{\PhRefoundMeanGSPlains}{$+$0.520}
\newcommand{\PhRefoundMeanGSPlainsBand}{[$-$1.539, $+$2.579]}
\newcommand{\PhRefoundMeanSAll}{65.6}
\newcommand{\PhRefoundMeanSAllBand}{[65.0, 66.1]}
\newcommand{\PhRefoundMeanSHighMtn}{69.7}
\newcommand{\PhRefoundMeanSHighMtnBand}{[68.6, 70.7]}
\newcommand{\PhRefoundMeanSHills}{62.3}
\newcommand{\PhRefoundMeanSHillsBand}{[61.2, 63.4]}
\newcommand{\PhRefoundMeanSLowMtn}{67.4}
\newcommand{\PhRefoundMeanSLowMtnBand}{[66.5, 68.2]}
\newcommand{\PhRefoundMeanSPlains}{65.3}
\newcommand{\PhRefoundMeanSPlainsBand}{[64.1, 66.6]}
\newcommand{\PhRefoundPerHourAAll}{1.36}
\newcommand{\PhRefoundPerHourAAllBand}{[1.33, 1.38]}
\newcommand{\PhRefoundPerHourAHighMtn}{1.30}
\newcommand{\PhRefoundPerHourAHighMtnBand}{[1.25, 1.34]}
\newcommand{\PhRefoundPerHourAHills}{1.52}
\newcommand{\PhRefoundPerHourAHillsBand}{[1.42, 1.62]}
\newcommand{\PhRefoundPerHourALowMtn}{1.26}
\newcommand{\PhRefoundPerHourALowMtnBand}{[1.23, 1.29]}
\newcommand{\PhRefoundPerHourAPlains}{1.71}
\newcommand{\PhRefoundPerHourAPlainsBand}{[1.62, 1.80]}
\newcommand{\PhRefoundPerHourASAll}{$-$0.0782}
\newcommand{\PhRefoundPerHourASAllBand}{[$-$0.1061, $-$0.0503]}
\newcommand{\PhRefoundPerHourASHighMtn}{$-$0.0429}
\newcommand{\PhRefoundPerHourASHighMtnBand}{[$-$0.0927, $+$0.0069]}
\newcommand{\PhRefoundPerHourASHills}{$+$0.0624}
\newcommand{\PhRefoundPerHourASHillsBand}{[$-$0.0378, $+$0.1626]}
\newcommand{\PhRefoundPerHourASLowMtn}{$-$0.0726}
\newcommand{\PhRefoundPerHourASLowMtnBand}{[$-$0.1098, $-$0.0353]}
\newcommand{\PhRefoundPerHourASPlains}{$-$0.0756}
\newcommand{\PhRefoundPerHourASPlainsBand}{[$-$0.1638, $+$0.0126]}
\newcommand{\PhRefoundPerHourGAAll}{$+$0.205}
\newcommand{\PhRefoundPerHourGAAllBand}{[$+$0.174, $+$0.236]}
\newcommand{\PhRefoundPerHourGAHighMtn}{$+$0.198}
\newcommand{\PhRefoundPerHourGAHighMtnBand}{[$+$0.134, $+$0.263]}
\newcommand{\PhRefoundPerHourGAHills}{$+$0.122}
\newcommand{\PhRefoundPerHourGAHillsBand}{[$+$0.018, $+$0.227]}
\newcommand{\PhRefoundPerHourGALowMtn}{$+$0.215}
\newcommand{\PhRefoundPerHourGALowMtnBand}{[$+$0.172, $+$0.258]}
\newcommand{\PhRefoundPerHourGAPlains}{$+$0.185}
\newcommand{\PhRefoundPerHourGAPlainsBand}{[$+$0.090, $+$0.279]}
\newcommand{\PhRefoundPerHourGAll}{1.56}
\newcommand{\PhRefoundPerHourGAllBand}{[1.53, 1.59]}
\newcommand{\PhRefoundPerHourGHighMtn}{1.50}
\newcommand{\PhRefoundPerHourGHighMtnBand}{[1.44, 1.55]}
\newcommand{\PhRefoundPerHourGHills}{1.64}
\newcommand{\PhRefoundPerHourGHillsBand}{[1.55, 1.73]}
\newcommand{\PhRefoundPerHourGLowMtn}{1.47}
\newcommand{\PhRefoundPerHourGLowMtnBand}{[1.43, 1.52]}
\newcommand{\PhRefoundPerHourGPlains}{1.90}
\newcommand{\PhRefoundPerHourGPlainsBand}{[1.81, 1.98]}
\newcommand{\PhRefoundPerHourGSAll}{$+$0.127}
\newcommand{\PhRefoundPerHourGSAllBand}{[$+$0.094, $+$0.160]}
\newcommand{\PhRefoundPerHourGSHighMtn}{$+$0.156}
\newcommand{\PhRefoundPerHourGSHighMtnBand}{[$+$0.099, $+$0.212]}
\newcommand{\PhRefoundPerHourGSHills}{$+$0.185}
\newcommand{\PhRefoundPerHourGSHillsBand}{[$+$0.094, $+$0.275]}
\newcommand{\PhRefoundPerHourGSLowMtn}{$+$0.142}
\newcommand{\PhRefoundPerHourGSLowMtnBand}{[$+$0.097, $+$0.188]}
\newcommand{\PhRefoundPerHourGSPlains}{$+$0.109}
\newcommand{\PhRefoundPerHourGSPlainsBand}{[$+$0.024, $+$0.194]}
\newcommand{\PhRefoundPerHourSAll}{1.43}
\newcommand{\PhRefoundPerHourSAllBand}{[1.41, 1.46]}
\newcommand{\PhRefoundPerHourSHighMtn}{1.34}
\newcommand{\PhRefoundPerHourSHighMtnBand}{[1.32, 1.37]}
\newcommand{\PhRefoundPerHourSHills}{1.45}
\newcommand{\PhRefoundPerHourSHillsBand}{[1.41, 1.49]}
\newcommand{\PhRefoundPerHourSLowMtn}{1.33}
\newcommand{\PhRefoundPerHourSLowMtnBand}{[1.31, 1.36]}
\newcommand{\PhRefoundPerHourSPlains}{1.79}
\newcommand{\PhRefoundPerHourSPlainsBand}{[1.74, 1.84]}
\newcommand{\PhRefoundShareAAll}{51.2}
\newcommand{\PhRefoundShareAAllBand}{[50.7, 51.7]}
\newcommand{\PhRefoundShareAHighMtn}{50.1}
\newcommand{\PhRefoundShareAHighMtnBand}{[49.2, 51.0]}
\newcommand{\PhRefoundShareAHills}{57.4}
\newcommand{\PhRefoundShareAHillsBand}{[55.6, 59.1]}
\newcommand{\PhRefoundShareALowMtn}{47.9}
\newcommand{\PhRefoundShareALowMtnBand}{[47.2, 48.5]}
\newcommand{\PhRefoundShareAPlains}{59.3}
\newcommand{\PhRefoundShareAPlainsBand}{[58.0, 60.7]}
\newcommand{\PhRefoundShareASAll}{$-$0.772}
\newcommand{\PhRefoundShareASAllBand}{[$-$1.321, $-$0.222]}
\newcommand{\PhRefoundShareASHighMtn}{$+$1.58}
\newcommand{\PhRefoundShareASHighMtnBand}{[$+$0.62, $+$2.55]}
\newcommand{\PhRefoundShareASHills}{$+$3.52}
\newcommand{\PhRefoundShareASHillsBand}{[$+$1.75, $+$5.28]}
\newcommand{\PhRefoundShareASLowMtn}{$-$1.27}
\newcommand{\PhRefoundShareASLowMtnBand}{[$-$2.01, $-$0.53]}
\newcommand{\PhRefoundShareASPlains}{$+$1.15}
\newcommand{\PhRefoundShareASPlainsBand}{[$-$0.19, $+$2.50]}
\newcommand{\PhRefoundShareGAAll}{$+$1.41}
\newcommand{\PhRefoundShareGAAllBand}{[$+$0.83, $+$1.99]}
\newcommand{\PhRefoundShareGAHighMtn}{$+$0.709}
\newcommand{\PhRefoundShareGAHighMtnBand}{[$-$0.487, $+$1.904]}
\newcommand{\PhRefoundShareGAHills}{$-$0.693}
\newcommand{\PhRefoundShareGAHillsBand}{[$-$2.690, $+$1.304]}
\newcommand{\PhRefoundShareGALowMtn}{$+$1.72}
\newcommand{\PhRefoundShareGALowMtnBand}{[$+$0.90, $+$2.54]}
\newcommand{\PhRefoundShareGAPlains}{$+$1.61}
\newcommand{\PhRefoundShareGAPlainsBand}{[$+$0.03, $+$3.19]}
\newcommand{\PhRefoundShareGAll}{52.6}
\newcommand{\PhRefoundShareGAllBand}{[52.1, 53.2]}
\newcommand{\PhRefoundShareGHighMtn}{50.8}
\newcommand{\PhRefoundShareGHighMtnBand}{[49.9, 51.8]}
\newcommand{\PhRefoundShareGHills}{56.7}
\newcommand{\PhRefoundShareGHillsBand}{[55.3, 58.0]}
\newcommand{\PhRefoundShareGLowMtn}{49.6}
\newcommand{\PhRefoundShareGLowMtnBand}{[48.9, 50.3]}
\newcommand{\PhRefoundShareGPlains}{60.9}
\newcommand{\PhRefoundShareGPlainsBand}{[59.8, 62.1]}
\newcommand{\PhRefoundShareGSAll}{$+$0.639}
\newcommand{\PhRefoundShareGSAllBand}{[$+$0.058, $+$1.221]}
\newcommand{\PhRefoundShareGSHighMtn}{$+$2.29}
\newcommand{\PhRefoundShareGSHighMtnBand}{[$+$1.26, $+$3.33]}
\newcommand{\PhRefoundShareGSHills}{$+$2.82}
\newcommand{\PhRefoundShareGSHillsBand}{[$+$1.33, $+$4.32]}
\newcommand{\PhRefoundShareGSLowMtn}{$+$0.453}
\newcommand{\PhRefoundShareGSLowMtnBand}{[$-$0.254, $+$1.160]}
\newcommand{\PhRefoundShareGSPlains}{$+$2.76}
\newcommand{\PhRefoundShareGSPlainsBand}{[$+$1.49, $+$4.04]}
\newcommand{\PhRefoundShareSAll}{52.0}
\newcommand{\PhRefoundShareSAllBand}{[51.6, 52.4]}
\newcommand{\PhRefoundShareSHighMtn}{48.5}
\newcommand{\PhRefoundShareSHighMtnBand}{[48.0, 49.0]}
\newcommand{\PhRefoundShareSHills}{53.8}
\newcommand{\PhRefoundShareSHillsBand}{[53.1, 54.6]}
\newcommand{\PhRefoundShareSLowMtn}{49.1}
\newcommand{\PhRefoundShareSLowMtnBand}{[48.7, 49.6]}
\newcommand{\PhRefoundShareSPlains}{58.2}
\newcommand{\PhRefoundShareSPlainsBand}{[57.5, 58.9]}
\newcommand{\PhSearchDropped}{4.6}
\newcommand{\PhSearchDroppedAll}{4.6}
\newcommand{\PhSearchDroppedHighMtn}{4.6}
\newcommand{\PhSearchDroppedHills}{5.2}
\newcommand{\PhSearchDroppedLowMtn}{4.6}
\newcommand{\PhSearchDroppedPlains}{4.0}
\newcommand{\PhSearchMeanAAll}{82.3}
\newcommand{\PhSearchMeanAAllBand}{[81.2, 83.3]}
\newcommand{\PhSearchMeanAHighMtn}{84.5}
\newcommand{\PhSearchMeanAHighMtnBand}{[83.0, 86.1]}
\newcommand{\PhSearchMeanAHills}{67.4}
\newcommand{\PhSearchMeanAHillsBand}{[65.2, 69.7]}
\newcommand{\PhSearchMeanALowMtn}{87.6}
\newcommand{\PhSearchMeanALowMtnBand}{[86.2, 88.9]}
\newcommand{\PhSearchMeanAPlains}{71.1}
\newcommand{\PhSearchMeanAPlainsBand}{[68.9, 73.2]}
\newcommand{\PhSearchMeanASAll}{$+$4.44}
\newcommand{\PhSearchMeanASAllBand}{[$+$3.38, $+$5.51]}
\newcommand{\PhSearchMeanASHighMtn}{$+$0.638}
\newcommand{\PhSearchMeanASHighMtnBand}{[$-$1.080, $+$2.355]}
\newcommand{\PhSearchMeanASHills}{$-$5.15}
\newcommand{\PhSearchMeanASHillsBand}{[$-$7.55, $-$2.76]}
\newcommand{\PhSearchMeanASLowMtn}{$+$4.77}
\newcommand{\PhSearchMeanASLowMtnBand}{[$+$3.29, $+$6.25]}
\newcommand{\PhSearchMeanASPlains}{$-$0.0117}
\newcommand{\PhSearchMeanASPlainsBand}{[$-$2.2207, $+$2.1973]}
\newcommand{\PhSearchMeanGAAll}{$-$5.33}
\newcommand{\PhSearchMeanGAAllBand}{[$-$6.34, $-$4.33]}
\newcommand{\PhSearchMeanGAHighMtn}{$-$5.82}
\newcommand{\PhSearchMeanGAHighMtnBand}{[$-$7.80, $-$3.85]}
\newcommand{\PhSearchMeanGAHills}{$-$1.33}
\newcommand{\PhSearchMeanGAHillsBand}{[$-$3.61, $+$0.95]}
\newcommand{\PhSearchMeanGALowMtn}{$-$5.43}
\newcommand{\PhSearchMeanGALowMtnBand}{[$-$7.01, $-$3.85]}
\newcommand{\PhSearchMeanGAPlains}{$-$3.93}
\newcommand{\PhSearchMeanGAPlainsBand}{[$-$6.11, $-$1.76]}
\newcommand{\PhSearchMeanGAll}{76.9}
\newcommand{\PhSearchMeanGAllBand}{[76.0, 77.9]}
\newcommand{\PhSearchMeanGHighMtn}{78.7}
\newcommand{\PhSearchMeanGHighMtnBand}{[77.0, 80.4]}
\newcommand{\PhSearchMeanGHills}{66.1}
\newcommand{\PhSearchMeanGHillsBand}{[64.0, 68.2]}
\newcommand{\PhSearchMeanGLowMtn}{82.1}
\newcommand{\PhSearchMeanGLowMtnBand}{[80.7, 83.6]}
\newcommand{\PhSearchMeanGPlains}{67.1}
\newcommand{\PhSearchMeanGPlainsBand}{[65.6, 68.7]}
\newcommand{\PhSearchMeanGSAll}{$-$0.892}
\newcommand{\PhSearchMeanGSAllBand}{[$-$1.866, $+$0.082]}
\newcommand{\PhSearchMeanGSHighMtn}{$-$5.18}
\newcommand{\PhSearchMeanGSHighMtnBand}{[$-$6.95, $-$3.42]}
\newcommand{\PhSearchMeanGSHills}{$-$6.48}
\newcommand{\PhSearchMeanGSHillsBand}{[$-$8.70, $-$4.27]}
\newcommand{\PhSearchMeanGSLowMtn}{$-$0.664}
\newcommand{\PhSearchMeanGSLowMtnBand}{[$-$2.136, $+$0.809]}
\newcommand{\PhSearchMeanGSPlains}{$-$3.95}
\newcommand{\PhSearchMeanGSPlainsBand}{[$-$5.55, $-$2.34]}
\newcommand{\PhSearchMeanSAll}{77.8}
\newcommand{\PhSearchMeanSAllBand}{[77.2, 78.4]}
\newcommand{\PhSearchMeanSHighMtn}{83.9}
\newcommand{\PhSearchMeanSHighMtnBand}{[83.0, 84.8]}
\newcommand{\PhSearchMeanSHills}{72.6}
\newcommand{\PhSearchMeanSHillsBand}{[71.4, 73.8]}
\newcommand{\PhSearchMeanSLowMtn}{82.8}
\newcommand{\PhSearchMeanSLowMtnBand}{[82.0, 83.6]}
\newcommand{\PhSearchMeanSPlains}{71.1}
\newcommand{\PhSearchMeanSPlainsBand}{[70.1, 72.1]}
\newcommand{\PhSearchPerHourAAll}{4.23}
\newcommand{\PhSearchPerHourAAllBand}{[4.17, 4.30]}
\newcommand{\PhSearchPerHourAHighMtn}{4.19}
\newcommand{\PhSearchPerHourAHighMtnBand}{[4.08, 4.29]}
\newcommand{\PhSearchPerHourAHills}{3.85}
\newcommand{\PhSearchPerHourAHillsBand}{[3.66, 4.04]}
\newcommand{\PhSearchPerHourALowMtn}{4.41}
\newcommand{\PhSearchPerHourALowMtnBand}{[4.32, 4.50]}
\newcommand{\PhSearchPerHourAPlains}{4.11}
\newcommand{\PhSearchPerHourAPlainsBand}{[3.97, 4.25]}
\newcommand{\PhSearchPerHourASAll}{$-$0.249}
\newcommand{\PhSearchPerHourASAllBand}{[$-$0.315, $-$0.182]}
\newcommand{\PhSearchPerHourASHighMtn}{$-$0.360}
\newcommand{\PhSearchPerHourASHighMtnBand}{[$-$0.469, $-$0.250]}
\newcommand{\PhSearchPerHourASHills}{$-$0.416}
\newcommand{\PhSearchPerHourASHillsBand}{[$-$0.608, $-$0.224]}
\newcommand{\PhSearchPerHourASLowMtn}{$-$0.240}
\newcommand{\PhSearchPerHourASLowMtnBand}{[$-$0.330, $-$0.149]}
\newcommand{\PhSearchPerHourASPlains}{$-$0.423}
\newcommand{\PhSearchPerHourASPlainsBand}{[$-$0.558, $-$0.287]}
\newcommand{\PhSearchPerHourGAAll}{$+$0.107}
\newcommand{\PhSearchPerHourGAAllBand}{[$+$0.052, $+$0.161]}
\newcommand{\PhSearchPerHourGAHighMtn}{$+$0.0994}
\newcommand{\PhSearchPerHourGAHighMtnBand}{[$-$0.0103, $+$0.2090]}
\newcommand{\PhSearchPerHourGAHills}{$+$0.140}
\newcommand{\PhSearchPerHourGAHillsBand}{[$-$0.035, $+$0.314]}
\newcommand{\PhSearchPerHourGALowMtn}{$+$0.112}
\newcommand{\PhSearchPerHourGALowMtnBand}{[$+$0.028, $+$0.197]}
\newcommand{\PhSearchPerHourGAPlains}{$+$0.102}
\newcommand{\PhSearchPerHourGAPlainsBand}{[$-$0.028, $+$0.232]}
\newcommand{\PhSearchPerHourGAll}{4.34}
\newcommand{\PhSearchPerHourGAllBand}{[4.26, 4.42]}
\newcommand{\PhSearchPerHourGHighMtn}{4.29}
\newcommand{\PhSearchPerHourGHighMtnBand}{[4.16, 4.41]}
\newcommand{\PhSearchPerHourGHills}{3.99}
\newcommand{\PhSearchPerHourGHillsBand}{[3.82, 4.16]}
\newcommand{\PhSearchPerHourGLowMtn}{4.52}
\newcommand{\PhSearchPerHourGLowMtnBand}{[4.41, 4.63]}
\newcommand{\PhSearchPerHourGPlains}{4.22}
\newcommand{\PhSearchPerHourGPlainsBand}{[4.06, 4.37]}
\newcommand{\PhSearchPerHourGSAll}{$-$0.142}
\newcommand{\PhSearchPerHourGSAllBand}{[$-$0.218, $-$0.065]}
\newcommand{\PhSearchPerHourGSHighMtn}{$-$0.260}
\newcommand{\PhSearchPerHourGSHighMtnBand}{[$-$0.382, $-$0.138]}
\newcommand{\PhSearchPerHourGSHills}{$-$0.277}
\newcommand{\PhSearchPerHourGSHillsBand}{[$-$0.446, $-$0.107]}
\newcommand{\PhSearchPerHourGSLowMtn}{$-$0.127}
\newcommand{\PhSearchPerHourGSLowMtnBand}{[$-$0.235, $-$0.019]}
\newcommand{\PhSearchPerHourGSPlains}{$-$0.320}
\newcommand{\PhSearchPerHourGSPlainsBand}{[$-$0.473, $-$0.168]}
\newcommand{\PhSearchPerHourSAll}{4.48}
\newcommand{\PhSearchPerHourSAllBand}{[4.43, 4.53]}
\newcommand{\PhSearchPerHourSHighMtn}{4.55}
\newcommand{\PhSearchPerHourSHighMtnBand}{[4.49, 4.61]}
\newcommand{\PhSearchPerHourSHills}{4.27}
\newcommand{\PhSearchPerHourSHillsBand}{[4.18, 4.36]}
\newcommand{\PhSearchPerHourSLowMtn}{4.65}
\newcommand{\PhSearchPerHourSLowMtnBand}{[4.59, 4.71]}
\newcommand{\PhSearchPerHourSPlains}{4.54}
\newcommand{\PhSearchPerHourSPlainsBand}{[4.42, 4.65]}
\newcommand{\PhShareClimbAAll}{30.5}
\newcommand{\PhShareClimbAAllBand}{[30.1, 30.9]}
\newcommand{\PhShareClimbAHighMtn}{27.5}
\newcommand{\PhShareClimbAHighMtnBand}{[27.0, 27.9]}
\newcommand{\PhShareClimbAHills}{38.8}
\newcommand{\PhShareClimbAHillsBand}{[37.9, 39.6]}
\newcommand{\PhShareClimbALowMtn}{27.4}
\newcommand{\PhShareClimbALowMtnBand}{[27.0, 27.9]}
\newcommand{\PhShareClimbAPlains}{40.0}
\newcommand{\PhShareClimbAPlainsBand}{[39.3, 40.7]}
\newcommand{\PhShareClimbASAll}{$-$1.14}
\newcommand{\PhShareClimbASAllBand}{[$-$1.50, $-$0.77]}
\newcommand{\PhShareClimbASHighMtn}{$-$0.455}
\newcommand{\PhShareClimbASHighMtnBand}{[$-$0.919, $+$0.008]}
\newcommand{\PhShareClimbASHills}{$+$3.81}
\newcommand{\PhShareClimbASHillsBand}{[$+$2.98, $+$4.64]}
\newcommand{\PhShareClimbASLowMtn}{$-$0.684}
\newcommand{\PhShareClimbASLowMtnBand}{[$-$1.111, $-$0.257]}
\newcommand{\PhShareClimbASPlains}{$+$2.07}
\newcommand{\PhShareClimbASPlainsBand}{[$+$1.38, $+$2.76]}
\newcommand{\PhShareClimbGAAll}{$+$4.45}
\newcommand{\PhShareClimbGAAllBand}{[$+$4.18, $+$4.72]}
\newcommand{\PhShareClimbGAHighMtn}{$+$5.05}
\newcommand{\PhShareClimbGAHighMtnBand}{[$+$4.56, $+$5.55]}
\newcommand{\PhShareClimbGAHills}{$+$3.44}
\newcommand{\PhShareClimbGAHillsBand}{[$+$2.69, $+$4.18]}
\newcommand{\PhShareClimbGALowMtn}{$+$4.00}
\newcommand{\PhShareClimbGALowMtnBand}{[$+$3.64, $+$4.36]}
\newcommand{\PhShareClimbGAPlains}{$+$3.42}
\newcommand{\PhShareClimbGAPlainsBand}{[$+$2.76, $+$4.07]}
\newcommand{\PhShareClimbGAll}{34.9}
\newcommand{\PhShareClimbGAllBand}{[34.5, 35.4]}
\newcommand{\PhShareClimbGHighMtn}{32.5}
\newcommand{\PhShareClimbGHighMtnBand}{[32.0, 33.1]}
\newcommand{\PhShareClimbGHills}{42.2}
\newcommand{\PhShareClimbGHillsBand}{[41.4, 43.0]}
\newcommand{\PhShareClimbGLowMtn}{31.4}
\newcommand{\PhShareClimbGLowMtnBand}{[30.9, 31.9]}
\newcommand{\PhShareClimbGPlains}{43.4}
\newcommand{\PhShareClimbGPlainsBand}{[42.8, 44.1]}
\newcommand{\PhShareClimbGSAll}{$+$3.32}
\newcommand{\PhShareClimbGSAllBand}{[$+$2.91, $+$3.72]}
\newcommand{\PhShareClimbGSHighMtn}{$+$4.60}
\newcommand{\PhShareClimbGSHighMtnBand}{[$+$4.08, $+$5.12]}
\newcommand{\PhShareClimbGSHills}{$+$7.25}
\newcommand{\PhShareClimbGSHillsBand}{[$+$6.47, $+$8.02]}
\newcommand{\PhShareClimbGSLowMtn}{$+$3.31}
\newcommand{\PhShareClimbGSLowMtnBand}{[$+$2.84, $+$3.79]}
\newcommand{\PhShareClimbGSPlains}{$+$5.49}
\newcommand{\PhShareClimbGSPlainsBand}{[$+$4.85, $+$6.13]}
\newcommand{\PhShareClimbSAll}{31.6}
\newcommand{\PhShareClimbSAllBand}{[31.3, 32.0]}
\newcommand{\PhShareClimbSHighMtn}{27.9}
\newcommand{\PhShareClimbSHighMtnBand}{[27.6, 28.3]}
\newcommand{\PhShareClimbSHills}{35.0}
\newcommand{\PhShareClimbSHillsBand}{[34.5, 35.5]}
\newcommand{\PhShareClimbSLowMtn}{28.1}
\newcommand{\PhShareClimbSLowMtnBand}{[27.7, 28.5]}
\newcommand{\PhShareClimbSPlains}{37.9}
\newcommand{\PhShareClimbSPlainsBand}{[37.5, 38.4]}
\newcommand{\PhShareExitAAll}{2.94}
\newcommand{\PhShareExitAAllBand}{[2.86, 3.01]}
\newcommand{\PhShareExitAHighMtn}{2.65}
\newcommand{\PhShareExitAHighMtnBand}{[2.54, 2.76]}
\newcommand{\PhShareExitAHills}{3.14}
\newcommand{\PhShareExitAHillsBand}{[2.95, 3.33]}
\newcommand{\PhShareExitALowMtn}{2.79}
\newcommand{\PhShareExitALowMtnBand}{[2.69, 2.88]}
\newcommand{\PhShareExitAPlains}{3.92}
\newcommand{\PhShareExitAPlainsBand}{[3.72, 4.12]}
\newcommand{\PhShareExitASAll}{$-$0.353}
\newcommand{\PhShareExitASAllBand}{[$-$0.428, $-$0.278]}
\newcommand{\PhShareExitASHighMtn}{$-$0.174}
\newcommand{\PhShareExitASHighMtnBand}{[$-$0.294, $-$0.054]}
\newcommand{\PhShareExitASHills}{$-$0.311}
\newcommand{\PhShareExitASHillsBand}{[$-$0.508, $-$0.114]}
\newcommand{\PhShareExitASLowMtn}{$-$0.212}
\newcommand{\PhShareExitASLowMtnBand}{[$-$0.314, $-$0.110]}
\newcommand{\PhShareExitASPlains}{$-$0.413}
\newcommand{\PhShareExitASPlainsBand}{[$-$0.622, $-$0.203]}
\newcommand{\PhShareExitGAAll}{$+$0.286}
\newcommand{\PhShareExitGAAllBand}{[$+$0.207, $+$0.366]}
\newcommand{\PhShareExitGAHighMtn}{$+$0.152}
\newcommand{\PhShareExitGAHighMtnBand}{[$+$0.010, $+$0.295]}
\newcommand{\PhShareExitGAHills}{$+$0.513}
\newcommand{\PhShareExitGAHillsBand}{[$+$0.282, $+$0.743]}
\newcommand{\PhShareExitGALowMtn}{$+$0.173}
\newcommand{\PhShareExitGALowMtnBand}{[$+$0.060, $+$0.286]}
\newcommand{\PhShareExitGAPlains}{$+$0.461}
\newcommand{\PhShareExitGAPlainsBand}{[$+$0.225, $+$0.697]}
\newcommand{\PhShareExitGAll}{3.22}
\newcommand{\PhShareExitGAllBand}{[3.14, 3.31]}
\newcommand{\PhShareExitGHighMtn}{2.80}
\newcommand{\PhShareExitGHighMtnBand}{[2.67, 2.93]}
\newcommand{\PhShareExitGHills}{3.65}
\newcommand{\PhShareExitGHillsBand}{[3.39, 3.91]}
\newcommand{\PhShareExitGLowMtn}{2.96}
\newcommand{\PhShareExitGLowMtnBand}{[2.85, 3.07]}
\newcommand{\PhShareExitGPlains}{4.38}
\newcommand{\PhShareExitGPlainsBand}{[4.16, 4.60]}
\newcommand{\PhShareExitGSAll}{$-$0.0668}
\newcommand{\PhShareExitGSAllBand}{[$-$0.1525, $+$0.0189]}
\newcommand{\PhShareExitGSHighMtn}{$-$0.0218}
\newcommand{\PhShareExitGSHighMtnBand}{[$-$0.1633, $+$0.1198]}
\newcommand{\PhShareExitGSHills}{$+$0.202}
\newcommand{\PhShareExitGSHillsBand}{[$-$0.049, $+$0.452]}
\newcommand{\PhShareExitGSLowMtn}{$-$0.0391}
\newcommand{\PhShareExitGSLowMtnBand}{[$-$0.1504, $+$0.0721]}
\newcommand{\PhShareExitGSPlains}{$+$0.0482}
\newcommand{\PhShareExitGSPlainsBand}{[$-$0.1763, $+$0.2727]}
\newcommand{\PhShareExitSAll}{3.29}
\newcommand{\PhShareExitSAllBand}{[3.23, 3.35]}
\newcommand{\PhShareExitSHighMtn}{2.82}
\newcommand{\PhShareExitSHighMtnBand}{[2.75, 2.89]}
\newcommand{\PhShareExitSHills}{3.45}
\newcommand{\PhShareExitSHillsBand}{[3.35, 3.55]}
\newcommand{\PhShareExitSLowMtn}{3.00}
\newcommand{\PhShareExitSLowMtnBand}{[2.94, 3.06]}
\newcommand{\PhShareExitSPlains}{4.33}
\newcommand{\PhShareExitSPlainsBand}{[4.20, 4.46]}
\newcommand{\PhShareGlideAAll}{56.9}
\newcommand{\PhShareGlideAAllBand}{[56.5, 57.3]}
\newcommand{\PhShareGlideAHighMtn}{60.0}
\newcommand{\PhShareGlideAHighMtnBand}{[59.5, 60.6]}
\newcommand{\PhShareGlideAHills}{50.9}
\newcommand{\PhShareGlideAHillsBand}{[50.0, 51.7]}
\newcommand{\PhShareGlideALowMtn}{59.1}
\newcommand{\PhShareGlideALowMtnBand}{[58.6, 59.5]}
\newcommand{\PhShareGlideAPlains}{48.0}
\newcommand{\PhShareGlideAPlainsBand}{[47.2, 48.8]}
\newcommand{\PhShareGlideASAll}{$+$1.50}
\newcommand{\PhShareGlideASAllBand}{[$+$1.14, $+$1.87]}
\newcommand{\PhShareGlideASHighMtn}{$+$1.39}
\newcommand{\PhShareGlideASHighMtnBand}{[$+$0.87, $+$1.91]}
\newcommand{\PhShareGlideASHills}{$-$2.11}
\newcommand{\PhShareGlideASHillsBand}{[$-$2.97, $-$1.25]}
\newcommand{\PhShareGlideASLowMtn}{$+$0.863}
\newcommand{\PhShareGlideASLowMtnBand}{[$+$0.396, $+$1.331]}
\newcommand{\PhShareGlideASPlains}{$-$0.821}
\newcommand{\PhShareGlideASPlainsBand}{[$-$1.596, $-$0.046]}
\newcommand{\PhShareGlideGAAll}{$-$4.34}
\newcommand{\PhShareGlideGAAllBand}{[$-$4.65, $-$4.03]}
\newcommand{\PhShareGlideGAHighMtn}{$-$4.75}
\newcommand{\PhShareGlideGAHighMtnBand}{[$-$5.38, $-$4.11]}
\newcommand{\PhShareGlideGAHills}{$-$4.06}
\newcommand{\PhShareGlideGAHillsBand}{[$-$4.92, $-$3.21]}
\newcommand{\PhShareGlideGALowMtn}{$-$3.76}
\newcommand{\PhShareGlideGALowMtnBand}{[$-$4.22, $-$3.30]}
\newcommand{\PhShareGlideGAPlains}{$-$3.62}
\newcommand{\PhShareGlideGAPlainsBand}{[$-$4.41, $-$2.83]}
\newcommand{\PhShareGlideGAll}{52.6}
\newcommand{\PhShareGlideGAllBand}{[52.1, 53.0]}
\newcommand{\PhShareGlideGHighMtn}{55.3}
\newcommand{\PhShareGlideGHighMtnBand}{[54.7, 55.9]}
\newcommand{\PhShareGlideGHills}{46.8}
\newcommand{\PhShareGlideGHillsBand}{[45.9, 47.7]}
\newcommand{\PhShareGlideGLowMtn}{55.3}
\newcommand{\PhShareGlideGLowMtnBand}{[54.7, 55.8]}
\newcommand{\PhShareGlideGPlains}{44.4}
\newcommand{\PhShareGlideGPlainsBand}{[43.5, 45.2]}
\newcommand{\PhShareGlideGSAll}{$-$2.83}
\newcommand{\PhShareGlideGSAllBand}{[$-$3.26, $-$2.41]}
\newcommand{\PhShareGlideGSHighMtn}{$-$3.35}
\newcommand{\PhShareGlideGSHighMtnBand}{[$-$3.98, $-$2.73]}
\newcommand{\PhShareGlideGSHills}{$-$6.17}
\newcommand{\PhShareGlideGSHillsBand}{[$-$7.06, $-$5.29]}
\newcommand{\PhShareGlideGSLowMtn}{$-$2.90}
\newcommand{\PhShareGlideGSLowMtnBand}{[$-$3.44, $-$2.35]}
\newcommand{\PhShareGlideGSPlains}{$-$4.44}
\newcommand{\PhShareGlideGSPlainsBand}{[$-$5.25, $-$3.63]}
\newcommand{\PhShareGlideSAll}{55.4}
\newcommand{\PhShareGlideSAllBand}{[55.1, 55.7]}
\newcommand{\PhShareGlideSHighMtn}{58.7}
\newcommand{\PhShareGlideSHighMtnBand}{[58.3, 59.0]}
\newcommand{\PhShareGlideSHills}{53.0}
\newcommand{\PhShareGlideSHillsBand}{[52.5, 53.5]}
\newcommand{\PhShareGlideSLowMtn}{58.2}
\newcommand{\PhShareGlideSLowMtnBand}{[57.8, 58.6]}
\newcommand{\PhShareGlideSPlains}{48.8}
\newcommand{\PhShareGlideSPlainsBand}{[48.3, 49.3]}
\newcommand{\PhShareSearchAAll}{9.68}
\newcommand{\PhShareSearchAAllBand}{[9.47, 9.89]}
\newcommand{\PhShareSearchAHighMtn}{9.83}
\newcommand{\PhShareSearchAHighMtnBand}{[9.50, 10.17]}
\newcommand{\PhShareSearchAHills}{7.21}
\newcommand{\PhShareSearchAHillsBand}{[6.76, 7.67]}
\newcommand{\PhShareSearchALowMtn}{10.7}
\newcommand{\PhShareSearchALowMtnBand}{[10.4, 11.0]}
\newcommand{\PhShareSearchAPlains}{8.12}
\newcommand{\PhShareSearchAPlainsBand}{[7.72, 8.52]}
\newcommand{\PhShareSearchASAll}{$-$0.0151}
\newcommand{\PhShareSearchASAllBand}{[$-$0.2309, $+$0.2006]}
\newcommand{\PhShareSearchASHighMtn}{$-$0.764}
\newcommand{\PhShareSearchASHighMtnBand}{[$-$1.116, $-$0.412]}
\newcommand{\PhShareSearchASHills}{$-$1.39}
\newcommand{\PhShareSearchASHillsBand}{[$-$1.85, $-$0.93]}
\newcommand{\PhShareSearchASLowMtn}{$+$0.0332}
\newcommand{\PhShareSearchASLowMtnBand}{[$-$0.2737, $+$0.3400]}
\newcommand{\PhShareSearchASPlains}{$-$0.835}
\newcommand{\PhShareSearchASPlainsBand}{[$-$1.231, $-$0.439]}
\newcommand{\PhShareSearchGAAll}{$-$0.399}
\newcommand{\PhShareSearchGAAllBand}{[$-$0.573, $-$0.225]}
\newcommand{\PhShareSearchGAHighMtn}{$-$0.460}
\newcommand{\PhShareSearchGAHighMtnBand}{[$-$0.781, $-$0.139]}
\newcommand{\PhShareSearchGAHills}{$+$0.114}
\newcommand{\PhShareSearchGAHillsBand}{[$-$0.285, $+$0.514]}
\newcommand{\PhShareSearchGALowMtn}{$-$0.409}
\newcommand{\PhShareSearchGALowMtnBand}{[$-$0.700, $-$0.118]}
\newcommand{\PhShareSearchGAPlains}{$-$0.259}
\newcommand{\PhShareSearchGAPlainsBand}{[$-$0.599, $+$0.081]}
\newcommand{\PhShareSearchGAll}{9.28}
\newcommand{\PhShareSearchGAllBand}{[9.04, 9.52]}
\newcommand{\PhShareSearchGHighMtn}{9.37}
\newcommand{\PhShareSearchGHighMtnBand}{[9.00, 9.75]}
\newcommand{\PhShareSearchGHills}{7.33}
\newcommand{\PhShareSearchGHillsBand}{[6.89, 7.76]}
\newcommand{\PhShareSearchGLowMtn}{10.3}
\newcommand{\PhShareSearchGLowMtnBand}{[10.0, 10.7]}
\newcommand{\PhShareSearchGPlains}{7.86}
\newcommand{\PhShareSearchGPlainsBand}{[7.47, 8.25]}
\newcommand{\PhShareSearchGSAll}{$-$0.414}
\newcommand{\PhShareSearchGSAllBand}{[$-$0.651, $-$0.176]}
\newcommand{\PhShareSearchGSHighMtn}{$-$1.22}
\newcommand{\PhShareSearchGSHighMtnBand}{[$-$1.60, $-$0.84]}
\newcommand{\PhShareSearchGSHills}{$-$1.28}
\newcommand{\PhShareSearchGSHillsBand}{[$-$1.73, $-$0.83]}
\newcommand{\PhShareSearchGSLowMtn}{$-$0.376}
\newcommand{\PhShareSearchGSLowMtnBand}{[$-$0.732, $-$0.020]}
\newcommand{\PhShareSearchGSPlains}{$-$1.09}
\newcommand{\PhShareSearchGSPlainsBand}{[$-$1.48, $-$0.71]}
\newcommand{\PhShareSearchSAll}{9.69}
\newcommand{\PhShareSearchSAllBand}{[9.54, 9.85]}
\newcommand{\PhShareSearchSHighMtn}{10.6}
\newcommand{\PhShareSearchSHighMtnBand}{[10.4, 10.8]}
\newcommand{\PhShareSearchSHills}{8.60}
\newcommand{\PhShareSearchSHillsBand}{[8.33, 8.87]}
\newcommand{\PhShareSearchSLowMtn}{10.7}
\newcommand{\PhShareSearchSLowMtnBand}{[10.5, 10.9]}
\newcommand{\PhShareSearchSPlains}{8.95}
\newcommand{\PhShareSearchSPlainsBand}{[8.67, 9.23]}
\newcommand{\PhSignAgree}{203}
\newcommand{\PhSignTotal}{230}
\newcommand{\PhSinkAAll}{96.0}
\newcommand{\PhSinkAAllBand}{[95.0, 97.0]}
\newcommand{\PhSinkAHighMtn}{99.5}
\newcommand{\PhSinkAHighMtnBand}{[98.1, 101.0]}
\newcommand{\PhSinkAHills}{102}
\newcommand{\PhSinkAHillsBand}{[100, 105]}
\newcommand{\PhSinkALowMtn}{94.3}
\newcommand{\PhSinkALowMtnBand}{[93.1, 95.5]}
\newcommand{\PhSinkAPlains}{90.0}
\newcommand{\PhSinkAPlainsBand}{[87.2, 92.8]}
\newcommand{\PhSinkASAll}{$-$7.08}
\newcommand{\PhSinkASAllBand}{[$-$8.02, $-$6.14]}
\newcommand{\PhSinkASHighMtn}{$-$4.34}
\newcommand{\PhSinkASHighMtnBand}{[$-$5.81, $-$2.86]}
\newcommand{\PhSinkASHills}{$-$6.03}
\newcommand{\PhSinkASHillsBand}{[$-$8.44, $-$3.61]}
\newcommand{\PhSinkASLowMtn}{$-$7.40}
\newcommand{\PhSinkASLowMtnBand}{[$-$8.63, $-$6.18]}
\newcommand{\PhSinkASPlains}{$-$8.53}
\newcommand{\PhSinkASPlainsBand}{[$-$10.87, $-$6.20]}
\newcommand{\PhSinkGAAll}{$-$11.6}
\newcommand{\PhSinkGAAllBand}{[$-$12.3, $-$10.9]}
\newcommand{\PhSinkGAHighMtn}{$-$12.7}
\newcommand{\PhSinkGAHighMtnBand}{[$-$14.1, $-$11.3]}
\newcommand{\PhSinkGAHills}{$-$5.89}
\newcommand{\PhSinkGAHillsBand}{[$-$7.78, $-$4.00]}
\newcommand{\PhSinkGALowMtn}{$-$14.1}
\newcommand{\PhSinkGALowMtnBand}{[$-$15.2, $-$13.0]}
\newcommand{\PhSinkGAPlains}{$-$6.60}
\newcommand{\PhSinkGAPlainsBand}{[$-$8.00, $-$5.21]}
\newcommand{\PhSinkGAll}{84.4}
\newcommand{\PhSinkGAllBand}{[83.3, 85.6]}
\newcommand{\PhSinkGHighMtn}{86.8}
\newcommand{\PhSinkGHighMtnBand}{[85.4, 88.3]}
\newcommand{\PhSinkGHills}{96.3}
\newcommand{\PhSinkGHillsBand}{[94.2, 98.4]}
\newcommand{\PhSinkGLowMtn}{80.2}
\newcommand{\PhSinkGLowMtnBand}{[78.7, 81.7]}
\newcommand{\PhSinkGPlains}{83.4}
\newcommand{\PhSinkGPlainsBand}{[80.8, 86.0]}
\newcommand{\PhSinkGSAll}{$-$18.7}
\newcommand{\PhSinkGSAllBand}{[$-$19.8, $-$17.5]}
\newcommand{\PhSinkGSHighMtn}{$-$17.0}
\newcommand{\PhSinkGSHighMtnBand}{[$-$18.6, $-$15.5]}
\newcommand{\PhSinkGSHills}{$-$11.9}
\newcommand{\PhSinkGSHillsBand}{[$-$14.0, $-$9.9]}
\newcommand{\PhSinkGSLowMtn}{$-$21.5}
\newcommand{\PhSinkGSLowMtnBand}{[$-$22.9, $-$20.0]}
\newcommand{\PhSinkGSPlains}{$-$15.1}
\newcommand{\PhSinkGSPlainsBand}{[$-$17.4, $-$12.9]}
\newcommand{\PhSinkSAll}{103}
\newcommand{\PhSinkSAllBand}{[102, 104]}
\newcommand{\PhSinkSHighMtn}{104}
\newcommand{\PhSinkSHighMtnBand}{[103, 105]}
\newcommand{\PhSinkSHills}{108}
\newcommand{\PhSinkSHillsBand}{[107, 109]}
\newcommand{\PhSinkSLowMtn}{102}
\newcommand{\PhSinkSLowMtnBand}{[101, 103]}
\newcommand{\PhSinkSPlains}{98.6}
\newcommand{\PhSinkSPlainsBand}{[97.1, 100.0]}
\newcommand{\PhSpeedAAll}{33.4}
\newcommand{\PhSpeedAAllBand}{[33.1, 33.8]}
\newcommand{\PhSpeedAHighMtn}{31.6}
\newcommand{\PhSpeedAHighMtnBand}{[31.3, 32.0]}
\newcommand{\PhSpeedAHills}{37.6}
\newcommand{\PhSpeedAHillsBand}{[36.8, 38.5]}
\newcommand{\PhSpeedALowMtn}{31.3}
\newcommand{\PhSpeedALowMtnBand}{[31.0, 31.6]}
\newcommand{\PhSpeedAPlains}{43.2}
\newcommand{\PhSpeedAPlainsBand}{[42.4, 44.1]}
\newcommand{\PhSpeedASAll}{$+$0.517}
\newcommand{\PhSpeedASAllBand}{[$+$0.211, $+$0.823]}
\newcommand{\PhSpeedASHighMtn}{$+$1.14}
\newcommand{\PhSpeedASHighMtnBand}{[$+$0.82, $+$1.45]}
\newcommand{\PhSpeedASHills}{$+$3.02}
\newcommand{\PhSpeedASHillsBand}{[$+$2.19, $+$3.85]}
\newcommand{\PhSpeedASLowMtn}{$+$1.09}
\newcommand{\PhSpeedASLowMtnBand}{[$+$0.80, $+$1.39]}
\newcommand{\PhSpeedASPlains}{$+$2.21}
\newcommand{\PhSpeedASPlainsBand}{[$+$1.41, $+$3.01]}
\newcommand{\PhSpeedGAAll}{$+$1.45}
\newcommand{\PhSpeedGAAllBand}{[$+$1.28, $+$1.61]}
\newcommand{\PhSpeedGAHighMtn}{$+$1.44}
\newcommand{\PhSpeedGAHighMtnBand}{[$+$1.17, $+$1.71]}
\newcommand{\PhSpeedGAHills}{$+$1.65}
\newcommand{\PhSpeedGAHillsBand}{[$+$1.09, $+$2.21]}
\newcommand{\PhSpeedGALowMtn}{$+$0.985}
\newcommand{\PhSpeedGALowMtnBand}{[$+$0.779, $+$1.192]}
\newcommand{\PhSpeedGAPlains}{$+$0.878}
\newcommand{\PhSpeedGAPlainsBand}{[$+$0.461, $+$1.295]}
\newcommand{\PhSpeedGAll}{34.9}
\newcommand{\PhSpeedGAllBand}{[34.5, 35.3]}
\newcommand{\PhSpeedGHighMtn}{33.1}
\newcommand{\PhSpeedGHighMtnBand}{[32.7, 33.4]}
\newcommand{\PhSpeedGHills}{39.3}
\newcommand{\PhSpeedGHillsBand}{[38.5, 40.1]}
\newcommand{\PhSpeedGLowMtn}{32.3}
\newcommand{\PhSpeedGLowMtnBand}{[32.0, 32.6]}
\newcommand{\PhSpeedGPlains}{44.1}
\newcommand{\PhSpeedGPlainsBand}{[43.3, 45.0]}
\newcommand{\PhSpeedGSAll}{$+$1.96}
\newcommand{\PhSpeedGSAllBand}{[$+$1.62, $+$2.31]}
\newcommand{\PhSpeedGSHighMtn}{$+$2.58}
\newcommand{\PhSpeedGSHighMtnBand}{[$+$2.24, $+$2.91]}
\newcommand{\PhSpeedGSHills}{$+$4.67}
\newcommand{\PhSpeedGSHillsBand}{[$+$3.92, $+$5.41]}
\newcommand{\PhSpeedGSLowMtn}{$+$2.08}
\newcommand{\PhSpeedGSLowMtnBand}{[$+$1.77, $+$2.39]}
\newcommand{\PhSpeedGSPlains}{$+$3.09}
\newcommand{\PhSpeedGSPlainsBand}{[$+$2.33, $+$3.84]}
\newcommand{\PhSpeedSAll}{32.9}
\newcommand{\PhSpeedSAllBand}{[32.7, 33.2]}
\newcommand{\PhSpeedSHighMtn}{30.5}
\newcommand{\PhSpeedSHighMtnBand}{[30.3, 30.7]}
\newcommand{\PhSpeedSHills}{34.6}
\newcommand{\PhSpeedSHillsBand}{[34.2, 35.1]}
\newcommand{\PhSpeedSLowMtn}{30.2}
\newcommand{\PhSpeedSLowMtnBand}{[30.0, 30.5]}
\newcommand{\PhSpeedSPlains}{41.0}
\newcommand{\PhSpeedSPlainsBand}{[40.6, 41.5]}
\newcommand{\PhVChordASAll}{$-$29.4}
\newcommand{\PhVChordASAllBand}{[$-$45.0, $-$13.8]}
\newcommand{\PhVChordASHighMtn}{$-$0.493}
\newcommand{\PhVChordASHighMtnBand}{[$-$8.567, $+$7.582]}
\newcommand{\PhVChordASHills}{$+$94.7}
\newcommand{\PhVChordASHillsBand}{[$+$62.6, $+$126.7]}
\newcommand{\PhVChordASLowMtn}{$-$15.8}
\newcommand{\PhVChordASLowMtnBand}{[$-$28.0, $-$3.7]}
\newcommand{\PhVChordASPlains}{$+$88.4}
\newcommand{\PhVChordASPlainsBand}{[$+$53.8, $+$122.9]}
\newcommand{\PhVChordGAAll}{$+$6.09}
\newcommand{\PhVChordGAAllBand}{[$-$1.68, $+$13.87]}
\newcommand{\PhVChordGAHighMtn}{$-$6.44}
\newcommand{\PhVChordGAHighMtnBand}{[$-$12.89, $+$0.01]}
\newcommand{\PhVChordGAHills}{$+$19.8}
\newcommand{\PhVChordGAHillsBand}{[$-$4.9, $+$44.5]}
\newcommand{\PhVChordGALowMtn}{$-$9.53}
\newcommand{\PhVChordGALowMtnBand}{[$-$16.25, $-$2.81]}
\newcommand{\PhVChordGAPlains}{$-$37.5}
\newcommand{\PhVChordGAPlainsBand}{[$-$62.4, $-$12.6]}
\newcommand{\PhVClimbMeanASAll}{$-$2.16}
\newcommand{\PhVClimbMeanASAllBand}{[$-$4.32, $+$0.00]}
\newcommand{\PhVClimbMeanASHighMtn}{$-$0.510}
\newcommand{\PhVClimbMeanASHighMtnBand}{[$-$3.201, $+$2.181]}
\newcommand{\PhVClimbMeanASHills}{$+$13.1}
\newcommand{\PhVClimbMeanASHillsBand}{[$+$7.1, $+$19.1]}
\newcommand{\PhVClimbMeanASLowMtn}{$-$0.0844}
\newcommand{\PhVClimbMeanASLowMtnBand}{[$-$2.7726, $+$2.6037]}
\newcommand{\PhVClimbMeanASPlains}{$+$10.2}
\newcommand{\PhVClimbMeanASPlainsBand}{[$+$5.8, $+$14.7]}
\newcommand{\PhVClimbMeanGAAll}{$+$1.89}
\newcommand{\PhVClimbMeanGAAllBand}{[$+$0.24, $+$3.55]}
\newcommand{\PhVClimbMeanGAHighMtn}{$+$5.24}
\newcommand{\PhVClimbMeanGAHighMtnBand}{[$+$2.49, $+$7.99]}
\newcommand{\PhVClimbMeanGAHills}{$+$3.06}
\newcommand{\PhVClimbMeanGAHillsBand}{[$-$2.46, $+$8.58]}
\newcommand{\PhVClimbMeanGALowMtn}{$+$0.394}
\newcommand{\PhVClimbMeanGALowMtnBand}{[$-$1.652, $+$2.441]}
\newcommand{\PhVClimbMeanGAPlains}{$-$7.52}
\newcommand{\PhVClimbMeanGAPlainsBand}{[$-$11.84, $-$3.20]}
\newcommand{\PhVCphASAll}{$-$0.205}
\newcommand{\PhVCphASAllBand}{[$-$0.257, $-$0.154]}
\newcommand{\PhVCphASHighMtn}{$-$0.108}
\newcommand{\PhVCphASHighMtnBand}{[$-$0.211, $-$0.004]}
\newcommand{\PhVCphASHills}{$+$0.167}
\newcommand{\PhVCphASHillsBand}{[$+$0.013, $+$0.321]}
\newcommand{\PhVCphASLowMtn}{$-$0.255}
\newcommand{\PhVCphASLowMtnBand}{[$-$0.330, $-$0.180]}
\newcommand{\PhVCphASPlains}{$-$0.0580}
\newcommand{\PhVCphASPlainsBand}{[$-$0.1732, $+$0.0572]}
\newcommand{\PhVCphGAAll}{$+$0.853}
\newcommand{\PhVCphGAAllBand}{[$+$0.796, $+$0.910]}
\newcommand{\PhVCphGAHighMtn}{$+$0.875}
\newcommand{\PhVCphGAHighMtnBand}{[$+$0.752, $+$0.998]}
\newcommand{\PhVCphGAHills}{$+$0.479}
\newcommand{\PhVCphGAHillsBand}{[$+$0.316, $+$0.642]}
\newcommand{\PhVCphGALowMtn}{$+$0.925}
\newcommand{\PhVCphGALowMtnBand}{[$+$0.840, $+$1.010]}
\newcommand{\PhVCphGAPlains}{$+$0.803}
\newcommand{\PhVCphGAPlainsBand}{[$+$0.668, $+$0.938]}
\newcommand{\PhVCtxExit}{26.7}
\newcommand{\PhVCtxExitAll}{26.7}
\newcommand{\PhVCtxExitHighMtn}{23.9}
\newcommand{\PhVCtxExitHills}{30.6}
\newcommand{\PhVCtxExitLowMtn}{24.3}
\newcommand{\PhVCtxExitPlains}{33.0}
\newcommand{\PhVCtxFailed}{26.0}
\newcommand{\PhVCtxFailedAll}{26.0}
\newcommand{\PhVCtxFailedHighMtn}{26.7}
\newcommand{\PhVCtxFailedHills}{23.4}
\newcommand{\PhVCtxFailedLowMtn}{28.3}
\newcommand{\PhVCtxFailedPlains}{21.1}
\newcommand{\PhVCtxFound}{22.9}
\newcommand{\PhVCtxFoundAll}{22.9}
\newcommand{\PhVCtxFoundHighMtn}{24.4}
\newcommand{\PhVCtxFoundHills}{21.0}
\newcommand{\PhVCtxFoundLowMtn}{23.8}
\newcommand{\PhVCtxFoundPlains}{20.3}
\newcommand{\PhVCtxRefound}{21.1}
\newcommand{\PhVCtxRefoundAll}{21.1}
\newcommand{\PhVCtxRefoundHighMtn}{21.9}
\newcommand{\PhVCtxRefoundHills}{20.8}
\newcommand{\PhVCtxRefoundLowMtn}{20.4}
\newcommand{\PhVCtxRefoundPlains}{22.8}
\newcommand{\PhVCycleASAll}{$+$0.294}
\newcommand{\PhVCycleASAllBand}{[$+$0.219, $+$0.369]}
\newcommand{\PhVCycleASHighMtn}{$+$0.163}
\newcommand{\PhVCycleASHighMtnBand}{[$+$0.005, $+$0.320]}
\newcommand{\PhVCycleASHills}{$-$0.223}
\newcommand{\PhVCycleASHillsBand}{[$-$0.425, $-$0.022]}
\newcommand{\PhVCycleASLowMtn}{$+$0.383}
\newcommand{\PhVCycleASLowMtnBand}{[$+$0.267, $+$0.498]}
\newcommand{\PhVCycleASPlains}{$+$0.0756}
\newcommand{\PhVCycleASPlainsBand}{[$-$0.0753, $+$0.2265]}
\newcommand{\PhVCycleGAAll}{$-$1.11}
\newcommand{\PhVCycleGAAllBand}{[$-$1.19, $-$1.04]}
\newcommand{\PhVCycleGAHighMtn}{$-$1.18}
\newcommand{\PhVCycleGAHighMtnBand}{[$-$1.35, $-$1.01]}
\newcommand{\PhVCycleGAHills}{$-$0.583}
\newcommand{\PhVCycleGAHillsBand}{[$-$0.786, $-$0.379]}
\newcommand{\PhVCycleGALowMtn}{$-$1.26}
\newcommand{\PhVCycleGALowMtnBand}{[$-$1.37, $-$1.14]}
\newcommand{\PhVCycleGAPlains}{$-$0.944}
\newcommand{\PhVCycleGAPlainsBand}{[$-$1.105, $-$0.782]}
\newcommand{\PhVDriftASAll}{$-$0.130}
\newcommand{\PhVDriftASAllBand}{[$-$0.197, $-$0.064]}
\newcommand{\PhVDriftASHighMtn}{$+$0.0041}
\newcommand{\PhVDriftASHighMtnBand}{[$-$0.0375, $+$0.0457]}
\newcommand{\PhVDriftASHills}{$+$0.240}
\newcommand{\PhVDriftASHillsBand}{[$+$0.119, $+$0.361]}
\newcommand{\PhVDriftASLowMtn}{$-$0.0981}
\newcommand{\PhVDriftASLowMtnBand}{[$-$0.1495, $-$0.0468]}
\newcommand{\PhVDriftASPlains}{$+$0.193}
\newcommand{\PhVDriftASPlainsBand}{[$+$0.059, $+$0.326]}
\newcommand{\PhVDriftGAAll}{$+$0.0021}
\newcommand{\PhVDriftGAAllBand}{[$-$0.0288, $+$0.0330]}
\newcommand{\PhVDriftGAHighMtn}{$-$0.112}
\newcommand{\PhVDriftGAHighMtnBand}{[$-$0.144, $-$0.080]}
\newcommand{\PhVDriftGAHills}{$+$0.0409}
\newcommand{\PhVDriftGAHillsBand}{[$-$0.0341, $+$0.1159]}
\newcommand{\PhVDriftGALowMtn}{$-$0.0658}
\newcommand{\PhVDriftGALowMtnBand}{[$-$0.0968, $-$0.0349]}
\newcommand{\PhVDriftGAPlains}{$-$0.0087}
\newcommand{\PhVDriftGAPlainsBand}{[$-$0.0723, $+$0.0549]}
\newcommand{\PhVDurMeanClimbASAll}{$-$4.20}
\newcommand{\PhVDurMeanClimbASAllBand}{[$-$6.25, $-$2.16]}
\newcommand{\PhVDurMeanClimbASHighMtn}{$+$1.37}
\newcommand{\PhVDurMeanClimbASHighMtnBand}{[$-$0.90, $+$3.64]}
\newcommand{\PhVDurMeanClimbASHills}{$+$9.60}
\newcommand{\PhVDurMeanClimbASHillsBand}{[$+$4.30, $+$14.91]}
\newcommand{\PhVDurMeanClimbASLowMtn}{$-$1.36}
\newcommand{\PhVDurMeanClimbASLowMtnBand}{[$-$3.83, $+$1.11]}
\newcommand{\PhVDurMeanClimbASPlains}{$+$5.97}
\newcommand{\PhVDurMeanClimbASPlainsBand}{[$+$2.07, $+$9.86]}
\newcommand{\PhVDurMeanClimbGAAll}{$+$3.47}
\newcommand{\PhVDurMeanClimbGAAllBand}{[$+$1.25, $+$5.69]}
\newcommand{\PhVDurMeanClimbGAHighMtn}{$+$1.08}
\newcommand{\PhVDurMeanClimbGAHighMtnBand}{[$-$1.48, $+$3.64]}
\newcommand{\PhVDurMeanClimbGAHills}{$+$3.28}
\newcommand{\PhVDurMeanClimbGAHillsBand}{[$-$1.89, $+$8.44]}
\newcommand{\PhVDurMeanClimbGALowMtn}{$-$1.16}
\newcommand{\PhVDurMeanClimbGALowMtnBand}{[$-$3.60, $+$1.27]}
\newcommand{\PhVDurMeanClimbGAPlains}{$-$9.14}
\newcommand{\PhVDurMeanClimbGAPlainsBand}{[$-$14.36, $-$3.92]}
\newcommand{\PhVDurMeanGlideASAll}{$+$28.5}
\newcommand{\PhVDurMeanGlideASAllBand}{[$+$24.9, $+$32.1]}
\newcommand{\PhVDurMeanGlideASHighMtn}{$+$18.9}
\newcommand{\PhVDurMeanGlideASHighMtnBand}{[$+$13.9, $+$23.8]}
\newcommand{\PhVDurMeanGlideASHills}{$+$8.36}
\newcommand{\PhVDurMeanGlideASHillsBand}{[$+$1.51, $+$15.22]}
\newcommand{\PhVDurMeanGlideASLowMtn}{$+$28.8}
\newcommand{\PhVDurMeanGlideASLowMtnBand}{[$+$23.8, $+$33.9]}
\newcommand{\PhVDurMeanGlideASPlains}{$+$11.4}
\newcommand{\PhVDurMeanGlideASPlainsBand}{[$+$4.7, $+$18.0]}
\newcommand{\PhVDurMeanGlideGAAll}{$-$33.9}
\newcommand{\PhVDurMeanGlideGAAllBand}{[$-$37.2, $-$30.5]}
\newcommand{\PhVDurMeanGlideGAHighMtn}{$-$24.9}
\newcommand{\PhVDurMeanGlideGAHighMtnBand}{[$-$31.4, $-$18.5]}
\newcommand{\PhVDurMeanGlideGAHills}{$-$20.8}
\newcommand{\PhVDurMeanGlideGAHillsBand}{[$-$28.6, $-$13.0]}
\newcommand{\PhVDurMeanGlideGALowMtn}{$-$30.7}
\newcommand{\PhVDurMeanGlideGALowMtnBand}{[$-$35.8, $-$25.6]}
\newcommand{\PhVDurMeanGlideGAPlains}{$-$32.3}
\newcommand{\PhVDurMeanGlideGAPlainsBand}{[$-$38.6, $-$26.0]}
\newcommand{\PhVDurMedianClimbASAll}{$-$2.00}
\newcommand{\PhVDurMedianClimbASAllBand}{[$-$3.76, $-$0.24]}
\newcommand{\PhVDurMedianClimbASHighMtn}{$+$2.00}
\newcommand{\PhVDurMedianClimbASHighMtnBand}{[$-$0.45, $+$4.45]}
\newcommand{\PhVDurMedianClimbASHills}{$+$10.0}
\newcommand{\PhVDurMedianClimbASHillsBand}{[$+$4.9, $+$15.1]}
\newcommand{\PhVDurMedianClimbASLowMtn}{$+$2.00}
\newcommand{\PhVDurMedianClimbASLowMtnBand}{[$-$0.22, $+$4.22]}
\newcommand{\PhVDurMedianClimbASPlains}{$+$4.00}
\newcommand{\PhVDurMedianClimbASPlainsBand}{[$-$0.08, $+$8.08]}
\newcommand{\PhVDurMedianClimbGAAll}{$+$5.00}
\newcommand{\PhVDurMedianClimbGAAllBand}{[$+$2.86, $+$7.14]}
\newcommand{\PhVDurMedianClimbGAHighMtn}{$+$2.00}
\newcommand{\PhVDurMedianClimbGAHighMtnBand}{[$-$0.84, $+$4.84]}
\newcommand{\PhVDurMedianClimbGAHills}{$+$8.00}
\newcommand{\PhVDurMedianClimbGAHillsBand}{[$+$2.54, $+$13.46]}
\newcommand{\PhVDurMedianClimbGALowMtn}{0}
\newcommand{\PhVDurMedianClimbGALowMtnBand}{[$-$2, $+$2]}
\newcommand{\PhVDurMedianClimbGAPlains}{$-$1.00}
\newcommand{\PhVDurMedianClimbGAPlainsBand}{[$-$6.61, $+$4.61]}
\newcommand{\PhVDurMedianGlideASAll}{$+$13.0}
\newcommand{\PhVDurMedianGlideASAllBand}{[$+$10.1, $+$15.9]}
\newcommand{\PhVDurMedianGlideASHighMtn}{$+$17.0}
\newcommand{\PhVDurMedianGlideASHighMtnBand}{[$+$12.3, $+$21.7]}
\newcommand{\PhVDurMedianGlideASHills}{$+$9.00}
\newcommand{\PhVDurMedianGlideASHillsBand}{[$+$2.65, $+$15.35]}
\newcommand{\PhVDurMedianGlideASLowMtn}{$+$12.0}
\newcommand{\PhVDurMedianGlideASLowMtnBand}{[$+$8.2, $+$15.8]}
\newcommand{\PhVDurMedianGlideASPlains}{$+$13.0}
\newcommand{\PhVDurMedianGlideASPlainsBand}{[$+$6.2, $+$19.8]}
\newcommand{\PhVDurMedianGlideGAAll}{$-$17.0}
\newcommand{\PhVDurMedianGlideGAAllBand}{[$-$19.9, $-$14.1]}
\newcommand{\PhVDurMedianGlideGAHighMtn}{$-$19.0}
\newcommand{\PhVDurMedianGlideGAHighMtnBand}{[$-$25.2, $-$12.8]}
\newcommand{\PhVDurMedianGlideGAHills}{$-$6.00}
\newcommand{\PhVDurMedianGlideGAHillsBand}{[$-$13.69, $+$1.69]}
\newcommand{\PhVDurMedianGlideGALowMtn}{$-$14.0}
\newcommand{\PhVDurMedianGlideGALowMtnBand}{[$-$18.1, $-$9.9]}
\newcommand{\PhVDurMedianGlideGAPlains}{$-$21.0}
\newcommand{\PhVDurMedianGlideGAPlainsBand}{[$-$27.1, $-$14.9]}
\newcommand{\PhVExitDurASAll}{0}
\newcommand{\PhVExitDurASAllBand}{[$-$1, $+$1]}
\newcommand{\PhVExitDurASHighMtn}{$-$1.00}
\newcommand{\PhVExitDurASHighMtnBand}{[$-$2.46, $+$0.46]}
\newcommand{\PhVExitDurASHills}{0}
\newcommand{\PhVExitDurASHillsBand}{[$-$2, $+$2]}
\newcommand{\PhVExitDurASLowMtn}{0}
\newcommand{\PhVExitDurASLowMtnBand}{[$-$1, $+$1]}
\newcommand{\PhVExitDurASPlains}{0}
\newcommand{\PhVExitDurASPlainsBand}{[$-$2, $+$2]}
\newcommand{\PhVExitDurGAAll}{$+$2.00}
\newcommand{\PhVExitDurGAAllBand}{[$+$1.00, $+$3.00]}
\newcommand{\PhVExitDurGAHighMtn}{$+$1.00}
\newcommand{\PhVExitDurGAHighMtnBand}{[$-$0.59, $+$2.59]}
\newcommand{\PhVExitDurGAHills}{$+$4.00}
\newcommand{\PhVExitDurGAHillsBand}{[$+$1.20, $+$6.80]}
\newcommand{\PhVExitDurGALowMtn}{$+$1.00}
\newcommand{\PhVExitDurGALowMtnBand}{[$-$0.29, $+$2.29]}
\newcommand{\PhVExitDurGAPlains}{$+$3.00}
\newcommand{\PhVExitDurGAPlainsBand}{[$+$1.09, $+$4.91]}
\newcommand{\PhVExitMeanASAll}{$+$2.01}
\newcommand{\PhVExitMeanASAllBand}{[$+$0.67, $+$3.35]}
\newcommand{\PhVExitMeanASHighMtn}{$+$0.0976}
\newcommand{\PhVExitMeanASHighMtnBand}{[$-$2.5147, $+$2.7098]}
\newcommand{\PhVExitMeanASHills}{$-$2.60}
\newcommand{\PhVExitMeanASHillsBand}{[$-$5.36, $+$0.17]}
\newcommand{\PhVExitMeanASLowMtn}{$+$2.15}
\newcommand{\PhVExitMeanASLowMtnBand}{[$-$0.02, $+$4.33]}
\newcommand{\PhVExitMeanASPlains}{$+$0.0201}
\newcommand{\PhVExitMeanASPlainsBand}{[$-$2.7739, $+$2.8140]}
\newcommand{\PhVExitMeanGAAll}{$-$1.58}
\newcommand{\PhVExitMeanGAAllBand}{[$-$3.03, $-$0.13]}
\newcommand{\PhVExitMeanGAHighMtn}{$-$3.76}
\newcommand{\PhVExitMeanGAHighMtnBand}{[$-$6.92, $-$0.60]}
\newcommand{\PhVExitMeanGAHills}{$+$3.32}
\newcommand{\PhVExitMeanGAHillsBand}{[$+$0.13, $+$6.51]}
\newcommand{\PhVExitMeanGALowMtn}{$-$3.42}
\newcommand{\PhVExitMeanGALowMtnBand}{[$-$5.70, $-$1.14]}
\newcommand{\PhVExitMeanGAPlains}{$+$2.26}
\newcommand{\PhVExitMeanGAPlainsBand}{[$-$0.87, $+$5.39]}
\newcommand{\PhVExitPerHourASAll}{$-$0.116}
\newcommand{\PhVExitPerHourASAllBand}{[$-$0.147, $-$0.085]}
\newcommand{\PhVExitPerHourASHighMtn}{$-$0.0340}
\newcommand{\PhVExitPerHourASHighMtnBand}{[$-$0.0808, $+$0.0129]}
\newcommand{\PhVExitPerHourASHills}{$-$0.0155}
\newcommand{\PhVExitPerHourASHillsBand}{[$-$0.1169, $+$0.0859]}
\newcommand{\PhVExitPerHourASLowMtn}{$-$0.109}
\newcommand{\PhVExitPerHourASLowMtnBand}{[$-$0.149, $-$0.069]}
\newcommand{\PhVExitPerHourASPlains}{$-$0.0737}
\newcommand{\PhVExitPerHourASPlainsBand}{[$-$0.1627, $+$0.0154]}
\newcommand{\PhVExitPerHourGAAll}{$+$0.114}
\newcommand{\PhVExitPerHourGAAllBand}{[$+$0.082, $+$0.145]}
\newcommand{\PhVExitPerHourGAHighMtn}{$+$0.0152}
\newcommand{\PhVExitPerHourGAHighMtnBand}{[$-$0.0374, $+$0.0678]}
\newcommand{\PhVExitPerHourGAHills}{$+$0.154}
\newcommand{\PhVExitPerHourGAHillsBand}{[$+$0.043, $+$0.265]}
\newcommand{\PhVExitPerHourGALowMtn}{$+$0.0851}
\newcommand{\PhVExitPerHourGALowMtnBand}{[$+$0.0430, $+$0.1273]}
\newcommand{\PhVExitPerHourGAPlains}{$+$0.219}
\newcommand{\PhVExitPerHourGAPlainsBand}{[$+$0.134, $+$0.304]}
\newcommand{\PhVFlightsAAll}{8\,161}
\newcommand{\PhVFlightsAHighMtn}{2\,038}
\newcommand{\PhVFlightsAHills}{837}
\newcommand{\PhVFlightsALowMtn}{3\,868}
\newcommand{\PhVFlightsAPlains}{1\,217}
\newcommand{\PhVFlightsGAll}{9\,818}
\newcommand{\PhVFlightsGHighMtn}{2\,273}
\newcommand{\PhVFlightsGHills}{1\,159}
\newcommand{\PhVFlightsGLowMtn}{4\,492}
\newcommand{\PhVFlightsGPlains}{1\,674}
\newcommand{\PhVFlightsSAll}{33\,933}
\newcommand{\PhVFlightsSHighMtn}{5\,590}
\newcommand{\PhVFlightsSHills}{6\,309}
\newcommand{\PhVFlightsSLowMtn}{12\,844}
\newcommand{\PhVFlightsSPlains}{6\,338}
\newcommand{\PhVFoundASAll}{$+$1.47}
\newcommand{\PhVFoundASAllBand}{[$+$0.81, $+$2.14]}
\newcommand{\PhVFoundASHighMtn}{$+$0.953}
\newcommand{\PhVFoundASHighMtnBand}{[$-$0.199, $+$2.106]}
\newcommand{\PhVFoundASHills}{$+$5.03}
\newcommand{\PhVFoundASHillsBand}{[$+$3.08, $+$6.97]}
\newcommand{\PhVFoundASLowMtn}{$+$1.24}
\newcommand{\PhVFoundASLowMtnBand}{[$+$0.25, $+$2.23]}
\newcommand{\PhVFoundASPlains}{$+$3.44}
\newcommand{\PhVFoundASPlainsBand}{[$+$1.52, $+$5.37]}
\newcommand{\PhVFoundGAAll}{$+$5.79}
\newcommand{\PhVFoundGAAllBand}{[$+$5.06, $+$6.53]}
\newcommand{\PhVFoundGAHighMtn}{$+$6.49}
\newcommand{\PhVFoundGAHighMtnBand}{[$+$5.13, $+$7.85]}
\newcommand{\PhVFoundGAHills}{$+$4.76}
\newcommand{\PhVFoundGAHillsBand}{[$+$2.78, $+$6.75]}
\newcommand{\PhVFoundGALowMtn}{$+$6.31}
\newcommand{\PhVFoundGALowMtnBand}{[$+$5.19, $+$7.43]}
\newcommand{\PhVFoundGAPlains}{$+$3.31}
\newcommand{\PhVFoundGAPlainsBand}{[$+$0.99, $+$5.63]}
\newcommand{\PhVGainASAll}{$+$8.35}
\newcommand{\PhVGainASAllBand}{[$+$5.20, $+$11.51]}
\newcommand{\PhVGainASHighMtn}{$+$3.08}
\newcommand{\PhVGainASHighMtnBand}{[$-$1.80, $+$7.97]}
\newcommand{\PhVGainASHills}{$+$17.9}
\newcommand{\PhVGainASHillsBand}{[$+$10.0, $+$25.8]}
\newcommand{\PhVGainASLowMtn}{$+$8.20}
\newcommand{\PhVGainASLowMtnBand}{[$+$3.79, $+$12.60]}
\newcommand{\PhVGainASPlains}{$+$8.86}
\newcommand{\PhVGainASPlainsBand}{[$+$3.18, $+$14.54]}
\newcommand{\PhVGainGAAll}{$+$1.10}
\newcommand{\PhVGainGAAllBand}{[$-$1.28, $+$3.48]}
\newcommand{\PhVGainGAHighMtn}{$+$9.98}
\newcommand{\PhVGainGAHighMtnBand}{[$+$5.59, $+$14.36]}
\newcommand{\PhVGainGAHills}{$+$2.92}
\newcommand{\PhVGainGAHillsBand}{[$-$4.66, $+$10.51]}
\newcommand{\PhVGainGALowMtn}{$+$0.906}
\newcommand{\PhVGainGALowMtnBand}{[$-$2.509, $+$4.321]}
\newcommand{\PhVGainGAPlains}{$-$12.7}
\newcommand{\PhVGainGAPlainsBand}{[$-$18.4, $-$7.0]}
\newcommand{\PhVGlideMeanASAll}{$+$20.5}
\newcommand{\PhVGlideMeanASAllBand}{[$+$17.0, $+$23.9]}
\newcommand{\PhVGlideMeanASHighMtn}{$+$14.8}
\newcommand{\PhVGlideMeanASHighMtnBand}{[$+$8.7, $+$20.9]}
\newcommand{\PhVGlideMeanASHills}{$+$0.809}
\newcommand{\PhVGlideMeanASHillsBand}{[$-$6.161, $+$7.778]}
\newcommand{\PhVGlideMeanASLowMtn}{$+$19.2}
\newcommand{\PhVGlideMeanASLowMtnBand}{[$+$14.5, $+$24.0]}
\newcommand{\PhVGlideMeanASPlains}{$+$9.09}
\newcommand{\PhVGlideMeanASPlainsBand}{[$+$2.53, $+$15.64]}
\newcommand{\PhVGlideMeanGAAll}{$-$28.6}
\newcommand{\PhVGlideMeanGAAllBand}{[$-$31.6, $-$25.6]}
\newcommand{\PhVGlideMeanGAHighMtn}{$-$29.7}
\newcommand{\PhVGlideMeanGAHighMtnBand}{[$-$36.3, $-$23.1]}
\newcommand{\PhVGlideMeanGAHills}{$-$17.6}
\newcommand{\PhVGlideMeanGAHillsBand}{[$-$25.0, $-$10.3]}
\newcommand{\PhVGlideMeanGALowMtn}{$-$25.9}
\newcommand{\PhVGlideMeanGALowMtnBand}{[$-$30.6, $-$21.2]}
\newcommand{\PhVGlideMeanGAPlains}{$-$30.6}
\newcommand{\PhVGlideMeanGAPlainsBand}{[$-$36.3, $-$24.9]}
\newcommand{\PhVGlidePerHourASAll}{$-$0.349}
\newcommand{\PhVGlidePerHourASAllBand}{[$-$0.403, $-$0.295]}
\newcommand{\PhVGlidePerHourASHighMtn}{$-$0.193}
\newcommand{\PhVGlidePerHourASHighMtnBand}{[$-$0.288, $-$0.098]}
\newcommand{\PhVGlidePerHourASHills}{$-$0.271}
\newcommand{\PhVGlidePerHourASHillsBand}{[$-$0.401, $-$0.141]}
\newcommand{\PhVGlidePerHourASLowMtn}{$-$0.332}
\newcommand{\PhVGlidePerHourASLowMtnBand}{[$-$0.407, $-$0.256]}
\newcommand{\PhVGlidePerHourASPlains}{$-$0.425}
\newcommand{\PhVGlidePerHourASPlainsBand}{[$-$0.563, $-$0.286]}
\newcommand{\PhVGlidePerHourGAAll}{$+$0.213}
\newcommand{\PhVGlidePerHourGAAllBand}{[$+$0.159, $+$0.267]}
\newcommand{\PhVGlidePerHourGAHighMtn}{$+$0.209}
\newcommand{\PhVGlidePerHourGAHighMtnBand}{[$+$0.099, $+$0.320]}
\newcommand{\PhVGlidePerHourGAHills}{$-$0.0060}
\newcommand{\PhVGlidePerHourGAHillsBand}{[$-$0.1473, $+$0.1354]}
\newcommand{\PhVGlidePerHourGALowMtn}{$+$0.188}
\newcommand{\PhVGlidePerHourGALowMtnBand}{[$+$0.109, $+$0.268]}
\newcommand{\PhVGlidePerHourGAPlains}{$+$0.394}
\newcommand{\PhVGlidePerHourGAPlainsBand}{[$+$0.255, $+$0.534]}
\newcommand{\PhVHoursAAll}{18\,621}
\newcommand{\PhVHoursAHighMtn}{5\,190}
\newcommand{\PhVHoursAHills}{1\,359}
\newcommand{\PhVHoursALowMtn}{9\,301}
\newcommand{\PhVHoursAPlains}{2\,309}
\newcommand{\PhVHoursGAll}{15\,292}
\newcommand{\PhVHoursGHighMtn}{3\,511}
\newcommand{\PhVHoursGHills}{1\,793}
\newcommand{\PhVHoursGLowMtn}{6\,958}
\newcommand{\PhVHoursGPlains}{2\,735}
\newcommand{\PhVHoursSAll}{77\,609}
\newcommand{\PhVHoursSHighMtn}{15\,846}
\newcommand{\PhVHoursSHills}{11\,732}
\newcommand{\PhVHoursSLowMtn}{30\,773}
\newcommand{\PhVHoursSPlains}{12\,532}
\newcommand{\PhVMuClimbASAll}{$-$0.559}
\newcommand{\PhVMuClimbASAllBand}{[$-$1.878, $+$0.759]}
\newcommand{\PhVMuClimbASHighMtn}{$-$0.155}
\newcommand{\PhVMuClimbASHighMtnBand}{[$-$0.965, $+$0.654]}
\newcommand{\PhVMuClimbASHills}{$+$0.134}
\newcommand{\PhVMuClimbASHillsBand}{[$-$2.136, $+$2.403]}
\newcommand{\PhVMuClimbASLowMtn}{$+$0.117}
\newcommand{\PhVMuClimbASLowMtnBand}{[$-$0.618, $+$0.853]}
\newcommand{\PhVMuClimbASPlains}{$-$1.51}
\newcommand{\PhVMuClimbASPlainsBand}{[$-$3.29, $+$0.27]}
\newcommand{\PhVMuClimbGAAll}{$+$1.35}
\newcommand{\PhVMuClimbGAAllBand}{[$-$0.83, $+$3.54]}
\newcommand{\PhVMuClimbGAHighMtn}{$+$0.0173}
\newcommand{\PhVMuClimbGAHighMtnBand}{[$-$1.0307, $+$1.0652]}
\newcommand{\PhVMuClimbGAHills}{$+$6.82}
\newcommand{\PhVMuClimbGAHillsBand}{[$+$0.35, $+$13.30]}
\newcommand{\PhVMuClimbGALowMtn}{$+$1.19}
\newcommand{\PhVMuClimbGALowMtnBand}{[$-$0.27, $+$2.65]}
\newcommand{\PhVMuClimbGAPlains}{$+$1.69}
\newcommand{\PhVMuClimbGAPlainsBand}{[$-$0.28, $+$3.67]}
\newcommand{\PhVMuGlideASAll}{$-$0.357}
\newcommand{\PhVMuGlideASAllBand}{[$-$0.752, $+$0.037]}
\newcommand{\PhVMuGlideASHighMtn}{$-$0.883}
\newcommand{\PhVMuGlideASHighMtnBand}{[$-$1.724, $-$0.042]}
\newcommand{\PhVMuGlideASHills}{$+$0.574}
\newcommand{\PhVMuGlideASHillsBand}{[$-$0.311, $+$1.458]}
\newcommand{\PhVMuGlideASLowMtn}{$-$0.355}
\newcommand{\PhVMuGlideASLowMtnBand}{[$-$0.766, $+$0.056]}
\newcommand{\PhVMuGlideASPlains}{$+$0.610}
\newcommand{\PhVMuGlideASPlainsBand}{[$-$0.298, $+$1.518]}
\newcommand{\PhVMuGlideGAAll}{$-$0.958}
\newcommand{\PhVMuGlideGAAllBand}{[$-$1.363, $-$0.552]}
\newcommand{\PhVMuGlideGAHighMtn}{$-$0.767}
\newcommand{\PhVMuGlideGAHighMtnBand}{[$-$1.522, $-$0.013]}
\newcommand{\PhVMuGlideGAHills}{$+$0.0527}
\newcommand{\PhVMuGlideGAHillsBand}{[$-$0.7629, $+$0.8683]}
\newcommand{\PhVMuGlideGALowMtn}{$-$0.857}
\newcommand{\PhVMuGlideGALowMtnBand}{[$-$1.291, $-$0.424]}
\newcommand{\PhVMuGlideGAPlains}{$-$0.606}
\newcommand{\PhVMuGlideGAPlainsBand}{[$-$1.992, $+$0.780]}
\newcommand{\PhVMuSearchASAll}{$+$0.190}
\newcommand{\PhVMuSearchASAllBand}{[$-$0.237, $+$0.617]}
\newcommand{\PhVMuSearchASHighMtn}{$-$0.0196}
\newcommand{\PhVMuSearchASHighMtnBand}{[$-$0.5073, $+$0.4682]}
\newcommand{\PhVMuSearchASHills}{$+$0.290}
\newcommand{\PhVMuSearchASHillsBand}{[$-$0.718, $+$1.299]}
\newcommand{\PhVMuSearchASLowMtn}{$+$0.110}
\newcommand{\PhVMuSearchASLowMtnBand}{[$-$0.357, $+$0.577]}
\newcommand{\PhVMuSearchASPlains}{$+$0.151}
\newcommand{\PhVMuSearchASPlainsBand}{[$-$0.733, $+$1.034]}
\newcommand{\PhVMuSearchGAAll}{$+$0.275}
\newcommand{\PhVMuSearchGAAllBand}{[$-$0.363, $+$0.913]}
\newcommand{\PhVMuSearchGAHighMtn}{$+$0.276}
\newcommand{\PhVMuSearchGAHighMtnBand}{[$-$0.403, $+$0.956]}
\newcommand{\PhVMuSearchGAHills}{$+$0.648}
\newcommand{\PhVMuSearchGAHillsBand}{[$-$6.313, $+$7.608]}
\newcommand{\PhVMuSearchGALowMtn}{$+$0.0354}
\newcommand{\PhVMuSearchGALowMtnBand}{[$-$0.6170, $+$0.6877]}
\newcommand{\PhVMuSearchGAPlains}{$+$0.786}
\newcommand{\PhVMuSearchGAPlainsBand}{[$-$0.329, $+$1.900]}
\newcommand{\PhVRateASAll}{$+$0.0620}
\newcommand{\PhVRateASAllBand}{[$+$0.0496, $+$0.0744]}
\newcommand{\PhVRateASHighMtn}{$+$0.0239}
\newcommand{\PhVRateASHighMtnBand}{[$+$0.0058, $+$0.0419]}
\newcommand{\PhVRateASHills}{$+$0.0194}
\newcommand{\PhVRateASHillsBand}{[$-$0.0075, $+$0.0463]}
\newcommand{\PhVRateASLowMtn}{$+$0.0521}
\newcommand{\PhVRateASLowMtnBand}{[$+$0.0359, $+$0.0683]}
\newcommand{\PhVRateASPlains}{$-$0.0046}
\newcommand{\PhVRateASPlainsBand}{[$-$0.0232, $+$0.0140]}
\newcommand{\PhVRateGAAll}{$-$0.0063}
\newcommand{\PhVRateGAAllBand}{[$-$0.0144, $+$0.0018]}
\newcommand{\PhVRateGAHighMtn}{$+$0.0197}
\newcommand{\PhVRateGAHighMtnBand}{[$+$0.0048, $+$0.0346]}
\newcommand{\PhVRateGAHills}{$-$0.0017}
\newcommand{\PhVRateGAHillsBand}{[$-$0.0217, $+$0.0182]}
\newcommand{\PhVRateGALowMtn}{$+$0.0027}
\newcommand{\PhVRateGALowMtnBand}{[$-$0.0092, $+$0.0145]}
\newcommand{\PhVRateGAPlains}{$-$0.0262}
\newcommand{\PhVRateGAPlainsBand}{[$-$0.0424, $-$0.0101]}
\newcommand{\PhVRefoundMeanASAll}{$+$2.34}
\newcommand{\PhVRefoundMeanASAllBand}{[$+$0.81, $+$3.86]}
\newcommand{\PhVRefoundMeanASHighMtn}{$+$1.98}
\newcommand{\PhVRefoundMeanASHighMtnBand}{[$-$1.06, $+$5.03]}
\newcommand{\PhVRefoundMeanASHills}{$-$4.91}
\newcommand{\PhVRefoundMeanASHillsBand}{[$-$8.85, $-$0.98]}
\newcommand{\PhVRefoundMeanASLowMtn}{$+$3.27}
\newcommand{\PhVRefoundMeanASLowMtnBand}{[$+$0.93, $+$5.62]}
\newcommand{\PhVRefoundMeanASPlains}{$-$0.612}
\newcommand{\PhVRefoundMeanASPlainsBand}{[$-$3.942, $+$2.718]}
\newcommand{\PhVRefoundMeanGAAll}{$+$0.0583}
\newcommand{\PhVRefoundMeanGAAllBand}{[$-$1.6709, $+$1.7876]}
\newcommand{\PhVRefoundMeanGAHighMtn}{$-$1.31}
\newcommand{\PhVRefoundMeanGAHighMtnBand}{[$-$4.93, $+$2.31]}
\newcommand{\PhVRefoundMeanGAHills}{$+$4.03}
\newcommand{\PhVRefoundMeanGAHillsBand}{[$+$0.15, $+$7.90]}
\newcommand{\PhVRefoundMeanGALowMtn}{$+$0.0108}
\newcommand{\PhVRefoundMeanGALowMtnBand}{[$-$2.7039, $+$2.7255]}
\newcommand{\PhVRefoundMeanGAPlains}{$+$0.372}
\newcommand{\PhVRefoundMeanGAPlainsBand}{[$-$3.362, $+$4.106]}
\newcommand{\PhVRefoundPerHourASAll}{$+$0.0127}
\newcommand{\PhVRefoundPerHourASAllBand}{[$-$0.0160, $+$0.0415]}
\newcommand{\PhVRefoundPerHourASHighMtn}{$-$0.0111}
\newcommand{\PhVRefoundPerHourASHighMtnBand}{[$-$0.0673, $+$0.0452]}
\newcommand{\PhVRefoundPerHourASHills}{$+$0.0687}
\newcommand{\PhVRefoundPerHourASHillsBand}{[$-$0.0197, $+$0.1570]}
\newcommand{\PhVRefoundPerHourASLowMtn}{$-$0.0261}
\newcommand{\PhVRefoundPerHourASLowMtnBand}{[$-$0.0661, $+$0.0139]}
\newcommand{\PhVRefoundPerHourASPlains}{$+$0.0504}
\newcommand{\PhVRefoundPerHourASPlainsBand}{[$-$0.0300, $+$0.1308]}
\newcommand{\PhVRefoundPerHourGAAll}{$+$0.109}
\newcommand{\PhVRefoundPerHourGAAllBand}{[$+$0.076, $+$0.141]}
\newcommand{\PhVRefoundPerHourGAHighMtn}{$+$0.0454}
\newcommand{\PhVRefoundPerHourGAHighMtnBand}{[$-$0.0232, $+$0.1141]}
\newcommand{\PhVRefoundPerHourGAHills}{$+$0.0417}
\newcommand{\PhVRefoundPerHourGAHillsBand}{[$-$0.0557, $+$0.1390]}
\newcommand{\PhVRefoundPerHourGALowMtn}{$+$0.153}
\newcommand{\PhVRefoundPerHourGALowMtnBand}{[$+$0.108, $+$0.199]}
\newcommand{\PhVRefoundPerHourGAPlains}{$+$0.119}
\newcommand{\PhVRefoundPerHourGAPlainsBand}{[$+$0.018, $+$0.221]}
\newcommand{\PhVRefoundShareASAll}{$+$0.722}
\newcommand{\PhVRefoundShareASAllBand}{[$+$0.066, $+$1.378]}
\newcommand{\PhVRefoundShareASHighMtn}{$+$1.72}
\newcommand{\PhVRefoundShareASHighMtnBand}{[$+$0.50, $+$2.94]}
\newcommand{\PhVRefoundShareASHills}{$+$5.43}
\newcommand{\PhVRefoundShareASHillsBand}{[$+$3.04, $+$7.81]}
\newcommand{\PhVRefoundShareASLowMtn}{$-$0.498}
\newcommand{\PhVRefoundShareASLowMtnBand}{[$-$1.405, $+$0.410]}
\newcommand{\PhVRefoundShareASPlains}{$+$2.80}
\newcommand{\PhVRefoundShareASPlainsBand}{[$+$0.99, $+$4.61]}
\newcommand{\PhVRefoundShareGAAll}{$+$1.05}
\newcommand{\PhVRefoundShareGAAllBand}{[$+$0.27, $+$1.82]}
\newcommand{\PhVRefoundShareGAHighMtn}{$-$1.12}
\newcommand{\PhVRefoundShareGAHighMtnBand}{[$-$2.72, $+$0.47]}
\newcommand{\PhVRefoundShareGAHills}{$-$1.35}
\newcommand{\PhVRefoundShareGAHillsBand}{[$-$4.07, $+$1.37]}
\newcommand{\PhVRefoundShareGALowMtn}{$+$1.80}
\newcommand{\PhVRefoundShareGALowMtnBand}{[$+$0.74, $+$2.85]}
\newcommand{\PhVRefoundShareGAPlains}{$+$2.46}
\newcommand{\PhVRefoundShareGAPlainsBand}{[$+$0.08, $+$4.84]}
\newcommand{\PhVSearchDropped}{3.3}
\newcommand{\PhVSearchDroppedAll}{3.3}
\newcommand{\PhVSearchDroppedHighMtn}{3.0}
\newcommand{\PhVSearchDroppedHills}{4.2}
\newcommand{\PhVSearchDroppedLowMtn}{3.2}
\newcommand{\PhVSearchDroppedPlains}{2.8}
\newcommand{\PhVSearchMeanASAll}{$+$4.33}
\newcommand{\PhVSearchMeanASAllBand}{[$+$3.05, $+$5.60]}
\newcommand{\PhVSearchMeanASHighMtn}{$+$0.0780}
\newcommand{\PhVSearchMeanASHighMtnBand}{[$-$2.1148, $+$2.2708]}
\newcommand{\PhVSearchMeanASHills}{$-$6.49}
\newcommand{\PhVSearchMeanASHillsBand}{[$-$9.59, $-$3.39]}
\newcommand{\PhVSearchMeanASLowMtn}{$+$4.35}
\newcommand{\PhVSearchMeanASLowMtnBand}{[$+$2.60, $+$6.10]}
\newcommand{\PhVSearchMeanASPlains}{$+$1.23}
\newcommand{\PhVSearchMeanASPlainsBand}{[$-$1.60, $+$4.06]}
\newcommand{\PhVSearchMeanGAAll}{$-$5.01}
\newcommand{\PhVSearchMeanGAAllBand}{[$-$6.26, $-$3.76]}
\newcommand{\PhVSearchMeanGAHighMtn}{$-$4.32}
\newcommand{\PhVSearchMeanGAHighMtnBand}{[$-$6.77, $-$1.87]}
\newcommand{\PhVSearchMeanGAHills}{$-$2.56}
\newcommand{\PhVSearchMeanGAHillsBand}{[$-$5.46, $+$0.34]}
\newcommand{\PhVSearchMeanGALowMtn}{$-$5.01}
\newcommand{\PhVSearchMeanGALowMtnBand}{[$-$6.97, $-$3.05]}
\newcommand{\PhVSearchMeanGAPlains}{$-$3.77}
\newcommand{\PhVSearchMeanGAPlainsBand}{[$-$6.40, $-$1.15]}
\newcommand{\PhVSearchPerHourASAll}{$-$0.106}
\newcommand{\PhVSearchPerHourASAllBand}{[$-$0.183, $-$0.029]}
\newcommand{\PhVSearchPerHourASHighMtn}{$-$0.244}
\newcommand{\PhVSearchPerHourASHighMtnBand}{[$-$0.370, $-$0.119]}
\newcommand{\PhVSearchPerHourASHills}{$-$0.411}
\newcommand{\PhVSearchPerHourASHillsBand}{[$-$0.597, $-$0.226]}
\newcommand{\PhVSearchPerHourASLowMtn}{$-$0.142}
\newcommand{\PhVSearchPerHourASLowMtnBand}{[$-$0.244, $-$0.041]}
\newcommand{\PhVSearchPerHourASPlains}{$-$0.215}
\newcommand{\PhVSearchPerHourASPlainsBand}{[$-$0.359, $-$0.071]}
\newcommand{\PhVSearchPerHourGAAll}{$-$0.0663}
\newcommand{\PhVSearchPerHourGAAllBand}{[$-$0.1271, $-$0.0054]}
\newcommand{\PhVSearchPerHourGAHighMtn}{$-$0.143}
\newcommand{\PhVSearchPerHourGAHighMtnBand}{[$-$0.264, $-$0.023]}
\newcommand{\PhVSearchPerHourGAHills}{$-$0.0256}
\newcommand{\PhVSearchPerHourGAHillsBand}{[$-$0.2050, $+$0.1538]}
\newcommand{\PhVSearchPerHourGALowMtn}{$-$0.0228}
\newcommand{\PhVSearchPerHourGALowMtnBand}{[$-$0.1170, $+$0.0714]}
\newcommand{\PhVSearchPerHourGAPlains}{$+$0.0074}
\newcommand{\PhVSearchPerHourGAPlainsBand}{[$-$0.1334, $+$0.1483]}
\newcommand{\PhVShareClimbASAll}{$-$1.39}
\newcommand{\PhVShareClimbASAllBand}{[$-$1.82, $-$0.96]}
\newcommand{\PhVShareClimbASHighMtn}{$-$0.561}
\newcommand{\PhVShareClimbASHighMtnBand}{[$-$1.059, $-$0.063]}
\newcommand{\PhVShareClimbASHills}{$+$3.38}
\newcommand{\PhVShareClimbASHillsBand}{[$+$2.42, $+$4.34]}
\newcommand{\PhVShareClimbASLowMtn}{$-$1.15}
\newcommand{\PhVShareClimbASLowMtnBand}{[$-$1.67, $-$0.63]}
\newcommand{\PhVShareClimbASPlains}{$+$1.58}
\newcommand{\PhVShareClimbASPlainsBand}{[$+$0.82, $+$2.34]}
\newcommand{\PhVShareClimbGAAll}{$+$4.50}
\newcommand{\PhVShareClimbGAAllBand}{[$+$4.20, $+$4.81]}
\newcommand{\PhVShareClimbGAHighMtn}{$+$4.87}
\newcommand{\PhVShareClimbGAHighMtnBand}{[$+$4.33, $+$5.40]}
\newcommand{\PhVShareClimbGAHills}{$+$3.38}
\newcommand{\PhVShareClimbGAHillsBand}{[$+$2.52, $+$4.24]}
\newcommand{\PhVShareClimbGALowMtn}{$+$4.19}
\newcommand{\PhVShareClimbGALowMtnBand}{[$+$3.76, $+$4.62]}
\newcommand{\PhVShareClimbGAPlains}{$+$3.28}
\newcommand{\PhVShareClimbGAPlainsBand}{[$+$2.53, $+$4.03]}
\newcommand{\PhVShareExitASAll}{$-$0.128}
\newcommand{\PhVShareExitASAllBand}{[$-$0.209, $-$0.046]}
\newcommand{\PhVShareExitASHighMtn}{$-$0.0591}
\newcommand{\PhVShareExitASHighMtnBand}{[$-$0.1887, $+$0.0705]}
\newcommand{\PhVShareExitASHills}{$-$0.140}
\newcommand{\PhVShareExitASHillsBand}{[$-$0.344, $+$0.065]}
\newcommand{\PhVShareExitASLowMtn}{$-$0.123}
\newcommand{\PhVShareExitASLowMtnBand}{[$-$0.239, $-$0.008]}
\newcommand{\PhVShareExitASPlains}{$-$0.129}
\newcommand{\PhVShareExitASPlainsBand}{[$-$0.340, $+$0.082]}
\newcommand{\PhVShareExitGAAll}{$+$0.141}
\newcommand{\PhVShareExitGAAllBand}{[$+$0.059, $+$0.224]}
\newcommand{\PhVShareExitGAHighMtn}{$-$0.105}
\newcommand{\PhVShareExitGAHighMtnBand}{[$-$0.249, $+$0.039]}
\newcommand{\PhVShareExitGAHills}{$+$0.404}
\newcommand{\PhVShareExitGAHillsBand}{[$+$0.157, $+$0.650]}
\newcommand{\PhVShareExitGALowMtn}{$+$0.0322}
\newcommand{\PhVShareExitGALowMtnBand}{[$-$0.0845, $+$0.1488]}
\newcommand{\PhVShareExitGAPlains}{$+$0.507}
\newcommand{\PhVShareExitGAPlainsBand}{[$+$0.283, $+$0.732]}
\newcommand{\PhVShareGlideASAll}{$+$1.30}
\newcommand{\PhVShareGlideASAllBand}{[$+$0.86, $+$1.73]}
\newcommand{\PhVShareGlideASHighMtn}{$+$1.19}
\newcommand{\PhVShareGlideASHighMtnBand}{[$+$0.60, $+$1.77]}
\newcommand{\PhVShareGlideASHills}{$-$1.84}
\newcommand{\PhVShareGlideASHillsBand}{[$-$2.84, $-$0.85]}
\newcommand{\PhVShareGlideASLowMtn}{$+$1.12}
\newcommand{\PhVShareGlideASLowMtnBand}{[$+$0.54, $+$1.71]}
\newcommand{\PhVShareGlideASPlains}{$-$1.16}
\newcommand{\PhVShareGlideASPlainsBand}{[$-$2.08, $-$0.24]}
\newcommand{\PhVShareGlideGAAll}{$-$3.98}
\newcommand{\PhVShareGlideGAAllBand}{[$-$4.34, $-$3.63]}
\newcommand{\PhVShareGlideGAHighMtn}{$-$3.98}
\newcommand{\PhVShareGlideGAHighMtnBand}{[$-$4.67, $-$3.29]}
\newcommand{\PhVShareGlideGAHills}{$-$3.51}
\newcommand{\PhVShareGlideGAHillsBand}{[$-$4.47, $-$2.55]}
\newcommand{\PhVShareGlideGALowMtn}{$-$3.62}
\newcommand{\PhVShareGlideGALowMtnBand}{[$-$4.14, $-$3.09]}
\newcommand{\PhVShareGlideGAPlains}{$-$3.46}
\newcommand{\PhVShareGlideGAPlainsBand}{[$-$4.31, $-$2.60]}
\newcommand{\PhVShareSearchASAll}{$+$0.218}
\newcommand{\PhVShareSearchASAllBand}{[$-$0.024, $+$0.461]}
\newcommand{\PhVShareSearchASHighMtn}{$-$0.569}
\newcommand{\PhVShareSearchASHighMtnBand}{[$-$0.960, $-$0.178]}
\newcommand{\PhVShareSearchASHills}{$-$1.40}
\newcommand{\PhVShareSearchASHillsBand}{[$-$1.88, $-$0.91]}
\newcommand{\PhVShareSearchASLowMtn}{$+$0.150}
\newcommand{\PhVShareSearchASLowMtnBand}{[$-$0.186, $+$0.486]}
\newcommand{\PhVShareSearchASPlains}{$-$0.300}
\newcommand{\PhVShareSearchASPlainsBand}{[$-$0.716, $+$0.116]}
\newcommand{\PhVShareSearchGAAll}{$-$0.661}
\newcommand{\PhVShareSearchGAAllBand}{[$-$0.848, $-$0.474]}
\newcommand{\PhVShareSearchGAHighMtn}{$-$0.780}
\newcommand{\PhVShareSearchGAHighMtnBand}{[$-$1.157, $-$0.403]}
\newcommand{\PhVShareSearchGAHills}{$-$0.270}
\newcommand{\PhVShareSearchGAHillsBand}{[$-$0.666, $+$0.127]}
\newcommand{\PhVShareSearchGALowMtn}{$-$0.606}
\newcommand{\PhVShareSearchGALowMtnBand}{[$-$0.929, $-$0.283]}
\newcommand{\PhVShareSearchGAPlains}{$-$0.333}
\newcommand{\PhVShareSearchGAPlainsBand}{[$-$0.686, $+$0.021]}
\newcommand{\PhVSinkASAll}{$-$6.07}
\newcommand{\PhVSinkASAllBand}{[$-$7.11, $-$5.02]}
\newcommand{\PhVSinkASHighMtn}{$-$4.39}
\newcommand{\PhVSinkASHighMtnBand}{[$-$6.04, $-$2.75]}
\newcommand{\PhVSinkASHills}{$-$5.30}
\newcommand{\PhVSinkASHillsBand}{[$-$7.88, $-$2.72]}
\newcommand{\PhVSinkASLowMtn}{$-$6.65}
\newcommand{\PhVSinkASLowMtnBand}{[$-$7.99, $-$5.32]}
\newcommand{\PhVSinkASPlains}{$-$6.37}
\newcommand{\PhVSinkASPlainsBand}{[$-$8.77, $-$3.97]}
\newcommand{\PhVSinkGAAll}{$-$12.8}
\newcommand{\PhVSinkGAAllBand}{[$-$13.6, $-$12.0]}
\newcommand{\PhVSinkGAHighMtn}{$-$14.4}
\newcommand{\PhVSinkGAHighMtnBand}{[$-$16.0, $-$12.9]}
\newcommand{\PhVSinkGAHills}{$-$7.05}
\newcommand{\PhVSinkGAHillsBand}{[$-$9.05, $-$5.05]}
\newcommand{\PhVSinkGALowMtn}{$-$15.5}
\newcommand{\PhVSinkGALowMtnBand}{[$-$16.7, $-$14.3]}
\newcommand{\PhVSinkGAPlains}{$-$7.13}
\newcommand{\PhVSinkGAPlainsBand}{[$-$8.67, $-$5.59]}
\newcommand{\PhVSpeedASAll}{$+$1.24}
\newcommand{\PhVSpeedASAllBand}{[$+$0.89, $+$1.58]}
\newcommand{\PhVSpeedASHighMtn}{$+$1.31}
\newcommand{\PhVSpeedASHighMtnBand}{[$+$0.95, $+$1.67]}
\newcommand{\PhVSpeedASHills}{$+$3.83}
\newcommand{\PhVSpeedASHillsBand}{[$+$2.94, $+$4.72]}
\newcommand{\PhVSpeedASLowMtn}{$+$1.54}
\newcommand{\PhVSpeedASLowMtnBand}{[$+$1.19, $+$1.89]}
\newcommand{\PhVSpeedASPlains}{$+$3.12}
\newcommand{\PhVSpeedASPlainsBand}{[$+$2.22, $+$4.01]}
\newcommand{\PhVSpeedGAAll}{$+$1.30}
\newcommand{\PhVSpeedGAAllBand}{[$+$1.12, $+$1.49]}
\newcommand{\PhVSpeedGAHighMtn}{$+$1.21}
\newcommand{\PhVSpeedGAHighMtnBand}{[$+$0.90, $+$1.51]}
\newcommand{\PhVSpeedGAHills}{$+$1.41}
\newcommand{\PhVSpeedGAHillsBand}{[$+$0.83, $+$1.99]}
\newcommand{\PhVSpeedGALowMtn}{$+$1.00}
\newcommand{\PhVSpeedGALowMtnBand}{[$+$0.78, $+$1.23]}
\newcommand{\PhVSpeedGAPlains}{$+$0.457}
\newcommand{\PhVSpeedGAPlainsBand}{[$-$0.027, $+$0.940]}
\newcommand{\PhVXminClimbAll}{1\,033}
\newcommand{\PhVXminClimbHighMtn}{654}
\newcommand{\PhVXminClimbHills}{993}
\newcommand{\PhVXminClimbLowMtn}{744}
\newcommand{\PhVXminClimbPlains}{908}
\newcommand{\PhVXminGlideAll}{1\,682}
\newcommand{\PhVXminGlideHighMtn}{1\,972}
\newcommand{\PhVXminGlideHills}{1\,003}
\newcommand{\PhVXminGlideLowMtn}{1\,632}
\newcommand{\PhVXminGlidePlains}{1\,130}
\newcommand{\PhVXminSearchAll}{654}
\newcommand{\PhVXminSearchHighMtn}{360}
\newcommand{\PhVXminSearchHills}{195}
\newcommand{\PhVXminSearchLowMtn}{495}
\newcommand{\PhVXminSearchPlains}{336}
\newcommand{\PhXminClimbAll}{1\,075}
\newcommand{\PhXminClimbHighMtn}{660}
\newcommand{\PhXminClimbHills}{1\,043}
\newcommand{\PhXminClimbLowMtn}{782}
\newcommand{\PhXminClimbPlains}{1\,118}
\newcommand{\PhXminGlideAll}{1\,682}
\newcommand{\PhXminGlideHighMtn}{1\,538}
\newcommand{\PhXminGlideHills}{1\,075}
\newcommand{\PhXminGlideLowMtn}{1\,665}
\newcommand{\PhXminGlidePlains}{1\,130}
\newcommand{\PhXminSearchAll}{510}
\newcommand{\PhXminSearchHighMtn}{574}
\newcommand{\PhXminSearchHills}{382}
\newcommand{\PhXminSearchLowMtn}{654}
\newcommand{\PhXminSearchPlains}{232}
\newcommand{\PieceMinMinutes}{21}
\newcommand{\PieceMinS}{1\,250}
\newcommand{\PilotShareS}{82}
\newcommand{\PlusFactor}{2.2}
\newcommand{\PlusLagHours}{3}
\newcommand{\PlusPieceHours}{6.7}
\newcommand{\PoolAMax}{1.29}
\newcommand{\PoolAMinusChainLong}{$+$0.049}
\newcommand{\PoolAMinusChainThousand}{$+$0.003}
\newcommand{\PoolAMinusChainVeryLong}{$+$0.122}
\newcommand{\PoolAMinusLong}{$+$0.044}
\newcommand{\PoolAMinusThousand}{$+$0.002}
\newcommand{\PoolAMinusVeryLong}{$+$0.107}
\newcommand{\PoolAboveA}{1\,500}
\newcommand{\PoolAboveG}{1\,800}
\newcommand{\PoolBelowEnd}{2\,800}
\newcommand{\PoolGMax}{1.05}
\newcommand{\PoolGMinusChainLong}{$+$0.052}
\newcommand{\PoolGMinusChainThousand}{$-$0.000}
\newcommand{\PoolGMinusChainVeryLong}{$+$0.158}
\newcommand{\PoolGMinusLong}{$+$0.047}
\newcommand{\PoolGMinusThousand}{$-$0.000}
\newcommand{\PoolGMinusVeryLong}{$+$0.144}
\newcommand{\PoolLeave}{870}
\newcommand{\PoolMsdCross}{6\,200}
\newcommand{\PoolRunsAwayA}{870}
\newcommand{\PoolRunsAwayG}{1\,500}
\newcommand{\PoolRunsAwayS}{3\,400}
\newcommand{\PoolSMax}{1.46}
\newcommand{\PoolSMinusChainLong}{$-$0.005}
\newcommand{\PoolSMinusChainThousand}{$-$0.010}
\newcommand{\PoolSMinusChainVeryLong}{$+$0.107}
\newcommand{\PoolSMinusLong}{$-$0.005}
\newcommand{\PoolSMinusThousand}{$-$0.009}
\newcommand{\PoolSMinusVeryLong}{$+$0.099}
\newcommand{\PoolSame}{710}
\newcommand{\PoolSteep}{4\,200}
\newcommand{\PreCAClimb}{$-$1.3}
\newcommand{\PreCAClimbBand}{[$-$1.8, $-$0.8]}
\newcommand{\PreCACph}{$-$0.36}
\newcommand{\PreCACphBand}{[$-$0.46, $-$0.26]}
\newcommand{\PreCAFound}{$-$2.5}
\newcommand{\PreCAFoundBand}{[$-$3.6, $-$1.3]}
\newcommand{\PreCAGlide}{$-$0.2}
\newcommand{\PreCAGlideBand}{[$-$0.5, $+$0.1]}
\newcommand{\PreCAHK}{$-$0.000}
\newcommand{\PreCAHKBand}{[$-$0.005, $+$0.004]}
\newcommand{\PreCARate}{$+$0.006}
\newcommand{\PreCARateBand}{[$-$0.007, $+$0.020]}
\newcommand{\PreCASearch}{$-$0.1}
\newcommand{\PreCASearchBand}{[$-$0.5, $+$0.3]}
\newcommand{\PreCASink}{$+$2}
\newcommand{\PreCASinkBand}{[$+$1, $+$3]}
\newcommand{\PreCAXc}{$+$0.9}
\newcommand{\PreCAXcBand}{[$+$0.5, $+$1.2]}
\newcommand{\PreCSClimb}{$-$1.1}
\newcommand{\PreCSClimbBand}{[$-$1.3, $-$0.9]}
\newcommand{\PreCSCph}{$-$0.44}
\newcommand{\PreCSCphBand}{[$-$0.49, $-$0.39]}
\newcommand{\PreCSFound}{$-$1.7}
\newcommand{\PreCSFoundBand}{[$-$2.2, $-$1.2]}
\newcommand{\PreCSGlide}{$+$0.3}
\newcommand{\PreCSGlideBand}{[$+$0.1, $+$0.4]}
\newcommand{\PreCSHK}{$+$0.002}
\newcommand{\PreCSHKBand}{[$-$0.000, $+$0.004]}
\newcommand{\PreCSRate}{$-$0.011}
\newcommand{\PreCSRateBand}{[$-$0.017, $-$0.006]}
\newcommand{\PreCSSearch}{$-$0.4}
\newcommand{\PreCSSearchBand}{[$-$0.6, $-$0.2]}
\newcommand{\PreCSSink}{$+$3}
\newcommand{\PreCSSinkBand}{[$+$2, $+$3]}
\newcommand{\PreCSXc}{$+$1.1}
\newcommand{\PreCSXcBand}{[$+$1.0, $+$1.2]}
\newcommand{\PreDwClimb}{$-$0.8}
\newcommand{\PreDwClimbBand}{[$-$1.4, $-$0.3]}
\newcommand{\PreDwCph}{$-$0.21}
\newcommand{\PreDwCphBand}{[$-$0.33, $-$0.10]}
\newcommand{\PreDwFound}{$-$0.6}
\newcommand{\PreDwFoundBand}{[$-$1.9, $+$0.7]}
\newcommand{\PreDwGlide}{$-$0.2}
\newcommand{\PreDwGlideBand}{[$-$0.4, $+$0.1]}
\newcommand{\PreDwHK}{$-$0.001}
\newcommand{\PreDwHKBand}{[$-$0.006, $+$0.004]}
\newcommand{\PreDwRate}{$+$0.015}
\newcommand{\PreDwRateBand}{[$-$0.001, $+$0.032]}
\newcommand{\PreDwSearch}{$+$0.5}
\newcommand{\PreDwSearchBand}{[$+$0.1, $+$0.9]}
\newcommand{\PreDwSink}{$+$4}
\newcommand{\PreDwSinkBand}{[$+$2, $+$5]}
\newcommand{\PreDwXc}{$+$0.0}
\newcommand{\PreDwXcBand}{[$-$0.3, $+$0.3]}
\newcommand{\PreEvents}{1\,595}
\newcommand{\PreEventsMSDCA}{5\,205}
\newcommand{\PreEventsMSDCS}{26\,175}
\newcommand{\PreEventsMSDDw}{1\,595}
\newcommand{\PreEventsMSDG}{1\,595}
\newcommand{\PreEventsMSDIn}{1\,595}
\newcommand{\PreEventsMSDOut}{1\,595}
\newcommand{\PreEventsMSDPP}{1\,818}
\newcommand{\PreEventsPhaseCA}{3\,678}
\newcommand{\PreEventsPhaseCS}{17\,651}
\newcommand{\PreEventsPhaseDw}{1\,150}
\newcommand{\PreEventsPhaseG}{1\,150}
\newcommand{\PreEventsPhaseIn}{1\,150}
\newcommand{\PreEventsPhaseOut}{1\,150}
\newcommand{\PreEventsPhasePP}{1\,298}
\newcommand{\PreFitEnd}{0.934}
\newcommand{\PreFitPost}{0.925}
\newcommand{\PreFitPre}{0.922}
\newcommand{\PreFitStart}{0.931}
\newcommand{\PreFlightsPhaseCA}{3\,289}
\newcommand{\PreFlightsPhaseCS}{17\,651}
\newcommand{\PreFlightsPhaseDw}{1\,113}
\newcommand{\PreFlightsPhaseG}{1\,113}
\newcommand{\PreFlightsPhaseIn}{1\,113}
\newcommand{\PreFlightsPhaseOut}{1\,113}
\newcommand{\PreFlightsPhasePP}{1\,252}
\newcommand{\PreGClimb}{$+$2.8}
\newcommand{\PreGClimbBand}{[$+$2.0, $+$3.6]}
\newcommand{\PreGCph}{$+$0.36}
\newcommand{\PreGCphBand}{[$+$0.22, $+$0.51]}
\newcommand{\PreGCross}{50}
\newcommand{\PreGFound}{$+$2.7}
\newcommand{\PreGFoundBand}{[$+$1.1, $+$4.4]}
\newcommand{\PreGGlide}{$+$1.3}
\newcommand{\PreGGlideBand}{[$+$0.9, $+$1.6]}
\newcommand{\PreGHK}{$+$0.026}
\newcommand{\PreGHKBand}{[$+$0.020, $+$0.032]}
\newcommand{\PreGRate}{$+$0.041}
\newcommand{\PreGRateBand}{[$+$0.020, $+$0.062]}
\newcommand{\PreGSearch}{$-$0.9}
\newcommand{\PreGSearchBand}{[$-$1.4, $-$0.4]}
\newcommand{\PreGSink}{$-$5}
\newcommand{\PreGSinkBand}{[$-$7, $-$4]}
\newcommand{\PreGXc}{$+$1.4}
\newcommand{\PreGXcBand}{[$+$1.0, $+$1.8]}
\newcommand{\PreHKEnd}{0.926}
\newcommand{\PreHKPost}{0.899}
\newcommand{\PreHKPre}{0.899}
\newcommand{\PreHKStart}{0.924}
\newcommand{\PreInClimb}{$+$2.0}
\newcommand{\PreInClimbBand}{[$+$1.0, $+$2.9]}
\newcommand{\PreInCph}{$+$0.15}
\newcommand{\PreInCphBand}{[$-$0.03, $+$0.33]}
\newcommand{\PreInFound}{$+$2.1}
\newcommand{\PreInFoundBand}{[$+$0.1, $+$4.2]}
\newcommand{\PreInGlide}{$+$1.1}
\newcommand{\PreInGlideBand}{[$+$0.6, $+$1.6]}
\newcommand{\PreInHK}{$+$0.025}
\newcommand{\PreInHKBand}{[$+$0.017, $+$0.034]}
\newcommand{\PreInRate}{$+$0.056}
\newcommand{\PreInRateBand}{[$+$0.031, $+$0.081]}
\newcommand{\PreInSearch}{$-$0.4}
\newcommand{\PreInSearchBand}{[$-$1.0, $+$0.3]}
\newcommand{\PreInSink}{$-$2}
\newcommand{\PreInSinkBand}{[$-$4, $+$0]}
\newcommand{\PreInXc}{$+$1.4}
\newcommand{\PreInXcBand}{[$+$0.8, $+$1.9]}
\newcommand{\PreLags}{24}
\newcommand{\PreMaxLag}{1\,840}
\newcommand{\PreMockCaliper}{month $\pm2$, hour $\pm2$, same class}
\newcommand{\PreMockDwClimb}{$-$1.3}
\newcommand{\PreMockDwClimbBand}{[$-$1.6, $-$1.0]}
\newcommand{\PreMockDwCph}{$-$0.26}
\newcommand{\PreMockDwCphBand}{[$-$0.31, $-$0.21]}
\newcommand{\PreMockDwFound}{$-$1.3}
\newcommand{\PreMockDwFoundBand}{[$-$2.0, $-$0.7]}
\newcommand{\PreMockDwGlide}{$-$0.1}
\newcommand{\PreMockDwGlideBand}{[$-$0.3, $+$0.0]}
\newcommand{\PreMockDwHK}{$+$0.001}
\newcommand{\PreMockDwHKBand}{[$-$0.001, $+$0.003]}
\newcommand{\PreMockDwRate}{$-$0.001}
\newcommand{\PreMockDwRateBand}{[$-$0.008, $+$0.006]}
\newcommand{\PreMockDwSearch}{$+$0.0}
\newcommand{\PreMockDwSearchBand}{[$-$0.2, $+$0.2]}
\newcommand{\PreMockDwSink}{$+$1}
\newcommand{\PreMockDwSinkBand}{[$+$0, $+$2]}
\newcommand{\PreMockDwXc}{$+$0.3}
\newcommand{\PreMockDwXcBand}{[$+$0.2, $+$0.5]}
\newcommand{\PreMockEvents}{7\,666}
\newcommand{\PreMockEventsPhase}{5\,580}
\newcommand{\PreMockFlights}{4\,716}
\newcommand{\PreMockGClimb}{$+$1.0}
\newcommand{\PreMockGClimbBand}{[$+$0.7, $+$1.3]}
\newcommand{\PreMockGCph}{$+$0.02}
\newcommand{\PreMockGCphBand}{[$-$0.04, $+$0.09]}
\newcommand{\PreMockGFound}{$+$0.7}
\newcommand{\PreMockGFoundBand}{[$-$0.1, $+$1.5]}
\newcommand{\PreMockGGlide}{$+$0.5}
\newcommand{\PreMockGGlideBand}{[$+$0.3, $+$0.7]}
\newcommand{\PreMockGHK}{$+$0.000}
\newcommand{\PreMockGHKBand}{[$-$0.002, $+$0.003]}
\newcommand{\PreMockGRate}{$+$0.031}
\newcommand{\PreMockGRateBand}{[$+$0.022, $+$0.040]}
\newcommand{\PreMockGSearch}{$-$0.2}
\newcommand{\PreMockGSearchBand}{[$-$0.5, $+$0.0]}
\newcommand{\PreMockGSink}{$-$1}
\newcommand{\PreMockGSinkBand}{[$-$1, $+$0]}
\newcommand{\PreMockGXc}{$+$0.2}
\newcommand{\PreMockGXcBand}{[$+$0.0, $+$0.4]}
\newcommand{\PreMockK}{5}
\newcommand{\PreMockPool}{10}
\newcommand{\PreMockRealAll}{1\,595}
\newcommand{\PreMockRealEvents}{1\,534}
\newcommand{\PreMockRealEventsPhase}{1\,121}
\newcommand{\PreMockSiteDwClimb}{$-$0.7}
\newcommand{\PreMockSiteDwClimbBand}{[$-$1.1, $-$0.4]}
\newcommand{\PreMockSiteDwCph}{$-$0.17}
\newcommand{\PreMockSiteDwCphBand}{[$-$0.24, $-$0.10]}
\newcommand{\PreMockSiteDwFound}{$-$0.3}
\newcommand{\PreMockSiteDwFoundBand}{[$-$1.1, $+$0.5]}
\newcommand{\PreMockSiteDwGlide}{$+$0.0}
\newcommand{\PreMockSiteDwGlideBand}{[$-$0.2, $+$0.2]}
\newcommand{\PreMockSiteDwHK}{$+$0.001}
\newcommand{\PreMockSiteDwHKBand}{[$-$0.002, $+$0.004]}
\newcommand{\PreMockSiteDwRate}{$+$0.007}
\newcommand{\PreMockSiteDwRateBand}{[$-$0.003, $+$0.018]}
\newcommand{\PreMockSiteDwSearch}{$-$0.0}
\newcommand{\PreMockSiteDwSearchBand}{[$-$0.2, $+$0.2]}
\newcommand{\PreMockSiteDwSink}{$+$3}
\newcommand{\PreMockSiteDwSinkBand}{[$+$2, $+$4]}
\newcommand{\PreMockSiteDwXc}{$+$0.4}
\newcommand{\PreMockSiteDwXcBand}{[$+$0.2, $+$0.6]}
\newcommand{\PreMockSiteEvents}{6\,166}
\newcommand{\PreMockSiteEventsPhase}{4\,268}
\newcommand{\PreMockSiteFlights}{3\,339}
\newcommand{\PreMockSiteGClimb}{$+$0.7}
\newcommand{\PreMockSiteGClimbBand}{[$+$0.3, $+$1.1]}
\newcommand{\PreMockSiteGCph}{$+$0.09}
\newcommand{\PreMockSiteGCphBand}{[$-$0.00, $+$0.17]}
\newcommand{\PreMockSiteGFound}{$+$0.5}
\newcommand{\PreMockSiteGFoundBand}{[$-$0.5, $+$1.5]}
\newcommand{\PreMockSiteGGlide}{$+$0.6}
\newcommand{\PreMockSiteGGlideBand}{[$+$0.4, $+$0.8]}
\newcommand{\PreMockSiteGHK}{$+$0.003}
\newcommand{\PreMockSiteGHKBand}{[$-$0.001, $+$0.006]}
\newcommand{\PreMockSiteGRate}{$+$0.029}
\newcommand{\PreMockSiteGRateBand}{[$+$0.017, $+$0.040]}
\newcommand{\PreMockSiteGSearch}{$-$0.1}
\newcommand{\PreMockSiteGSearchBand}{[$-$0.4, $+$0.2]}
\newcommand{\PreMockSiteGSink}{$+$0}
\newcommand{\PreMockSiteGSinkBand}{[$-$1, $+$1]}
\newcommand{\PreMockSiteGXc}{$+$0.1}
\newcommand{\PreMockSiteGXcBand}{[$-$0.2, $+$0.4]}
\newcommand{\PreMockSiteRealAll}{1\,595}
\newcommand{\PreMockSiteRealEvents}{1\,382}
\newcommand{\PreMockSiteRealEventsPhase}{974}
\newcommand{\PreMockSubsetGClimb}{$+$2.8}
\newcommand{\PreMockSubsetGClimbBand}{[$+$2.0, $+$3.6]}
\newcommand{\PreMockSubsetGCph}{$+$0.37}
\newcommand{\PreMockSubsetGCphBand}{[$+$0.22, $+$0.52]}
\newcommand{\PreMockSubsetGFound}{$+$2.8}
\newcommand{\PreMockSubsetGFoundBand}{[$+$1.1, $+$4.5]}
\newcommand{\PreMockSubsetGGlide}{$+$1.2}
\newcommand{\PreMockSubsetGGlideBand}{[$+$0.9, $+$1.6]}
\newcommand{\PreMockSubsetGHK}{$+$0.028}
\newcommand{\PreMockSubsetGHKBand}{[$+$0.021, $+$0.034]}
\newcommand{\PreMockSubsetGRate}{$+$0.043}
\newcommand{\PreMockSubsetGRateBand}{[$+$0.022, $+$0.064]}
\newcommand{\PreMockSubsetGSearch}{$-$0.9}
\newcommand{\PreMockSubsetGSearchBand}{[$-$1.4, $-$0.4]}
\newcommand{\PreMockSubsetGSink}{$-$5}
\newcommand{\PreMockSubsetGSinkBand}{[$-$7, $-$4]}
\newcommand{\PreMockSubsetGXc}{$+$1.4}
\newcommand{\PreMockSubsetGXcBand}{[$+$1.0, $+$1.8]}
\newcommand{\PreMockSubsetSiteGClimb}{$+$2.9}
\newcommand{\PreMockSubsetSiteGClimbBand}{[$+$2.1, $+$3.8]}
\newcommand{\PreMockSubsetSiteGCph}{$+$0.34}
\newcommand{\PreMockSubsetSiteGCphBand}{[$+$0.19, $+$0.50]}
\newcommand{\PreMockSubsetSiteGFound}{$+$2.9}
\newcommand{\PreMockSubsetSiteGFoundBand}{[$+$1.0, $+$4.7]}
\newcommand{\PreMockSubsetSiteGGlide}{$+$1.2}
\newcommand{\PreMockSubsetSiteGGlideBand}{[$+$0.8, $+$1.6]}
\newcommand{\PreMockSubsetSiteGHK}{$+$0.029}
\newcommand{\PreMockSubsetSiteGHKBand}{[$+$0.023, $+$0.036]}
\newcommand{\PreMockSubsetSiteGRate}{$+$0.048}
\newcommand{\PreMockSubsetSiteGRateBand}{[$+$0.025, $+$0.071]}
\newcommand{\PreMockSubsetSiteGSearch}{$-$1.2}
\newcommand{\PreMockSubsetSiteGSearchBand}{[$-$1.7, $-$0.6]}
\newcommand{\PreMockSubsetSiteGSink}{$-$5}
\newcommand{\PreMockSubsetSiteGSinkBand}{[$-$7, $-$3]}
\newcommand{\PreMockSubsetSiteGXc}{$+$1.5}
\newcommand{\PreMockSubsetSiteGXcBand}{[$+$1.0, $+$1.9]}
\newcommand{\PreNetGClimb}{$+$1.8}
\newcommand{\PreNetGClimbBand}{[$+$1.0, $+$2.7]}
\newcommand{\PreNetGCph}{$+$0.35}
\newcommand{\PreNetGCphBand}{[$+$0.19, $+$0.51]}
\newcommand{\PreNetGFound}{$+$2.1}
\newcommand{\PreNetGFoundBand}{[$+$0.3, $+$4.0]}
\newcommand{\PreNetGGlide}{$+$0.7}
\newcommand{\PreNetGGlideBand}{[$+$0.3, $+$1.2]}
\newcommand{\PreNetGHK}{$+$0.027}
\newcommand{\PreNetGHKBand}{[$+$0.021, $+$0.034]}
\newcommand{\PreNetGRate}{$+$0.012}
\newcommand{\PreNetGRateBand}{[$-$0.011, $+$0.035]}
\newcommand{\PreNetGSearch}{$-$0.7}
\newcommand{\PreNetGSearchBand}{[$-$1.3, $-$0.2]}
\newcommand{\PreNetGSink}{$-$5}
\newcommand{\PreNetGSinkBand}{[$-$7, $-$3]}
\newcommand{\PreNetGXc}{$+$1.2}
\newcommand{\PreNetGXcBand}{[$+$0.7, $+$1.6]}
\newcommand{\PreNetSiteGClimb}{$+$2.2}
\newcommand{\PreNetSiteGClimbBand}{[$+$1.2, $+$3.1]}
\newcommand{\PreNetSiteGCph}{$+$0.26}
\newcommand{\PreNetSiteGCphBand}{[$+$0.08, $+$0.43]}
\newcommand{\PreNetSiteGFound}{$+$2.3}
\newcommand{\PreNetSiteGFoundBand}{[$+$0.2, $+$4.5]}
\newcommand{\PreNetSiteGGlide}{$+$0.6}
\newcommand{\PreNetSiteGGlideBand}{[$+$0.1, $+$1.1]}
\newcommand{\PreNetSiteGHK}{$+$0.026}
\newcommand{\PreNetSiteGHKBand}{[$+$0.019, $+$0.034]}
\newcommand{\PreNetSiteGRate}{$+$0.019}
\newcommand{\PreNetSiteGRateBand}{[$-$0.006, $+$0.045]}
\newcommand{\PreNetSiteGSearch}{$-$1.1}
\newcommand{\PreNetSiteGSearchBand}{[$-$1.7, $-$0.5]}
\newcommand{\PreNetSiteGSink}{$-$6}
\newcommand{\PreNetSiteGSinkBand}{[$-$8, $-$4]}
\newcommand{\PreNetSiteGXc}{$+$1.4}
\newcommand{\PreNetSiteGXcBand}{[$+$0.8, $+$1.9]}
\newcommand{\PreOrigins}{56}
\newcommand{\PreOutClimb}{$+$3.6}
\newcommand{\PreOutClimbBand}{[$+$2.7, $+$4.6]}
\newcommand{\PreOutCph}{$+$0.58}
\newcommand{\PreOutCphBand}{[$+$0.39, $+$0.76]}
\newcommand{\PreOutFound}{$+$3.3}
\newcommand{\PreOutFoundBand}{[$+$1.1, $+$5.5]}
\newcommand{\PreOutGlide}{$+$1.4}
\newcommand{\PreOutGlideBand}{[$+$1.0, $+$1.8]}
\newcommand{\PreOutHK}{$+$0.027}
\newcommand{\PreOutHKBand}{[$+$0.019, $+$0.035]}
\newcommand{\PreOutRate}{$+$0.026}
\newcommand{\PreOutRateBand}{[$-$0.002, $+$0.054]}
\newcommand{\PreOutSearch}{$-$1.3}
\newcommand{\PreOutSearchBand}{[$-$2.0, $-$0.7]}
\newcommand{\PreOutSink}{$-$9}
\newcommand{\PreOutSinkBand}{[$-$11, $-$7]}
\newcommand{\PreOutXc}{$+$1.3}
\newcommand{\PreOutXcBand}{[$+$0.9, $+$1.8]}
\newcommand{\PrePPClimb}{$-$2.2}
\newcommand{\PrePPClimbBand}{[$-$3.0, $-$1.4]}
\newcommand{\PrePPCph}{$-$0.61}
\newcommand{\PrePPCphBand}{[$-$0.78, $-$0.44]}
\newcommand{\PrePPFound}{$-$2.5}
\newcommand{\PrePPFoundBand}{[$-$4.3, $-$0.6]}
\newcommand{\PrePPGlide}{$-$0.2}
\newcommand{\PrePPGlideBand}{[$-$0.7, $+$0.3]}
\newcommand{\PrePPHK}{$+$0.002}
\newcommand{\PrePPHKBand}{[$-$0.006, $+$0.009]}
\newcommand{\PrePPRate}{$+$0.023}
\newcommand{\PrePPRateBand}{[$-$0.005, $+$0.051]}
\newcommand{\PrePPSearch}{$+$0.3}
\newcommand{\PrePPSearchBand}{[$-$0.4, $+$1.0]}
\newcommand{\PrePPSink}{$+$9}
\newcommand{\PrePPSinkBand}{[$+$6, $+$11]}
\newcommand{\PrePPXc}{$+$0.6}
\newcommand{\PrePPXcBand}{[$+$0.1, $+$1.2]}
\newcommand{\PreVsGapClimb}{$-$0.2}
\newcommand{\PreVsGapCph}{$-$0.37}
\newcommand{\PreVsGapFound}{$+$0.4}
\newcommand{\PreVsGapGlide}{$+$2.2}
\newcommand{\PreVsGapHK}{$+$0.016}
\newcommand{\PreVsGapRate}{$+$0.085}
\newcommand{\PreVsGapSearch}{$-$1.7}
\newcommand{\PreVsGapSink}{$-$11}
\newcommand{\PreVsInsideClimb}{yes}
\newcommand{\PreVsInsideCph}{no}
\newcommand{\PreVsInsideFound}{yes}
\newcommand{\PreVsInsideGlide}{no}
\newcommand{\PreVsInsideHK}{no}
\newcommand{\PreVsInsideRate}{no}
\newcommand{\PreVsInsideSearch}{no}
\newcommand{\PreVsInsideSink}{no}
\newcommand{\PreVsSignClimb}{yes}
\newcommand{\PreVsSignCph}{yes}
\newcommand{\PreVsSignFound}{yes}
\newcommand{\PreVsSignGlide}{no}
\newcommand{\PreVsSignHK}{yes}
\newcommand{\PreVsSignRate}{no}
\newcommand{\PreVsSignSearch}{no}
\newcommand{\PreVsSignSink}{no}
\newcommand{\PreWinAClimb}{$+$3.8}
\newcommand{\PreWinAClimbBand}{[$+$3.0, $+$4.6]}
\newcommand{\PreWinADIn}{$-$0.010}
\newcommand{\PreWinADOut}{$+$0.008}
\newcommand{\PreWinAEvents}{3\,549}
\newcommand{\PreWinAHK}{$-$0.001}
\newcommand{\PreWinAHKBand}{[$-$0.006, $+$0.004]}
\newcommand{\PreWinALen}{1\,260}
\newcommand{\PreWinAPeak}{$+$0.017}
\newcommand{\PreWinAPeakBand}{[$+$0.013, $+$0.021]}
\newcommand{\PreWinAPeakLag}{390}
\newcommand{\PreWinBClimb}{$+$2.5}
\newcommand{\PreWinBClimbBand}{[$+$1.7, $+$3.2]}
\newcommand{\PreWinBDIn}{$+$0.020}
\newcommand{\PreWinBDOut}{$+$0.026}
\newcommand{\PreWinBEvents}{2\,567}
\newcommand{\PreWinBHK}{$+$0.023}
\newcommand{\PreWinBHKBand}{[$+$0.019, $+$0.027]}
\newcommand{\PreWinBLen}{1\,800}
\newcommand{\PreWinBPeak}{$+$0.024}
\newcommand{\PreWinBPeakBand}{[$+$0.020, $+$0.028]}
\newcommand{\PreWinBPeakLag}{580}
\newcommand{\PreWinCClimb}{$+$2.8}
\newcommand{\PreWinCClimbBand}{[$+$2.0, $+$3.6]}
\newcommand{\PreWinCDIn}{$+$0.025}
\newcommand{\PreWinCDOut}{$+$0.027}
\newcommand{\PreWinCEvents}{1\,595}
\newcommand{\PreWinCHK}{$+$0.026}
\newcommand{\PreWinCHKBand}{[$+$0.020, $+$0.032]}
\newcommand{\PreWinCHQ}{$+$0.007}
\newcommand{\PreWinCHQBand}{[$-$0.002, $+$0.016]}
\newcommand{\PreWinCLen}{2\,400}
\newcommand{\PreWinCPeak}{$+$0.025}
\newcommand{\PreWinCPeakBand}{[$+$0.019, $+$0.031]}
\newcommand{\PreWinCPeakLag}{710}
\newcommand{\PreWinDClimb}{$+$2.8}
\newcommand{\PreWinDClimbBand}{[$+$1.8, $+$3.7]}
\newcommand{\PreWinDDIn}{$+$0.020}
\newcommand{\PreWinDDOut}{$+$0.030}
\newcommand{\PreWinDEvents}{695}
\newcommand{\PreWinDHK}{$+$0.025}
\newcommand{\PreWinDHKBand}{[$+$0.018, $+$0.032]}
\newcommand{\PreWinDHQ}{$+$0.023}
\newcommand{\PreWinDHQBand}{[$+$0.015, $+$0.031]}
\newcommand{\PreWinDLen}{3\,000}
\newcommand{\PreWinDPeak}{$+$0.024}
\newcommand{\PreWinDPeakBand}{[$+$0.017, $+$0.031]}
\newcommand{\PreWinDPeakLag}{710}
\newcommand{\PreWindow}{2\,400}
\newcommand{\PreWindowsEnd}{1\,595}
\newcommand{\PreWindowsPost}{1\,595}
\newcommand{\PreWindowsPre}{1\,595}
\newcommand{\PreWindowsStart}{1\,595}
\newcommand{\RawClimbMedianA}{137}
\newcommand{\RawClimbMedianG}{142}
\newcommand{\RawClimbMedianS}{140}
\newcommand{\RawFirstLag}{10}
\newcommand{\RawFitLongA}{0.929}
\newcommand{\RawFitLongABand}{[0.924, 0.935]}
\newcommand{\RawFitLongG}{0.943}
\newcommand{\RawFitLongGBand}{[0.930, 0.956]}
\newcommand{\RawFitLongS}{0.923}
\newcommand{\RawFitLongSBand}{[0.920, 0.927]}
\newcommand{\RawFitOneA}{0.923}
\newcommand{\RawFitOneABand}{[0.922, 0.925]}
\newcommand{\RawFitOneG}{0.924}
\newcommand{\RawFitOneGBand}{[0.922, 0.925]}
\newcommand{\RawFitOneS}{0.921}
\newcommand{\RawFitOneSBand}{[0.920, 0.922]}
\newcommand{\RawFitTwoA}{0.920}
\newcommand{\RawFitTwoABand}{[0.918, 0.922]}
\newcommand{\RawFitTwoG}{0.924}
\newcommand{\RawFitTwoGBand}{[0.922, 0.926]}
\newcommand{\RawFitTwoS}{0.915}
\newcommand{\RawFitTwoSBand}{[0.914, 0.916]}
\newcommand{\RawFlightsLongA}{602}
\newcommand{\RawFlightsLongG}{115}
\newcommand{\RawFlightsLongS}{4\,309}
\newcommand{\RawFlightsOneA}{8\,991}
\newcommand{\RawFlightsOneG}{11\,077}
\newcommand{\RawFlightsOneS}{43\,977}
\newcommand{\RawFlightsTwoA}{6\,235}
\newcommand{\RawFlightsTwoG}{6\,297}
\newcommand{\RawFlightsTwoS}{33\,613}
\newcommand{\RawGapClimb}{$+$3.3}
\newcommand{\RawGapClimbAS}{$-$1.4}
\newcommand{\RawGapClimbASBand}{[$-$1.9, $-$0.9]}
\newcommand{\RawGapClimbBand}{[$+$2.8, $+$3.8]}
\newcommand{\RawGapClimbGA}{$+$4.7}
\newcommand{\RawGapClimbGABand}{[$+$4.4, $+$5.0]}
\newcommand{\RawGapCph}{$+$0.72}
\newcommand{\RawGapCphAS}{$-$0.17}
\newcommand{\RawGapCphASBand}{[$-$0.22, $-$0.12]}
\newcommand{\RawGapCphBand}{[$+$0.65, $+$0.78]}
\newcommand{\RawGapCphGA}{$+$0.89}
\newcommand{\RawGapCphGABand}{[$+$0.82, $+$0.95]}
\newcommand{\RawGapCross}{170}
\newcommand{\RawGapFirst}{$-$0.006}
\newcommand{\RawGapFirstAS}{$+$0.005}
\newcommand{\RawGapFirstASBand}{[$+$0.003, $+$0.006]}
\newcommand{\RawGapFirstBand}{[$-$0.007, $-$0.004]}
\newcommand{\RawGapFirstGA}{$-$0.011}
\newcommand{\RawGapFirstGABand}{[$-$0.012, $-$0.009]}
\newcommand{\RawGapFound}{$+$6.0}
\newcommand{\RawGapFoundAS}{$+$1.6}
\newcommand{\RawGapFoundASBand}{[$+$1.0, $+$2.1]}
\newcommand{\RawGapFoundBand}{[$+$5.5, $+$6.6]}
\newcommand{\RawGapFoundGA}{$+$4.5}
\newcommand{\RawGapFoundGABand}{[$+$3.8, $+$5.1]}
\newcommand{\RawGapGlide}{$+$2.1}
\newcommand{\RawGapGlideAS}{$+$0.7}
\newcommand{\RawGapGlideASBand}{[$+$0.4, $+$1.1]}
\newcommand{\RawGapGlideBand}{[$+$1.7, $+$2.5]}
\newcommand{\RawGapGlideGA}{$+$1.3}
\newcommand{\RawGapGlideGABand}{[$+$1.1, $+$1.6]}
\newcommand{\RawGapHK}{$+$0.018}
\newcommand{\RawGapHKAS}{$-$0.003}
\newcommand{\RawGapHKASBand}{[$-$0.006, $-$0.000]}
\newcommand{\RawGapHKBand}{[$+$0.015, $+$0.021]}
\newcommand{\RawGapHKGA}{$+$0.021}
\newcommand{\RawGapHKGABand}{[$+$0.019, $+$0.024]}
\newcommand{\RawGapHQ}{$+$0.017}
\newcommand{\RawGapHQAS}{$-$0.003}
\newcommand{\RawGapHQASBand}{[$-$0.007, $+$0.001]}
\newcommand{\RawGapHQBand}{[$+$0.013, $+$0.021]}
\newcommand{\RawGapHQGA}{$+$0.020}
\newcommand{\RawGapHQGABand}{[$+$0.016, $+$0.024]}
\newcommand{\RawGapRate}{$+$0.044}
\newcommand{\RawGapRateAS}{$+$0.052}
\newcommand{\RawGapRateASBand}{[$+$0.039, $+$0.065]}
\newcommand{\RawGapRateBand}{[$+$0.030, $+$0.057]}
\newcommand{\RawGapRateGA}{$-$0.009}
\newcommand{\RawGapRateGABand}{[$-$0.017, $-$0.001]}
\newcommand{\RawGapSearch}{$-$0.9}
\newcommand{\RawGapSearchAS}{$+$0.2}
\newcommand{\RawGapSearchASBand}{[$-$0.1, $+$0.6]}
\newcommand{\RawGapSearchBand}{[$-$1.2, $-$0.6]}
\newcommand{\RawGapSearchGA}{$-$1.1}
\newcommand{\RawGapSearchGABand}{[$-$1.4, $-$0.8]}
\newcommand{\RawGapSig}{210}
\newcommand{\RawGapSink}{$-$18}
\newcommand{\RawGapSinkAS}{$-$6}
\newcommand{\RawGapSinkASBand}{[$-$7, $-$5]}
\newcommand{\RawGapSinkBand}{[$-$19, $-$17]}
\newcommand{\RawGapSinkGA}{$-$12}
\newcommand{\RawGapSinkGABand}{[$-$13, $-$11]}
\newcommand{\RawGapXc}{$+$1.4}
\newcommand{\RawGapXcAS}{$+$1.0}
\newcommand{\RawGapXcASBand}{[$+$0.6, $+$1.3]}
\newcommand{\RawGapXcBand}{[$+$1.0, $+$1.7]}
\newcommand{\RawGapXcGA}{$+$0.4}
\newcommand{\RawGapXcGABand}{[$+$0.1, $+$0.7]}
\newcommand{\RawGlideMedianA}{205}
\newcommand{\RawGlideMedianG}{188}
\newcommand{\RawGlideMedianS}{194}
\newcommand{\RawHKCutLongS}{0.909}
\newcommand{\RawHKCutTwoS}{0.895}
\newcommand{\RawHKLongA}{0.917}
\newcommand{\RawHKLongABand}{[0.910, 0.924]}
\newcommand{\RawHKLongG}{0.936}
\newcommand{\RawHKLongGBand}{[0.920, 0.952]}
\newcommand{\RawHKLongS}{0.912}
\newcommand{\RawHKLongSBand}{[0.908, 0.916]}
\newcommand{\RawHKOneA}{0.894}
\newcommand{\RawHKOneABand}{[0.891, 0.898]}
\newcommand{\RawHKOneG}{0.915}
\newcommand{\RawHKOneGBand}{[0.913, 0.918]}
\newcommand{\RawHKOneS}{0.897}
\newcommand{\RawHKOneSBand}{[0.895, 0.899]}
\newcommand{\RawHKTwoA}{0.898}
\newcommand{\RawHKTwoABand}{[0.894, 0.902]}
\newcommand{\RawHKTwoG}{0.918}
\newcommand{\RawHKTwoGBand}{[0.915, 0.922]}
\newcommand{\RawHKTwoS}{0.898}
\newcommand{\RawHKTwoSBand}{[0.896, 0.900]}
\newcommand{\RawHMeanGap}{0.002}
\newcommand{\RawHQLongA}{0.926}
\newcommand{\RawHQLongABand}{[0.919, 0.934]}
\newcommand{\RawHQLongG}{0.948}
\newcommand{\RawHQLongGBand}{[0.932, 0.965]}
\newcommand{\RawHQLongS}{0.922}
\newcommand{\RawHQLongSBand}{[0.917, 0.927]}
\newcommand{\RawHQTwoA}{0.903}
\newcommand{\RawHQTwoABand}{[0.898, 0.908]}
\newcommand{\RawHQTwoG}{0.924}
\newcommand{\RawHQTwoGBand}{[0.919, 0.928]}
\newcommand{\RawHQTwoS}{0.907}
\newcommand{\RawHQTwoSBand}{[0.904, 0.910]}
\newcommand{\RawHlocRangeHigh}{0.049}
\newcommand{\RawHlocRangeLong}{0.085}
\newcommand{\RawHlocRangeLow}{0.043}
\newcommand{\RawNonZeroCount}{7}
\newcommand{\RawPhaseFlightsA}{6\,397}
\newcommand{\RawPhaseFlightsG}{7\,896}
\newcommand{\RawPhaseFlightsS}{27\,665}
\newcommand{\RawPhaseHoursA}{13\,898}
\newcommand{\RawPhaseHoursG}{12\,338}
\newcommand{\RawPhaseHoursS}{60\,882}
\newcommand{\RawPiecesLongA}{602}
\newcommand{\RawPiecesLongG}{115}
\newcommand{\RawPiecesLongS}{4\,311}
\newcommand{\RawPiecesOneA}{13\,464}
\newcommand{\RawPiecesOneG}{14\,401}
\newcommand{\RawPiecesOneS}{46\,239}
\newcommand{\RawPiecesTwoA}{7\,507}
\newcommand{\RawPiecesTwoG}{7\,243}
\newcommand{\RawPiecesTwoS}{34\,269}
\newcommand{\RawRunsClimbA}{87\,378}
\newcommand{\RawRunsClimbG}{87\,000}
\newcommand{\RawRunsClimbS}{397\,341}
\newcommand{\RawRunsGlideA}{93\,424}
\newcommand{\RawRunsGlideG}{86\,766}
\newcommand{\RawRunsGlideS}{428\,324}
\newcommand{\RawRunsSearchA}{29\,692}
\newcommand{\RawRunsSearchG}{28\,476}
\newcommand{\RawRunsSearchS}{132\,029}
\newcommand{\RawShortestLong}{14\,312.5}
\newcommand{\RawShortestOne}{1\,250}
\newcommand{\RawShortestTwo}{3\,450}
\newcommand{\RawStatClimbA}{31.6}
\newcommand{\RawStatClimbABand}{[31.0, 32.1]}
\newcommand{\RawStatClimbG}{36.3}
\newcommand{\RawStatClimbGBand}{[35.7, 36.8]}
\newcommand{\RawStatClimbS}{32.9}
\newcommand{\RawStatClimbSBand}{[32.6, 33.3]}
\newcommand{\RawStatCphA}{6.36}
\newcommand{\RawStatCphABand}{[6.30, 6.41]}
\newcommand{\RawStatCphG}{7.24}
\newcommand{\RawStatCphGBand}{[7.18, 7.31]}
\newcommand{\RawStatCphS}{6.53}
\newcommand{\RawStatCphSBand}{[6.50, 6.56]}
\newcommand{\RawStatFoundA}{42.1}
\newcommand{\RawStatFoundABand}{[41.6, 42.6]}
\newcommand{\RawStatFoundG}{46.6}
\newcommand{\RawStatFoundGBand}{[46.0, 47.1]}
\newcommand{\RawStatFoundS}{40.5}
\newcommand{\RawStatFoundSBand}{[40.2, 40.9]}
\newcommand{\RawStatGlideA}{33.9}
\newcommand{\RawStatGlideABand}{[33.5, 34.3]}
\newcommand{\RawStatGlideG}{35.2}
\newcommand{\RawStatGlideGBand}{[34.8, 35.7]}
\newcommand{\RawStatGlideS}{33.2}
\newcommand{\RawStatGlideSBand}{[32.9, 33.4]}
\newcommand{\RawStatRateA}{1.15}
\newcommand{\RawStatRateABand}{[1.14, 1.17]}
\newcommand{\RawStatRateG}{1.14}
\newcommand{\RawStatRateGBand}{[1.13, 1.16]}
\newcommand{\RawStatRateS}{1.10}
\newcommand{\RawStatRateSBand}{[1.09, 1.11]}
\newcommand{\RawStatSearchA}{11.6}
\newcommand{\RawStatSearchABand}{[11.3, 12.0]}
\newcommand{\RawStatSearchG}{10.5}
\newcommand{\RawStatSearchGBand}{[10.2, 10.8]}
\newcommand{\RawStatSearchS}{11.4}
\newcommand{\RawStatSearchSBand}{[11.2, 11.6]}
\newcommand{\RawStatSinkA}{94}
\newcommand{\RawStatSinkABand}{[93, 96]}
\newcommand{\RawStatSinkG}{82}
\newcommand{\RawStatSinkGBand}{[81, 84]}
\newcommand{\RawStatSinkS}{100}
\newcommand{\RawStatSinkSBand}{[100, 101]}
\newcommand{\RawStatXcA}{21.5}
\newcommand{\RawStatXcABand}{[21.0, 21.9]}
\newcommand{\RawStatXcG}{21.9}
\newcommand{\RawStatXcGBand}{[21.5, 22.3]}
\newcommand{\RawStatXcS}{20.5}
\newcommand{\RawStatXcSBand}{[20.2, 20.7]}
\newcommand{\RawSumClimbA}{32.0}
\newcommand{\RawSumClimbG}{36.4}
\newcommand{\RawSumClimbS}{34.2}
\newcommand{\RawSumCphA}{6.35}
\newcommand{\RawSumCphG}{7.12}
\newcommand{\RawSumCphS}{6.59}
\newcommand{\RawSumFoundA}{44.8}
\newcommand{\RawSumFoundG}{47.2}
\newcommand{\RawSumFoundS}{43.1}
\newcommand{\RawSumGlideA}{34.9}
\newcommand{\RawSumGlideG}{36.1}
\newcommand{\RawSumGlideS}{34.2}
\newcommand{\RawSumRateA}{1.23}
\newcommand{\RawSumRateG}{1.16}
\newcommand{\RawSumRateS}{1.16}
\newcommand{\RawSumSearchA}{10.3}
\newcommand{\RawSumSearchG}{9.9}
\newcommand{\RawSumSearchS}{10.3}
\newcommand{\RawSumSinkA}{89}
\newcommand{\RawSumSinkG}{82}
\newcommand{\RawSumSinkS}{96}
\newcommand{\RawSumXcA}{22.3}
\newcommand{\RawSumXcG}{22.7}
\newcommand{\RawSumXcS}{21.3}
\newcommand{\RawTauLong}{11\,450}
\newcommand{\RawTauOne}{1\,000}
\newcommand{\RawTauTwo}{2\,760}
\newcommand{\RawXcFlightsA}{5\,463}
\newcommand{\RawXcFlightsG}{7\,823}
\newcommand{\RawXcFlightsS}{26\,336}
\newcommand{\RefitAgreeHang}{98.1}
\newcommand{\RefitAgreeHighMtn}{97.9}
\newcommand{\RefitAgreeHills}{97.1}
\newcommand{\RefitAgreeLowMtn}{97.7}
\newcommand{\RefitAgreePara}{97.7}
\newcommand{\RefitAgreePlains}{96.3}
\newcommand{\RefitCadenceHangFive}{98.2}
\newcommand{\RefitCadenceParaFive}{97.3}
\newcommand{\RefitCadenceParaTwo}{98.9}
\newcommand{\RefitFlightsHang}{1\,666}
\newcommand{\RefitFlightsPara}{104\,170}
\newcommand{\RefitHangVsPara}{99.999}
\newcommand{\RefitHeldHang}{349}
\newcommand{\RefitHeldPara}{377}
\newcommand{\RefitIterMax}{19}
\newcommand{\RefitLLBestParaFix}{$-$0.702}
\newcommand{\RefitLLHisParaFix}{$-$0.890}
\newcommand{\RefitLLOursParaFix}{$-$0.757}
\newcommand{\RefitPoolPara}{120\,496}
\newcommand{\RefitSearchPerHourHis}{5.35}
\newcommand{\RefitSearchPerHourOurs}{6.19}
\newcommand{\RefitSearchShareHis}{11.7}
\newcommand{\RefitSearchShareOurs}{13.4}
\newcommand{\RefitStarts}{21}
\newcommand{\RefitTrainHang}{1\,396}
\newcommand{\RefitTrainPara}{1\,510}
\newcommand{\RouteASAllThousand}{$+$0.001}
\newcommand{\RouteASAllThousandBand}{[$-$0.002, $+$0.004]}
\newcommand{\RouteASAllThreeTwenty}{$+$0.006}
\newcommand{\RouteASAllThreeTwentyBand}{[$+$0.004, $+$0.007]}
\newcommand{\RouteASHighMtnThousand}{$+$0.029}
\newcommand{\RouteASHighMtnThousandBand}{[$+$0.024, $+$0.034]}
\newcommand{\RouteASHighMtnThreeTwenty}{$+$0.016}
\newcommand{\RouteASHighMtnThreeTwentyBand}{[$+$0.014, $+$0.018]}
\newcommand{\RouteASHillsThousand}{$+$0.014}
\newcommand{\RouteASHillsThousandBand}{[$+$0.008, $+$0.019]}
\newcommand{\RouteASHillsThreeTwenty}{$+$0.010}
\newcommand{\RouteASHillsThreeTwentyBand}{[$+$0.007, $+$0.014]}
\newcommand{\RouteASLowMtnThousand}{$+$0.016}
\newcommand{\RouteASLowMtnThousandBand}{[$+$0.012, $+$0.020]}
\newcommand{\RouteASLowMtnThreeTwenty}{$+$0.013}
\newcommand{\RouteASLowMtnThreeTwentyBand}{[$+$0.012, $+$0.015]}
\newcommand{\RouteASPlainsThousand}{$+$0.005}
\newcommand{\RouteASPlainsThousandBand}{[$+$0.003, $+$0.007]}
\newcommand{\RouteASPlainsThreeTwenty}{$+$0.002}
\newcommand{\RouteASPlainsThreeTwentyBand}{[$+$0.001, $+$0.004]}
\newcommand{\RouteCheckGAMax}{0.005}
\newcommand{\RouteCheckMax}{0.012}
\newcommand{\RouteDs}{100}
\newcommand{\RouteGAAllThousand}{$+$0.023}
\newcommand{\RouteGAAllThousandBand}{[$+$0.020, $+$0.025]}
\newcommand{\RouteGAAllThreeTwenty}{$+$0.012}
\newcommand{\RouteGAAllThreeTwentyBand}{[$+$0.011, $+$0.013]}
\newcommand{\RouteGAHighMtnThousand}{$+$0.028}
\newcommand{\RouteGAHighMtnThousandBand}{[$+$0.023, $+$0.032]}
\newcommand{\RouteGAHighMtnThreeTwenty}{$+$0.015}
\newcommand{\RouteGAHighMtnThreeTwentyBand}{[$+$0.013, $+$0.017]}
\newcommand{\RouteGAHillsThousand}{$+$0.012}
\newcommand{\RouteGAHillsThousandBand}{[$+$0.007, $+$0.017]}
\newcommand{\RouteGAHillsThreeTwenty}{$+$0.006}
\newcommand{\RouteGAHillsThreeTwentyBand}{[$+$0.003, $+$0.009]}
\newcommand{\RouteGALowMtnThousand}{$+$0.022}
\newcommand{\RouteGALowMtnThousandBand}{[$+$0.019, $+$0.026]}
\newcommand{\RouteGALowMtnThreeTwenty}{$+$0.013}
\newcommand{\RouteGALowMtnThreeTwentyBand}{[$+$0.012, $+$0.015]}
\newcommand{\RouteGAPlainsThousand}{$+$0.004}
\newcommand{\RouteGAPlainsThousandBand}{[$+$0.003, $+$0.006]}
\newcommand{\RouteGAPlainsThreeTwenty}{$+$0.002}
\newcommand{\RouteGAPlainsThreeTwentyBand}{[$+$0.000, $+$0.003]}
\newcommand{\RouteGSAllThousand}{$+$0.023}
\newcommand{\RouteGSAllThousandBand}{[$+$0.020, $+$0.026]}
\newcommand{\RouteGSAllThreeTwenty}{$+$0.018}
\newcommand{\RouteGSAllThreeTwentyBand}{[$+$0.017, $+$0.020]}
\newcommand{\RouteGSHighMtnThousand}{$+$0.057}
\newcommand{\RouteGSHighMtnThousandBand}{[$+$0.052, $+$0.061]}
\newcommand{\RouteGSHighMtnThreeTwenty}{$+$0.031}
\newcommand{\RouteGSHighMtnThreeTwentyBand}{[$+$0.029, $+$0.033]}
\newcommand{\RouteGSHillsThousand}{$+$0.026}
\newcommand{\RouteGSHillsThousandBand}{[$+$0.021, $+$0.031]}
\newcommand{\RouteGSHillsThreeTwenty}{$+$0.016}
\newcommand{\RouteGSHillsThreeTwentyBand}{[$+$0.013, $+$0.019]}
\newcommand{\RouteGSLowMtnThousand}{$+$0.039}
\newcommand{\RouteGSLowMtnThousandBand}{[$+$0.034, $+$0.043]}
\newcommand{\RouteGSLowMtnThreeTwenty}{$+$0.027}
\newcommand{\RouteGSLowMtnThreeTwentyBand}{[$+$0.025, $+$0.029]}
\newcommand{\RouteGSPlainsThousand}{$+$0.009}
\newcommand{\RouteGSPlainsThousandBand}{[$+$0.007, $+$0.011]}
\newcommand{\RouteGSPlainsThreeTwenty}{$+$0.004}
\newcommand{\RouteGSPlainsThreeTwentyBand}{[$+$0.002, $+$0.006]}
\newcommand{\RouteKeepHighMtn}{84}
\newcommand{\RouteKeepHills}{94}
\newcommand{\RouteKeepLowMtn}{73}
\newcommand{\RouteKeepMax}{94}
\newcommand{\RouteKeepMin}{55}
\newcommand{\RouteKeepPlains}{55}
\newcommand{\RouteLagThousand}{870}
\newcommand{\RouteLagThreeTwenty}{260}
\newcommand{\RouteMinusASAllThousand}{$-$0.005}
\newcommand{\RouteMinusASAllThousandBand}{[$-$0.006, $-$0.003]}
\newcommand{\RouteMinusASAllThreeTwenty}{$-$0.003}
\newcommand{\RouteMinusASAllThreeTwentyBand}{[$-$0.004, $-$0.002]}
\newcommand{\RouteMinusASHighMtnThousand}{$-$0.006}
\newcommand{\RouteMinusASHighMtnThousandBand}{[$-$0.008, $-$0.004]}
\newcommand{\RouteMinusASHighMtnThreeTwenty}{$-$0.004}
\newcommand{\RouteMinusASHighMtnThreeTwentyBand}{[$-$0.005, $-$0.002]}
\newcommand{\RouteMinusASHillsThousand}{$-$0.001}
\newcommand{\RouteMinusASHillsThousandBand}{[$-$0.005, $+$0.002]}
\newcommand{\RouteMinusASHillsThreeTwenty}{$+$0.000}
\newcommand{\RouteMinusASHillsThreeTwentyBand}{[$-$0.003, $+$0.003]}
\newcommand{\RouteMinusASLowMtnThousand}{$-$0.005}
\newcommand{\RouteMinusASLowMtnThousandBand}{[$-$0.007, $-$0.003]}
\newcommand{\RouteMinusASLowMtnThreeTwenty}{$-$0.001}
\newcommand{\RouteMinusASLowMtnThreeTwentyBand}{[$-$0.003, $-$0.000]}
\newcommand{\RouteMinusASPlainsThousand}{$+$0.000}
\newcommand{\RouteMinusASPlainsThousandBand}{[$-$0.002, $+$0.002]}
\newcommand{\RouteMinusASPlainsThreeTwenty}{$+$0.002}
\newcommand{\RouteMinusASPlainsThreeTwentyBand}{[$+$0.000, $+$0.004]}
\newcommand{\RouteMinusGAAllThousand}{$+$0.005}
\newcommand{\RouteMinusGAAllThousandBand}{[$+$0.004, $+$0.006]}
\newcommand{\RouteMinusGAAllThreeTwenty}{$-$0.001}
\newcommand{\RouteMinusGAAllThreeTwentyBand}{[$-$0.002, $-$0.000]}
\newcommand{\RouteMinusGAHighMtnThousand}{$+$0.005}
\newcommand{\RouteMinusGAHighMtnThousandBand}{[$+$0.003, $+$0.008]}
\newcommand{\RouteMinusGAHighMtnThreeTwenty}{$-$0.000}
\newcommand{\RouteMinusGAHighMtnThreeTwentyBand}{[$-$0.002, $+$0.002]}
\newcommand{\RouteMinusGAHillsThousand}{$+$0.001}
\newcommand{\RouteMinusGAHillsThousandBand}{[$-$0.003, $+$0.005]}
\newcommand{\RouteMinusGAHillsThreeTwenty}{$-$0.005}
\newcommand{\RouteMinusGAHillsThreeTwentyBand}{[$-$0.008, $-$0.002]}
\newcommand{\RouteMinusGALowMtnThousand}{$+$0.008}
\newcommand{\RouteMinusGALowMtnThousandBand}{[$+$0.006, $+$0.010]}
\newcommand{\RouteMinusGALowMtnThreeTwenty}{$+$0.000}
\newcommand{\RouteMinusGALowMtnThreeTwentyBand}{[$-$0.001, $+$0.001]}
\newcommand{\RouteMinusGAPlainsThousand}{$+$0.004}
\newcommand{\RouteMinusGAPlainsThousandBand}{[$+$0.001, $+$0.006]}
\newcommand{\RouteMinusGAPlainsThreeTwenty}{$-$0.003}
\newcommand{\RouteMinusGAPlainsThreeTwentyBand}{[$-$0.005, $-$0.001]}
\newcommand{\RouteMinusGSAllThousand}{$+$0.000}
\newcommand{\RouteMinusGSAllThousandBand}{[$-$0.001, $+$0.002]}
\newcommand{\RouteMinusGSAllThreeTwenty}{$-$0.004}
\newcommand{\RouteMinusGSAllThreeTwentyBand}{[$-$0.005, $-$0.003]}
\newcommand{\RouteMinusGSHighMtnThousand}{$-$0.001}
\newcommand{\RouteMinusGSHighMtnThousandBand}{[$-$0.003, $+$0.002]}
\newcommand{\RouteMinusGSHighMtnThreeTwenty}{$-$0.004}
\newcommand{\RouteMinusGSHighMtnThreeTwentyBand}{[$-$0.005, $-$0.002]}
\newcommand{\RouteMinusGSHillsThousand}{$-$0.000}
\newcommand{\RouteMinusGSHillsThousandBand}{[$-$0.004, $+$0.003]}
\newcommand{\RouteMinusGSHillsThreeTwenty}{$-$0.005}
\newcommand{\RouteMinusGSHillsThreeTwentyBand}{[$-$0.007, $-$0.002]}
\newcommand{\RouteMinusGSLowMtnThousand}{$+$0.003}
\newcommand{\RouteMinusGSLowMtnThousandBand}{[$+$0.001, $+$0.005]}
\newcommand{\RouteMinusGSLowMtnThreeTwenty}{$-$0.001}
\newcommand{\RouteMinusGSLowMtnThreeTwentyBand}{[$-$0.003, $-$0.000]}
\newcommand{\RouteMinusGSPlainsThousand}{$+$0.004}
\newcommand{\RouteMinusGSPlainsThousandBand}{[$+$0.001, $+$0.006]}
\newcommand{\RouteMinusGSPlainsThreeTwenty}{$-$0.001}
\newcommand{\RouteMinusGSPlainsThreeTwentyBand}{[$-$0.003, $+$0.001]}
\newcommand{\RoutePathSpeed}{9.9}
\newcommand{\RouteSpeed}{7.9}
\newcommand{\RouteTimeLagThousand}{1\,098}
\newcommand{\RouteTimeLagThreeTwenty}{328}
\newcommand{\RouteWaitHighMtn}{16}
\newcommand{\RouteWaitHills}{6}
\newcommand{\RouteWaitLowMtn}{27}
\newcommand{\RouteWaitPlains}{45}
\newcommand{\StrictShare}{26}
\newcommand{\Stricter}{0.5}
\newcommand{\Support}{0.8}
\newcommand{\SupportOriginsPercent}{20}
\newcommand{\TailAboveMax}{522}
\newcommand{\TailAboveMin}{15}
\newcommand{\TailBinPercent}{1}
\newcommand{\TailCandidates}{80}
\newcommand{\TailFineBelow}{200}
\newcommand{\TailMinRuns}{50}
\newcommand{\TailMuMax}{16.3}
\newcommand{\TailMuMin}{3.2}
\newcommand{\TailQuantileHi}{99.99}
\newcommand{\TailQuantileLo}{50}
\newcommand{\TakeoffHighFlights}{34\,333}
\newcommand{\TakeoffHighKeep}{80}
\newcommand{\TakeoffHighKeepN}{27\,415}
\newcommand{\TakeoffHighMedian}{0.24}
\newcommand{\TakeoffHighStrict}{13}
\newcommand{\TakeoffHighStrictN}{4\,307}
\newcommand{\TakeoffLowFlights}{29\,958}
\newcommand{\TakeoffLowKeep}{87}
\newcommand{\TakeoffLowKeepN}{26\,201}
\newcommand{\TakeoffLowMedian}{0.45}
\newcommand{\TakeoffLowStrict}{42}
\newcommand{\TakeoffLowStrictN}{12\,511}
\newcommand{\ToyAvgFirst}{12.6}
\newcommand{\ToyCgmB}{2}
\newcommand{\ToyCgmC}{6}
\newcommand{\ToyCgmP}{8}
\newcommand{\ToyCgmPig}{16}
\newcommand{\ToyCgmV}{4}
\newcommand{\ToyDa}{10}
\newcommand{\ToyDb}{20}
\newcommand{\ToyDc}{30}
\newcommand{\ToyDin}{3}
\newcommand{\ToyDout}{1}
\newcommand{\ToyDw}{1}
\newcommand{\ToyG}{2}
\newcommand{\ToyH}{0.54}
\newcommand{\ToyMeanShare}{13.9}
\newcommand{\ToyMsdOne}{15.33}
\newcommand{\ToyMsdTwo}{32.5}
\newcommand{\ToyNa}{1}
\newcommand{\ToyNb}{5}
\newcommand{\ToyNc}{2}
\newcommand{\ToyPlAlpha}{0.9}
\newcommand{\ToyPlBeta}{0.3}
\newcommand{\ToyPlChained}{0.900}
\newcommand{\ToyPlCohorts}{2\,000, 5\,000, 10\,000, 20\,000, 50\,000}
\newcommand{\ToyPlJump}{0.959}
\newcommand{\ToyPlMean}{0.929}
\newcommand{\ToyPlMeanMax}{0.954}
\newcommand{\ToyPlPerlag}{0.900}
\newcommand{\ToyPlPool}{0.967}
\newcommand{\ToyPlTau}{10\,000}
\newcommand{\ToyPlTmax}{100\,000}
\newcommand{\ToyPlTop}{0.975}
\newcommand{\ToyPlTzero}{10\,000}
\newcommand{\ToyPos}{0, 3, 4, 10}
\newcommand{\ToyRatio}{13.3}
\newcommand{\ToyRwA}{5}
\newcommand{\ToyRwB}{7}
\newcommand{\ToyRwCa}{2}
\newcommand{\ToyRwCb}{4}
\newcommand{\ToyRwCc}{6}
\newcommand{\ToyRwControl}{4.5}
\newcommand{\ToyRwGroup}{6.0}
\newcommand{\ToyRwPooled}{4.0}
\newcommand{\ToyShapleyF}{0.01}
\newcommand{\ToyShapleyNone}{0.46}
\newcommand{\ToyShapleyR}{0.30}
\newcommand{\ToyShapleyT}{0.05}
\newcommand{\ToyShapleyW}{0.10}
\newcommand{\ToySumD}{60}
\newcommand{\ToySumN}{8}
\newcommand{\ToySurvAtThirty}{0.25}
\newcommand{\ToySurvAtTwenty}{0.75}
\newcommand{\ToySurvDurations}{10, 30, 30, 100}
\newcommand{\ToyTa}{100}
\newcommand{\ToyTb}{400}
\newcommand{\ToyTc}{500}
\newcommand{\WeightASAllThousand}{$-$0.003}
\newcommand{\WeightASAllThousandBand}{[$-$0.007, $-$0.000]}
\newcommand{\WeightASHighMtnThousand}{$+$0.018}
\newcommand{\WeightASHighMtnThousandBand}{[$+$0.013, $+$0.023]}
\newcommand{\WeightASHillsThousand}{$+$0.013}
\newcommand{\WeightASHillsThousandBand}{[$+$0.008, $+$0.019]}
\newcommand{\WeightASLowMtnThousand}{$+$0.007}
\newcommand{\WeightASLowMtnThousandBand}{[$+$0.003, $+$0.011]}
\newcommand{\WeightASPlainsThousand}{$+$0.005}
\newcommand{\WeightASPlainsThousandBand}{[$+$0.002, $+$0.008]}
\newcommand{\WeightGAAllThousand}{$+$0.027}
\newcommand{\WeightGAAllThousandBand}{[$+$0.025, $+$0.029]}
\newcommand{\WeightGAHighMtnThousand}{$+$0.035}
\newcommand{\WeightGAHighMtnThousandBand}{[$+$0.030, $+$0.040]}
\newcommand{\WeightGAHillsThousand}{$+$0.012}
\newcommand{\WeightGAHillsThousandBand}{[$+$0.006, $+$0.017]}
\newcommand{\WeightGALowMtnThousand}{$+$0.034}
\newcommand{\WeightGALowMtnThousandBand}{[$+$0.030, $+$0.038]}
\newcommand{\WeightGAPlainsThousand}{$+$0.008}
\newcommand{\WeightGAPlainsThousandBand}{[$+$0.005, $+$0.011]}
\newcommand{\WeightGSAllThousand}{$+$0.023}
\newcommand{\WeightGSAllThousandBand}{[$+$0.020, $+$0.027]}
\newcommand{\WeightGSHighMtnThousand}{$+$0.053}
\newcommand{\WeightGSHighMtnThousandBand}{[$+$0.048, $+$0.058]}
\newcommand{\WeightGSHillsThousand}{$+$0.025}
\newcommand{\WeightGSHillsThousandBand}{[$+$0.020, $+$0.030]}
\newcommand{\WeightGSLowMtnThousand}{$+$0.041}
\newcommand{\WeightGSLowMtnThousandBand}{[$+$0.037, $+$0.045]}
\newcommand{\WeightGSPlainsThousand}{$+$0.013}
\newcommand{\WeightGSPlainsThousandBand}{[$+$0.009, $+$0.016]}
\newcommand{\WeightMaxShift}{0.006}
\newcommand{\WeightMinusASAllThousand}{$-$0.000}
\newcommand{\WeightMinusASAllThousandBand}{[$-$0.001, $+$0.001]}
\newcommand{\WeightMinusASHighMtnThousand}{$-$0.004}
\newcommand{\WeightMinusASHighMtnThousandBand}{[$-$0.006, $-$0.002]}
\newcommand{\WeightMinusASHillsThousand}{$+$0.001}
\newcommand{\WeightMinusASHillsThousandBand}{[$-$0.001, $+$0.002]}
\newcommand{\WeightMinusASLowMtnThousand}{$-$0.004}
\newcommand{\WeightMinusASLowMtnThousandBand}{[$-$0.005, $-$0.003]}
\newcommand{\WeightMinusASPlainsThousand}{$-$0.001}
\newcommand{\WeightMinusASPlainsThousandBand}{[$-$0.001, $+$0.000]}
\newcommand{\WeightMinusGAAllThousand}{$-$0.001}
\newcommand{\WeightMinusGAAllThousandBand}{[$-$0.001, $+$0.000]}
\newcommand{\WeightMinusGAHighMtnThousand}{$+$0.002}
\newcommand{\WeightMinusGAHighMtnThousandBand}{[$+$0.000, $+$0.004]}
\newcommand{\WeightMinusGAHillsThousand}{$-$0.001}
\newcommand{\WeightMinusGAHillsThousandBand}{[$-$0.003, $+$0.001]}
\newcommand{\WeightMinusGALowMtnThousand}{$+$0.003}
\newcommand{\WeightMinusGALowMtnThousandBand}{[$+$0.001, $+$0.004]}
\newcommand{\WeightMinusGAPlainsThousand}{$+$0.000}
\newcommand{\WeightMinusGAPlainsThousandBand}{[$-$0.001, $+$0.002]}
\newcommand{\WeightMinusGSAllThousand}{$-$0.001}
\newcommand{\WeightMinusGSAllThousandBand}{[$-$0.001, $+$0.000]}
\newcommand{\WeightMinusGSHighMtnThousand}{$-$0.002}
\newcommand{\WeightMinusGSHighMtnThousandBand}{[$-$0.003, $-$0.000]}
\newcommand{\WeightMinusGSHillsThousand}{$+$0.000}
\newcommand{\WeightMinusGSHillsThousandBand}{[$-$0.001, $+$0.002]}
\newcommand{\WeightMinusGSLowMtnThousand}{$-$0.001}
\newcommand{\WeightMinusGSLowMtnThousandBand}{[$-$0.002, $-$0.000]}
\newcommand{\WeightMinusGSPlainsThousand}{$-$0.000}
\newcommand{\WeightMinusGSPlainsThousandBand}{[$-$0.001, $+$0.001]}
\newcommand{\WeightPiecesPerFlightA}{1.52}
\newcommand{\WeightPiecesPerFlightG}{1.30}
\newcommand{\WeightPiecesPerFlightS}{1.05}
\newcommand{\WhenAllClimbEnd}{33.6}
\newcommand{\WhenAllClimbEndBand}{[33.0, 34.1]}
\newcommand{\WhenAllClimbPost}{30.5}
\newcommand{\WhenAllClimbPostBand}{[29.9, 31.1]}
\newcommand{\WhenAllClimbPre}{32.9}
\newcommand{\WhenAllClimbPreBand}{[32.2, 33.5]}
\newcommand{\WhenAllClimbStart}{35.7}
\newcommand{\WhenAllClimbStartBand}{[35.0, 36.3]}
\newcommand{\WhenAllControl}{$+$0.005}
\newcommand{\WhenAllControlBand}{[$+$0.004, $+$0.007]}
\newcommand{\WhenAllEndMinusStart}{$+$0.004}
\newcommand{\WhenAllEndMinusStartBand}{[$+$0.002, $+$0.007]}
\newcommand{\WhenAllEvents}{4\,591}
\newcommand{\WhenAllLongControl}{$+$0.001}
\newcommand{\WhenAllLongControlBand}{[$-$0.001, $+$0.003]}
\newcommand{\WhenAllLongEndMinusStart}{$+$0.007}
\newcommand{\WhenAllLongEndMinusStartBand}{[$+$0.003, $+$0.010]}
\newcommand{\WhenAllLongPostMinusEnd}{$-$0.017}
\newcommand{\WhenAllLongPostMinusEndBand}{[$-$0.021, $-$0.013]}
\newcommand{\WhenAllLongStartMinusPre}{$+$0.008}
\newcommand{\WhenAllLongStartMinusPreBand}{[$+$0.004, $+$0.011]}
\newcommand{\WhenAllPostMinusEnd}{$-$0.002}
\newcommand{\WhenAllPostMinusEndBand}{[$-$0.004, $+$0.000]}
\newcommand{\WhenAllStartMinusPre}{$+$0.002}
\newcommand{\WhenAllStartMinusPreBand}{[$-$0.001, $+$0.004]}
\newcommand{\WhenHighMtnClimbEnd}{31.3}
\newcommand{\WhenHighMtnClimbEndBand}{[30.4, 32.2]}
\newcommand{\WhenHighMtnClimbPost}{28.1}
\newcommand{\WhenHighMtnClimbPostBand}{[27.3, 29.0]}
\newcommand{\WhenHighMtnClimbPre}{32.3}
\newcommand{\WhenHighMtnClimbPreBand}{[31.4, 33.2]}
\newcommand{\WhenHighMtnClimbStart}{34.5}
\newcommand{\WhenHighMtnClimbStartBand}{[33.5, 35.5]}
\newcommand{\WhenHighMtnControl}{$+$0.002}
\newcommand{\WhenHighMtnControlBand}{[$-$0.001, $+$0.005]}
\newcommand{\WhenHighMtnEndMinusStart}{$+$0.006}
\newcommand{\WhenHighMtnEndMinusStartBand}{[$+$0.001, $+$0.010]}
\newcommand{\WhenHighMtnEvents}{1\,299}
\newcommand{\WhenHighMtnLongControl}{$-$0.003}
\newcommand{\WhenHighMtnLongControlBand}{[$-$0.008, $+$0.003]}
\newcommand{\WhenHighMtnLongEndMinusStart}{$+$0.001}
\newcommand{\WhenHighMtnLongEndMinusStartBand}{[$-$0.006, $+$0.008]}
\newcommand{\WhenHighMtnLongPostMinusEnd}{$-$0.017}
\newcommand{\WhenHighMtnLongPostMinusEndBand}{[$-$0.025, $-$0.008]}
\newcommand{\WhenHighMtnLongStartMinusPre}{$+$0.011}
\newcommand{\WhenHighMtnLongStartMinusPreBand}{[$+$0.003, $+$0.019]}
\newcommand{\WhenHighMtnPostMinusEnd}{$-$0.004}
\newcommand{\WhenHighMtnPostMinusEndBand}{[$-$0.008, $+$0.001]}
\newcommand{\WhenHighMtnStartMinusPre}{$+$0.001}
\newcommand{\WhenHighMtnStartMinusPreBand}{[$-$0.004, $+$0.006]}
\newcommand{\WhenHillsClimbEnd}{41.2}
\newcommand{\WhenHillsClimbEndBand}{[39.4, 43.0]}
\newcommand{\WhenHillsClimbPost}{40.9}
\newcommand{\WhenHillsClimbPostBand}{[39.0, 42.8]}
\newcommand{\WhenHillsClimbPre}{42.9}
\newcommand{\WhenHillsClimbPreBand}{[41.2, 44.5]}
\newcommand{\WhenHillsClimbStart}{46.3}
\newcommand{\WhenHillsClimbStartBand}{[44.5, 48.1]}
\newcommand{\WhenHillsControl}{$+$0.013}
\newcommand{\WhenHillsControlBand}{[$+$0.010, $+$0.017]}
\newcommand{\WhenHillsEndMinusStart}{$+$0.012}
\newcommand{\WhenHillsEndMinusStartBand}{[$+$0.005, $+$0.019]}
\newcommand{\WhenHillsEvents}{327}
\newcommand{\WhenHillsLongControl}{$+$0.009}
\newcommand{\WhenHillsLongControlBand}{[$+$0.004, $+$0.014]}
\newcommand{\WhenHillsLongEndMinusStart}{$+$0.013}
\newcommand{\WhenHillsLongEndMinusStartBand}{[$+$0.003, $+$0.023]}
\newcommand{\WhenHillsLongPostMinusEnd}{$-$0.013}
\newcommand{\WhenHillsLongPostMinusEndBand}{[$-$0.023, $-$0.003]}
\newcommand{\WhenHillsLongStartMinusPre}{$+$0.005}
\newcommand{\WhenHillsLongStartMinusPreBand}{[$-$0.006, $+$0.015]}
\newcommand{\WhenHillsPostMinusEnd}{$-$0.006}
\newcommand{\WhenHillsPostMinusEndBand}{[$-$0.013, $+$0.001]}
\newcommand{\WhenHillsStartMinusPre}{$+$0.000}
\newcommand{\WhenHillsStartMinusPreBand}{[$-$0.007, $+$0.008]}
\newcommand{\WhenLongHi}{1\,000}
\newcommand{\WhenLongLo}{320}
\newcommand{\WhenLowMtnClimbEnd}{31.8}
\newcommand{\WhenLowMtnClimbEndBand}{[31.1, 32.5]}
\newcommand{\WhenLowMtnClimbPost}{27.8}
\newcommand{\WhenLowMtnClimbPostBand}{[27.1, 28.5]}
\newcommand{\WhenLowMtnClimbPre}{29.5}
\newcommand{\WhenLowMtnClimbPreBand}{[28.8, 30.3]}
\newcommand{\WhenLowMtnClimbStart}{33.1}
\newcommand{\WhenLowMtnClimbStartBand}{[32.3, 33.8]}
\newcommand{\WhenLowMtnControl}{$+$0.005}
\newcommand{\WhenLowMtnControlBand}{[$+$0.003, $+$0.007]}
\newcommand{\WhenLowMtnEndMinusStart}{$+$0.005}
\newcommand{\WhenLowMtnEndMinusStartBand}{[$+$0.002, $+$0.009]}
\newcommand{\WhenLowMtnEvents}{2\,426}
\newcommand{\WhenLowMtnLongControl}{$-$0.001}
\newcommand{\WhenLowMtnLongControlBand}{[$-$0.005, $+$0.003]}
\newcommand{\WhenLowMtnLongEndMinusStart}{$+$0.011}
\newcommand{\WhenLowMtnLongEndMinusStartBand}{[$+$0.006, $+$0.017]}
\newcommand{\WhenLowMtnLongPostMinusEnd}{$-$0.020}
\newcommand{\WhenLowMtnLongPostMinusEndBand}{[$-$0.026, $-$0.014]}
\newcommand{\WhenLowMtnLongStartMinusPre}{$+$0.006}
\newcommand{\WhenLowMtnLongStartMinusPreBand}{[$-$0.001, $+$0.012]}
\newcommand{\WhenLowMtnPostMinusEnd}{$-$0.001}
\newcommand{\WhenLowMtnPostMinusEndBand}{[$-$0.005, $+$0.002]}
\newcommand{\WhenLowMtnStartMinusPre}{$+$0.000}
\newcommand{\WhenLowMtnStartMinusPreBand}{[$-$0.004, $+$0.004]}
\newcommand{\WhenMockAllEndMinusStart}{$+$0.003}
\newcommand{\WhenMockAllEndMinusStartBand}{[$+$0.001, $+$0.004]}
\newcommand{\WhenMockAllPostMinusEnd}{$-$0.001}
\newcommand{\WhenMockAllPostMinusEndBand}{[$-$0.003, $+$0.000]}
\newcommand{\WhenMockAllStartMinusPre}{$+$0.000}
\newcommand{\WhenMockAllStartMinusPreBand}{[$-$0.001, $+$0.002]}
\newcommand{\WhenMockEvents}{22\,199}
\newcommand{\WhenMockHighMtnEndMinusStart}{$+$0.002}
\newcommand{\WhenMockHighMtnEndMinusStartBand}{[$-$0.001, $+$0.005]}
\newcommand{\WhenMockHighMtnPostMinusEnd}{$-$0.001}
\newcommand{\WhenMockHighMtnPostMinusEndBand}{[$-$0.004, $+$0.002]}
\newcommand{\WhenMockHighMtnStartMinusPre}{$+$0.000}
\newcommand{\WhenMockHighMtnStartMinusPreBand}{[$-$0.003, $+$0.003]}
\newcommand{\WhenMockHillsEndMinusStart}{$+$0.005}
\newcommand{\WhenMockHillsEndMinusStartBand}{[$-$0.000, $+$0.010]}
\newcommand{\WhenMockHillsPostMinusEnd}{$+$0.002}
\newcommand{\WhenMockHillsPostMinusEndBand}{[$-$0.004, $+$0.007]}
\newcommand{\WhenMockHillsStartMinusPre}{$-$0.001}
\newcommand{\WhenMockHillsStartMinusPreBand}{[$-$0.006, $+$0.005]}
\newcommand{\WhenMockLowMtnEndMinusStart}{$+$0.002}
\newcommand{\WhenMockLowMtnEndMinusStartBand}{[$+$0.000, $+$0.005]}
\newcommand{\WhenMockLowMtnPostMinusEnd}{$-$0.001}
\newcommand{\WhenMockLowMtnPostMinusEndBand}{[$-$0.004, $+$0.001]}
\newcommand{\WhenMockLowMtnStartMinusPre}{$+$0.001}
\newcommand{\WhenMockLowMtnStartMinusPreBand}{[$-$0.001, $+$0.003]}
\newcommand{\WhenMockPlainsEndMinusStart}{$+$0.005}
\newcommand{\WhenMockPlainsEndMinusStartBand}{[$-$0.001, $+$0.010]}
\newcommand{\WhenMockPlainsPostMinusEnd}{$-$0.004}
\newcommand{\WhenMockPlainsPostMinusEndBand}{[$-$0.009, $+$0.000]}
\newcommand{\WhenMockPlainsStartMinusPre}{$-$0.000}
\newcommand{\WhenMockPlainsStartMinusPreBand}{[$-$0.005, $+$0.005]}
\newcommand{\WhenMockRealEvents}{4\,441}
\newcommand{\WhenPlainsClimbEnd}{43.9}
\newcommand{\WhenPlainsClimbEndBand}{[42.3, 45.5]}
\newcommand{\WhenPlainsClimbPost}{43.2}
\newcommand{\WhenPlainsClimbPostBand}{[41.9, 44.6]}
\newcommand{\WhenPlainsClimbPre}{44.9}
\newcommand{\WhenPlainsClimbPreBand}{[43.5, 46.2]}
\newcommand{\WhenPlainsClimbStart}{45.8}
\newcommand{\WhenPlainsClimbStartBand}{[44.1, 47.5]}
\newcommand{\WhenPlainsControl}{$+$0.005}
\newcommand{\WhenPlainsControlBand}{[$+$0.003, $+$0.007]}
\newcommand{\WhenPlainsEndMinusStart}{$-$0.002}
\newcommand{\WhenPlainsEndMinusStartBand}{[$-$0.007, $+$0.003]}
\newcommand{\WhenPlainsEvents}{433}
\newcommand{\WhenPlainsLongControl}{$+$0.001}
\newcommand{\WhenPlainsLongControlBand}{[$-$0.002, $+$0.004]}
\newcommand{\WhenPlainsLongEndMinusStart}{$+$0.001}
\newcommand{\WhenPlainsLongEndMinusStartBand}{[$-$0.004, $+$0.006]}
\newcommand{\WhenPlainsLongPostMinusEnd}{$-$0.007}
\newcommand{\WhenPlainsLongPostMinusEndBand}{[$-$0.011, $-$0.002]}
\newcommand{\WhenPlainsLongStartMinusPre}{$+$0.002}
\newcommand{\WhenPlainsLongStartMinusPreBand}{[$-$0.003, $+$0.007]}
\newcommand{\WhenPlainsPostMinusEnd}{$+$0.000}
\newcommand{\WhenPlainsPostMinusEndBand}{[$-$0.004, $+$0.004]}
\newcommand{\WhenPlainsStartMinusPre}{$+$0.004}
\newcommand{\WhenPlainsStartMinusPreBand}{[$-$0.001, $+$0.008]}
\newcommand{\WhenRealMinusMockAllEndMinusStart}{$+$0.004}
\newcommand{\WhenRealMinusMockAllEndMinusStartBand}{[$+$0.000, $+$0.008]}
\newcommand{\WhenRealMinusMockAllPostMinusEnd}{$-$0.016}
\newcommand{\WhenRealMinusMockAllPostMinusEndBand}{[$-$0.020, $-$0.012]}
\newcommand{\WhenRealMinusMockAllStartMinusPre}{$+$0.007}
\newcommand{\WhenRealMinusMockAllStartMinusPreBand}{[$+$0.003, $+$0.012]}
\newcommand{\WhenRealMinusMockHighMtnEndMinusStart}{$-$0.000}
\newcommand{\WhenRealMinusMockHighMtnEndMinusStartBand}{[$-$0.008, $+$0.008]}
\newcommand{\WhenRealMinusMockHighMtnPostMinusEnd}{$-$0.016}
\newcommand{\WhenRealMinusMockHighMtnPostMinusEndBand}{[$-$0.025, $-$0.006]}
\newcommand{\WhenRealMinusMockHighMtnStartMinusPre}{$+$0.010}
\newcommand{\WhenRealMinusMockHighMtnStartMinusPreBand}{[$+$0.002, $+$0.018]}
\newcommand{\WhenRealMinusMockHillsEndMinusStart}{$+$0.009}
\newcommand{\WhenRealMinusMockHillsEndMinusStartBand}{[$-$0.003, $+$0.020]}
\newcommand{\WhenRealMinusMockHillsPostMinusEnd}{$-$0.014}
\newcommand{\WhenRealMinusMockHillsPostMinusEndBand}{[$-$0.025, $-$0.003]}
\newcommand{\WhenRealMinusMockHillsStartMinusPre}{$+$0.003}
\newcommand{\WhenRealMinusMockHillsStartMinusPreBand}{[$-$0.010, $+$0.016]}
\newcommand{\WhenRealMinusMockLowMtnEndMinusStart}{$+$0.009}
\newcommand{\WhenRealMinusMockLowMtnEndMinusStartBand}{[$+$0.004, $+$0.015]}
\newcommand{\WhenRealMinusMockLowMtnPostMinusEnd}{$-$0.019}
\newcommand{\WhenRealMinusMockLowMtnPostMinusEndBand}{[$-$0.026, $-$0.013]}
\newcommand{\WhenRealMinusMockLowMtnStartMinusPre}{$+$0.006}
\newcommand{\WhenRealMinusMockLowMtnStartMinusPreBand}{[$-$0.001, $+$0.012]}
\newcommand{\WhenRealMinusMockPlainsEndMinusStart}{$-$0.005}
\newcommand{\WhenRealMinusMockPlainsEndMinusStartBand}{[$-$0.013, $+$0.003]}
\newcommand{\WhenRealMinusMockPlainsPostMinusEnd}{$-$0.002}
\newcommand{\WhenRealMinusMockPlainsPostMinusEndBand}{[$-$0.008, $+$0.005]}
\newcommand{\WhenRealMinusMockPlainsStartMinusPre}{$+$0.004}
\newcommand{\WhenRealMinusMockPlainsStartMinusPreBand}{[$-$0.004, $+$0.012]}
\newcommand{\WhenWindow}{1\,260}
\newcommand{\WindowLagsMax}{7}
\newcommand{\WindowLagsMin}{5}
\newcommand{\tauStar}{38.1}


% One block per step: Q, Why, Tool, Result, Observation, Deduction, Take home; Hypothesis and Comment where needed.
% Colors: questions in purple (whole text); answer labels in green; the rest black.
\definecolor{question}{HTML}{7B2D8E}
\definecolor{problem}{HTML}{C0392B}
\definecolor{answer}{HTML}{1E7B34}
\newcommand{\Q}{\item[\textcolor{question}{\textbf{Q.}}]\color{question}}
\newcommand{\Res}{\item[\textcolor{answer}{\textbf{Result.}}]\color{black}}
\newcommand{\Home}{\item[\textcolor{answer}{\textbf{Take home.}}]\color{black}}
\newcommand{\Obs}{\item[\textcolor{answer}{\textbf{Observation.}}]\color{black}}
\newcommand{\Hyp}{\item[\textcolor{question}{\textbf{Hypothesis.}}]\color{question}}
\newcommand{\Ded}{\item[\textcolor{answer}{\textbf{Deduction.}}]\color{black}}
\newcommand{\Com}{\item[\textcolor{black}{\textbf{Comment.}}]\color{black}}
\newcommand{\Why}{\item[\textcolor{black}{\textbf{Why.}}]\color{black}}
\newcommand{\Tool}{\item[\textcolor{black}{\textbf{Tool.}}]\color{black}}
\newenvironment{step}{\begin{description}[leftmargin=0pt, itemsep=7pt, topsep=8pt]}{\end{description}}
% Method box: how a statistic is built and why, apart from the observations. Figures inside use \captionof.
\newtcolorbox{method}[1]{breakable, enhanced, colback=black!3, colframe=black!35, boxrule=0.5pt, arc=2pt,
  left=7pt, right=7pt, top=5pt, bottom=5pt, fonttitle=\bfseries, coltitle=black, colbacktitle=black!10,
  title={Method: #1}}
% Problem box: what goes wrong with an obvious choice.
\newtcolorbox{problembox}[1]{breakable, enhanced, colback=problem!4, colframe=problem!70, boxrule=0.6pt, arc=2pt,
  left=7pt, right=7pt, top=5pt, bottom=5pt, fonttitle=\bfseries, coltitle=problem!80!black,
  colbacktitle=problem!12, title={Problem: #1}}

\title{\textbf{Group time and solo time}}
\author{}
\date{\today\\ paragliders, declared open flights logged every \DtMax~s or faster, grouping run \texttt{d120\_eh500\_ez300\_th45\_g120\_q0.1}}

\begin{document}
\maketitle

\section{The question}\label{question}

Do paragliders move differently when they fly in a group? We answer on one property of the motion, its
persistence: how fast the displacement grows with time, read on the exponent $\alpha$ of the mean squared
displacement (Sec.~\ref{measures}).

\begin{method}{three populations and one identity}
  Three populations of flight time:
  \begin{itemize}
  \item $\Sl$: the solo flights, never in a group;
  \item $\A$: the solo time of the flights that have group time;
  \item $\G$: their group time.
  \end{itemize}
  For any statistic $X$ measured with the same estimator on the three,
  \[ X_\G-X_\Sl=(X_\G-X_\A)+(X_\A-X_\Sl). \]
  \begin{itemize}
  \item $\G-\A$ compares two moments of the \emph{same flights}: same pilot, wing, day, site and flight shape.
    What changes is flying with others, and the moment of the flight. It is the group effect proper, up to the
    moment (Sec.~\ref{ch:when}) and the direction (Sec.~\ref{which-way}).
  \item $\A-\Sl$ compares \emph{different flights}, both alone: who flies, where, when, which flight. It is the
    selection of the flights that have group time, plus any carry-over of the group on the solo time around it.
  \item $\G-\Sl$ is the raw total.
  \end{itemize}
  $\G-\A$ already holds the flight fixed, so the controls of the chapters below that compare flights (the ground
  class, the opening, the same pilot) act mostly on $\A-\Sl$; the timeline of the flight acts on $\G-\A$.
\end{method}

\textbf{How to read a step.} Every step answers one question in a fixed order:
\begin{itemize}
\item \textbf{Q.} the question;
\item \textbf{Why.} why we ask it and what the answer would change; from the second chapter on, it follows from
  the take home of the chapter before;
\item \textbf{Tool.} the analysis, and why this one;
\item \textbf{Result.} the figures, tables and numbers;
\item \textbf{Observation.} what we see in them;
\item \textbf{Deduction.} what follows;
\item \textbf{Take home.} what we keep.
\end{itemize}
A \textbf{Hypothesis} states what we expect before the result, and a \textbf{Comment} adds what the step does not
need; both appear only where they carry something. How a statistic is built goes in a \emph{method} box, what goes
wrong with an obvious choice in a \emph{problem} box, apart from the observations.

\textbf{The chapters.} Is group time more persistent (Sec.~\ref{ch:straighter})? Is it the group or the flights
(Sec.~\ref{ch:identity})? Where (Sec.~\ref{ground}), on which flights (Sec.~\ref{ch:flights}), for whom
(Sec.~\ref{ch:who}), when in the flight (Sec.~\ref{ch:when})? How: the dynamics (Sec.~\ref{ch:dynamics}) and the
geometry (Sec.~\ref{ch:geometry}) of the motion. Which way round (Sec.~\ref{which-way})? What we found
(Sec.~\ref{ch:found}).

\textbf{Data.} Paragliders, flights declared open on the FFVL catalogue (free distance or distance with 1--3
points), logged every \DtMax~s or faster, so that every flight has phases (glide, search, climb,
Sec.~\ref{phases}); those logged more slowly (\DtDroppedPercent\%) are dropped. \DtFlights{} of them give at least
one piece (Sec.~\ref{definitions}) and enter the statistics of $\alpha$; the tables of the opening count every
declared open flight, with or without pieces.

\textbf{Group time.} It comes from the grouping run of the group-identification note, whose name gives its
parameters: two pilots are in contact when one is on the path of the other within $\delta=120$~s, 500~m
horizontally and 300~m vertically (headings compared only in glide), for at least the duration $\tau^*$ at which
pairs of pilots shifted in time explain at most a share $q=0.1$ of the contacts. A pilot flies in a group while
in contact with at least one other. The run is built on every flight: a partner logged more slowly still makes a
flight's time group time, but only the tracks of the flights kept are analysed.

\section{Measures}\label{measures}

\subsection{Definitions}\label{definitions}

\begin{description}[leftmargin=0pt, itemsep=7pt]
\item[S] Flights never in a group.
\item[G] Group time: the stretches of group flights spent in a group. \textbf{The unit is the stretch,
  not the flight}: a flight with three stretches in a group gives three pieces and counts three times.
\item[A] The solo time of the same group flights (same pilots and days); the unit is again the stretch.
\item[Segment] Loggers sometimes miss fixes. The cleaning of the tracks (old repository soaring-anomalous-transport) splits a track
  where two consecutive fixes are more than $\min(10\,\mathrm{d}t,\ \max(20~\mathrm{s},\ 2\,\mathrm{d}t))$
  apart, with $\mathrm{d}t$ the logging interval of the flight. Each part is a segment.
\item[Piece] The part of one segment that lies in one population: for $\Sl$ a whole segment, for $\G$ a
  stretch in a group, for $\A$ a stretch alone. A piece never crosses a gap. A flight can give several
  pieces: several segments, or several stretches in and out of a group.

  \textbf{Our choice: only the pieces of at least \PieceMinS~s are kept} ($=1\,000~\mathrm{s}/\Support$, so
  that at the lag 1\,000~s at least \SupportOriginsPercent\% of the piece gives origins). Every shorter piece is dropped: short
  segments, and the stretches in a group or alone shorter than about \PieceMinMinutes~min. The statistics of $\alpha$ and of
  the phases are on these pieces. Why:
  \begin{itemize}
  \item \emph{One population up to 1\,000~s.} Every lag from 10~s to 1\,000~s sees the same pieces, so along
    these two decades $\alpha$ changes with the dynamics only, never with the pieces kept. Without the
    threshold the shortest lags would also take pieces of a few minutes, and the population would change
    along the curve exactly where it is now fixed.
  \item \emph{Short stretches are mostly edges.} A stretch in a group of a few minutes is mostly its entry and
    exit, where the label of the group is least reliable; a short segment lies between gaps of the logger.
  \end{itemize}
  The price: $\G$ describes the stretches in a group that last, not the brief encounters.
\item[TA-MSD] Time-averaged MSD of one piece $p$, in three steps.
  \begin{enumerate}
  \item Positions: the fixes, linearly interpolated on a grid of $\Delta t=\GridS$~s, give the horizontal
    position $x_j=(E_j,N_j)$ east and north, $j=0,\dots,M-1$; duration $T_p=(M-1)\Delta t$.
  \item At lag $\tau=k\Delta t$, every origin $j$ with $j+k\le M-1$ gives one squared displacement
    $|x_{j+k}-x_j|^2$; there are $M-k$ origins. Their mean is the TA-MSD:
    \[ \overline{\delta^2_p}(\tau) = \frac{1}{M-k}\sum_{j=0}^{M-k-1}|x_{j+k}-x_j|^2 , \]
    the discrete form of $\frac{1}{T_p-\tau}\int_0^{T_p-\tau}|x(t+\tau)-x(t)|^2\,dt$.
  \item Repeated at every lag of a \textbf{fixed grid}, not at every multiple of $\Delta t$: one curve
    $\overline{\delta^2_p}(\tau)$ per piece. A piece of the cohort lasts at least $\tau_\mathrm{max}/\Support$, so
    even at $\tau_\mathrm{max}$ at least \SupportOriginsPercent\% of its length gives origins.
  \end{enumerate}
  \emph{The grid of lags.} Geometric: each lag is about \LagRatio{} times the previous one, then rounded to a
  multiple of $\Delta t=\GridS$~s,
  \[ \frac{\tau_{i+1}}{\tau_i}\simeq10^{\LagDecades/\LagSteps}\approx\LagRatio \qquad\text{(\LagSteps{} steps per factor \LagFactor, about \LagsPerDecade{}
     lags per decade)}, \]
  from \LagMin~s to \LagMax~s. The \LagCount{} lags, in seconds:
  \begin{center}\small
  \LagList.
  \end{center}
  Below 200~s the rounding to \GridS~s makes the steps larger than \LagRatio.
\item[MSD] At each lag $\tau$, the mean of the $N$ values $\overline{\delta^2_p}(\tau)$ of the pieces of a
  fixed cohort, one weight per piece:
  \[ \mathrm{MSD}(\tau) = \frac{1}{N}\sum_{p}\overline{\delta^2_p}(\tau). \]
  Weights do not depend on $\tau$. In the raw comparison the MSD does not depend on the unit, pieces or
  flights within a few thousandths (Appendix~\ref{app:weights}).
\item[$\boldsymbol{\alpha(\tau)}$] Local exponent, the numerical form of
  $\alpha(\tau)=\dfrac{1}{2}\dfrac{d\log\mathrm{MSD}}{d\log\tau}$. For a power law $\mathrm{MSD}\propto\tau^{2\alpha}$
  it is the exponent.
  \begin{enumerate}
  \item \emph{Lags}: the grid $\tau_k$ of the TA-MSD, consecutive lags in the ratio $\approx\LagRatio$.
  \item \emph{Window}: the lags within \HalfDecade{} decades of $\tau$,
    $W(\tau)=\{k:\ |\log_{10}\tau_k-\log_{10}\tau|\le\HalfDecade\}$, that is $\tau_k\in[\tau^-,\tau^+]$ with
    $\tau^\pm=\tau\cdot10^{\pm\HalfDecade}$. It holds \WindowLagsMin{} to \WindowLagsMax{} lags over most of the grid; with fewer than 3 (the
    ends of the grid), the 3 nearest.
  \item \emph{Fit}: with $u_k=\log\tau_k$, $y_k=\log\mathrm{MSD}(\tau_k)$ and the means $\bar u,\bar y$ over $W$,
    the least-squares line through the points $(u_k,y_k)$, $k\in W(\tau)$, equal weight per lag, has half-slope
    \[ \alpha(\tau) = \frac{1}{2}\,\frac{\sum_{k\in W}(u_k-\bar u)(y_k-\bar y)}{\sum_{k\in W}(u_k-\bar u)^2}. \]
  \end{enumerate}

\item[$\boldsymbol{H}$] Half the least-squares slope of the straight line fitted to $\log\mathrm{MSD}$ against
  $\log\tau$ over the lags $\tau_1,\dots,\tau_K$ of the grid, from \LagMin~s (the shortest lag of the \GridS-s grid) to
  the $\tau_\mathrm{max}$ of the cohort, one weight per lag. So $H$ depends on the cohort through its range. With $u_k=\log\tau_k$ and
  $y_k=\log\mathrm{MSD}(\tau_k)$:
  \[ H = \frac{1}{2}\,\frac{\sum_k (u_k-\bar u)(y_k-\bar y)}{\sum_k (u_k-\bar u)^2}. \]

  \emph{Link to $\alpha$.} Take the lags dense and evenly spaced in $u=\log\tau$ over $[u_0,u_1]$, with
  $L=u_1-u_0$ and $\bar u$ the middle. Integrating by parts, $H$ is a weighted mean of $\alpha$, with a
  parabolic weight that peaks at the middle of the range and vanishes at its ends:
  \[ H = \int_{u_0}^{u_1} w(u)\,\alpha(u)\,du, \qquad w(u) = \frac{6}{L^3}\Big[\big(\tfrac{L}{2}\big)^2-(u-\bar u)^2\Big],
  \qquad \int w = 1. \]
  Our lags are geometric, rounded to the \GridS-s grid, so this holds approximately. For a pure power law
  $\alpha$ is constant and $H=\alpha$.
\end{description}

\subsection{Bands}\label{bands}

Pieces are not independent: those of one pilot share skill and wing, those of one block share the day and
the air. A \emph{block} is the set of flights joined by a group episode or by the same site and day.
\begin{itemize}
\item Pilots and blocks cross, so three bootstraps of \NBoot{} draws each resample with replacement the
  pilots, the blocks and the cells (pilot $\times$ block); in a draw each piece weighs the count of its
  pilot, block or cell.
\item With $V_p$, $V_b$, $V_c$ the variances of an estimate over the three sets of draws (Cameron, Gelbach
  and Miller~\cite{cameron2011}):
  \[ V = \max(V_p + V_b - V_c,\ V_p,\ V_b), \qquad \text{band} = \text{estimate} \pm \BandZ\sqrt{V}\quad(90\%). \]
\item \emph{Shared draws.} $\Sl$, $\A$ and $\G$ are computed in the same draws, so every gap and every difference
  has its own band, from its values over the draws; "paired" draws are the same for two estimators of one
  population.
\item Code: \texttt{twoway.py}; \texttt{coverage.py} checks on synthetic data with pilot and block effects that the
  band of nominal 90\% covers the true value in \CovNewMin--\CovNewMax\% of the datasets.
\end{itemize}


\subsection{Fixed cohort}\label{fixed-cohort}
\begin{step}
\Q Should a lag-resolved statistic use every flight, with fewer flights at longer lags?
  \textcolor{answer}{\textbf{No}}
\Why The TA-MSD of a piece $p$ at lag $\tau$ averages over the origins $t_j=j\Delta t\in[0,T_p-\tau]$. For
  enough origins, and so a reliable estimate, we ask $T_p-\tau\ge (1-\Support)\,T_p$, that is
  $T_p\ge\tau/\Support$, our choice. So the longer
  the lag, the fewer the pieces, and an exponent read across lags may also read which pieces leave.
\end{step}

\begin{problembox}{the population changes with the lag}
  If every lag uses all the pieces it can, the population changes with $\tau$. From $\tau_1$ to $\tau_2$
  the short pieces drop out. Let
  \begin{itemize}
  \item $a$: MSD at $\tau_1$ of all the pieces present at $\tau_1$;
  \item $b$: MSD at $\tau_1$ of only the pieces still present at $\tau_2$;
  \item $c$: MSD at $\tau_2$ of these same pieces;
  \item $H$: exponent of the pieces present at both lags, $2H = \log(c/b)/\log(\tau_2/\tau_1)$;
  \item $\hat H_\mathrm{pool}$: exponent read from the pool, $2\hat H_\mathrm{pool} = \log(c/a)/\log(\tau_2/\tau_1)$.
  \end{itemize}
  Since $c/a = (c/b)(b/a)$,
  \[ 2\hat H_\mathrm{pool} = 2H + \frac{\log(b/a)}{\log(\tau_2/\tau_1)}. \]
  The extra term is zero only if $b=a$, namely when the dropped pieces had the same MSD as those that stay.\\

  \emph{The same for $\alpha$.} $\alpha(\tau)$ is also a slope, over the lags $\tau_k$ of its window
  $[\tau^-,\tau^+]$ (Definitions). For each lag $\tau_k$ of the window let
  \begin{itemize}
  \item $a_k$: MSD at the lag $\tau_k$ of all the pieces present at $\tau_k$ (the pool);
  \item $b_k$: MSD at the same lag $\tau_k$ of only the pieces present at $\tau^+$, the longest lag of the window.
  \end{itemize}
  Then
  $\log\mathrm{MSD}_\mathrm{pool}(\tau_k)=\log b_k+\log(a_k/b_k)$, and since the least-squares slope is linear in
  the points,
  \[ \hat\alpha_\mathrm{pool}(\tau) = \alpha(\tau) + \frac12\,\frac{\sum_k(u_k-\bar u)\,(r_k-\bar r)}{\sum_k(u_k-\bar u)^2},
     \qquad u_k=\log \tau_k,\quad r_k=\log(a_k/b_k). \]
  At $\tau^+$ the two sets coincide and $r=0$. Below, the pool also holds the shorter pieces: if they have a
  lower MSD ($a_k<b_k$), $r_k$ is negative and rises to 0 across the window, its slope is positive and
  $\hat\alpha_\mathrm{pool}>\alpha$. With the two lags $\tau^-,\tau^+$ only, the extra term is
  $\tfrac12\log(b/a)/\log(\tau^+/\tau^-)$, the formula of $H$ above.
\end{problembox}

\begin{step}
\Tool Nested fixed cohorts, each on its own lag range. If the length of the pieces does not select, the
  cohorts give the same $\alpha$ at the lags they share, within the bands. Two checks:
  \begin{itemize}
  \item the whole curves $\alpha(\tau)$ of the cohorts, overlaid: do they coincide where they overlap
    (Fig.~\ref{fig:cohorts})?
  \item one shared lag, $\Hk$, against the $\tau_\mathrm{max}$ of the cohort: does it stay flat
    (Fig.~\ref{fig:tradeoff} left)?
  \end{itemize}
\Res Cohorts to 1\,000~s, 2\,760~s, 1~h, 2~h, 3~h; MSD one weight per piece (Fig.~\ref{fig:cohorts}), on
  \textbf{every declared open flight} (free distance or distance with 1--3 points on the FFVL catalogue), with
  no other condition.
  \begin{figure}[H]
  \centering
  \includegraphics[width=\textwidth]{open_declared_dt5/cohorts_msd.pdf}
  \caption{Top: $\mathrm{MSD}(\tau)$ of each cohort; middle: $\mathrm{MSD}(\tau)/\tau$; bottom: its $\alpha(\tau)$. Red: the whole dataset, each lag on every piece that reaches it (no fixed cohort). Nested cohorts on every declared open flight
  (no condition on the opening), one weight per piece; 90\% bands. Each curve stops at the $\tau_\mathrm{max}$ of its cohort.}
  \label{fig:cohorts}
  \end{figure}
\Obs How the statistics change with the cohort:
  \begin{itemize}
  \item \emph{Pieces}: a longer $\tau_\mathrm{max}$ explores longer lags but keeps fewer pieces ($\G$: from
    \CohortPiecesFirst{} to \CohortPiecesLast, Fig.~\ref{fig:tradeoff} right).
  \item \emph{Level}: from the cohort to 1~h on, each longer cohort has a higher curve $\alpha(\tau)$ than the one
    before, in the three populations, at every shared lag but at most one. Below 1~h the cohorts are close:
    $\alpha(1\,000~\mathrm{s})$ grows with $\tau_\mathrm{max}$ for $\A$ and $\G$, while for $\Sl$ it first goes
    down, from the cohort to 1\,000~s to the one to 2\,760~s, and then up (Fig.~\ref{fig:tradeoff} left). From the
    cohort to 1~h to the one to 2~h, $\alpha(1\,000~\mathrm{s})$ rises by \AlphaStepHourMin{} to
    \AlphaStepHourMax, and the next step is of the same order, against at most \AlphaStepBeforeMax{} between the
    shorter cohorts.
  \item \emph{Shape}: from 2~h on the end of the curve changes: a second peak, then a descent
    (Fig.~\ref{fig:cohorts}).
  \item \emph{Contrasts} (Fig.~\ref{fig:gaps}): up to \GapSpreadLag~s the cohort moves the gaps less than $\alpha$
    itself. Among the cohorts that reach beyond, the gaps differ at most by \GapSpreadLowGS{} ($\G-\Sl$),
    \GapSpreadLowAS{} ($\A-\Sl$) and \GapSpreadLowGA{} ($\G-\A$) up to \GapSpreadLag~s; above it by
    \GapSpreadHighGS, \GapSpreadHighAS{} and \GapSpreadHighGA.
  \item \emph{Whole dataset} (red, Fig.~\ref{fig:cohorts}): up to about \PoolSame~s every piece reaches the lag and
    the red curve is the cohort to 1\,000~s. From about \PoolLeave~s it leaves the band of the cohort to 1~h. In
    $\Sl$ it first falls below it, up to about \PoolBelowEnd~s; in $\A$ it rises above it from about
    \PoolAboveA~s, in $\G$ from about \PoolAboveG~s. From about \PoolSteep~s it is above every cohort in the three
    populations and climbs steeply, up to $\alpha=\PoolSMax$ ($\Sl$),
    \PoolAMax{} ($\A$), \PoolGMax{} ($\G$), above the ballistic 1, while the cohorts slowly descend. In the
    top row its MSD lies below the longer cohorts and joins each of them at its $\tau_\mathrm{max}$: it is not
    higher, it grows faster, and in log--log it bends upward where the cohorts are straight.
  \end{itemize}
  \begin{figure}[H]
  \centering
  \includegraphics[width=\textwidth]{open_declared_dt5/cohorts_tradeoff.pdf}
  \caption{Left: $\Hk$ against the $\tau_\mathrm{max}$ of the cohort (pieces of at least $\tau_\mathrm{max}/\Support$);
  90\% bands. Right: pieces of each cohort. Every declared open flight (no condition on the opening).}
  \label{fig:tradeoff}
  \end{figure}
  \begin{figure}[H]
  \centering
  \includegraphics[width=\textwidth]{open_declared_dt5/cohorts_gaps_msd.pdf}
  \caption{Top: the ratio of the MSDs of the two populations in each cohort, whose half log-slope is the gap;
  bottom: the gaps $\Delta\alpha(\tau)$ between the populations in each cohort, one weight per piece; 90\%
  bands from the shared two-way draws. Every declared open flight (no condition on the opening).}
  \label{fig:gaps}
  \end{figure}
\Ded\mbox{}
  \begin{itemize}
  \item \textbf{The choice of the cohort matters for the gaps $\G-\Sl$, $\G-\A$, $\A-\Sl$ from about
    \GapSpreadLag~s on.}
  \item \textbf{The red curve is the extra term of the Problem}, $\log(b/a)/\log(\tau_2/\tau_1)$. Where the
    pieces that drop out have the lower MSD ($b>a$) the term is positive and the pool climbs above every cohort;
    in $\Sl$ between \PoolLeave{} and \PoolBelowEnd~s the solo flights that drop out are a little straighter than
    those that stay ($b<a$), and the pool falls below. Its rise from \PoolSteep~s and its $\alpha>1$ say how
    fast the short flights leave, not how the flights move. \textbf{So the pool cannot be used.}
  \item Any single cohort is a compromise. Carrying five curves per population is not practical either.
  \end{itemize}
\Home\mbox{}
  \begin{itemize}
  \item \textbf{From the cohort to 1~h on, the longer flights are straighter at every shared lag}
    $\Longrightarrow$ a longer cohort selects straighter flights.
  \item \textbf{Never the pool for a lag-resolved statistic.} Without a fixed cohort the short pieces leave at
    long lags, the remaining ones have a higher MSD ($b>a$), and the extra term
    $\log(b/a)/\log(\tau_2/\tau_1)>0$ raises $\hat\alpha_\mathrm{pool}$ above 1 only because the short flights
    left (Fig.~\ref{fig:cohorts}, red).
  \end{itemize}
\end{step}

\subsection{One curve per population}\label{which-cohort}

\begin{step}
\Q How do we get \textbf{one curve per population over every lag}, using at each lag the pieces that reach it?
\Why The step above rules out the pool, and any single cohort is a compromise between lags and pieces.
\Tool The chained MSD: what we want of it, what goes wrong with the obvious choices, how it is built, and the proof
  on a toy (the boxes below); then its MSD and its exponent on the data.
\end{step}

\begin{method}{what we want}
  A statistic computed as a function of the lag, on pieces that must reach that lag (the MSD, $\alpha$, and so
  any distribution or correlation of the increments), should have four properties:
  \begin{enumerate}[label=(\arabic*)]
  \item \emph{Every lag}, from 10~s to the longest flights. The long lags exist only in the long flights, so
    these must enter.
  \item \emph{At each lag, as many pieces as possible}: the short flights too, at the lags they reach.
  \item \emph{No composition term}: each increment of $\log\mathrm{MSD}$ along the lags of one fit of $\alpha$
    must compare one set of pieces with itself, or the slope also measures which pieces leave
    (Sec.~\ref{fixed-cohort}). Which pieces reach the long lags (the selection) is another matter, which no
    method removes (Comment of Sec.~\ref{which-cohort}).
  \item \emph{One curve per population.} Five cohorts give five curves for each of $\Sl$, $\A$, $\G$ and for each
    gap, and above $10^3$~s the gaps depend on the cohort (Sec.~\ref{fixed-cohort}): no single message.
  \end{enumerate}
  The pool meets (1), (2), (4) and fails (3). A single fixed cohort meets (3) and fails (1) if short or (2) if
  long; all of them together fail (4). The chained MSD below meets all four.

  Statistics without a lag, such as the glide speed or the climb rate, are a mean over one set of pieces: the
  choice of the flights changes their population, but no set changes along an axis, and the problem does not
  arise.
\end{method}

\begin{problembox}{the pool and the mean of the cohorts bias the slope}
  \textbf{The pool.} At each lag, every piece that reaches it:
  \[ \mathrm{MSD}_\mathrm{pool}(\tau_k)=\frac{1}{|P(\tau_k)|}\sum_{p\in P(\tau_k)}\overline{\delta^2_p}(\tau_k),\qquad
     P(\tau)=\{p:\ T_p\ge\tau/\Support\}. \]
  One curve over every lag, on every piece. But along the window of a fit, $\tau_k\in[\tau^-,\tau^+]$, the set
  shrinks, $P(\tau^+)\subseteq P(\tau_k)\subseteq P(\tau^-)$, and the slope picks up the composition term of
  Sec.~\ref{fixed-cohort}:
  \[ \hat\alpha_\mathrm{pool}(\tau)=\alpha(\tau)+\frac12\,\frac{\sum_k(u_k-\bar u)(r_k-\bar r)}{\sum_k(u_k-\bar u)^2},
     \qquad r_k=\log\frac{a_k}{b_k}, \]
  with $a_k$ the MSD of the pool at $\tau_k$, $b_k$ that of the pieces of $P(\tau^+)$ at the same lag, and
  $\alpha(\tau)$ the exponent of these pieces. Where the shorter pieces have the lower MSD, $r_k<0$ rises to 0 at
  $\tau^+$, and \textbf{the pool is steeper than the pieces it is made of}. What it inflates is the growth of the
  MSD with the lag: within the window the pool still holds shorter pieces, so its level lies below that of the
  pieces of $P(\tau^+)$ and climbs toward it at $\tau^+$. In the
  data this is the red curve of Fig.~\ref{fig:cohorts} (Sec.~\ref{fixed-cohort}, Observation).

  \textbf{The mean of the cohort curves.} With nested cohorts $c$ to $\tau_c$ and any weights $w_c$ (equal, or the
  number of pieces $n_c$), only the cohorts that reach the lag enter:
  \[ \overline{\mathrm{MSD}}(\tau_k)=\frac{\sum_{c:\,\tau_c\ge\tau_k}w_c\,\mathrm{MSD}_c(\tau_k)}{\sum_{c:\,\tau_c\ge\tau_k}w_c}. \]
  A cohort leaves at its $\tau_c$, in analogy with a piece that leaves the pool. With $w_c=n_c$ the mean is a
  pool in which each piece weighs the number of cohorts that hold it and reach the lag,
  \[ \sum_{c:\,\tau_c\ge\tau_k}n_c\,\mathrm{MSD}_c(\tau_k)=\sum_p m_p(\tau_k)\,\overline{\delta^2_p}(\tau_k),\qquad
     m_p(\tau_k)=\#\{c\ni p:\ \tau_c\ge\tau_k\}, \]
  so the same composition term appears in every window where a cohort ends. In the toy below the mean of the
  cohorts gives \ToyPlMean{} at \ToyPlTau~s and up to \ToyPlMeanMax, against $\alpha=\ToyPlAlpha$ for every
  piece (Fig.~\ref{fig:chained-scheme}, purple). Averaging the curves $\alpha_c(\tau)$ instead removes the term
  from each fit, but the cohorts in the mean still change with $\tau$, and the curve jumps where one ends.
\end{problembox}

\begin{method}{the chained MSD}
  Start from every piece at the shortest lag, then move one lag at a time, each step on the pieces that reach
  both of its lags. With $M_S(\tau)$ the mean of $\overline{\delta^2_p}(\tau)$ over a set $S$ of pieces,
  \[ \mathrm{MSD}^*(\tau_1)=M_{P(\tau_1)}(\tau_1),\qquad
     \mathrm{MSD}^*(\tau_{k+1})=\mathrm{MSD}^*(\tau_k)\,\frac{M_{P(\tau_{k+1})}(\tau_{k+1})}{M_{P(\tau_{k+1})}(\tau_k)}. \]
  The pieces of $P(\tau_{k+1})$ reach $\tau_k$ too, so each factor compares two lags on one set. It is the
  product-limit estimator of Kaplan and Meier~\cite{kaplan1958}, where the units of a factor are the patients
  still followed at both times, and the chain index of Marshall~\cite{marshall1887}, where they are the goods
  sold in both periods.
  \begin{itemize}
  \item \emph{Every lag}: the chain runs as far as the pool, one step at a time.
  \item \emph{Below 1\,000~s it is the mean of every piece.} Every piece of the cache lasts at least \PieceMinS~s
    ($=1\,000/\Support$), so up to 1\,000~s every step runs on all of them, the factors telescope and
    $\mathrm{MSD}^*=M_{P(\tau_1)}$, the pool.
  \item \emph{The level beyond.} From 1\,000~s on, each step grows as the pieces that reach it, while the level
    started from every piece. So at long lags $\mathrm{MSD}^*$ is the MSD that every piece would have if it grew
    as the long ones. Kaplan and Meier make the same assumption: leaving does not depend on what is measured.
    Here it fails, since the longest flights are straighter (Sec.~\ref{fixed-cohort}). Beyond 1\,000~s the slope
    of $\mathrm{MSD}^*$ is measured, its level is a construction.
  \end{itemize}
\end{method}

\begin{step}
\Res \textbf{The MSD} (Fig.~\ref{fig:chained-msd}, every declared open flight), log--log: the slope of a curve
  is $2\alpha$.
  \begin{itemize}
  \item Up to 1\,000~s the chained MSD and the pool coincide: every piece reaches every lag.
  \item Above, in $\A$ and $\G$ the pool lies above $\mathrm{MSD}^*$, more and more with the lag: at each step the
    pieces that leave had the lower MSD ($b>a$). In $\Sl$ it lies slightly below up to about \PoolMsdCross~s, since
    the short solo flights that leave are a little straighter (Sec.~\ref{fixed-cohort}), and above beyond.
  \item At the longest lags the pool bends upward, while $\mathrm{MSD}^*$ stays close to a line.
  \end{itemize}
  \begin{figure}[H]
  \centering
  \includegraphics[width=\textwidth]{open_declared_dt5/chained_msd.pdf}
  \caption{$\mathrm{MSD}^*(\tau)$ (color) and the pool (red, each lag on every piece that reaches it). Every
  declared open flight, one weight per piece; 90\% bands.}
  \label{fig:chained-msd}
  \end{figure}
\end{step}

\begin{method}{the exponent $\alpha^*$}
  $\alpha^*(\tau)$ is half the slope of the line fitted to $\log\mathrm{MSD}^*(\tau_k)$ against $\log\tau_k$ over
  the window $\tau_k\in[\tau^-,\tau^+]$, $\tau^\pm=\tau\cdot10^{\pm\HalfDecade}$ (Definitions), the fit of every
  curve of this note.
  \begin{itemize}
  \item \emph{No composition term.} Each increment of $\log\mathrm{MSD}^*$ is on one set,
    \[ \log\mathrm{MSD}^*(\tau_{k+1})-\log\mathrm{MSD}^*(\tau_k)=\log\frac{M_{P(\tau_{k+1})}(\tau_{k+1})}{M_{P(\tau_{k+1})}(\tau_k)}, \]
    and the least-squares slope over the window is a mean of the slopes of its steps, with positive weights.
    So $\alpha^*$ averages exponents of fixed sets: no step compares two different sets.
  \item \emph{The pool is the chain plus the pieces that leave.} The increment of the pool,
    $\log M_{P(\tau_{k+1})}(\tau_{k+1})-\log M_{P(\tau_k)}(\tau_k)$, is that of the chain plus
    $\log\big(M_{P(\tau_{k+1})}(\tau_k)/M_{P(\tau_k)}(\tau_k)\big)$, the $\log(b/a)$ of Sec.~\ref{fixed-cohort}
    at that step. Summing the steps,
    \[ \log\frac{\mathrm{MSD}_\mathrm{pool}(\tau_k)}{\mathrm{MSD}^*(\tau_k)}=\sum_{j<k}\log\frac{M_{P(\tau_{j+1})}(\tau_j)}{M_{P(\tau_j)}(\tau_j)}, \]
    and $\hat\alpha_\mathrm{pool}-\alpha^*$ is half the slope of this sum over the window: the composition term
    alone.
  \item \emph{The four properties}: every lag (1), at each step every piece that reaches it (2), one set per
    step (3), one curve (4).
  \item \emph{What it does not remove}: the selection. At long lags a step runs on the long flights only
    (Comment below).
  \end{itemize}
\end{method}

\begin{method}{the proof on a toy}
  \textbf{The toy.} Every piece has the same exponent $\alpha$ and a prefactor that grows with its duration,
  so the longer pieces are straighter at every scale, as in the data (Sec.~\ref{fixed-cohort}):
  \[ \overline{\delta^2_p}(\tau)=c_p\,\tau^{2\alpha},\qquad c_p=\Big(\frac{T_p}{T_0}\Big)^{\beta},\qquad
     \alpha=\ToyPlAlpha,\quad \beta=\ToyPlBeta,\quad T_0=\ToyPlTzero~\mathrm{s}, \]
  with the durations $T_p$ log-uniform on $[T_\mathrm{min},T_\mathrm{max}]=[\PieceMinS,\ \ToyPlTmax]$~s, so
  that, as in the data, every piece reaches 1\,000~s. The right answer is known: $\alpha$ at every lag.

  \textbf{The two MSDs.} The mean over a set $S$ of pieces factorizes: $M_S(\tau)=C(S)\,\tau^{2\alpha}$,
  with $C(S)$ the mean of $c_p$ over $S$. In each factor of the chain $C$ cancels,
  \[ \frac{M_{P(\tau_{k+1})}(\tau_{k+1})}{M_{P(\tau_{k+1})}(\tau_k)}=\Big(\frac{\tau_{k+1}}{\tau_k}\Big)^{2\alpha}
     \qquad\Longrightarrow\qquad \mathrm{MSD}^*(\tau)=C\big(P(\tau_1)\big)\,\tau^{2\alpha}, \]
  while $\mathrm{MSD}_\mathrm{pool}(\tau)=C\big(P(\tau)\big)\,\tau^{2\alpha}$, so
  \[ \frac{\mathrm{MSD}_\mathrm{pool}(\tau)}{\mathrm{MSD}^*(\tau)}=\frac{C\big(P(\tau)\big)}{C\big(P(\tau_1)\big)}\ge1, \]
  equal to 1 up to 1\,000~s, where every piece counts. Above, the pool lies above the chain, since it has lost
  the short pieces, which have a smaller $c_p$ (Fig.~\ref{fig:chained-scheme}, bottom left).

  \textbf{The chained slope is exact.} $\log\mathrm{MSD}^*(\tau)=2\alpha\log\tau+\text{const}$ at every lag: the
  points lie on one line of slope $2\alpha$, and the fit gives $\hat\alpha^*=\alpha$ in every window
  (\ToyPlChained{} in the figure). This holds for any durations and prefactors: each factor compares one set with
  itself.

  \textbf{The pool slope is not.} $C(P(\tau))$ grows with $\tau$, as the short pieces leave. With
  $x=\tau/\Support$ the shortest duration that reaches $\tau$, for $x\ge T_\mathrm{min}$,
  \[ C\big(P(\tau)\big)=\frac{1}{\ln(T_\mathrm{max}/x)}\int_x^{T_\mathrm{max}}\Big(\frac{T}{T_0}\Big)^{\beta}\frac{dT}{T}
     =\frac{(T_\mathrm{max}/T_0)^\beta-(x/T_0)^\beta}{\beta\,\ln(T_\mathrm{max}/x)}. \]
  Since $\log\mathrm{MSD}_\mathrm{pool}(\tau)=2\alpha\log\tau+\log C(P(\tau))$ and $d\log x=d\log\tau$,
  \[ \frac{d\log C}{d\log x}=\frac{1}{\ln(T_\mathrm{max}/x)}-\frac{\beta\,x^\beta}{T_\mathrm{max}^\beta-x^\beta}
     =\beta\,g(y),\qquad g(y)=\frac1y-\frac1{e^y-1},\qquad y=\beta\ln\frac{T_\mathrm{max}}{x}, \]
  \[ \alpha_\mathrm{pool}(\tau)=\frac12\,\frac{d\log\mathrm{MSD}_\mathrm{pool}}{d\log\tau}=\alpha+\frac{\beta}{2}\,g(y). \]
  Below $\Support\,T_\mathrm{min}=1\,000$~s every piece counts, $C$ is constant and $\alpha_\mathrm{pool}=\alpha$.
  Above, $g$ grows from $g(\beta\ln(T_\mathrm{max}/T_\mathrm{min}))$ to $\tfrac12$ as $y\to0$: $\alpha_\mathrm{pool}$
  jumps to \ToyPlJump{} and rises toward $\alpha+\beta/4=\ToyPlTop$ as $\tau\to\Support\,T_\mathrm{max}$, where
  only the longest pieces are left. So the pool has no single slope: it is steeper than every piece, and more so
  at longer lags, although every piece has the exponent $\alpha$. Over the window of \ToyPlTau~s the fit gives
  $\hat\alpha_\mathrm{pool}=\ToyPlPool$, the formula of the problem box; the fits over the windows follow the local
  curve (Fig.~\ref{fig:chained-scheme}, bottom right).

  \textbf{The mean of the cohorts is not exact either.} Cohorts to \ToyPlCohorts~s, one weight per cohort. While
  every cohort reaches the whole window the mean is over a fixed set of cohorts, a sum of power laws with the
  same exponent, and the fit gives $\alpha$. Where a cohort ends inside the window it leaves the mean: the fit
  gives \ToyPlMean{} at \ToyPlTau~s and up to \ToyPlMeanMax, irregular along $\tau$, since it depends on where the
  ends of the cohorts fall in the window.

  The toy proves the mechanism, not its size: with $\beta=\ToyPlBeta$ the pool stays below \ToyPlTop, while in the
  data it goes above 1 (Result below).

  \begin{center}
  \includegraphics[width=0.9\textwidth]{chained_scheme.pdf}
  \captionof{figure}{The toy. Top: which pieces (rows, short to long) enter each lag of the window of
  \ToyPlTau~s in the pool (left) and each step of the chain (right). Bottom left: the MSDs of that window and
  their fits; the chained points lie on a line of slope $2\alpha$, the pool points above it and on a steeper
  line, the mean of the cohorts (purple) on a steeper line too. Bottom right: $\alpha(\tau)$ from the fit over
  every window; line, the local exponent of the pool $\alpha+\frac{\beta}{2}g(y)$; dotted, the exponent of every
  piece.}
  \label{fig:chained-scheme}
  \end{center}
\end{method}

\begin{step}
\Res \textbf{The exponent} (Fig.~\ref{fig:chained-pool}): $\alpha^*$ and the half-slopes of the pool, with the
  same fit.
  \begin{itemize}
  \item \textbf{The pool is not an option}: pool minus $\alpha^*$ is above 0, band included, from \PoolRunsAwayA~s
    in $\A$, \PoolRunsAwayG~s in $\G$ and \PoolRunsAwayS~s in $\Sl$, and its $\alpha$ runs away, up to \PoolSMax{} ($\Sl$),
    \PoolAMax{} ($\A$), \PoolGMax{} ($\G$), above the ballistic $\alpha=1$, which no cohort reaches.
    Pool minus $\alpha^*$ at 6\,220~s: \PoolSMinusChainVeryLong{} ($\Sl$), \PoolAMinusChainVeryLong{} ($\A$),
    \PoolGMinusChainVeryLong{} ($\G$), the composition term alone.
  \item $\alpha^*$ stays within \ChainSMinusCohortMax{} ($\Sl$), \ChainAMinusCohortMax{} ($\A$),
    \ChainGMinusCohortMax{} ($\G$) of the cohort to 1~h, at every lag up to \ChainSolidEnd~s, with bands from
    the paired draws.
  \end{itemize}
  \begin{figure}[H]
  \centering
  \includegraphics[width=\textwidth]{open_declared_dt5/chained_pool.pdf}
  \caption{$\alpha^*(\tau)$ (color), the pool (red) and the cohort to 1~h (dotted), every declared open
  flight. 90\% bands.}
  \label{fig:chained-pool}
  \end{figure}
\Com \textbf{What it does not solve: the selection.} The window of $\tau$ ends with the step into $\tau^+$,
  on the pieces of at least $\tau^+/\Support$~s: every piece at 10~s (\ChainGFirstPieces{} in $\G$), only the
  long flights at the end (\ChainGLastPieces{} pieces at \ChainGLastLag~s in $\G$). Since the longest flights are
  straighter at every scale (Sec.~\ref{fixed-cohort}), a change of $\alpha^*$ at long lags may come from the
  flights that change, not from the dynamics. No method removes it: the lag $\tau^+$ exists only in the pieces of
  at least $\tau^+/\Support$~s, about $\PlusFactor\,\tau$, so at \PlusLagHours~h only the pieces of about \PlusPieceHours~h count.
\Ded \textbf{Solid and dashed.} Solid while the window ends within \ChainSelFree~s ($\tau\le\ChainSolidEnd$~s):
  there every step runs on at least the pieces of the cohort to 1~h, and the selection is measured small
  (Result above), against $\G-\A=\ChainGAThousand$ at 1\,000~s. Dashed beyond: the long flights only.
\Home\mbox{}
  \begin{itemize}
  \item \textbf{One curve per population}: $\alpha^*$, the half-slope of the chained MSD: at each step every
    piece that reaches it, no composition term.
  \item Not the pool: using all the pieces at each lag, $\alpha$ goes above 1 at long lags.
  \item Solid up to \ChainSolidEnd~s, where the selection is small; dashed beyond, the long flights only.
  \item Against the cohort to 1~h, $\G-\A$ hardly moves (\ChainGAThousand{} against \GapAllGAThousand{} at
    1\,000~s); $\A-\Sl$ does, at short lags (\ChainASTen{} against \GapAllASTen{} at 10~s).
  \end{itemize}
\end{step}

\subsection{Ground class}\label{ground-class}

\begin{method}{the ground class of a flight}
  \begin{itemize}
  \item \emph{Ground}: the IGN BD ALTI v2 terrain model, 25~m, averaged over blocks of \GroundCell~m. It is the
    bare ground under the pilot, whatever the altitude of the flight.
  \item \emph{Points}: the positions of the flight every 10~s, without the first and the last \GroundTrim\% of
    its duration, where the pilot reaches the zone of the flight from the takeoff and leaves it toward the
    landing.
  \item \emph{Ground of the flight}: the median of the ground under these points. Points off the model are
    dropped; a flight with less than \GroundMinOnDem\% of its points on it has no class (\GroundUnknown\% of
    the flights).
  \item \emph{Classes}: plains $<\GroundEdgeLow$~m $\le$ hills $<\GroundEdgeMid$~m $\le$ low mountains
    $<\GroundEdgeHigh$~m $\le$ high mountains.
  \end{itemize}
  The trim changes the class of \GroundChanged\% of the flights. One class per flight: its group time and its
  solo time take the class of the flight. In each class, $\alpha^*$ and its gaps with the two-way bands
  (Sec.~\ref{which-cohort}).
\end{method}

\subsection{Opening}\label{open-flights}\label{opening}

\begin{method}{the opening of a flight}
  With $x_0$, $x_\mathrm{end}$ the first and last fix and $\ell$ the path flown,
  \[ \mathcal{O} = \frac{|x_\mathrm{end}-x_0|}{\ell}\in[0,1]: \]
  1 for a straight flight, near 0 for a flight that comes back where it started. It describes the shape of the
  whole flight.

  \emph{Declared open is not open.} Every flight of the story is declared open on the FFVL catalogue (free
  distance or distance with 1--3 points), yet the median $\mathcal O$ is \MedianEvery{} (Fig.~\ref{fig:opening},
  left), and the long lags rest on the most folded flights: median \MedianLongest{} and \BelowLongest\% below
  \MinOpening{} (against \BelowEvery\%) in the cohort to \DeclTauLongest~s. A flight that comes back moves away
  and then back, so its displacement stops growing, and at long lags $\alpha$ falls.

  \emph{With what it goes.} With $z_0$ the altitude of the takeoff and $T$ the duration of the flight, $\mathcal O$
  is lower in the mountains, where relief and airspace bend the route
  (median \TakeoffHighMedian{} for takeoffs at $z_0\ge\MountainM$~m, \TakeoffLowMedian{} below), and in the long
  flights (median \MedianEvery{} on every flight, \MedianFixed{} in the cohort to \DeclTauFixed~s, \MedianLongest{}
  to \DeclTauLongest~s). A threshold $\mathcal O>\theta$ would keep these shares of the flights:
  \begin{center}\footnotesize
  \begin{tabular}{@{}l c c c c@{}}
  & $z_0\ge\MountainM$~m & $z_0<\MountainM$~m & $T\ge\DurationS$~s & $T<\DurationS$~s \\
  \midrule
  $\theta=\Stricter$ & \TakeoffHighStrict\% & \TakeoffLowStrict\% & \DurationLongStrict\% & \DurationShortStrict\% \\
   & (\TakeoffHighStrictN) & (\TakeoffLowStrictN) & (\DurationLongStrictN) & (\DurationShortStrictN) \\
  $\theta=\MinOpening$ & \TakeoffHighKeep\% & \TakeoffLowKeep\% & \DurationLongKeep\% & \DurationShortKeep\% \\
   & (\TakeoffHighKeepN) & (\TakeoffLowKeepN) & (\DurationLongKeepN) & (\DurationShortKeepN) \\
  \midrule
  declared open flights & \TakeoffHighFlights & \TakeoffLowFlights & \DurationLongFlights & \DurationShortFlights \\
  \end{tabular}
  \end{center}

  \textbf{Our choice: no threshold.} A threshold on $\mathcal O$ is a condition on the shape of the paths, not a
  cleaning of the data: it selects the region (the mountains drop first), the length (the long flights drop first)
  and the straightness itself, the outcome we measure. So every declared open flight is kept, and $\mathcal O$ is a
  property of the flight: whether the group effect depends on it is a question of the analysis
  (Sec.~\ref{ch:flights}).
  \begin{center}
  \centering
  \includegraphics[width=0.85\textwidth]{open_declared_dt5/opening.pdf}
  \captionof{figure}{Opening of the declared open flights. Left: every flight and the cohorts of the long lags. Right:
  takeoffs below and above \MountainM~m. Vertical line: \MinOpening.}
  \label{fig:opening}
  \end{center}
\end{method}

\subsection{Phases}\label{phases}

\begin{method}{phases and their statistics}
  \emph{Labels.} Every fix is labelled transition (glide), search or climb by the three-state hidden Markov
  model of Vilpellet, Darmon and Benzaquen~\cite{vilpellet2026}, refitted on our flights
  (Appendix~\ref{app:phases}). Tracks logged every 2 to \DtMax~s are splined to 1~s first. A run is a maximal
  sequence of fixes with one label; it belongs to the population of the piece that holds its middle, whole
  (runs whose middle lies in no piece are dropped), since cutting it where a flight enters or leaves a group would shorten the durations.

  \emph{What a search is.} The model labels search whatever is neither a climb nor a straight glide: a pilot
  turning without rising. The same label covers four moments of the cycle, about a quarter each
  (\PhCtxFound\%, \PhCtxFailed\%, \PhCtxRecentre\%, \PhCtxExit\% of all the searches), told apart only by the runs
  before and after it:
  \begin{center}\footnotesize
  \begin{tabular}{@{}l l l@{}}
  \toprule
  sequence & moment & in the statistics \\
  \midrule
  glide $\to$ search $\to$ climb & a search that finds a thermal & search \\
  glide $\to$ search $\to$ glide & a search that fails & search \\
  climb $\to$ search $\to$ climb & the thermal lost and found again & part of the climb \\
  climb $\to$ search $\to$ glide & leaving the thermal & exit, apart \\
  \bottomrule
  \end{tabular}
  \end{center}
  So a search is what follows a glide; two climbs joined by a re-centring are one climb; the exit is kept out of
  the searches and of the glides, and only its duration is read. Searches at the edge of a segment, whose
  neighbour is unknown (\PhSearchDropped\%), are left out. The model alone cannot tell these moments apart,
  since the pilot turns without rising in all four; the sequence does it directly, with no new training. Vilpellet's
  labels show the same four moments in shares within \PhCtxMaxDiff{} points.

  \emph{Statistics}, one value per flight and population (the part of the flight in it), then averaged over the
  flights:
  \begin{itemize}
  \item glide: share of time; glide speed; height lost per km;
  \item search: share of time; searches per hour; share of searches that find a thermal, found over found
    and failed;
  \item exit: its duration;
  \item climb: climbs per hour; climb rate; share of time; per climb, the height gained (vertical) and the
    horizontal distance from its first to its last fix (its chord); the drift speed, chord over duration, a proxy
    of the wind at the height of the thermals.
  \end{itemize}
  Each statistic is a ratio of a numerator $n$ and a denominator $d$ (for the climb rate, the height gained and
  the time in climb). Average first, ratio after: for each unit $u$, the part of a flight in one population, of
  duration $T_u$, the rates $n_u/T_u$ and $d_u/T_u$; then
  \[ X=\frac{\sum_u w_u\,n_u/T_u}{\sum_u w_u\,d_u/T_u}, \]
  with $w_u=1$ for the value and the bootstrap counts in the draws. Each unit weighs by its rates, not by its
  length. They answer how the motion is flown, where
  $\alpha$ only says how persistent it is.
\end{method}

\subsection{Tails of the durations}\label{tails}

\begin{method}{the tail exponent and its threshold}
  \emph{Why.} The means of the phase statistics say how a phase is flown on average; the tails of the durations
  say how often a glide, a search or a climb is very long, as in the distributions of durations of Vilpellet,
  Darmon and Benzaquen~\cite{vilpellet2026}.

  The durations $x$ of glides, of searches that end in a climb and of climbs: survival $P(X>x)$ and, above a
  threshold $x_\mathrm{min}$, the Hill estimate of the tail exponent of $P(X>x)\sim x^{-\mu}$,
  \[ \hat\mu=\frac{n}{\sum_{i:\,x_i\ge x_\mathrm{min}}\ln(x_i/x_\mathrm{min})}, \]
  over the $n$ durations above $x_\mathrm{min}$.

  \emph{The threshold.} For each curve, the $x_\mathrm{min}$ that minimises the Kolmogorov--Smirnov distance
  between the data above it and the fitted power law~\cite{clauset2009}. We read $\mu$ of $\Sl$, $\A$ and $\G$
  above one threshold per phase, the largest of their three thresholds, re-chosen in every bootstrap draw so
  that the band includes it: above it every curve follows its power law, and the three exponents describe the
  same durations. Separate thresholds would give each tail its own best estimate, over different ranges of
  durations, and their difference would mix the range with the population. One threshold for the three also
  keeps the identity of Sec.~\ref{question} exact for $\mu$. The per-curve thresholds and the Hill plots
  $\hat\mu(x_\mathrm{min})$ are the check (Appendix~\ref{app:tails}).
\end{method}

\subsection{The route}\label{route}

\begin{method}{the route}
  The path with each climb replaced by its chord, from its first to its last fix, and parametrised by the arc
  length: the waiting in the climbs is gone, the shape of the path is left. Its exponent $\alpha_\mathrm{route}$
  is read on the MSD in arc length as $\alpha$ is in time, as in our study of the ridges (xc-lab). In arc length
  with the circles kept, the exponent equals the one in time within \RouteCheckMax{} (Sec.~\ref{ch:geometry}): a
  paraglider waits while circling, so only the chord removes the waiting.
\end{method}

\section{Is group time more persistent?}\label{ch:straighter}

\begin{step}
\Q Does a pilot's motion persist more while flying in a group than when flying alone?
\Why It is the starting fact: without it there is nothing to explain.
\Tool $\alpha^*(\tau)$ of $\G$ and $\Sl$ at every lag, and the gap $\Delta\alpha^*_{\G-\Sl}(\tau)$, every flight
  (Sec.~\ref{which-cohort}).
\Res The curves $\alpha^*$ and their gaps (Fig.~\ref{fig:chained}); $\A$ is drawn already, for the next chapter.
  \begin{figure}[H]
  \centering
  \includegraphics[width=\textwidth]{open_declared_dt5/chained.pdf}
  \caption{Left: $\alpha^*(\tau)$ of $\Sl$, $\A$, $\G$; solid up to \ChainSolidEnd~s, dashed beyond (long
  flights only); dotted, the cohort to 1~h. Right: the gaps. Every declared open flight, 90\% bands from the
  shared two-way draws.}
  \label{fig:chained}
  \end{figure}
\Obs $\G-\Sl$ is about zero up to 100~s (\ChainGSTen{} at 10~s, \ChainGSHundred{} at 100~s), grows from about
  \ChainGSGrow~s, is \ChainGSThreeTwenty{} at 320~s and \ChainGSThousand{} at 1\,000~s, and stays at that level up to the end
  of the solid range, \ChainSolidEnd~s.
\Ded The gap is real at the lags of a stretch of flight, hundreds to thousands of seconds, and absent at the lags of
  a turn. Below 100~s there is no net gap; whether it hides two terms that cancel is the next question.
\Home \textbf{Yes, from a few hundred seconds on}: group time is more persistent than solo flights, by
  \ChainGSThousand{} at 1\,000~s; below 100~s there is no net gap.
\end{step}

\section{The group or the flights?}\label{ch:identity}

\begin{step}
\Q Is the gap a property of the moment spent in a group, or of the flights that have group time?
\Why Chapter~\ref{ch:straighter} found $\G-\Sl>0$ from a few hundred seconds on and no net gap below 100~s. But
  $\G-\Sl$ mixes two things: flying with others, and which pilots, days and places have group time. They call for
  different explanations.
\Tool The identity of Sec.~\ref{question}: $\G-\A$ (same flights) and $\A-\Sl$ (different flights), at every lag,
  with the shared draws (Fig.~\ref{fig:identity}).
\Res At 10~s: $\G-\A=\ChainGATen$, $\A-\Sl=\ChainASTen$. At 1\,000~s: $\G-\A=\ChainGAThousand$,
  $\A-\Sl=\ChainASThousand$.
  \begin{figure}[H]
  \centering
  \includegraphics[width=0.6\textwidth]{open_declared_dt5/identity.pdf}
  \caption{$\G-\Sl$ and its two terms, $\G-\A$ and $\A-\Sl$, of $\alpha^*$; solid up to \ChainSolidEnd~s, dashed
  beyond; 90\% bands from the shared two-way draws. Every flight.}
  \label{fig:identity}
  \end{figure}
\Obs\mbox{}
  \begin{itemize}
  \item Opposite shapes. $\G-\A$ is negative at the short lags and grows from about \ChainGAGrow~s; $\A-\Sl$ is positive
    at the short lags and turns slightly negative near \ChainASNegative~s.
  \item Below 100~s the two terms cancel, and $\G-\Sl\approx0$. From about $10^3$~s the whole of $\G-\Sl$ is
    $\G-\A$.
  \end{itemize}
\Ded\mbox{}
  \begin{itemize}
  \item \textbf{At long lags the gap belongs to the moment in a group}: it lives within the same flights.
  \item The absence of a gap below 100~s is a cancellation: within the same flights group time is \emph{less}
    persistent there, and the flights with group time are more persistent than solo flights.
  \item $\A-\Sl$ compares different flights, and its sign changes with the lag. Different flights fly over
    different ground, which changes $\alpha^*$ itself: before reading $\A-\Sl$, the ground.
  \end{itemize}
\Home On every flight, the group effect proper, $\G-\A$, carries $\G-\Sl$ from about $10^3$~s on. $\A-\Sl$ is
  small and changes sign with the lag: a comparison between flights, to be read where they fly.
\end{step}

\section{Where? Plains and mountains}\label{ground}

\begin{step}
\Q Does the group effect depend on the ground under the flight?
\Why Chapter~\ref{ch:identity} left $\A-\Sl$ to be read where the flights fly. The ground is both a moderator
  and a confounder: groups fly more in the mountains (\GroundMountainG\% of the flights of $\G$, \GroundMountainS\%
  of $\Sl$), and the relief changes $\alpha^*$ itself.
\Tool The ground class of each flight (Sec.~\ref{ground-class}); in each class $\alpha^*$ and its gaps
  (Sec.~\ref{which-cohort}).
\Res Flights and gaps at 1\,000~s per class (Table~\ref{tab:ground}); the gaps at every lag
  (Fig.~\ref{fig:ground-gaps}) and $\alpha^*$ (Fig.~\ref{fig:ground-alpha}).
  \begin{table}[H]
  \centering\footnotesize
  \begin{tabular}{@{}l r r r r r r@{}}
\toprule
 & \multicolumn{3}{c}{flights} & \multicolumn{3}{c}{gap at 1\,000~s} \\
\cmidrule(lr){2-4}\cmidrule(lr){5-7}
ground & $\Sl$ & $\A$ & $\G$ & $\G-\Sl$ & $\G-\A$ & $\A-\Sl$ \\
\midrule
plains & 9\,763 & 1\,548 & 2\,124 & $+$0.013 [$+$0.010, $+$0.016] & $+$0.007 [$+$0.005, $+$0.010] & $+$0.005 [$+$0.002, $+$0.008] \\
hills & 9\,815 & 1\,120 & 1\,550 & $+$0.025 [$+$0.020, $+$0.030] & $+$0.012 [$+$0.007, $+$0.018] & $+$0.012 [$+$0.007, $+$0.018] \\
low mountains & 18\,902 & 5\,277 & 6\,160 & $+$0.042 [$+$0.038, $+$0.046] & $+$0.031 [$+$0.027, $+$0.035] & $+$0.011 [$+$0.007, $+$0.015] \\
high mountains & 8\,028 & 2\,800 & 3\,112 & $+$0.055 [$+$0.050, $+$0.059] & $+$0.033 [$+$0.028, $+$0.038] & $+$0.022 [$+$0.017, $+$0.027] \\
\bottomrule
\end{tabular}

  \caption{Flights of each population at the first lag and gaps of $\alpha^*$ at 1\,000~s per ground class,
  90\% bands. Every declared open flight with a class.}
  \label{tab:ground}
  \end{table}
  \begin{figure}[H]
  \centering
  \includegraphics[width=\textwidth]{open_declared_dt5/ground_gaps.pdf}
  \caption{Gaps $\Delta\alpha^*(\tau)$ in each ground class (color) and on every flight (gray); solid up to
  \ChainSolidEnd~s, dashed beyond (long flights only); 90\% bands from the shared two-way draws.}
  \label{fig:ground-gaps}
  \end{figure}
  \begin{figure}[H]
  \centering
  \includegraphics[width=\textwidth]{open_declared_dt5/ground_alpha.pdf}
  \caption{$\alpha^*(\tau)$ of $\Sl$, $\A$, $\G$ in each ground class; solid up to \ChainSolidEnd~s, dashed
  beyond; 90\% bands.}
  \label{fig:ground-alpha}
  \end{figure}
\Obs\mbox{}
  \begin{itemize}
  \item \emph{$\G-\Sl$ grows with the relief}: at 1\,000~s \GroundPlainsGSThousand{} in the plains,
    \GroundHillsGSThousand{} in the hills, \GroundLowMtnGSThousand{} in the low and \GroundHighMtnGSThousand{} in
    the high mountains. $\G-\A$ the same way, from \GroundPlainsGAThousand{} to \GroundHighMtnGAThousand. In the
    plains $\G-\A$ stays negative or unresolved up to about \GroundPlainsGAPositive~s.
  \item \emph{$\A-\Sl$ is positive in every class}: at 1\,000~s \GroundPlainsASThousand, \GroundHillsASThousand,
    \GroundLowMtnASThousand, \GroundHighMtnASThousand; on every flight it is \ChainASThousand.
  \item \emph{$\alpha^*$ changes regime with the ground} (Fig.~\ref{fig:ground-alpha}): in the plains and the
    hills it rises again beyond about 500~s for the three populations; in the mountains it falls from about
    100~s on, and $\G$ falls last.
  \end{itemize}
\Ded\mbox{}
  \begin{itemize}
  \item \textbf{In absolute terms the group effect grows with the relief}: in the mountains $\G-\Sl$ is several
    times its value in the plains. But $\alpha^*$ itself is lower there: the next step compares the classes on
    a common scale.
  \item \textbf{The negative $\A-\Sl$ on every flight is a composition effect} (Simpson's paradox). $\Sl$ flies
    more in the lowlands than $\A$ (\GroundLowlandS\% of its flights in plains or hills, against
    \GroundLowlandA\%), where $\alpha^*$ is higher for every population. Mixing the classes puts $\Sl$ above
    $\A$, while within each class $\A$ is above $\Sl$.
  \item The curves on every flight mix two regimes of $\alpha^*$: a comparison of the populations has to hold
    the ground class fixed.
  \end{itemize}
\Home\mbox{}
  \begin{itemize}
  \item In absolute terms the group effect grows with the relief: small in the plains, largest in the high
    mountains (normalised, see the next step).
  \item The solo time of group flights is straighter than solo flights in every class; the opposite sign on
    every flight comes from where they fly.
  \item Compare the populations at fixed ground class.
  \end{itemize}
\end{step}


\begin{step}
\Q Is the effect larger in the mountains, or the same effect on another scale?
\Why $\alpha^*$ itself differs between the classes, so the same $\Delta\alpha^*$ is a different share of what a
  solo flight lacks to be ballistic. To compare the classes we divide each gap by $1-\alpha^*_b$, our choice of
  normalisation, used only to compare ground classes.
\Tool The normalised gaps $(\alpha^*_a-\alpha^*_b)/(1-\alpha^*_b)$ per class, computed in every draw
  (Fig.~\ref{fig:ground-rel}); their peak and the lag of the peak.
\Res Peak of the normalised $\G-\Sl$: \GroundPlainsPeakGS{} at \GroundPlainsPeakGSLag~s (plains),
  \GroundHillsPeakGS{} at \GroundHillsPeakGSLag~s (hills), \GroundLowMtnPeakGS{} at \GroundLowMtnPeakGSLag~s (low
  mountains), \GroundHighMtnPeakGS{} at \GroundHighMtnPeakGSLag~s (high mountains). Of $\G-\A$:
  \GroundPlainsPeakGA, \GroundHillsPeakGA, \GroundLowMtnPeakGA, \GroundHighMtnPeakGA.
  \begin{figure}[H]
  \centering
  \includegraphics[width=\textwidth]{open_declared_dt5/ground_relgaps.pdf}
  \caption{Normalised gaps $(\alpha^*_a-\alpha^*_b)/(1-\alpha^*_b)$ per ground class; solid up to \ChainSolidEnd~s,
  dashed beyond; 90\% bands from the shared two-way draws.}
  \label{fig:ground-rel}
  \end{figure}
\Obs\mbox{}
  \begin{itemize}
  \item The peaks of the normalised $\G-\A$ are close in the four classes (\GroundPlainsPeakGA{} to
    \GroundHighMtnPeakGA); those of $\G-\Sl$ are of the same order, the hills highest, within wide bands.
  \item The lag of the peak of $\G-\Sl$ is shorter in the mountains (\GroundLowMtnPeakGSLag{} and
    \GroundHighMtnPeakGSLag~s) than in the plains and the hills, where it lies at the last solid lag
    (\GroundPlainsPeakGSLag{} and \GroundHillsPeakGSLag~s): there it is a lower bound.
  \end{itemize}
\Ded \textbf{Normalised, the effect has about the same size on every ground; the ground sets the scale at which
  it shows}: about $10^3$~s or less in the mountains, at least the end of the solid range in the plains and the
  hills. The growth of the absolute gap with the relief is mostly the lower $\alpha^*$ of the mountains.
\Com Peaks over the solid lags only. In the plains and the hills the peak of $\G-\Sl$ rests on the longer flights.
\Home Same effect, different scale.
\end{step}

\begin{step}
\Q Does the effect peak at the scale of one glide and climb cycle?
\Why The step above found the effect at a shorter scale in the mountains. A cycle is the natural unit of progress
  of a soaring flight; if the cycle is shorter in the mountains, it would explain that scale.
\Tool The mean length of a cycle, 60/climbs per hour, of $\Sl$, $\A$, $\G$ per class (Sec.~\ref{phases}), against
  the lag of the peak (Table~\ref{tab:cycle}).
\Res\mbox{}
  \begin{table}[H]
  \centering\footnotesize
  \begin{tabular}{@{}l r r r r r@{}}
\toprule
 & \multicolumn{3}{c}{cycle (min)} & \multicolumn{2}{c}{peak lag (s)} \\
\cmidrule(lr){2-4}\cmidrule(l){5-6}
ground & $\Sl$ & $\A$ & $\G$ & $\G-\A$ & $\G-\Sl$ \\
\midrule
plains & 11.3 & 11.3 & 10.2 & 1\,500 & 2\,250 \\
hills & 11.2 & 10.7 & 10.1 & 2\,250 & 2\,250 \\
low mountains & 11.4 & 11.7 & 10.2 & 710 & 710 \\
high mountains & 11.6 & 11.8 & 10.1 & 870 & 1\,230 \\
all & 11.3 & 11.5 & 10.2 & 1\,230 & 870 \\
\bottomrule
\end{tabular}

  \caption{Mean length of a glide and climb cycle per population and ground class, and the lag of the peak of
  the normalised gaps (solid lags).}
  \label{tab:cycle}
  \end{table}
\Obs The cycle lasts between \PhCycleGHighMtn{} and \PhCycleAHighMtn~min in every class, \PhCycleShortMin{} to
  \PhCycleShortMax~min shorter in $\G$ than in $\A$; the peaks lie at \PhCycleRatioMin{} to \PhCycleRatioMax{} cycles and move between the classes while the cycle does not.
\Ded \textbf{No}: the scale of the effect is not set by the cycle. Why it is shorter in the mountains stays open
  (Sec.~\ref{ch:found}).
\Home The ground sets the scale of the effect, the cycle does not; from here on, a comparison made on every
  flight is checked at fixed ground class.
\end{step}

\section{Which flights? The opening}\label{ch:flights}

\subsection{A first look: where group time is straighter}\label{first-look}

\begin{step}
\Q Does flying in a group add straightness whatever the geometry of the path?
\Why Chapter~\ref{ground} found the effect on every ground, at a scale set by the ground. The shape of the whole
  flight is another property of the flights compared: most declared open flights come back toward their start
  (Sec.~\ref{opening}), and a group may straighten some shapes and not others. A threshold on $\mathcal O$ is the
  obvious way to look, and where we started.
\Tool The nested cohorts (Sec.~\ref{fixed-cohort}) of every declared open flight and of those with
  $\mathcal O>\MinOpening$ and $\mathcal O>\Stricter$: their $\alpha(\tau)$ and the gaps; the numbers in the cohort
  to \GapCohortTau~s.
\Res Fig.~\ref{fig:declared-cohorts} and \ref{fig:gaps-opening}. At 1\,000~s, $\G-\Sl$ is \GapAllGSThousand{} (a),
  \GapMinGSThousand{} (b) and \GapStrictGSThousand{} (c); $\G-\A$ is \GapAllGAThousand, \GapMinGAThousand{} and
  \GapStrictGAThousand. $\A-\Sl$ in (a): \GapAllASTen{} at 10~s, \GapAllASThousand{} at 1\,000~s, \GapAllASLong{}
  at \OpenTauFixed~s.
  \begin{figure}[H]
  \centering
  \textbf{(a)} every declared open flight\\
  \includegraphics[width=\textwidth]{open_declared_dt5/cohorts.pdf}\\[2pt]
  \textbf{(b)} $\mathcal O>\MinOpening$\\
  \includegraphics[width=\textwidth]{open_declared_dt5/cohorts_min0.1.pdf}\\[2pt]
  \textbf{(c)} $\mathcal O>\Stricter$\\
  \includegraphics[width=\textwidth]{open_declared_dt5/cohorts_min0.5.pdf}
  \caption{$\alpha(\tau)$ of the nested cohorts (Sec.~\ref{fixed-cohort}) on the declared open flights: (a) all of
  them, with no condition on the opening; (b) opening above \MinOpening; (c) above \Stricter. One weight per
  piece, 90\% bands, same axes. Each curve stops at the $\tau_\mathrm{max}$ of its cohort.}
  \label{fig:declared-cohorts}
  \end{figure}
  \begin{figure}[H]
  \centering
  \textbf{(a)} every declared open flight\\
  \includegraphics[width=\textwidth]{open_declared_dt5/cohorts_gaps.pdf}\\[2pt]
  \textbf{(b)} $\mathcal O>\MinOpening$\\
  \includegraphics[width=\textwidth]{open_declared_dt5/cohorts_gaps_min0.1.pdf}\\[2pt]
  \textbf{(c)} $\mathcal O>\Stricter$\\
  \includegraphics[width=\textwidth]{open_declared_dt5/cohorts_gaps_min0.5.pdf}
  \caption{Gaps $\Delta\alpha(\tau)$ in the nested cohorts at the three thresholds on the opening; 90\% bands
  from the shared two-way draws, same axes.}
  \label{fig:gaps-opening}
  \end{figure}
\Obs\mbox{}
  \begin{itemize}
  \item In (a) and (b) the curves of $\G$ sit above those of $\Sl$ and $\A$ and resist the fall of $\alpha$ at long
    lags; $\G$ is straighter than $\A$ from $\tau\approx\GapAllGAFrom$~s in (a) and
    \GapMinGAFrom~s in (b) on.
  \item The gaps shrink as the threshold rises. In (c) the curves are close, and the gaps small but still resolved
    from about \GapStrictGSFrom~s ($\G-\Sl$) and \GapStrictGAFrom~s ($\G-\A$) on.
  \item $\A-\Sl$ has the opposite shape: largest at 10~s, about zero at \OpenTauFixed~s, where
    $\G-\Sl\approx\G-\A$.
  \end{itemize}
\Ded \textbf{At first sight the group effect depends on the geometry of the path}: the more open the
  flights kept, the smaller it is. Two readings: the flights that fly in a group have another shape (next step), or
  the threshold changes which flights are compared (last step of this chapter, Sec.~\ref{opening-ground}).
\Home On every flight, the more open the flights kept, the smaller the group effect; the next two steps ask why.
\end{step}

\subsection{The flight or the stretch?}\label{flight-or-stretch}

\begin{step}
\Q Is $\G-\Sl$ a property of the flights that fly in a group, or of their group time?
\Why The first look found the effect shrinking as the flights kept get more open. If the flights that fly in a
  group had another shape, $\G-\Sl$ would only say which flights fly in a group.
\Hyp \textbf{The flight.} Group flights are more open as a whole, and more open flights have a higher $\alpha$
  (Sec.~\ref{open-flights}): then $\G-\Sl$ would only say which flights fly in a group.
\Tool Opening of solo flights and of group flights (entered a group at least once), every declared open
  flight (Fig.~\ref{fig:opening-kind}).
  \begin{figure}[H]
  \centering
  \includegraphics[width=0.55\textwidth]{open_declared_dt5/opening_kind.pdf}
  \caption{Opening of the declared open solo and group flights. Solid: \MinOpening; dotted: \Stricter.}
  \label{fig:opening-kind}
  \end{figure}
\Res Median opening of solo and group flights, on every flight and by takeoff altitude and duration:
  \begin{center}
  \begin{tabular}{@{}l c c@{}}
  median $\mathcal O$ & solo flights & group flights \\
  \midrule
  every flight & \KindSoloMedian & \KindGroupMedian \\
  $z_0<\MountainM$~m & \KindSoloMedianLow & \KindGroupMedianLow \\
  $z_0\ge\MountainM$~m & \KindSoloMedianHigh & \KindGroupMedianHigh \\
  $T<\DurationS$~s & \KindSoloMedianShort & \KindGroupMedianShort \\
  $T\ge\DurationS$~s & \KindSoloMedianLong & \KindGroupMedianLong \\
  \midrule
  flights & \KindSoloFlights & \KindGroupFlights \\
  median duration (s) & \KindSoloDuration & \KindGroupDuration \\
  takeoffs at $z_0\ge\MountainM$~m & \KindSoloMountain\% & \KindGroupMountain\% \\
  \end{tabular}
  \end{center}
\Obs Group flights are never more open: slightly less in every stratum, as open among the flights of at least
  \DurationS~s.
\Ded\mbox{}
  \begin{itemize}
  \item \textbf{The hypothesis of the flight fails}: group flights are not more open, they are slightly less.
  \item Three facts point inside the flight:
    \begin{enumerate}
    \item group flights are not more open than solo flights (Fig.~\ref{fig:opening-kind});
    \item $\G-\A>0$, and $\G$ and $\A$ come from the same flights, so the opening of the flight cannot
      explain it;
    \item on every flight $\A-\Sl$ is small at the long lags (Sec.~\ref{ch:identity}): there, outside its group
      time a group flight is about as straight as a solo flight.
    \end{enumerate}
  \item \textbf{So group time looks like the stretch where the flight gets straighter}, in flights that, as a
    whole, are not more open. The direct test, the opening of each piece, is still to do (Comment).
  \end{itemize}
\Com $\mathcal O$ and $\alpha$ measure different things: $\mathcal O$ the whole path, from end to end; $\alpha$
  the straightness of the stretches up to $\tau_\mathrm{max}$. A flight can have straight legs and come back.
\Com Why the gaps shrink at $\mathcal O>\Stricter$, three readings not separated here:
  \begin{itemize}
  \item ceiling: $\alpha$ is already near 1, little room left;
  \item in a straight flight the stretches alone are straight too;
  \item fewer pieces.
  \end{itemize}
\Com Left out:
  \begin{itemize}
  \item The opening of each piece ($\G$, $\A$, $\Sl$): the direct test that $\G$ is the straighter stretch;
    it needs the path length of the pieces, not in the cache.
  \item \OpeningAboveOne{} solo flights have $\mathcal O>1$ (path shorter than the chord): a data error to check.
  \end{itemize}
\Home\mbox{}
  \begin{itemize}
  \item The group effect is not a straightness added whatever the path: it shrinks as the flights kept get
    more open.
  \item The shape of the flights that fly in a group does not explain it: they are never more open than solo
    flights.
  \item It points to the group time: the straighter stretch of its flight. At long lags the whole
    gap $\G-\Sl$ is $\G-\A$.
  \end{itemize}
\end{step}

\subsection{The opening or the ground?}\label{opening-ground}

\begin{step}
\Q Is the effect of the opening an effect of the ground?
\Why The first look was on every flight, but a threshold on $\mathcal O$ keeps the lowland flights
  (Sec.~\ref{opening}) and the effect is smaller there (Sec.~\ref{ground}): the comparison has to be made at fixed
  ground class. And the identity says where the opening can act: $\mathcal O$ is a property of the whole
  flight, the same for $\G$ and $\A$ of one flight.
\Tool The gaps per ground class and threshold on $\mathcal O$ (Table~\ref{tab:ground-opening},
  Fig.~\ref{fig:ground-opening}); the opening of $\Sl$, $\A$, $\G$ per class (Table~\ref{tab:opening-class}).
  Caution: $\mathcal O$ and $\alpha$ both measure how straight a flight is, so a threshold on $\mathcal O$ selects
  on the outcome and compresses every gap.
\Res\mbox{}
  \begin{table}[H]
  \centering\footnotesize
  \begin{tabular}{@{}l r r r r r@{}}
\toprule
 & \multicolumn{3}{c}{median $\mathcal O$} & \multicolumn{2}{c}{share of the flights (\%)} \\
\cmidrule(lr){2-4}\cmidrule(l){5-6}
ground & $\Sl$ & $\A$ & $\G$ & all & $\mathcal O>\Stricter$ \\
\midrule
plains & 0.57 & 0.61 & 0.60 & 20.0 & 50.3 \\
hills & 0.44 & 0.49 & 0.50 & 19.1 & 27.4 \\
low mountains & 0.23 & 0.18 & 0.19 & 42.2 & 16.9 \\
high mountains & 0.20 & 0.19 & 0.20 & 18.7 & 5.4 \\
\bottomrule
\end{tabular}

  \caption{Median opening of $\Sl$, $\A$, $\G$ per ground class, and the share of each class among every flight
  and among $\mathcal O>\Stricter$.}
  \label{tab:opening-class}
  \end{table}
  \begin{table}[H]
  \centering\footnotesize
  \begin{tabular}{@{}l l l l l@{}}
\toprule
 & & \multicolumn{3}{c}{opening threshold} \\
\cmidrule(l){3-5}
 ground & gap & none & $>0.1$ & $>0.5$ \\
\midrule
plains & $\G-\A$ & $+$0.007 [$+$0.005, $+$0.010] & $+$0.007 [$+$0.004, $+$0.010] & $+$0.007 [$+$0.004, $+$0.010] \\
 & $\A-\Sl$ & $+$0.005 [$+$0.002, $+$0.008] & $+$0.005 [$+$0.002, $+$0.008] & $+$0.001 [$-$0.002, $+$0.004] \\
\addlinespace
hills & $\G-\A$ & $+$0.012 [$+$0.007, $+$0.018] & $+$0.010 [$+$0.005, $+$0.016] & $+$0.006 [$+$0.001, $+$0.012] \\
 & $\A-\Sl$ & $+$0.012 [$+$0.007, $+$0.018] & $+$0.009 [$+$0.004, $+$0.014] & $+$0.000 [$-$0.004, $+$0.004] \\
\addlinespace
low mountains & $\G-\A$ & $+$0.031 [$+$0.027, $+$0.035] & $+$0.027 [$+$0.022, $+$0.031] & $+$0.009 [$+$0.001, $+$0.017] \\
 & $\A-\Sl$ & $+$0.011 [$+$0.007, $+$0.015] & $+$0.003 [$-$0.001, $+$0.008] & $-$0.004 [$-$0.010, $+$0.002] \\
\addlinespace
high mountains & $\G-\A$ & $+$0.033 [$+$0.028, $+$0.038] & $+$0.036 [$+$0.030, $+$0.041] & $+$0.017 [$+$0.009, $+$0.026] \\
 & $\A-\Sl$ & $+$0.022 [$+$0.017, $+$0.027] & $+$0.014 [$+$0.008, $+$0.019] & $-$0.001 [$-$0.008, $+$0.005] \\
\bottomrule
\end{tabular}

  \caption{$\G-\A$ and $\A-\Sl$ of $\alpha^*$ at 1\,000~s per ground class and threshold on $\mathcal O$, 90\% bands.}
  \label{tab:ground-opening}
  \end{table}
  \begin{figure}[H]
  \centering
  \includegraphics[width=\textwidth]{open_declared_dt5/ground_opening_rel.pdf}
  \caption{Normalised gaps $(\alpha^*_a-\alpha^*_b)/(1-\alpha^*_b)$ per ground class (columns) and threshold on
  $\mathcal O$ (colours): rows $\G-\Sl$, $\G-\A$, $\A-\Sl$. Solid up to \ChainSolidEnd~s, dashed beyond; 90\% bands.}
  \label{fig:ground-opening}
  \end{figure}
\Obs\mbox{}
  \begin{itemize}
  \item A threshold on $\mathcal O$ keeps the lowlands: \OpenLowlandAll\% of the flights are in plains or hills,
    \OpenLowlandStrict\% of those with $\mathcal O>\Stricter$.
  \item Within a class, $\G$ is about as open as $\Sl$, a little more in the plains and the hills, a little less
    in the low mountains (Table~\ref{tab:opening-class}). $\G$ and $\A$ fly more in the mountains than $\Sl$
    (\GroundMountainG\% and \GroundMountainA\% of their flights, against \GroundMountainS\%).
  \item $\A-\Sl$ falls to about zero at $\mathcal O>\Stricter$ in every class; normalised, slightly negative in the
    low mountains.
  \item $\G-\A$ stays positive in every class. In absolute terms it shrinks at $\mathcal O>\Stricter$ in the hills
    and the mountains; normalised by $1-\alpha^*_\A$ it changes little, within wide bands in the low mountains
    (Fig.~\ref{fig:ground-opening}).
  \end{itemize}
\Ded\mbox{}
  \begin{itemize}
  \item \textbf{The opening acts on the comparison between flights}: it selects which flights are compared and
    removes $\A-\Sl$. Within a flight, the group effect survives every threshold; what it loses in absolute terms
    is the room left by a more persistent solo time.
  \item Within a class, group time is about as open as solo flights: on every flight, group flights look less
    open mostly because they fly more over the mountains (by ground class; by takeoff altitude the shares are
    close).
  \item So the shrinking of the effect with the opening, in the first step, is mostly the ground and the
    selection of the flights. Within a class no dependence on the geometry of the path is detected, though a
    threshold that selects on the outcome could hide one.
  \end{itemize}
\Home Within a flight, no dependence of the group effect on the shape of the flight is detected; the opening,
  like the ground, acts mostly on $\A-\Sl$, which compares flights.
\end{step}

\section{Who?}\label{ch:who}

\begin{step}
\Q Does $\A-\Sl$ come from the pilots?
\Why Chapters~\ref{ground} and~\ref{ch:flights} left $\A-\Sl$ positive in every class, a difference between
  flights. The pilots who fly in groups may be other pilots: more experienced, on other wings. If so, part of $\G-\Sl$ is who flies, not the group.
\Tool $\A-\Sl'$, with $\Sl'$ the solo flights of the pilots who also have group flights (\PilotShareS\% of
  $\Sl$), next to $\A-\Sl$, per ground class (Table~\ref{tab:pilot}, Fig.~\ref{fig:pilot}). $\G-\A$ already holds
  the pilot fixed.
\Res\mbox{}
  \begin{table}[H]
  \centering\footnotesize
  \begin{tabular}{@{}l r r l l@{}}
\toprule
ground & $\Sl$ & $\Sl'$ & $\A-\Sl$ & $\A-\Sl'$ \\
\midrule
plains & 9\,763 & 8\,232 & $+$0.005 [$+$0.002, $+$0.008] & $+$0.005 [$+$0.001, $+$0.008] \\
hills & 9\,815 & 8\,072 & $+$0.012 [$+$0.007, $+$0.018] & $+$0.011 [$+$0.005, $+$0.017] \\
low mountains & 18\,902 & 15\,351 & $+$0.011 [$+$0.007, $+$0.015] & $+$0.010 [$+$0.005, $+$0.014] \\
high mountains & 8\,028 & 6\,590 & $+$0.022 [$+$0.017, $+$0.027] & $+$0.019 [$+$0.014, $+$0.024] \\
\midrule
all & 50\,189 & 41\,280 & $-$0.003 [$-$0.006, $-$0.000] & $-$0.006 [$-$0.009, $-$0.002] \\
\bottomrule
\end{tabular}

  \caption{Solo flights of every pilot ($\Sl$) and of the pilots with group flights ($\Sl'$); $\A-\Sl$ and
  $\A-\Sl'$ of $\alpha^*$ at 1\,000~s, 90\% bands, per ground class.}
  \label{tab:pilot}
  \end{table}
  \begin{figure}[H]
  \centering
  \includegraphics[width=\textwidth]{open_declared_dt5/pilot.pdf}
  \caption{$\A-\Sl$ and $\A-\Sl'$ of $\alpha^*$ per ground class, and $\G-\Sl'$; solid up to \ChainSolidEnd~s, dashed beyond; 90\%
  bands.}
  \label{fig:pilot}
  \end{figure}
\Obs $\A-\Sl'$ follows $\A-\Sl$ in every class, within the bands, equal or a little smaller (Table~\ref{tab:pilot}).
  The pilots with group flights fly most of the solo flights.
\Ded \textbf{$\A-\Sl$ does not come from which pilots fly, to first order}: the pilots with group flights, alone on a
  flight without group time, are less persistent than alone on a flight with group time. $\Sl'$ restricts the
  pilots but does not pair them: each pilot weighs by its solo flights in $\Sl'$ and by its group flights in $\A$.
  What is left: when and where within a class (day, site), and a carry-over of the group on the solo time around it (Sec.~\ref{ch:when}).
\Home Not who flies, to first order: $\A-\Sl$ belongs to the flights with group time more than to their
  pilots.
\end{step}

\section{When in the flight?}\label{ch:when}

\begin{step}
\Q Is group time more persistent than the solo time just before and just after it?
\Why Chapters~\ref{ground}--\ref{ch:who} read $\A-\Sl$, the comparison between flights. $\G-\A$, the group effect
  proper, compares the group time with any solo time of the same flight; but the moment matters: the start of a
  flight, the hour, the conditions that change. The timeline is also the first test of the direction
  (Sec.~\ref{which-way}).
\Tool For every stretch of group time, four windows of \WhenWindow~s of the same flight (\texttt{prepost.py}):
  \emph{pre}, the window before it, if a solo stretch at least that long precedes it; \emph{start} and \emph{end},
  its first and last window, which overlap when the stretch is shorter than two windows; \emph{post}, the window
  after it, if a solo stretch at least that long follows it. An event is a stretch with the four windows. Windows
  of one length need no chain: every window reaches every lag. Control: two consecutive solo windows of the same
  flights, with no change of state. Per ground class.
\Res Windows of \WhenWindow~s, the exponent fitted over the lags from 10~s to 1\,000~s of each window
  (Table~\ref{tab:prepost}, Fig.~\ref{fig:prepost}); \WhenAllEvents{} events with the four windows. Every flight:
  start $-$ pre $=\WhenAllStartMinusPre$ \WhenAllStartMinusPreBand, end $-$ start $=\WhenAllEndMinusStart$, post $-$
  end $=\WhenAllPostMinusEnd$ \WhenAllPostMinusEndBand; two consecutive solo windows of group flights differ by
  \WhenAllControl.
  \begin{table}[H]
  \centering
  \resizebox{\textwidth}{!}{\begin{tabular}{@{}l r r r r r@{}}
\toprule
ground & events & start $-$ pre & end $-$ start & post $-$ end & control \\
\midrule
plains & 433 & $+$0.004 [$-$0.001, $+$0.008] & $-$0.002 [$-$0.007, $+$0.003] & $+$0.000 [$-$0.004, $+$0.004] & $+$0.005 [$+$0.003, $+$0.007] \\
hills & 327 & $+$0.000 [$-$0.007, $+$0.008] & $+$0.012 [$+$0.005, $+$0.019] & $-$0.006 [$-$0.013, $+$0.001] & $+$0.013 [$+$0.010, $+$0.017] \\
low mountains & 2\,426 & $+$0.000 [$-$0.004, $+$0.004] & $+$0.005 [$+$0.002, $+$0.009] & $-$0.001 [$-$0.005, $+$0.002] & $+$0.005 [$+$0.003, $+$0.007] \\
high mountains & 1\,299 & $+$0.001 [$-$0.004, $+$0.006] & $+$0.006 [$+$0.001, $+$0.010] & $-$0.004 [$-$0.008, $+$0.001] & $+$0.002 [$-$0.001, $+$0.005] \\
all & 4\,591 & $+$0.002 [$-$0.001, $+$0.004] & $+$0.004 [$+$0.002, $+$0.007] & $-$0.002 [$-$0.004, $+$0.000] & $+$0.005 [$+$0.004, $+$0.007] \\
\bottomrule
\end{tabular}
}
  \caption{Events and differences of the exponent between the windows of the timeline, per ground class, and the
  control (two consecutive solo windows of group flights), 90\% bands.}
  \label{tab:prepost}
  \end{table}
  \begin{figure}[H]
  \centering
  \includegraphics[width=\textwidth]{open_declared_dt5/prepost_ground.pdf}
  \caption{The exponent of the windows before, at the start, at the end and after each stretch of group time, per
  ground class, 90\% bands.}
  \label{fig:prepost}
  \end{figure}
\Obs Over 10 to 1\,000~s the four windows have the same exponent, within differences as large as the drift
  between two consecutive solo windows.
\Ded The fit over 10 to 1\,000~s sees no change. It weighs the middle of the range most, around 100~s, where
  $\G-\A\approx0$: negative below about \ChainGAGrow~s and positive above (Sec.~\ref{ch:identity}), so the two cancel.
\Home Over the whole range the timeline is flat; it has to be read where the effect lives.
\end{step}

\begin{step}
\Q Read on the long lags, does the persistence change when the pilot joins and leaves the group?
\Why The step above found the timeline flat over 10 to 1\,000~s, a range where $\G-\A$ changes sign.
\Tool The same windows, the exponent fitted over \WhenLongLo{} to \WhenLongHi~s only; the shares of time in glide,
  search and climb of each window (our labels), as the control of the composition: groups may form in climbs, so
  the start window may hold a climb by construction.
\Res Every flight: start $-$ pre $=\WhenAllLongStartMinusPre$ \WhenAllLongStartMinusPreBand, end $-$ start
  $=\WhenAllLongEndMinusStart$ \WhenAllLongEndMinusStartBand, post $-$ end $=\WhenAllLongPostMinusEnd$
  \WhenAllLongPostMinusEndBand; the control of two consecutive solo windows of group flights is
  \WhenAllLongControl{} \WhenAllLongControlBand. Time in climb: \WhenAllClimbPre\% before, \WhenAllClimbStart\% at the
  start, \WhenAllClimbEnd\% at the end, \WhenAllClimbPost\% after (Table~\ref{tab:prepost-long},
  Fig.~\ref{fig:prepost-long}, \ref{fig:prepost-phases}).
  \begin{table}[H]
  \centering
  \resizebox{\textwidth}{!}{\begin{tabular}{@{}l r r r r r r r r@{}}
\toprule
 & \multicolumn{4}{c}{exponent over 320--1000~s} & \multicolumn{4}{c}{time in climb (\%)} \\
\cmidrule(lr){2-5}\cmidrule(l){6-9}
ground & start $-$ pre & end $-$ start & post $-$ end & control & pre & start & end & post \\
\midrule
plains & $+$0.002 & $+$0.001 & $-$0.007$^\ast$ & $+$0.001 & 44.9 & 45.8 & 43.9 & 43.2 \\
hills & $+$0.005 & $+$0.013$^\ast$ & $-$0.013$^\ast$ & $+$0.009$^\ast$ & 42.9 & 46.3 & 41.2 & 40.9 \\
low mountains & $+$0.006 & $+$0.011$^\ast$ & $-$0.020$^\ast$ & $-$0.001 & 29.5 & 33.1 & 31.8 & 27.8 \\
high mountains & $+$0.011$^\ast$ & $+$0.001 & $-$0.017$^\ast$ & $-$0.003 & 32.3 & 34.5 & 31.3 & 28.1 \\
all & $+$0.008$^\ast$ & $+$0.007$^\ast$ & $-$0.017$^\ast$ & $+$0.001 & 32.9 & 35.7 & 33.6 & 30.5 \\
\bottomrule
\end{tabular}
}
  \caption{Differences of the exponent fitted over \WhenLongLo{} to \WhenLongHi~s between the windows of the
  timeline and the control, and the share of time in climb of each window, per ground class; $^*$: the 90\% band
  excludes 0.}
  \label{tab:prepost-long}
  \end{table}
  \begin{figure}[H]
  \centering
  \includegraphics[width=\textwidth]{open_declared_dt5/prepost_ground_long.pdf}
  \caption{The exponent over \WhenLongLo{} to \WhenLongHi~s of the windows before, at the start, at the end and after
  each stretch of group time, per ground class, with the control differences; 90\% bands.}
  \label{fig:prepost-long}
  \end{figure}
  \begin{figure}[H]
  \centering
  \includegraphics[width=\textwidth]{open_declared_dt5/prepost_phases.pdf}
  \caption{Shares of time in glide, search and climb of the windows of the timeline, per ground class (our labels).}
  \label{fig:prepost-phases}
  \end{figure}
\Obs\mbox{}
  \begin{itemize}
  \item On every flight the exponent rises into the group time, a little more by its end, and falls after it, back
    to about the level before; the fall is far outside the control.
  \item Per class (Table~\ref{tab:prepost-long}), the fall after the group is resolved in the hills and the
    mountains; the rise at the entry only in the high mountains; the rise from start to end in the hills and the
    low mountains, where in the hills the control drifts as much.
  \item The start window holds more climb than the window before, and the window after holds less.
  \end{itemize}
\Ded\mbox{}
  \begin{itemize}
  \item \textbf{On every flight the persistence rises at the entry into the group and falls at the exit}, on the
    same flights, between adjacent windows of \WhenWindow~s. Per class the fall is resolved in the hills and the
    mountains, the rise only in the high mountains, neither in the plains.
  \item It is not the composition, if a climb lowers the exponent as circling does: more climb at the start would
    lower it, which rises; less climb after would raise it, which falls.
  \end{itemize}
\Home Group time is more persistent than the solo time right before and right after it: the effect belongs to the
  time in the group more than to the flight, if the moment of the flight does not make it (next step).
\end{step}

\begin{step}
\Q Is the change at the entry and the exit made by the group, or by the moment of the flight?
\Why The window before the group lies earlier in the flight and the window after it later. If the exponent were
  low after the takeoff and before the landing, it would rise from pre to start and fall from end to post with no
  group at all.
\Tool Mock events (\texttt{prepost.py}): for every event, solo flights, never in a group, in which the four windows
  fit at the same flight-clock times inside one solo stretch, under calipers of month, local hour and wing class;
  measured like the real windows, the mocks of one event weighing one event. The steps of the mocks show what the
  moment of the flight does alone; real $-$ mock is what is left for the group, with the shared draws.
  \WhenMockEvents{} mock events for \WhenMockRealEvents{} real events.
\Res Table~\ref{tab:prepost-mock}. Every flight, mock: start $-$ pre $=\WhenMockAllStartMinusPre$
  \WhenMockAllStartMinusPreBand, end $-$ start $=\WhenMockAllEndMinusStart$ \WhenMockAllEndMinusStartBand, post $-$
  end $=\WhenMockAllPostMinusEnd$ \WhenMockAllPostMinusEndBand. Real $-$ mock: \WhenRealMinusMockAllStartMinusPre{}
  \WhenRealMinusMockAllStartMinusPreBand, \WhenRealMinusMockAllEndMinusStart{}
  \WhenRealMinusMockAllEndMinusStartBand, \WhenRealMinusMockAllPostMinusEnd{}
  \WhenRealMinusMockAllPostMinusEndBand.
  \begin{table}[H]
  \centering
  \resizebox{\textwidth}{!}{\begin{tabular}{@{}l r r r r r r r r r@{}}
\toprule
 & \multicolumn{3}{c}{real} & \multicolumn{3}{c}{mock} & \multicolumn{3}{c}{real $-$ mock} \\
\cmidrule(lr){2-4}\cmidrule(lr){5-7}\cmidrule(l){8-10}
ground & start $-$ pre & end $-$ start & post $-$ end & start $-$ pre & end $-$ start & post $-$ end & start $-$ pre & end $-$ start & post $-$ end \\
\midrule
plains & $+$0.004 & $-$0.000 & $-$0.006$^\ast$ & $-$0.000 & $+$0.005 & $-$0.004 & $+$0.004 & $-$0.005 & $-$0.002 \\
hills & $+$0.002 & $+$0.014$^\ast$ & $-$0.012$^\ast$ & $-$0.001 & $+$0.005 & $+$0.002 & $+$0.003 & $+$0.009 & $-$0.014$^\ast$ \\
low mountains & $+$0.007$^\ast$ & $+$0.011$^\ast$ & $-$0.021$^\ast$ & $+$0.001 & $+$0.002$^\ast$ & $-$0.001 & $+$0.006 & $+$0.009$^\ast$ & $-$0.019$^\ast$ \\
high mountains & $+$0.010$^\ast$ & $+$0.002 & $-$0.017$^\ast$ & $+$0.000 & $+$0.002 & $-$0.001 & $+$0.010$^\ast$ & $-$0.000 & $-$0.016$^\ast$ \\
all & $+$0.008$^\ast$ & $+$0.007$^\ast$ & $-$0.017$^\ast$ & $+$0.000 & $+$0.003$^\ast$ & $-$0.001 & $+$0.007$^\ast$ & $+$0.004$^\ast$ & $-$0.016$^\ast$ \\
\bottomrule
\end{tabular}
}
  \caption{Steps of the exponent fitted over \WhenLongLo{} to \WhenLongHi~s between the windows of the timeline: real
  events, mock events in solo flights at the same flight-clock times, and real $-$ mock, per ground class;
  $^*$: the 90\% band excludes 0.}
  \label{tab:prepost-mock}
  \end{table}
\Obs\mbox{}
  \begin{itemize}
  \item In the solo flights the moment alone moves the exponent little: no rise at the entry, no fall at the exit,
    a small rise from start to end.
  \item Real $-$ mock: the fall after the group stays on every flight and in the hills and the mountains; the rise
    at the entry stays on every flight and in the high mountains; the rise from start to end shrinks, resolved on
    every flight and in the low mountains only.
  \end{itemize}
\Ded \textbf{The moment of the flight does not make the timeline}: the fall at the exit and the rise at the entry
  belong to the group. Part of the rise within the group time is the drift of the exponent along the flight.
\Home On every flight the persistence rises at the entry and falls at the exit, beyond what the moment of the flight
  does. Per class the fall is resolved in the hills and the mountains, the rise only in the high mountains, neither
  in the plains.
\end{step}

\subsection{The solo time around the group}\label{carry}

\begin{step}
\Q Does the group carry over to the solo time around it?
\Why Chapter~\ref{ch:who} left two sources of $\A-\Sl$ within a class: the day and the site, and a carry-over. Alone
  but near the group, a pilot may keep some of its benefit. The windows of the timeline can measure it.
\Tool $\A(d)-\Sl$: windows of \WhenWindow~s in the solo time of group flights, binned by their distance $d$ from the
  nearest group time of the same flight, against windows of solo flights; the exponent over \WhenLongLo{} to
  \WhenLongHi~s, per ground class. A carry-over makes $\A(d)-\Sl$ largest near the group.
\Res Table~\ref{tab:carry}. In the high mountains \CarryHighMtnNear{} at \CarryNearLabel{} and \CarryHighMtnFar{} at
  \CarryFarLabel; in the low mountains \CarryLowMtnNear{} and \CarryLowMtnFar.
  \begin{table}[H]
  \centering
  \resizebox{\textwidth}{!}{\begin{tabular}{@{}l r rrrr@{}}
\toprule
ground & $\Sl$ & 0--20~min & 20--40~min & 40--80~min & $\ge80$~min \\
\midrule
\multicolumn{6}{@{}l}{$\A(d)-\Sl$, exponent over 320--1000~s} \\
plains & 0.937 [0.936, 0.939] & $-$0.000 [$-$0.004, $+$0.004] & $+$0.004 [$+$0.001, $+$0.008] & $+$0.001 [$-$0.004, $+$0.006] & $+$0.012 [$+$0.008, $+$0.017] \\
hills & 0.900 [0.897, 0.903] & $+$0.010 [$+$0.005, $+$0.016] & $+$0.012 [$+$0.005, $+$0.018] & $+$0.014 [$+$0.006, $+$0.023] & $+$0.024 [$+$0.011, $+$0.037] \\
low mountains & 0.858 [0.856, 0.860] & $+$0.013 [$+$0.008, $+$0.018] & $+$0.009 [$+$0.004, $+$0.014] & $+$0.023 [$+$0.018, $+$0.028] & $+$0.027 [$+$0.021, $+$0.033] \\
high mountains & 0.850 [0.847, 0.852] & $+$0.020 [$+$0.014, $+$0.026] & $+$0.018 [$+$0.012, $+$0.024] & $+$0.026 [$+$0.020, $+$0.032] & $+$0.026 [$+$0.019, $+$0.032] \\
all & 0.887 [0.884, 0.889] & $+$0.001 [$-$0.003, $+$0.004] & $+$0.001 [$-$0.003, $+$0.004] & $+$0.006 [$+$0.003, $+$0.010] & $+$0.010 [$+$0.006, $+$0.015] \\
\midrule
\multicolumn{6}{@{}l}{windows} \\
plains & 15\,779 & 889 & 872 & 553 & 418 \\
hills & 16\,554 & 637 & 687 & 460 & 174 \\
low mountains & 34\,372 & 3\,467 & 3\,533 & 2\,759 & 2\,002 \\
high mountains & 15\,286 & 1\,802 & 1\,851 & 1\,414 & 1\,199 \\
all & 88\,654 & 6\,965 & 7\,121 & 5\,306 & 3\,894 \\
\bottomrule
\end{tabular}
}
  \caption{$\A(d)-\Sl$ of the exponent over \WhenLongLo{} to \WhenLongHi~s, by the distance $d$ of the solo window
  from the nearest group time of its flight, per ground class, 90\% bands; numbers of windows.}
  \label{tab:carry}
  \end{table}
\Obs $\A(d)-\Sl$ does not fall as the solo time moves away from the group, in any class; it stays or grows.
\Ded \textbf{No sign of a carry-over} at these lags. The solo time far from the group comes from long solo
  stretches, a selection of the flight, so the growth itself is not read.
\Home No carry-over is detected over \WhenLongLo{} to \WhenLongHi~s. What is left of $\A-\Sl$ within a class is the
  day and the site, and at the short lags the phases (next chapter).
\end{step}

\section{How: the dynamics}\label{ch:dynamics}

\begin{step}
\Q Which phases of the flight does the group change?
\Why Chapter~\ref{ch:when} found that the persistence belongs to the time in the group. What changes in the way
  the pilot flies? $\alpha^*$ in time mixes the shape of the path and the way it is flown; the phases say the
  second: how long the pilot glides, searches and climbs, how fast and how well.
\Hyp Groups form mostly in the climbs, where the pilots circle together ($\G-\A<0$ at short lags,
  Sec.~\ref{ch:identity}), so the climbs go to $\G$ and $\A$ keeps more glides; and a group finds the climbs sooner
  (shorter searches, fewer weak climbs).
\Tool The phase statistics and the tails of the durations (Sec.~\ref{phases}, \ref{tails}), for $\G-\A$ and
  $\A-\Sl$, per ground class.
\Com The composition of the phases is to $\G-\A$ what the ground is to $\A-\Sl$: the grouping catches pilots in the
  same thermal, so group time may hold more climbs by construction. Two steps: the composition, then the
  performance within each phase.
\Res \textbf{Composition} (Fig.~\ref{fig:phases-composition}, Appendix~\ref{app:phases-table}): group time holds
  \PhShareClimbGAAll{} points more climb than the solo time of the same flights \PhShareClimbGAAllBand, and
  \PhCphGAAll{} more climbs per hour \PhCphGAAllBand; a glide and climb cycle lasts \PhCycleGAll~min in $\G$,
  \PhCycleAAll~min in $\A$.
  \textbf{Performance} (Table~\ref{tab:phases-main}, Fig.~\ref{fig:phases-performance}): in group time the glide is
  faster by \PhSpeedGAAll~km/h \PhSpeedGAAllBand{} and the height lost per km changes by \PhSinkGAAll~m
  \PhSinkGAAllBand; the searches find a thermal \PhFoundGAAll{} points more often \PhFoundGAAllBand; the climb rate
  is not higher (\PhRateGAAll~m/s \PhRateGAAllBand). \textbf{Tails} (Fig.~\ref{fig:phases-tails}): above
  \PhXminGlideAll~s the glides of $\G$ have a heavier tail on every flight, $\mu_\G-\mu_\A=\PhMuGlideGAAll$
  \PhMuGlideGAAllBand, resolved in the mountains only; searches and climbs differ within the bands.
  \begin{figure}[H]
  \centering
  \includegraphics[width=\textwidth]{open_declared_dt5/phases_composition.pdf}
  \caption{Composition of the phases: $\G-\A$ and $\A-\Sl$ per ground class, 90\% bands from the shared two-way
  draws. Our labels.}
  \label{fig:phases-composition}
  \end{figure}
  \begin{table}[H]
  \centering
  \resizebox{\textwidth}{!}{\begin{tabular}{@{}l l l l l l l @{}}
\toprule
 statistic & gap & plains & hills & low mountains & high mountains & all \\
\midrule
climbs per hour & $\G-\A$ & $+$0.593 [$+$0.487, $+$0.698] & $+$0.338 [$+$0.225, $+$0.451] & $+$0.761 [$+$0.702, $+$0.820] & $+$0.826 [$+$0.747, $+$0.906] & $+$0.708 [$+$0.667, $+$0.748] \\
 & $\A-\Sl$ & $+$0.0112 [$-$0.0717, $+$0.0940] & $+$0.211 [$+$0.115, $+$0.306] & $-$0.150 [$-$0.203, $-$0.096] & $-$0.0986 [$-$0.1718, $-$0.0255] & $-$0.118 [$-$0.155, $-$0.081] \\
\addlinespace
height lost per km of glide (m/km) & $\G-\A$ & $-$6.56 [$-$7.97, $-$5.16] & $-$5.86 [$-$7.74, $-$3.98] & $-$14.1 [$-$15.2, $-$12.9] & $-$12.7 [$-$14.1, $-$11.4] & $-$11.6 [$-$12.3, $-$10.9] \\
 & $\A-\Sl$ & $-$8.55 [$-$10.89, $-$6.21] & $-$6.05 [$-$8.46, $-$3.65] & $-$7.40 [$-$8.63, $-$6.18] & $-$4.31 [$-$5.79, $-$2.83] & $-$7.07 [$-$8.01, $-$6.14] \\
\addlinespace
searches ending in a climb (\%) & $\G-\A$ & $+$3.50 [$+$1.53, $+$5.47] & $+$5.05 [$+$3.01, $+$7.10] & $+$5.62 [$+$4.68, $+$6.57] & $+$7.16 [$+$5.97, $+$8.35] & $+$5.66 [$+$5.01, $+$6.32] \\
 & $\A-\Sl$ & $+$2.17 [$+$0.52, $+$3.83] & $+$3.94 [$+$2.16, $+$5.73] & $+$1.97 [$+$1.14, $+$2.79] & $+$0.367 [$-$0.688, $+$1.423] & $+$1.43 [$+$0.88, $+$1.98] \\
\addlinespace
climb rate (m/s) & $\G-\A$ & $-$0.0331 [$-$0.0507, $-$0.0155] & $-$0.0102 [$-$0.0301, $+$0.0096] & $-$0.0030 [$-$0.0153, $+$0.0093] & $+$0.0201 [$+$0.0042, $+$0.0359] & $-$0.0103 [$-$0.0188, $-$0.0018] \\
 & $\A-\Sl$ & $+$0.0107 [$-$0.0079, $+$0.0294] & $+$0.0313 [$+$0.0047, $+$0.0579] & $+$0.0453 [$+$0.0288, $+$0.0618] & $+$0.0157 [$-$0.0022, $+$0.0335] & $+$0.0630 [$+$0.0513, $+$0.0747] \\
\addlinespace
drift speed of a climb (m/s) & $\G-\A$ & $-$0.0010 [$-$0.0593, $+$0.0574] & $+$0.0451 [$-$0.0263, $+$0.1165] & $-$0.0646 [$-$0.0922, $-$0.0369] & $-$0.0980 [$-$0.1289, $-$0.0671] & $+$0.0103 [$-$0.0181, $+$0.0388] \\
 & $\A-\Sl$ & $+$0.145 [$+$0.014, $+$0.276] & $+$0.186 [$+$0.071, $+$0.302] & $-$0.159 [$-$0.207, $-$0.111] & $-$0.0022 [$-$0.0415, $+$0.0370] & $-$0.249 [$-$0.313, $-$0.185] \\
\addlinespace
searches per hour & $\G-\A$ & $-$0.0748 [$-$0.1672, $+$0.0175] & $+$0.0198 [$-$0.0991, $+$0.1386] & $-$0.106 [$-$0.173, $-$0.039] & $-$0.101 [$-$0.176, $-$0.025] & $-$0.0988 [$-$0.1416, $-$0.0561] \\
 & $\A-\Sl$ & $-$0.361 [$-$0.448, $-$0.274] & $-$0.483 [$-$0.613, $-$0.353] & $-$0.163 [$-$0.238, $-$0.089] & $-$0.315 [$-$0.401, $-$0.229] & $-$0.172 [$-$0.225, $-$0.119] \\
\addlinespace
mean exit duration (s) & $\G-\A$ & $+$2.06 [$-$0.51, $+$4.63] & $+$2.58 [$-$0.18, $+$5.35] & $-$1.79 [$-$3.82, $+$0.25] & $-$1.76 [$-$4.26, $+$0.74] & $-$0.547 [$-$1.721, $+$0.626] \\
 & $\A-\Sl$ & $-$1.96 [$-$4.21, $+$0.29] & $-$2.94 [$-$5.48, $-$0.40] & $+$1.91 [$+$0.20, $+$3.61] & $-$0.180 [$-$2.300, $+$1.940] & $+$0.756 [$-$0.325, $+$1.836] \\
\bottomrule
\end{tabular}
}
  \caption{The non-redundant phase statistics (\texttt{redundancy.py}): $\G-\A$ and $\A-\Sl$ per ground class, 90\%
  bands. Our labels; the full table in Appendix~\ref{app:phases-table}.}
  \label{tab:phases-main}
  \end{table}
  \begin{figure}[H]
  \centering
  \includegraphics[width=\textwidth]{open_declared_dt5/phases_performance.pdf}
  \caption{Performance within the phases: $\G-\A$ and $\A-\Sl$ per ground class, 90\% bands.}
  \label{fig:phases-performance}
  \end{figure}
  \begin{figure}[H]
  \centering
  \includegraphics[width=\textwidth]{open_declared_dt5/phases_tails.pdf}
  \caption{Survival of the durations of glides, searches that end in a climb and climbs, and $\mu$ above the
  common threshold of each phase (Sec.~\ref{tails}), 90\% bands with the threshold re-chosen in every draw.}
  \label{fig:phases-tails}
  \end{figure}
\Obs\mbox{}
  \begin{itemize}
  \item The hypothesis of the composition holds: group time is richer in climbs and climbs more often, in every
    ground class, consistent with $\G-\A<0$ at the short lags, where circling dominates.
  \item Yet group time is more persistent at long lags. More climbs alone would lower $\alpha^*$ there: the gain
    comes from the other phases.
  \item In a group the pilot glides faster and more efficiently, and searches succeed more often; the climb is not
    stronger. The height saved per km is larger in the mountains (Table~\ref{tab:phases-main}).
  \item In the mountains the durations of the long glides decay more slowly in a group, on tens to hundreds of
    glides above the threshold.
  \item $\A-\Sl$: the solo time of group flights climbs stronger than solo flights in the hills and the low
    mountains (Table~\ref{tab:phases-main}).
  \end{itemize}
\Ded\mbox{}
  \begin{itemize}
  \item \textbf{The group effect lives in the search and the glide, not in the climb}: the group climbs about as
    fast as a pilot alone, but finds the climbs more reliably, spends its cycles faster and glides better between
    them. A more regular progress would be a more persistent one, although the group circles more: how much of the
    effect is timing, Sec.~\ref{ch:geometry} measures.
  \item $\A-\Sl$ is a stronger climb even alone, on pilots who also fly solo flights (Sec.~\ref{ch:who}): it points
    to the day and the site more than to the pilot.
  \end{itemize}
\Com The link from the phases to $\alpha^*$ is qualitative here: no formula maps the phase statistics on the
  exponent. Sec.~\ref{ch:geometry} separates shape and timing on the route.
\Home The group finds the climbs more reliably and glides better between them; it does not climb faster.
\end{step}

\begin{step}
\Q \textbf{The wind test}: why does group time glide faster and lose less height per km?
\Why The step above found both, against the polar of the wing, which makes a faster glide lose more height in the
  same air. A tailwind would give both without any group effect.
\Hyp \textbf{Why the group glides faster.} The speed-to-fly of MacCready grows with the climb expected next: a group
  that makes the next climb more certain lets the pilot fly faster, as birds in a flock. But in the same air a
  faster glide loses more height per km (the polar of the wing), and group time is faster \emph{and} loses less.
  So the air of the glide is better too: lines of lift shared by the group, the pilot ahead showing where the air
  rises, or a tailwind.
\Tool The drift speed of the climbs (chord over duration, a proxy of the wind at the height of the thermals) and
  the height gained per climb, for $\G-\A$ and $\A-\Sl$ per class (Appendix~\ref{app:phases-table}).
\Res Drift speed, $\G-\A$: \PhDriftGAAll~m/s on every flight, \PhDriftGALowMtn{} and \PhDriftGAHighMtn~m/s in the low
  and high mountains, where the height saved per km is largest. Height gained per climb, $\G-\A$: \PhGainGAAll~m on every
  flight. $\A-\Sl$: \PhGainASAll~m per climb. Horizontal distance per climb, $\G-\A$: \PhChordGAAll~m
  \PhChordGAAllBand.
\Obs Group time drifts no faster than the solo time of the same flights, and slower in the mountains.
\Ded \textbf{The wind at the height of the thermals does not explain the faster and flatter glide of group
  time}. The drift is a speed, not the component along the glide: a tailwind along the glide line is not
  excluded. What is left is the speed-to-fly of a more certain next climb together with better air in glide. On
  every flight the group does not climb higher either; the solo time of group flights does, a sign of better days
  (Sec.~\ref{ch:who}).
\Home Not the wind of the thermals: most likely a more certain next climb and better air in glide.
\end{step}

\begin{step}
\Q How does the group use the thermals?
\Why The group climbs more often, at about the same rate, and glides better: what may change is how long it stays
  in a thermal, and so how high it leaves.
\Tool The height gained per climb, next to the climbs per hour, the climb rate and the height lost per km, $\G-\A$
  per class; the duration of the exit from a thermal.
\Res $\G-\A$ per class (Appendix~\ref{app:phases-table}):
  \begin{center}\footnotesize
  \begin{tabular}{@{}l c c c c@{}}
  \toprule
   & climbs per hour & height per climb (m) & climb rate (m/s) & height lost (m/km) \\
  \midrule
  plains & \PhCphGAPlains & \PhGainGAPlains & \PhRateGAPlains & \PhSinkGAPlains \\
  hills & \PhCphGAHills & \PhGainGAHills & \PhRateGAHills & \PhSinkGAHills \\
  low mountains & \PhCphGALowMtn & \PhGainGALowMtn & \PhRateGALowMtn & \PhSinkGALowMtn \\
  high mountains & \PhCphGAHighMtn & \PhGainGAHighMtn & \PhRateGAHighMtn & \PhSinkGAHighMtn \\
  \bottomrule
  \end{tabular}
  \end{center}
  A climb of the solo time of group flights gains \PhGainAAll~m.
\Obs\mbox{}
  \begin{itemize}
  \item In the plains group time climbs more often and \emph{less high} per climb; in the high mountains more often
    and \emph{higher}; in the hills and the low mountains the height per climb does not change within the bands.
  \item The climb rate moves by a few hundredths of m/s at most: lower in the plains, higher in the high mountains.
    The glide loses less height everywhere.
  \item In the high mountains the extra height belongs to group time: $\A-\Sl$ is compatible with zero there, while
    $\G-\A$ and $\G-\Sl$ are not.
  \item The exit from a thermal lasts a little longer in a group (median $\G-\A=\PhExitDurGAAll$~s
    \PhExitDurGAAllBand); its mean does not move (Appendix~\ref{app:phases-table}). The group does not leave
    faster.
  \end{itemize}
\Ded Two readings, both of a more certain next climb, still hypotheses:
  \begin{enumerate}
  \item \emph{Plains}: the group leaves the thermal earlier, when staying is no longer worth it, since the next
    climb is more certain: the race strategy of MacCready, taking only the good part of a thermal.
  \item \emph{High mountains}: height is margin for the crossings and the ridges, and the group often leaves
    together, so the first to arrive keep climbing while they wait.
  \end{enumerate}
\Com Small effects: \PhGainGAPlains{} and \PhGainGAHighMtn~m on a climb of \PhGainAAll~m, a few percent; the signs
  are clear, the readings are hypotheses. The test: when the pilot leaves a thermal against the time its climb rate falls, and in the
  mountains whether those who wait are the first to arrive.
\Home \textbf{The group does not find stronger thermals; it changes how it uses them.} In the plains it climbs less
  high per climb, in the high mountains higher, and everywhere it then glides better; the strategies behind are
  hypotheses.
\end{step}

\section{How: the geometry}\label{ch:geometry}

\begin{step}
\Q Is the group effect in the shape of the route or in the way the route is flown?
\Why Chapter~\ref{ch:dynamics} found the gains in the search and the glide, but no formula maps the phases on
  $\alpha^*$. The route measures the split directly.
\Hyp $\alpha^*$ is a statistic in time, and the MSD at $\tau$ mixes two things:
  \begin{enumerate}
  \item \textbf{Geometry}: how much the route turns at the scale covered in $\tau$ (detours, zig-zags between
    thermals, turning back). A group that marks the thermals ahead lets its pilots deviate less.
  \item \textbf{Dynamics}: how the route is flown in time. In a climb the pilot circles and the displacement
    waits while the time runs; the alternation of climbs and glides is an intermittency that lowers $\alpha$
    at the lags of a climb, on a straight route too. A group that finds the climbs sooner (shorter searches,
    fewer weak climbs) progresses more regularly on the same route.
  \end{enumerate}
  In the plains the thermals are spread more evenly and the ground constrains the route little: we expect
  mostly (2).
\Tool\mbox{}
  \begin{itemize}
  \item \textbf{The route}: each climb of the phases of Vilpellet replaced by its chord, the path in arc
    length, so that the waiting is gone and only the shape is left (as in the ridges work). Read in arc length,
    the route is flown as at constant speed: the clock is gone. $\alpha$ of $\Sl$,
    $\A$, $\G$ on the route, in each ground class. If $\G-\A$ stays on the route, the group flies straighter
    paths (geometry); if it vanishes, it comes from the way the route is flown (intermittency). The gap in time
    minus the gap on the route is the share of the dynamics.
  \item \textbf{The check}: the exponent in arc length with the circles kept must equal $\alpha^*$ in time, since a
    paraglider waits while circling.
  \item Arc grid of \RouteDs~m; the lags in arc length are placed on the scale of time through the median speed of
    $\Sl$ along its route, \RouteSpeed~m/s, the same for the three populations: an approximate mapping. Same chain, same
    window, same shared draws as in time.
  \end{itemize}
\Res The check holds: with the circles kept, the exponent in arc length stays within \RouteCheckMax{} of
  $\alpha^*$ in time over the solid range (Fig.~\ref{fig:route-check}): the changes of speed along glides and
  searches matter little, and the timing, time minus route, is the waiting in the circles. On the route (Fig.~\ref{fig:route-gaps}),
  $\G-\A$ is \RouteGAAllThreeTwenty{} at the scale of 320~s and \RouteGAAllThousand{} at that of 1\,000~s, every
  flight; time minus route, \RouteMinusGAAllThreeTwenty{} and \RouteMinusGAAllThousand. Per class at 1\,000~s:
  \RouteGAPlainsThousand{} (plains), \RouteGAHillsThousand{} (hills), \RouteGALowMtnThousand{} (low mountains),
  \RouteGAHighMtnThousand{} (high mountains); time minus route in the plains \RouteMinusGAPlainsThousand.
  \begin{figure}[H]
  \centering
  \includegraphics[width=\textwidth]{open_declared_dt5/route_gaps.pdf}
  \caption{$\G-\A$ and $\A-\Sl$ on the route (climbs replaced by their chords, arc length) and in time, per ground
  class, on the scale of time through the median route speed of $\Sl$; 90\% bands from the shared draws.}
  \label{fig:route-gaps}
  \end{figure}
  \begin{figure}[H]
  \centering
  \includegraphics[width=0.75\textwidth]{open_declared_dt5/route_check.pdf}
  \caption{Check: the exponent in arc length with the circles kept against $\alpha^*$ in time, per population.}
  \label{fig:route-check}
  \end{figure}
\Obs\mbox{}
  \begin{itemize}
  \item The group gap survives on the route almost whole: in time it is larger by a few thousandths at 1\,000~s
    and by nothing at 320~s.
  \item In the plains the route keeps about half of the small gap in time.
  \item $\A-\Sl$ on the route follows $\A-\Sl$ in time in each class; on every flight it stays near zero on the
    route, while in time it turns negative near \ChainASNegative~s.
  \end{itemize}
\Ded\mbox{}
  \begin{itemize}
  \item \textbf{The group effect is mostly in the shape of the path}: once the waiting in the climbs is gone, group
    time still follows a straighter route than the solo time of the same flights. The share of the timing is
    small; relative to the gap it is largest in the plains, where we expected it.
  \item This agrees with the phases: the route keeps the searches and the glides, where the group gains
    (Sec.~\ref{ch:dynamics}); fewer failed searches and straighter glide lines make a straighter route.
  \end{itemize}
\Com The route keeps the searches as flown: separating the searches from the glides on the route is the next
  step of the geometry, as in the ridges work.
\Home Group time follows a straighter route than the solo time of the same flights; the timing adds little,
  relatively more in the plains.
\end{step}

\section{Which way round?}\label{which-way}

\begin{step}
\Q Does the group make the stretch straighter, or does a straight stretch make the group?
\Why Every chapter above found an association: group time is more persistent than the solo time of the same flights,
  rises at the entry and falls at the exit, and follows a straighter route. Both directions predict all of it.
\Hyp Two options:
  \begin{enumerate}
  \item \textbf{group $\to$ straight}: the group finds, selects or holds straighter paths (shared search of
    the climbs, following);
  \item \textbf{straight $\to$ group}: on a straight leg the paths of pilots converge on the same corridor,
    so they meet, and they stay together while the leg lasts.
  \end{enumerate}
\Tool The long-lag timeline of Sec.~\ref{ch:when}. If a straight leg makes the group (2), the pilots converge on it,
  so the window before the group already lies on the straight leg and is as persistent as the start; and the group
  splits when the path bends, so the window after is less persistent. If the group makes the path straight (1), the
  persistence rises only at the entry and falls at the exit.
\Res Before $\to$ start: \WhenAllLongStartMinusPre{} \WhenAllLongStartMinusPreBand; end $\to$ after:
  \WhenAllLongPostMinusEnd{} \WhenAllLongPostMinusEndBand; the level after is about the level before. Per class the
  rise at the entry is resolved only in the high mountains (Table~\ref{tab:prepost-long}). Both steps survive the
  mock events in solo flights (Table~\ref{tab:prepost-mock}).
\Obs The window before is less persistent than the start; the window after falls back.
\Ded\mbox{}
  \begin{itemize}
  \item \textbf{The entry leans to (1)}: on every flight and in the high mountains, a straight leg of the length of a
    window (\WhenWindow~s) does not precede the group. A shorter straight leg, just before the meeting, would not
    show.
  \item The exit does not separate them: both readings predict the fall after the group.
  \end{itemize}
\Com Tests still to do:
  \begin{itemize}
  \item \emph{Timing at a shorter scale}: does the straightness rise before the entry into the group (2) or only
    after it (1)? Does the group split when the path bends (2)?
  \item \emph{Same place, same time}: solo pilots on the same leg in the same hour. As straight as $\G$: the
    straightness belongs to the leg (2, or the conditions). Less straight: the group adds it (1).
  \item \emph{How the group formed}: groups that leave the same takeoff together did not form on a straight
    leg. If their group time is as much straighter, (1).
  \end{itemize}
\Home The data lean to group $\to$ persistent at the entry, on every flight and in the high mountains; the direction
  is not settled.
\end{step}

\section{What we found}\label{ch:found}

\begin{step}
\Home \textbf{What the data say.}
  \begin{enumerate}
  \item Group time is more persistent than solo flights from a few hundred seconds on: $\G-\Sl\approx0$ up to
    100~s, \ChainGSThousand{} at 1\,000~s (Sec.~\ref{ch:straighter}).
  \item At the long lags the gap is the group effect proper, $\G-\A$, within the same flights; $\A-\Sl$, between
    flights, is small and changes sign with the lag (Sec.~\ref{ch:identity}).
  \item The ground: in absolute terms the effect grows with the relief; normalised, it has about the same size on
    every ground and shows at a shorter scale in the mountains. The negative $\A-\Sl$ on every flight is a Simpson's
    paradox of the ground (Sec.~\ref{ground}).
  \item The opening selects which flights are compared and acts on $\A-\Sl$; within a flight the group effect
    survives every threshold. Group flights look less open mostly because they fly over the mountains
    (Sec.~\ref{ch:flights}).
  \item $\A-\Sl$ does not come from which pilots fly, to first order (Sec.~\ref{ch:who}), nor from a carry-over of
    the group on the solo time around it (Sec.~\ref{carry}).
  \item On the long lags, on every flight, the persistence rises when the pilot joins the group and falls when the
    group ends, against the composition of the phases and beyond what the moment of the flight does in solo
    flights. Per class the fall is resolved in the hills and the mountains, the rise only in the high mountains,
    neither in the plains (Sec.~\ref{ch:when}).
  \item In a group the pilot climbs more often but not faster; the searches end in a climb more often; the glides
    are faster and lose less height, and the wind does not explain it. It changes how the thermals are used: less
    height per climb in the plains, more in the high mountains (Sec.~\ref{ch:dynamics}).
  \item The effect is mostly in the shape of the route: it survives once the waiting in the climbs is removed; the
    timing adds little, relatively more in the plains (Sec.~\ref{ch:geometry}).
  \item The entry into the group leans to group $\to$ persistent; the direction is not settled (Sec.~\ref{which-way}).
  \end{enumerate}
  In one line: \textbf{the group time of a paraglider is straighter than its solo time, on every ground, at a scale
  set by the ground; in a group the pilot finds the climbs more reliably and glides better between them. Which way
  round is not settled.}
\Com \textbf{What the data do not say.}
  \begin{itemize}
  \item The direction at a finer scale than a window of \WhenWindow~s.
  \item How the group shows the better lines: shared information on the climbs, following the pilot ahead, or the
    better air that a group finds; and why the faster glide does not cost height.
  \item Why the scale of the effect is shorter in the mountains: the glide and climb cycle does not explain it
    (Sec.~\ref{ground}).
  \item Where $\A-\Sl$ comes from, once the pilot and the carry-over are excluded: the day, the site.
  \item Beyond \ChainSolidEnd~s: the curves rest on the long flights only, and the level of the chained MSD is a
    construction.
  \end{itemize}
\Com \textbf{Tools still missing.}
  \begin{itemize}
  \item The timing of the entry at a shorter scale, and the groups that leave the same takeoff together (direction).
  \item Within a group, who leads and who follows, and whether the follower's route is straighter than the
    leader's (how the group shows the lines).
  \item The geometry of the route split between searches and glides, with the tools of the ridges work (box
    dimension, crests) (shape against timing beyond the route).
  \item The same site and day for $\A-\Sl$; the weather beyond the drift of the climbs.
  \end{itemize}
\end{step}

\appendix
\section{One weight per piece or per flight}\label{app:weights}

\begin{step}
\Q Do the flights split into many stretches drive the MSD?
\Why With one weight per piece, a flight with $n_f$ pieces weighs $n_f$. $\G$ and $\A$ have more than one
  piece per flight on average, $\Sl$ almost one: \WeightPiecesPerFlightS{} ($\Sl$), \WeightPiecesPerFlightA{}
  ($\A$), \WeightPiecesPerFlightG{} ($\G$).
\Tool The chained MSD with one weight per flight: in every step of the chain the pieces of one flight share a
  total weight 1. The gaps per ground class against those with one weight per piece, with the shared draws
  (Fig.~\ref{fig:weights}).
\Res Per flight minus per piece at 1\,000~s, every flight: $\G-\Sl$ \WeightMinusGSAllThousand, $\G-\A$
  \WeightMinusGAAllThousand, $\A-\Sl$ \WeightMinusASAllThousand; in the high mountains $\G-\A$
  \WeightMinusGAHighMtnThousand{} \WeightMinusGAHighMtnThousandBand{} and $\A-\Sl$ \WeightMinusASHighMtnThousand{}
  \WeightMinusASHighMtnThousandBand. The largest shift over the solid lags is \WeightMaxShift.
  \begin{figure}[H]
  \centering
  \includegraphics[width=\textwidth]{open_declared_dt5/weights.pdf}
  \caption{$\G-\Sl$, $\G-\A$, $\A-\Sl$ with one weight per piece (gray) and per flight (colour) per ground class, and
  their differences; 90\% bands from the shared draws.}
  \label{fig:weights}
  \end{figure}
\Obs The gaps do not move on every flight; in the high mountains, where $\A$ has the most pieces per flight,
  $\A-\Sl$ falls and $\G-\A$ rises by a few thousandths.
\Ded \textbf{No}: the conclusions do not depend on the weight.
\Home One weight per piece stays, the choice of the text.
\end{step}

\section{The per-lag cohort}\label{app:perlag}

Another way to one curve per population, kept for reference: each window of $\alpha$ on its own fixed set,
the pieces that reach its longest lag. It meets the four properties of Sec.~\ref{which-cohort} too, with
fewer pieces than the chained MSD at the lags of a window below $\tau^+$.

\begin{method}{the per-lag cohort, first the MSD}
  Each window of $\alpha(\tau)$, $\tau_k\in[\tau^-,\tau^+]$ with $\tau^\pm=\tau\cdot10^{\pm\HalfDecade}$
  (Definitions), gets its own fixed set: the pieces that reach its longest lag $\tau^+$, hence every lag of the
  window,
  \[ \mathrm{MSD}_\tau(\tau_k)=\frac{1}{|P(\tau^+)|}\sum_{p\in P(\tau^+)}\overline{\delta^2_p}(\tau_k),
     \qquad \tau_k\in[\tau^-,\tau^+]. \]
  The set depends on the window, through $\tau^+$, and not on the lag (Fig.~\ref{fig:perlag-scheme}, top right).
  In the pool it is the other way round: it depends on the lag and not on the window (top left).

  \textbf{One MSD curve per window.} A lag $\tau_k$ falls in several windows. In the pool its MSD is the same in
  all of them. In the per-lag cohort it changes with the window, since $\tau^+$ does. The lag \PerLagExample~s, for
  instance:
  \begin{center}\small
  \begin{tabular}{rrr}
\toprule
window of $\tau$ & longest grid lag $\le\tau^+$ & pieces of $\mathrm{MSD}_\tau(2\,250~\mathrm{s})$ \\
\midrule
1\,840~s & 2\,760~s & $T_p\ge3\,450$~s \\
2\,250~s & 3\,380~s & $T_p\ge4\,225$~s \\
2\,760~s & 4\,150~s & $T_p\ge5\,188$~s \\
\midrule
pool, any window & & $T_p\ge2\,812$~s \\
\bottomrule
\end{tabular}

  \end{center}
  Three windows, three sets, three values of the MSD at the same lag. So the per-lag cohort gives one short MSD
  curve per window, from $\tau^-$ to $\tau^+$, which meets the pool at $\tau^+$, where the two sets coincide.

  \textbf{Below 1\,000~s they coincide.} Every piece of the cache lasts at least \PieceMinS~s
  ($=1\,000/\Support$), so $P(\tau)$ holds all of them up to 1\,000~s: pool and per-lag cohort coincide while
  the window ends within 1\,000~s ($\tau\le\PerLagCoincideEnd$~s).
\end{method}

\begin{step}
\Res \textbf{The MSD of the two choices} (Fig.~\ref{fig:perlag-msd}, every declared open flight), log--log: the
  slope of a curve is $2\alpha$.
  \begin{itemize}
  \item Each per-lag window meets the pool at its last lag $\tau^+$, where the two sets coincide. To its left
    the distance between the two curves is $r_k$.
  \item In $\A$ and $\G$ the windows lie on or above the pool to their left: the pieces that reach $\tau^+$ have a
    higher MSD than the pool, which also holds the shorter ones ($b_k>a_k$). So the pool is steeper than the
    windows it meets. In $\Sl$ the same holds from a few thousand seconds; before, the windows lie slightly
    below.
  \item At the longest lags the pool bends upward, steeper than the windows it meets, and ends above the
    ballistic slope $2$.
  \end{itemize}
  \begin{figure}[H]
  \centering
  \includegraphics[width=\textwidth]{open_declared_dt5/perlag_msd.pdf}
  \caption{$\mathrm{MSD}(\tau)$ of the pool (gray, each lag on its own pieces) and of the per-lag cohort
  (color, one window every third lag from 1\,000~s, each on the pieces that reach its last lag). Every declared
  open flight, one weight per piece.}
  \label{fig:perlag-msd}
  \end{figure}
\end{step}

\begin{method}{the per-lag cohort, then $\alpha$}
  $\alpha(\tau)$ is half the slope of the line fitted to $\log\mathrm{MSD}_\tau(\tau_k)$ against $\log\tau_k$ over
  the window (Definitions). One window, one short MSD curve, one $\alpha$: the exponents make one curve over
  every lag.
  \begin{itemize}
  \item \emph{It follows the cohorts.} $P(\tau^+)$ is the fixed cohort to $\tau^+$, so at each $\tau$ the per-lag
    $\alpha$ is the $\alpha$ of that cohort: of the fixed cohorts that fit the window, the one with the most
    pieces.
  \item \emph{It has no composition term}: each fit runs on one set. The pool and the per-lag cohort share the
    fit and meet at $\tau^+$, so the formula of the problem box holds exactly at every $\tau$, with
    $\alpha(\tau)=\alpha_\mathrm{per}(\tau)$:
    \[ \hat\alpha_\mathrm{pool}(\tau)-\alpha_\mathrm{per}(\tau)=\frac12\,\frac{\sum_k(u_k-\bar u)(r_k-\bar r)}{\sum_k(u_k-\bar u)^2}. \]
    The difference of the two curves is the composition term alone.
  \item \emph{The four properties}: every lag that a cohort reaches (1), the largest set in each window (2), one
    fixed set per fit (3), one curve (4).
  \item \emph{What it does not remove}: the selection. At long $\tau$ only the long flights reach $\tau^+$
    (Comment below).
  \end{itemize}
\end{method}

\begin{method}{the per-lag cohort on the toy}
  The toy of Sec.~\ref{which-cohort}, $\overline{\delta^2_p}(\tau)=c_p\,\tau^{2\alpha}$. Over one window,
  \[ \mathrm{MSD}_\mathrm{per}(\tau_k)=C\big(P(\tau^+)\big)\,\tau_k^{2\alpha},\qquad
     \frac{\mathrm{MSD}_\mathrm{pool}(\tau_k)}{\mathrm{MSD}_\mathrm{per}(\tau_k)}
     =\frac{C\big(P(\tau_k)\big)}{C\big(P(\tau^+)\big)}=e^{r_k}\le1, \]
  with equality at $\tau_k=\tau^+$: to the left of $\tau^+$ the pool lies below the per-lag curve, since it
  still holds the shorter pieces. $C(P(\tau^+))$ is the same at every lag of the window, so the per-lag points
  lie on one line of slope $2\alpha$ and the fit gives $\hat\alpha=\alpha$ in every window (\ToyPlPerlag{} in the
  figure), for any durations and prefactors.
  \begin{center}
  \includegraphics[width=0.9\textwidth]{perlag_scheme.pdf}
  \captionof{figure}{The toy with the per-lag cohort. Top: which pieces (rows, short to long) enter the MSD of
  each lag of the window of \ToyPlTau~s. Bottom left: the MSDs of that window and their fits; the per-lag points
  lie on a line of slope $2\alpha$, the pool points below it and on a steeper line, the mean of the cohorts
  (purple) on a steeper line too. Bottom right: $\alpha(\tau)$ from the fit over every window; line, the local
  exponent of the pool $\alpha+\frac{\beta}{2}g(y)$; dotted, the exponent of every piece.}
  \label{fig:perlag-scheme}
  \end{center}
\end{method}

\begin{step}
\Res \textbf{The exponent} (Fig.~\ref{fig:perlag-pool}): the half-slopes of the curves of
  Fig.~\ref{fig:perlag-msd}, with the same fit for both.
  \begin{itemize}
  \item \textbf{The pool is not an option}: pool minus per-lag is above 0, band included, from \PerLagPoolRunsAwayA~s
    in $\A$, \PerLagPoolRunsAwayG~s in $\G$ and \PerLagPoolRunsAwayS~s in $\Sl$, and its $\alpha$ runs away, up to \PoolSMax{} ($\Sl$),
    \PoolAMax{} ($\A$), \PoolGMax{} ($\G$), above the ballistic $\alpha=1$, which no cohort reaches.
    Pool minus per-lag at 6\,220~s: \PoolSMinusVeryLong{} ($\Sl$), \PoolAMinusVeryLong{} ($\A$),
    \PoolGMinusVeryLong{} ($\G$). It is the composition term alone.
  \item The per-lag curve stays within \PerLagSMinusCohortMax{} ($\Sl$), \PerLagAMinusCohortMax{} ($\A$),
    \PerLagGMinusCohortMax{} ($\G$) of the cohort to 1~h, at every lag up to \PerLagSolidEnd~s, with bands from
    the paired draws.
  \end{itemize}
  \begin{figure}[H]
  \centering
  \includegraphics[width=\textwidth]{open_declared_dt5/perlag_pool.pdf}
  \caption{$\alpha(\tau)$ of the per-lag cohort (color), of the pool (gray) and of the cohort to 1~h
  (dotted), every declared open flight. 90\% bands.}
  \label{fig:perlag-pool}
  \end{figure}
\Com \textbf{What it does not solve: the selection.} At $\tau$ the curve describes the pieces of at least
  $\tau^+/\Support$~s: every piece at 10~s (\PerLagGFirstPieces{} in $\G$), only the long flights at the end
  (\PerLagGLastPieces{} pieces at \PerLagGLastLag~s in $\G$). Since the longest flights are straighter at every
  scale (Sec.~\ref{fixed-cohort}), a rise of $\alpha$ at long lags may come from the flights that change, not
  from the dynamics. No method removes it: the lag $\tau$ needs pieces of at least
  $\tau^+/\Support=\tau\cdot10^{\HalfDecade}/\Support\approx\PlusFactor\,\tau$, so at \PlusLagHours~h only the pieces of about
  \PlusPieceHours~h count.
\Ded \textbf{Solid and dashed.} Solid while the window ends within \PerLagSelFree~s ($\tau\le\PerLagSolidEnd$~s):
  there the pieces are at least those of the cohort to 1~h, and the selection is measured small (at most
  \PerLagAMinusCohortMax{} in $\alpha$, against $\G-\A=\PerLagGAThousand$ at 1\,000~s). Dashed beyond: the long
  flights only.
\Res The per-lag curves and their gaps (Fig.~\ref{fig:perlag}). $\G-\A$ is \PerLagGAThousand{} at 1\,000~s,
  against \GapAllGAThousand{} in the cohort to 1~h. $\A-\Sl$ at 10~s is \PerLagASTen, against \GapAllASTen{} in the
  cohort to 1~h: the pieces shorter than 1~h count for $\A-\Sl$.
  \begin{figure}[H]
  \centering
  \includegraphics[width=\textwidth]{open_declared_dt5/perlag.pdf}
  \caption{Left: $\alpha(\tau)$ of $\Sl$, $\A$, $\G$, each lag on its per-lag cohort; solid up to
  \PerLagSolidEnd~s, dashed beyond (long flights only); dotted, the cohort to 1~h. Right: the gaps. Every
  declared open flight, 90\% bands from the shared two-way draws.}
  \label{fig:perlag}
  \end{figure}
\end{step}

\begin{step}
\Res \textbf{Against the chained MSD} (Fig.~\ref{fig:perlag-chained}). Up to about 1~h the two exponents
  agree, far within the gaps between the populations. Beyond, $\alpha^*$ lies below the per-lag $\alpha$, more
  so at longer lags: every step of the chain runs on the pieces that reach it, while the per-lag window keeps
  only those that reach $\tau^+$, longer and straighter (Sec.~\ref{fixed-cohort}). The chain also reaches the
  last lag of the pool.
  \begin{figure}[H]
  \centering
  \includegraphics[width=\textwidth]{open_declared_dt5/perlag_chained.pdf}
  \caption{Top: $\alpha^*$ (color), the per-lag $\alpha$ (dashed) and the pool (red). Bottom: $\alpha^*$ minus
  the per-lag $\alpha$, 90\% bands from the paired draws. Every declared open flight.}
  \label{fig:perlag-chained}
  \end{figure}
\end{step}

\section{Phase labels}\label{app:phases}

\begin{step}
\Q How are the phases labelled, and how much does our refit change Vilpellet's labels?
\Tool His three-state model (three yes/no features over 30~s, Viterbi decoding, 90-fix majority vote),
  refitted on 1-Hz flights by Baum--Welch with the transition matrix estimated, several starts, held-out
  likelihood, one model per discipline; tracks at 2--\DtMax~s splined to 1~s (cadence check); agreement with his
  labels per state and ground class; the change of the grouping on a sample of days.
\Res\mbox{}
  \begin{itemize}
  \item \emph{Training.} Paragliders: \RefitTrainPara{} flights logged at 1~Hz (and \RefitHeldPara{} held out),
    drawn in equal quotas over the ground classes from \RefitPoolPara; hang gliders: \RefitTrainHang{} (and
    \RefitHeldHang{} held out). \RefitStarts{} starts per discipline: his parameters, random ones and perturbed
    ones. Convergence in at most \RefitIterMax{} iterations.
  \item \emph{Which optimum.} The starts end in different local optima. The best held-out likelihood
    (\RefitLLBestParaFix{} per fix, against \RefitLLHisParaFix{} with his parameters) turns the search into a
    partly straight and rising state. \textbf{We keep the best optimum whose three states keep his meaning}
    (\RefitLLOursParaFix{} per fix): climb rises and turns one way (both in at least 80\% of its windows, straight
    in at most 20\%), transition is straight (at least 80\%), search is neither. Several starts reach it with the
    same parameters. The labels serve because they mean a phase of the flight; a better fit that redefines the
    search would not answer our questions.
  \item \emph{Hang gliders}: their own fit converges and keeps the meaning; its labels and those of the
    paraglider fit agree on \RefitHangVsPara\% of the fixes.
  \item \emph{Cadence}: thinned to 2~s and splined back, \RefitCadenceParaTwo\% of the paraglider labels are
    unchanged; at 5~s \RefitCadenceParaFive\% (hang gliders \RefitCadenceHangFive\%).
  \item \emph{Against his labels}: \RefitAgreePara\% of the paraglider fixes agree (hang gliders
    \RefitAgreeHang\%), from \RefitAgreePlains\% in the plains to \RefitAgreeHighMtn\% in the high mountains. Our
    model changes phase more often: \RefitSearchPerHourOurs{} searches per hour against \RefitSearchPerHourHis,
    \RefitSearchShareOurs\% of the time in search against \RefitSearchShareHis\%, at the expense of the glide. These
    count every run labelled search; the statistics of the text count fewer, since the context rule of
    Sec.~\ref{phases} moves the re-centrings into the climbs and the exits apart.
  \item \emph{The statistics of Sec.~\ref{ch:dynamics} with his labels} (on the 1-Hz flights his labels cover): the
    sign of the gaps $\G-\A$ and $\A-\Sl$ agrees in \PhSignAgree{} of \PhSignTotal{} cases.
  \item \emph{Grouping.} The phases enter it only through the heading condition, checked when both pilots glide
    and both flights are 1~Hz. On \GroupCheckDays{} busy days the contacts and the groups recomputed with our labels,
    everything else fixed, change the flights with group time by \GroupCheckFlightsChange\% and the total group time
    by \GroupCheckTimeChange\%; \GroupCheckShared\% of the group seconds are common to the two (over their union).
    The grouping of the run stands.
  \end{itemize}
\end{step}

\section{Phase statistics, full table}\label{app:phases-table}

\begin{table}[H]
\centering
\adjustbox{max width=\textwidth, max totalheight=0.85\textheight}{\begin{tabular}{@{}l r r r l l l@{}}
\toprule
statistic & $\Sl$ & $\A$ & $\G$ & $\G-\A$ & $\A-\Sl$ & $\G-\Sl$ \\
\midrule
\multicolumn{7}{@{}l}{\textit{plains}: 9\,634, 1\,525, 2\,097 flights} \\
time in glide (\%) & 48.8 & 48.0 & 44.4 & $-$3.62 [$-$4.41, $-$2.83] & $-$0.821 [$-$1.596, $-$0.046] & $-$4.44 [$-$5.25, $-$3.63] \\
time in search (\%) & 8.95 & 8.12 & 7.86 & $-$0.259 [$-$0.599, $+$0.081] & $-$0.835 [$-$1.231, $-$0.439] & $-$1.09 [$-$1.48, $-$0.71] \\
time in climb (\%) & 37.9 & 40.0 & 43.4 & $+$3.42 [$+$2.76, $+$4.07] & $+$2.07 [$+$1.38, $+$2.76] & $+$5.49 [$+$4.85, $+$6.13] \\
climbs per hour & 7.09 & 7.00 & 7.80 & $+$0.798 [$+$0.669, $+$0.927] & $-$0.0933 [$-$0.2128, $+$0.0262] & $+$0.704 [$+$0.591, $+$0.818] \\
glide speed (km/h) & 41.0 & 43.2 & 44.1 & $+$0.878 [$+$0.461, $+$1.295] & $+$2.21 [$+$1.41, $+$3.01] & $+$3.09 [$+$2.33, $+$3.84] \\
height lost per km of glide (m/km) & 98.6 & 90.0 & 83.4 & $-$6.60 [$-$8.00, $-$5.21] & $-$8.53 [$-$10.87, $-$6.20] & $-$15.1 [$-$17.4, $-$12.9] \\
searches ending in a climb (\%) & 67.7 & 70.1 & 73.8 & $+$3.69 [$+$2.12, $+$5.27] & $+$2.41 [$+$1.00, $+$3.82] & $+$6.10 [$+$5.05, $+$7.15] \\
climb rate (m/s) & 0.970 & 0.980 & 0.948 & $-$0.0317 [$-$0.0474, $-$0.0160] & $+$0.0098 [$-$0.0080, $+$0.0276] & $-$0.0219 [$-$0.0379, $-$0.0059] \\
distance gained per climb (m) & 879 & 971 & 947 & $-$24.7 [$-$46.2, $-$3.1] & $+$92.3 [$+$62.5, $+$122.1] & $+$67.6 [$+$42.4, $+$92.9] \\
height gained per climb (m) & 186 & 202 & 190 & $-$12.1 [$-$17.3, $-$6.9] & $+$15.7 [$+$10.7, $+$20.7] & $+$3.64 [$-$1.22, $+$8.50] \\
drift speed of a climb (m/s) & 4.58 & 4.72 & 4.73 & $+$0.0148 [$-$0.0423, $+$0.0719] & $+$0.134 [$+$0.007, $+$0.261] & $+$0.148 [$+$0.035, $+$0.262] \\
time in exit (\%) & 4.33 & 3.92 & 4.38 & $+$0.461 [$+$0.225, $+$0.697] & $-$0.413 [$-$0.622, $-$0.203] & $+$0.0482 [$-$0.1763, $+$0.2727] \\
searches per hour & 4.54 & 4.11 & 4.22 & $+$0.102 [$-$0.028, $+$0.232] & $-$0.423 [$-$0.558, $-$0.287] & $-$0.320 [$-$0.473, $-$0.168] \\
exits per hour & 2.20 & 2.05 & 2.23 & $+$0.174 [$+$0.094, $+$0.255] & $-$0.152 [$-$0.231, $-$0.073] & $+$0.0223 [$-$0.0599, $+$0.1044] \\
mean exit duration (s) & 70.7 & 68.7 & 70.8 & $+$2.07 [$-$0.50, $+$4.63] & $-$1.99 [$-$4.25, $+$0.26] & $+$0.0724 [$-$1.8207, $+$1.9655] \\
median exit duration (s) & 55.0 & 54.0 & 57.0 & $+$3.00 [$+$0.98, $+$5.02] & $-$1.00 [$-$2.57, $+$0.57] & $+$2.00 [$+$0.44, $+$3.56] \\
glide and climb cycle (min) & 8.46 & 8.57 & 7.69 & $-$0.877 [$-$1.023, $-$0.731] & $+$0.113 [$-$0.033, $+$0.258] & $-$0.764 [$-$0.880, $-$0.648] \\
\midrule
\multicolumn{7}{@{}l}{\textit{hills}: 9\,731, 1\,113, 1\,543 flights} \\
time in glide (\%) & 53.0 & 50.9 & 46.8 & $-$4.06 [$-$4.92, $-$3.21] & $-$2.11 [$-$2.97, $-$1.25] & $-$6.17 [$-$7.06, $-$5.29] \\
time in search (\%) & 8.60 & 7.21 & 7.33 & $+$0.114 [$-$0.285, $+$0.514] & $-$1.39 [$-$1.85, $-$0.93] & $-$1.28 [$-$1.73, $-$0.83] \\
time in climb (\%) & 35.0 & 38.8 & 42.2 & $+$3.44 [$+$2.69, $+$4.18] & $+$3.81 [$+$2.98, $+$4.64] & $+$7.25 [$+$6.47, $+$8.02] \\
climbs per hour & 6.83 & 7.10 & 7.56 & $+$0.460 [$+$0.301, $+$0.619] & $+$0.268 [$+$0.122, $+$0.414] & $+$0.728 [$+$0.589, $+$0.867] \\
glide speed (km/h) & 34.6 & 37.6 & 39.3 & $+$1.65 [$+$1.09, $+$2.21] & $+$3.02 [$+$2.19, $+$3.85] & $+$4.67 [$+$3.92, $+$5.41] \\
height lost per km of glide (m/km) & 108 & 102 & 96.3 & $-$5.89 [$-$7.78, $-$4.00] & $-$6.03 [$-$8.44, $-$3.61] & $-$11.9 [$-$14.0, $-$9.9] \\
searches ending in a climb (\%) & 63.3 & 68.7 & 72.5 & $+$3.81 [$+$2.20, $+$5.42] & $+$5.37 [$+$3.88, $+$6.86] & $+$9.18 [$+$7.77, $+$10.59] \\
climb rate (m/s) & 1.06 & 1.09 & 1.08 & $-$0.0049 [$-$0.0233, $+$0.0136] & $+$0.0262 [$+$0.0010, $+$0.0514] & $+$0.0213 [$-$0.0030, $+$0.0456] \\
distance gained per climb (m) & 644 & 723 & 753 & $+$29.8 [$+$7.2, $+$52.4] & $+$79.6 [$+$51.9, $+$107.3] & $+$109 [$+$84, $+$135] \\
height gained per climb (m) & 195 & 214 & 219 & $+$4.61 [$-$2.09, $+$11.32] & $+$19.3 [$+$12.4, $+$26.3] & $+$24.0 [$+$17.0, $+$30.9] \\
drift speed of a climb (m/s) & 3.51 & 3.68 & 3.73 & $+$0.0539 [$-$0.0164, $+$0.1243] & $+$0.166 [$+$0.054, $+$0.278] & $+$0.220 [$+$0.119, $+$0.321] \\
time in exit (\%) & 3.45 & 3.14 & 3.65 & $+$0.513 [$+$0.282, $+$0.743] & $-$0.311 [$-$0.508, $-$0.114] & $+$0.202 [$-$0.049, $+$0.452] \\
searches per hour & 4.27 & 3.85 & 3.99 & $+$0.140 [$-$0.035, $+$0.314] & $-$0.416 [$-$0.608, $-$0.224] & $-$0.277 [$-$0.446, $-$0.107] \\
exits per hour & 1.92 & 1.83 & 2.04 & $+$0.210 [$+$0.113, $+$0.307] & $-$0.0869 [$-$0.1734, $-$0.0004] & $+$0.123 [$+$0.032, $+$0.215] \\
mean exit duration (s) & 64.8 & 61.8 & 64.5 & $+$2.70 [$-$0.05, $+$5.45] & $-$3.04 [$-$5.58, $-$0.51] & $-$0.347 [$-$2.951, $+$2.257] \\
median exit duration (s) & 50.0 & 50.0 & 52.0 & $+$2.00 [$-$0.04, $+$4.04] & 0 [$-$2, $+$2] & $+$2.00 [$+$0.15, $+$3.85] \\
glide and climb cycle (min) & 8.78 & 8.45 & 7.94 & $-$0.514 [$-$0.694, $-$0.334] & $-$0.332 [$-$0.508, $-$0.156] & $-$0.846 [$-$0.999, $-$0.694] \\
\midrule
\multicolumn{7}{@{}l}{\textit{low mountains}: 18\,796, 5\,241, 6\,120 flights} \\
time in glide (\%) & 58.2 & 59.1 & 55.3 & $-$3.76 [$-$4.22, $-$3.30] & $+$0.863 [$+$0.396, $+$1.331] & $-$2.90 [$-$3.44, $-$2.35] \\
time in search (\%) & 10.7 & 10.7 & 10.3 & $-$0.409 [$-$0.700, $-$0.118] & $+$0.0332 [$-$0.2737, $+$0.3400] & $-$0.376 [$-$0.732, $-$0.020] \\
time in climb (\%) & 28.1 & 27.4 & 31.4 & $+$4.00 [$+$3.64, $+$4.36] & $-$0.684 [$-$1.111, $-$0.257] & $+$3.31 [$+$2.84, $+$3.79] \\
climbs per hour & 6.59 & 6.37 & 7.34 & $+$0.969 [$+$0.891, $+$1.047] & $-$0.222 [$-$0.293, $-$0.152] & $+$0.747 [$+$0.671, $+$0.823] \\
glide speed (km/h) & 30.2 & 31.3 & 32.3 & $+$0.985 [$+$0.779, $+$1.192] & $+$1.09 [$+$0.80, $+$1.39] & $+$2.08 [$+$1.77, $+$2.39] \\
height lost per km of glide (m/km) & 102 & 94.3 & 80.2 & $-$14.1 [$-$15.2, $-$13.0] & $-$7.40 [$-$8.63, $-$6.18] & $-$21.5 [$-$22.9, $-$20.0] \\
searches ending in a climb (\%) & 58.3 & 59.7 & 65.7 & $+$6.08 [$+$5.19, $+$6.96] & $+$1.36 [$+$0.54, $+$2.19] & $+$7.44 [$+$6.66, $+$8.22] \\
climb rate (m/s) & 1.18 & 1.23 & 1.23 & $+$0.0012 [$-$0.0095, $+$0.0120] & $+$0.0517 [$+$0.0364, $+$0.0669] & $+$0.0529 [$+$0.0379, $+$0.0679] \\
distance gained per climb (m) & 396 & 377 & 367 & $-$10.2 [$-$16.1, $-$4.4] & $-$18.5 [$-$29.0, $-$8.1] & $-$28.8 [$-$39.1, $-$18.4] \\
height gained per climb (m) & 180 & 189 & 189 & $-$0.480 [$-$3.537, $+$2.578] & $+$9.84 [$+$5.83, $+$13.85] & $+$9.36 [$+$5.34, $+$13.38] \\
drift speed of a climb (m/s) & 2.59 & 2.45 & 2.39 & $-$0.0580 [$-$0.0849, $-$0.0311] & $-$0.147 [$-$0.193, $-$0.101] & $-$0.205 [$-$0.252, $-$0.158] \\
time in exit (\%) & 3.00 & 2.79 & 2.96 & $+$0.173 [$+$0.060, $+$0.286] & $-$0.212 [$-$0.314, $-$0.110] & $-$0.0391 [$-$0.1504, $+$0.0721] \\
searches per hour & 4.65 & 4.41 & 4.52 & $+$0.112 [$+$0.028, $+$0.197] & $-$0.240 [$-$0.330, $-$0.149] & $-$0.127 [$-$0.235, $-$0.019] \\
exits per hour & 1.57 & 1.42 & 1.55 & $+$0.127 [$+$0.090, $+$0.165] & $-$0.151 [$-$0.185, $-$0.116] & $-$0.0233 [$-$0.0642, $+$0.0176] \\
mean exit duration (s) & 68.6 & 70.5 & 68.7 & $-$1.77 [$-$3.80, $+$0.27] & $+$1.89 [$+$0.18, $+$3.59] & $+$0.121 [$-$1.519, $+$1.762] \\
median exit duration (s) & 50.0 & 50.0 & 51.0 & $+$1.00 [$-$0.35, $+$2.35] & 0 [$-$1, $+$1] & $+$1.00 [$-$0.32, $+$2.32] \\
glide and climb cycle (min) & 9.11 & 9.42 & 8.18 & $-$1.25 [$-$1.35, $-$1.14] & $+$0.318 [$+$0.215, $+$0.421] & $-$0.927 [$-$1.017, $-$0.837] \\
\midrule
\multicolumn{7}{@{}l}{\textit{high mountains}: 7\,952, 2\,786, 3\,097 flights} \\
time in glide (\%) & 58.7 & 60.0 & 55.3 & $-$4.75 [$-$5.38, $-$4.11] & $+$1.39 [$+$0.87, $+$1.91] & $-$3.35 [$-$3.98, $-$2.73] \\
time in search (\%) & 10.6 & 9.83 & 9.37 & $-$0.460 [$-$0.781, $-$0.139] & $-$0.764 [$-$1.116, $-$0.412] & $-$1.22 [$-$1.60, $-$0.84] \\
time in climb (\%) & 27.9 & 27.5 & 32.5 & $+$5.05 [$+$4.56, $+$5.55] & $-$0.455 [$-$0.919, $+$0.008] & $+$4.60 [$+$4.08, $+$5.12] \\
climbs per hour & 6.53 & 6.40 & 7.40 & $+$1.00 [$+$0.89, $+$1.11] & $-$0.130 [$-$0.227, $-$0.034] & $+$0.871 [$+$0.771, $+$0.970] \\
glide speed (km/h) & 30.5 & 31.6 & 33.1 & $+$1.44 [$+$1.17, $+$1.71] & $+$1.14 [$+$0.82, $+$1.45] & $+$2.58 [$+$2.24, $+$2.91] \\
height lost per km of glide (m/km) & 104 & 99.5 & 86.8 & $-$12.7 [$-$14.1, $-$11.3] & $-$4.34 [$-$5.81, $-$2.86] & $-$17.0 [$-$18.6, $-$15.5] \\
searches ending in a climb (\%) & 60.8 & 61.8 & 68.7 & $+$6.83 [$+$5.75, $+$7.92] & $+$1.09 [$+$0.07, $+$2.11] & $+$7.92 [$+$6.97, $+$8.87] \\
climb rate (m/s) & 1.28 & 1.30 & 1.31 & $+$0.0161 [$+$0.0024, $+$0.0298] & $+$0.0205 [$+$0.0045, $+$0.0365] & $+$0.0366 [$+$0.0197, $+$0.0535] \\
distance gained per climb (m) & 346 & 347 & 338 & $-$8.86 [$-$14.79, $-$2.93] & $+$0.737 [$-$6.732, $+$8.206] & $-$8.12 [$-$15.44, $-$0.80] \\
height gained per climb (m) & 195 & 199 & 207 & $+$7.46 [$+$3.44, $+$11.48] & $+$4.16 [$-$0.13, $+$8.45] & $+$11.6 [$+$7.0, $+$16.3] \\
drift speed of a climb (m/s) & 2.26 & 2.26 & 2.15 & $-$0.111 [$-$0.141, $-$0.081] & $-$0.0069 [$-$0.0453, $+$0.0315] & $-$0.118 [$-$0.153, $-$0.082] \\
time in exit (\%) & 2.82 & 2.65 & 2.80 & $+$0.152 [$+$0.010, $+$0.295] & $-$0.174 [$-$0.294, $-$0.054] & $-$0.0218 [$-$0.1633, $+$0.1198] \\
searches per hour & 4.55 & 4.19 & 4.29 & $+$0.0994 [$-$0.0103, $+$0.2090] & $-$0.360 [$-$0.469, $-$0.250] & $-$0.260 [$-$0.382, $-$0.138] \\
exits per hour & 1.47 & 1.38 & 1.50 & $+$0.117 [$+$0.067, $+$0.167] & $-$0.0865 [$-$0.1301, $-$0.0429] & $+$0.0307 [$-$0.0194, $+$0.0809] \\
mean exit duration (s) & 69.1 & 68.9 & 67.1 & $-$1.72 [$-$4.25, $+$0.80] & $-$0.212 [$-$2.331, $+$1.907] & $-$1.94 [$-$4.13, $+$0.25] \\
median exit duration (s) & 50.0 & 49.0 & 50.0 & $+$1.00 [$-$0.60, $+$2.60] & $-$1.00 [$-$2.38, $+$0.38] & 0 [$-$1, $+$1] \\
glide and climb cycle (min) & 9.19 & 9.38 & 8.11 & $-$1.27 [$-$1.41, $-$1.12] & $+$0.187 [$+$0.047, $+$0.328] & $-$1.08 [$-$1.20, $-$0.97] \\
\midrule
\multicolumn{7}{@{}l}{\textit{all}: 49\,781, 10\,916, 13\,133 flights} \\
time in glide (\%) & 55.4 & 56.9 & 52.6 & $-$4.34 [$-$4.65, $-$4.03] & $+$1.50 [$+$1.14, $+$1.87] & $-$2.83 [$-$3.26, $-$2.41] \\
time in search (\%) & 9.69 & 9.68 & 9.28 & $-$0.399 [$-$0.573, $-$0.225] & $-$0.0151 [$-$0.2309, $+$0.2006] & $-$0.414 [$-$0.651, $-$0.176] \\
time in climb (\%) & 31.6 & 30.5 & 34.9 & $+$4.45 [$+$4.18, $+$4.72] & $-$1.14 [$-$1.50, $-$0.77] & $+$3.32 [$+$2.91, $+$3.72] \\
climbs per hour & 6.75 & 6.56 & 7.46 & $+$0.907 [$+$0.854, $+$0.959] & $-$0.199 [$-$0.249, $-$0.150] & $+$0.707 [$+$0.653, $+$0.761] \\
glide speed (km/h) & 32.9 & 33.4 & 34.9 & $+$1.45 [$+$1.28, $+$1.61] & $+$0.517 [$+$0.211, $+$0.823] & $+$1.96 [$+$1.62, $+$2.31] \\
height lost per km of glide (m/km) & 103 & 96.0 & 84.4 & $-$11.6 [$-$12.3, $-$10.9] & $-$7.08 [$-$8.02, $-$6.14] & $-$18.7 [$-$19.8, $-$17.5] \\
searches ending in a climb (\%) & 61.6 & 62.6 & 68.4 & $+$5.80 [$+$5.19, $+$6.40] & $+$0.991 [$+$0.447, $+$1.535] & $+$6.79 [$+$6.24, $+$7.34] \\
climb rate (m/s) & 1.11 & 1.18 & 1.17 & $-$0.0089 [$-$0.0164, $-$0.0014] & $+$0.0701 [$+$0.0588, $+$0.0814] & $+$0.0612 [$+$0.0492, $+$0.0732] \\
distance gained per climb (m) & 540 & 501 & 507 & $+$6.01 [$-$0.71, $+$12.73] & $-$39.6 [$-$53.2, $-$26.0] & $-$33.6 [$-$48.2, $-$19.1] \\
height gained per climb (m) & 186 & 197 & 197 & $+$0.136 [$-$2.009, $+$2.282] & $+$11.0 [$+$8.2, $+$13.8] & $+$11.2 [$+$8.2, $+$14.1] \\
drift speed of a climb (m/s) & 3.23 & 3.00 & 3.01 & $+$0.0109 [$-$0.0164, $+$0.0382] & $-$0.226 [$-$0.287, $-$0.164] & $-$0.215 [$-$0.280, $-$0.149] \\
time in exit (\%) & 3.29 & 2.94 & 3.22 & $+$0.286 [$+$0.207, $+$0.366] & $-$0.353 [$-$0.428, $-$0.278] & $-$0.0668 [$-$0.1525, $+$0.0189] \\
searches per hour & 4.48 & 4.23 & 4.34 & $+$0.107 [$+$0.052, $+$0.161] & $-$0.249 [$-$0.315, $-$0.182] & $-$0.142 [$-$0.218, $-$0.065] \\
exits per hour & 1.75 & 1.54 & 1.71 & $+$0.163 [$+$0.135, $+$0.192] & $-$0.204 [$-$0.232, $-$0.177] & $-$0.0407 [$-$0.0743, $-$0.0072] \\
mean exit duration (s) & 67.7 & 68.5 & 67.9 & $-$0.522 [$-$1.696, $+$0.653] & $+$0.729 [$-$0.355, $+$1.813] & $+$0.208 [$-$0.784, $+$1.199] \\
median exit duration (s) & 50.0 & 50.0 & 52.0 & $+$2.00 [$+$1.16, $+$2.84] & 0 [$-$1, $+$1] & $+$2.00 [$+$1.26, $+$2.74] \\
glide and climb cycle (min) & 8.88 & 9.15 & 8.04 & $-$1.11 [$-$1.18, $-$1.05] & $+$0.270 [$+$0.202, $+$0.338] & $-$0.842 [$-$0.904, $-$0.780] \\
\bottomrule
\end{tabular}
}
\caption{Every phase statistic: $\G-\A$ and $\A-\Sl$ per ground class, 90\% bands. Our labels.}
\end{table}

\section{Thresholds of the tails}\label{app:tails}

\begin{step}
\Q Do the tail exponents depend on the threshold?
\Tool The Kolmogorov--Smirnov threshold of each curve and the Hill plots $\hat\mu(x_\mathrm{min})$ of $\Sl$,
  $\A$, $\G$ per phase. The scan of the threshold runs on the durations binned at 1~s below \TailFineBelow~s and \TailBinPercent\% above,
  over \TailCandidates{} candidates between the percentiles \TailQuantileLo{} and \TailQuantileHi{} of the durations, with at least
  \TailMinRuns{} durations above.
\Res\mbox{}
  \begin{table}[H]
  \centering
  \resizebox{\textwidth}{!}{\begin{tabular}{@{}l l r l l l l l@{}}
\toprule
ground & phase & $x_\mathrm{min}$ (s) & $\Sl$ & $\A$ & $\G$ & $\G-\A$ & $\A-\Sl$ \\
\midrule
plains & glide & 1\,130 & 5.89 [5.05, 6.73] & 5.26 [4.33, 6.20] & 5.39 [4.09, 6.69] & $+$0.127 [$-$1.232, $+$1.485] & $-$0.632 [$-$1.479, $+$0.215] \\
 & search & 232 & 3.42 [2.48, 4.36] & 3.40 [1.73, 5.07] & 3.81 [2.45, 5.17] & $+$0.407 [$-$0.636, $+$1.450] & $-$0.0191 [$-$1.0263, $+$0.9882] \\
 & climb & 1\,118 & 10.5 [7.7, 13.4] & 9.77 [7.41, 12.14] & 8.00 [6.27, 9.73] & $-$1.78 [$-$3.77, $+$0.22] & $-$0.768 [$-$2.620, $+$1.084] \\
\addlinespace
hills & glide & 1\,075 & 4.37 [4.06, 4.68] & 5.12 [4.19, 6.05] & 5.36 [4.28, 6.44] & $+$0.240 [$-$0.773, $+$1.253] & $+$0.749 [$-$0.041, $+$1.538] \\
 & search & 382 & 3.75 [3.07, 4.43] & 4.31 [2.62, 6.01] & 4.42 [2.95, 5.89] & $+$0.106 [$-$1.575, $+$1.788] & $+$0.562 [$-$0.767, $+$1.890] \\
 & climb & 1\,043 & 9.54 [7.00, 12.08] & 8.38 [5.88, 10.88] & 16.3 [9.0, 23.6] & $+$7.91 [$+$1.43, $+$14.39] & $-$1.15 [$-$3.40, $+$1.09] \\
\addlinespace
low mountains & glide & 1\,665 & 4.63 [4.31, 4.95] & 4.34 [4.02, 4.67] & 3.48 [3.14, 3.83] & $-$0.861 [$-$1.304, $-$0.419] & $-$0.286 [$-$0.683, $+$0.111] \\
 & search & 654 & 3.92 [3.57, 4.28] & 4.35 [3.78, 4.92] & 4.33 [3.75, 4.91] & $-$0.0209 [$-$0.7458, $+$0.7040] & $+$0.427 [$-$0.092, $+$0.947] \\
 & climb & 782 & 7.13 [5.22, 9.04] & 6.94 [4.96, 8.92] & 8.27 [5.40, 11.14] & $+$1.33 [$-$0.40, $+$3.06] & $-$0.194 [$-$1.098, $+$0.711] \\
\addlinespace
high mountains & glide & 1\,538 & 4.59 [4.21, 4.98] & 4.38 [3.97, 4.80] & 3.19 [2.77, 3.60] & $-$1.20 [$-$1.82, $-$0.58] & $-$0.208 [$-$0.746, $+$0.330] \\
 & search & 574 & 3.98 [3.42, 4.53] & 3.84 [3.22, 4.46] & 4.70 [3.68, 5.71] & $+$0.858 [$+$0.125, $+$1.592] & $-$0.138 [$-$0.619, $+$0.342] \\
 & climb & 660 & 6.47 [5.48, 7.46] & 6.49 [5.06, 7.93] & 6.39 [5.19, 7.58] & $-$0.107 [$-$1.141, $+$0.926] & $+$0.0260 [$-$0.7933, $+$0.8452] \\
\addlinespace
all & glide & 1\,682 & 4.73 [4.46, 4.99] & 4.34 [4.02, 4.66] & 3.46 [3.14, 3.77] & $-$0.883 [$-$1.311, $-$0.455] & $-$0.388 [$-$0.792, $+$0.016] \\
 & search & 510 & 3.80 [3.54, 4.06] & 3.93 [3.53, 4.33] & 4.18 [3.64, 4.72] & $+$0.250 [$-$0.368, $+$0.867] & $+$0.131 [$-$0.257, $+$0.520] \\
 & climb & 1\,075 & 9.84 [7.84, 11.84] & 10.1 [7.0, 13.2] & 9.50 [7.41, 11.60] & $-$0.600 [$-$3.547, $+$2.348] & $+$0.258 [$-$1.824, $+$2.341] \\
\bottomrule
\end{tabular}
}
  \caption{Common threshold per phase and ground class, $\mu$ of $\Sl$, $\A$, $\G$ and the gaps, 90\% bands with the
  threshold re-chosen in every draw.}
  \end{table}
  \begin{figure}[H]
  \centering
  \includegraphics[width=\textwidth]{open_declared_dt5/phases_hill.pdf}
  \caption{Hill plots $\hat\mu(x_\mathrm{min})$ of $\Sl$, $\A$, $\G$ per phase, with the per-curve thresholds.}
  \end{figure}
\Obs Above the thresholds the tails are thin: \TailAboveMin{} to \TailAboveMax{} durations of $\G$, and $\mu$ between about \TailMuMin{} and \TailMuMax,
  steep tails. The climb bands are wide.
\end{step}

\begin{thebibliography}{5}\small
\bibitem{cameron2011} A. C. Cameron, J. B. Gelbach, D. L. Miller. Robust inference with multiway clustering.
  \emph{Journal of Business \& Economic Statistics} 29, 238--249 (2011).
\bibitem{kaplan1958} E. L. Kaplan, P. Meier. Nonparametric estimation from incomplete observations.
\emph{Journal of the American Statistical Association} 53(282):457--481, 1958. doi:10.1080/01621459.1958.10501452.
\bibitem{marshall1887} A. Marshall. Remedies for fluctuations of general prices. \emph{Contemporary Review}
51:355--375, 1887.
\bibitem{clauset2009} A. Clauset, C. R. Shalizi, M. E. J. Newman. Power-law distributions in empirical data.
\emph{SIAM Review} 51(4):661--703, 2009.
\bibitem{vilpellet2026} J. Vilpellet, A. Darmon, M. Benzaquen. From random walks to thermal rides: universal
anomalous transport in soaring flights. arXiv:2601.01293, 2026.
\end{thebibliography}

\end{document}
